\input{preamble}

% OK, start here.
%
\begin{document}

\title{Variétés}


\maketitle

\phantomsection
\label{section-phantom}

\tableofcontents

\section{Introduction}
\label{section-introduction}

\noindent
Dans ce chapitre, nous commençons à étudier les variétés et, plus
généralement, les schémas sur un corps. Une référence fondamentale est \cite{EGA}.








\section{Notation}
\label{section-notation}

\noindent
Tout au long de ce chapitre, nous utilisons la lettre $k$ pour désigner le corps de base.










\section{Variétés}
\label{section-varieties}

\noindent
Dans le projet Champs, nous adoptons la définition suivante
d'une variété.

\begin{definition}
\label{definition-variety}
Soit $k$ un corps. Une {\it variété} est un schéma $X$ sur $k$
qui est intègre et dont le morphisme structural
$X \to \Spec(k)$ est séparé et de type fini.
\end{definition}

\noindent
Cette définition présente l'inconvénient suivant. Soit $k'/k$
une extension de corps et soit $X$ une variété sur $k$. Le changement
de base $X_{k'} = X \times_{\Spec(k)} \Spec(k')$ n'est alors
pas nécessairement une variété sur $k'$. Ce phénomène sera étudié
en détail, dans un caunere plus général, dans les sections suivantes.
Le produit de deux variétés n'est pas nécessairement d variété non
plus (il s'agit en fait du même phénomène).
En voici un exemple.


\begin{example}
\label{example-product-not-a-variety}
Soit $k = \mathbf{Q}$. Posons $X = \Spec(\mathbf{Q}(i))$
et $Y = \Spec(\mathbf{Q}(i))$. Alors le produit
$X \times_{\Spec(k)} Y$ des variétés $X$ et $Y$
n'est pas d variété, car il est réductible. (Il est
isomorphe à la réunion disjointe de deux copies de $X$.)
\end{example}

\noindent
Toutefois, si le corps de base est algébriquement clos, le produit de
variétés est une variété. Cela résulte du chapitre sur l'algèbre, où
l'on traite des situations bien plus générales. Nous en donnons également
ici une démonstration directe et simple.

\begin{lemma}
\label{lemma-product-varieties}
\begin{slogan}
Les produits de variétés sont des variétés sur des corps algébriquement clos.
\end{slogan}
Soit $k$ un corps algébriquement clos.
Soient $X$ et $Y$ des variétés sur $k$.
Alors $X \times_{\Spec(k)} Y$ est d variété sur $k$.
\end{lemma}

\begin{proof}
Le morphisme $X \times_{\Spec(k)} Y \to \Spec(k)$ est de type
fini et séparé : c'est le composé des morphismes
$X \times_{\Spec(k)} Y \to Y \to \Spec(k)$, qui sont séparés
et de type fini ; voir Morphismes, lemmes
\ref{morphisms-lemma-base-change-finite-type} et
\ref{morphisms-lemma-composition-finite-type}, ainsi que
Schémas, lemme \ref{schemes-lemma-separated-permanence}.
Il reste à montrer que $X \times_{\Spec(k)} Y$ est intègre.
Soient $X = \bigcup_{i = 1, \ldots, n} U_i$ et
$Y = \bigcup_{j = 1, \ldots, m} V_j$ des recouvrements
ouverts affines finis. Il suffit que chacun des schémas
$U_i \times_{\Spec(k)} V_j$ soit intègre, d'après
Propriétés, lemmes \ref{properties-lemma-characterize-reduced},
\ref{properties-lemma-characterize-irreducible} et
\ref{properties-lemma-characterize-integral}.
On est ainsi ramené au cas affine.

\medskip\noindent
Le cas affine équivaut à l'énoncé algébrique suivant : si $A$ et
$B$ sont des anneaux intègres et des $k$-algèbres de type fini, alors $A
\otimes_k B$ est un anneau intègre. Raisonnons par l'absurde et supposons que
$$
(\sum\nolimits_{i = 1, \ldots, n} a_i \otimes b_i)
(\sum\nolimits_{j = 1, \ldots, m} c_j \otimes une_j) = 0
$$
dans $A \otimes_k B$, les deux facteurs étant non nuls dans $A
\otimes_k B$. On peut supposer que $b_1, \ldots, b_n$ sont
linéairement indépendants sur $k$ dans $B$, et de même pour
$d_1, \ldots, d_m$. On peut aussi supposer que $a_1$ et $c_1$
sont non nuls dans $A$. Ainsi, $D(a_1c_1) \subset \Spec(A)$ est
non vide. D'après le théorème des zéros de Hqubert
(Algèbre, théorème \ref{algebra-theorem-nullstellensatz}), il existe
un idéal maximal $\mathfrak m \in D(a_1c_1)$ tel que
$A/\mathfrak m = k$, puisque $k$ est algébriquement clos. Notons
$\overline{a}_i$ et $\overline{c}_j$ les classes de $a_i$ et $c_j$
dans $A/\mathfrak m = k$. L'égalité précédente devient
$$
(\sum\nolimits_{i = 1, \ldots, n} \overline{a}_i b_i)
(\sum\nolimits_{j = 1, \ldots, m} \overline{c}_j une_j) = 0
$$
ce qui contredit $\mathfrak m \in D(a_1c_1)$, l'indépendance linéaire
des famqules $b_1, \ldots, b_n$ et $une_1, \ldots, une_m$, ainsi que
le fait que $B$ est un anneau intègre.
\end{proof}






\section{Variétés et applications rationnelles}
\label{section-varieties-rational-maps}

\noindent
Soit $k$ un corps et soient $X$ et $Y$ des variétés sur $k$.
Par {\it application rationnelle de variétés de $X$ vers $Y$}, nous entendrons
une application rationnelle sur $\Spec(k)$ du schéma $X$ vers le schéma $Y$,
au sens de Morphismes, définition \ref{morphisms-definition-rational-map}.
Selon l'usage, l'expression «~application rationnelle de variétés~» ne
mentionne pas le corps de base (commun) des variétés, bien que, pour des
schémas généraux, nous distinguions les applications rationnelles des
applications rationnelles sur une base donnée.

\medskip\noindent
Le titre de cette section fait référence au théorème fondamental suivant.

\begin{theorem}
\label{theorem-varieties-rational-maps}
Soit $k$ un corps. La catégorie des variétés
et des applications rationnelles dominantes est équivalente
à la catégorie des extensions de corps de type fini $K/k$.
\end{theorem}

\begin{proof}
Soient $X$ et $Y$ des variétés de points génériques respectifs
$x \in X$ et $y \in Y$. Rappelons que les applications rationnelles
dominantes de $X$ vers $Y$ sont précisément les applications
rationnelles qui envoient $x$ sur $y$ (Morphismes,
définition \ref{morphisms-definition-dominant-rational} et discussion
qui suit). Toute application rationnelle dominante $X
\supset U \to Y$ induit donc un homomorphisme de corps de fonctions
$$
k(Y) = \kappa(y) = \mathcal{O}_{Y, y}
\longrightarrow
\mathcal{O}_{X, x} = \kappa(x) = k(X)
$$
Inversement, un tel homomorphisme de $k$-algèbres (automatiquement
local, puisque sa source et son but sont des corps) détermine de
manière unique d application rationnelle dominante, d'après Morphismes,
lemme \ref{morphisms-lemma-rational-map-finite-presentation}. On obtient ainsi
un foncteur pleinement fidèle. Il reste à voir que toute extension de
corps de type fini $K/k$ appartient à l'image essentielle. Comme $K/k$
est de type fini, qu existe d $k$-algèbre de type fini $A \subset K$
dont $K$ est le corps des fractions. Alors $X = \Spec(A)$ est d variété
de corps de fonctions $K$.
\end{proof}

\noindent
Soit $k$ un corps et soient $X$ et $Y$ des variétés sur $k$.
Nous dirons que {\it $X$ et $Y$ sont des variétés birationnelles}
si $X$ et $Y$ sont birationnels sur $\Spec(k)$, au sens de
Morphismes, définition \ref{morphisms-definition-birational}.
Selon l'usage, l'expression «~variétés birationnelles~» ne mentionne pas
le corps de base (commun) des variétés, bien que, pour des schémas
irréductibles généraux, nous distinguions le fait d'être birationnels
du fait de l'être sur une base donnée.

\begin{lemma}
\label{lemma-birational-varieties}
Soient $X$ et $Y$ des variétés sur un corps $k$.
Les conditions suivantes sont équivalentes :
\begin{enumerate}
\item $X$ et $Y$ sont des variétés birationnelles,
\item les corps de fonctions $k(X)$ et $k(Y)$ sont isomorphes,
\item il existe des ouverts non vides de $X$ et de $Y$ qui sont
isomorphes comme variétés,
\item il existe un ouvert $U \subset X$ et un morphisme
birationnel de variétés $U \to Y$.
\end{enumerate}
\end{lemma}

\begin{proof}
C'est un cas particulier de Morphismes, lemme
\ref{morphisms-lemma-criterion-birational-finite-presentation}.
\end{proof}







\section{Changement de corps et anneaux locaux}
\label{section-local-rings}

\noindent
Voici quelques résultats préliminaires sur le comportement des
anneaux locaux par extension du corps de base.

\begin{lemma}
\label{lemma-change-fields-flat}
Soit $K/k$ d extension de corps. Soit $X$ un schéma sur $k$
et posons $Y = X_K$. Si $y \in Y$ a pour image $x \in X$, alors
\begin{enumerate}
\item $\mathcal{O}_{X, x} \to \mathcal{O}_{Y, y}$ est un
homomorphisme local fidèlement plat,
\item si $\mathfrak p_0 = \Ker(\kappa(x) \otimes_k K \to \kappa(y))$,
alors $\kappa(y) = \kappa(\mathfrak p_0)$,
\item $\mathcal{O}_{Y, y} =
(\mathcal{O}_{X, x} \otimes_k K)_\mathfrak p$, où
$\mathfrak p \subset \mathcal{O}_{X, x} \otimes_k K$ est l'image
réciproque de $\mathfrak p_0$,
\item on a
$\mathcal{O}_{Y, y}/\mathfrak m_x\mathcal{O}_{Y, y} =
(\kappa(x) \otimes_k K)_{\mathfrak p_0}$.
\end{enumerate}
\end{lemma}

\begin{proof}
On peut supposer $X = \Spec(A)$ affine. Alors
$Y = \Spec(A \otimes_k K)$. Puisque $K$ est plat sur $k$,
$A \to A \otimes_k K$ est plat. Il en va donc de même de $Y \to X$,
et la première assertion résulte aussi d'Algèbre, lemme
\ref{algebra-lemma-local-flat-ff}. La deuxième résulte de
Schémas, lemme \ref{schemes-lemma-points-fibre-product}. Le point
$y$ correspond à un idéal premier $\mathfrak q \subset A \otimes_k K$
et $x$ à $\mathfrak r = A \cap \mathfrak q$. Alors $\mathfrak p_0$
est le noyau de l'homomorphisme induit
$\kappa(\mathfrak r) \otimes_k K \to \kappa(\mathfrak q)$.
L'homomorphisme d'anneaux locaux est
$$
A_\mathfrak r \longrightarrow (A \otimes_k K)_\mathfrak q
$$
Il se factorise par
$A_\mathfrak r \otimes_k K =
(A \otimes_k K)_{\mathfrak r}$ sous la forme
$$
A_\mathfrak r \longrightarrow A_\mathfrak r \otimes_k K
\longrightarrow (A \otimes_k K)_\mathfrak q
$$
où la deuxième flèche est la localisation en un certain idéal
premier. Celui-ci est l'image réciproque de $\mathfrak p_0$
(nous omettons les détails), ce qui prouve (3). Pour (4), on utilise
(3) et l'exactitude de la localisation et du foncteur $- \otimes_k K$.
\end{proof}

\begin{lemma}
\label{lemma-change-fields-algebraic-dim}
Avec les notations du lemme \ref{lemma-change-fields-flat},
si $X$ est localement de type fini sur $k$, alors
$$
\dim(\mathcal{O}_{Y, y}/\mathfrak m_x\mathcal{O}_{Y, y}) =
\text{trdeg}_k(\kappa(x)) - \text{trdeg}_K(\kappa(y)) =
\dim(\mathcal{O}_{Y, y}) - \dim(\mathcal{O}_{X, x})
$$
\end{lemma}

\begin{proof}
Il s'agit d'une reformulation d'Algèbre, lemme
\ref{algebra-lemma-inequalities-under-field-extension}.
\end{proof}

\begin{lemma}
\label{lemma-change-fields-algebraic-unramified}
Avec les notations du lemme
\ref{lemma-change-fields-flat}, supposons $X$ localement de type fini
sur $k$, $\dim(\mathcal{O}_{X, x}) = \dim(\mathcal{O}_{Y, y})$
et $\kappa(x) \otimes_k K$ réduit (par exemple si
$\kappa(x)/k$ est séparable ou si $K/k$ est séparable). Alors
$\mathfrak m_x \mathcal{O}_{Y, y} = \mathfrak m_y$.
\end{lemma}

\begin{proof}
(L'assertion entre parenthèses résulte d'Algèbre, lemme
\ref{algebra-lemma-separable-extension-preserves-reducedness}.)
En combinant les lemmes \ref{lemma-change-fields-flat} et
\ref{lemma-change-fields-algebraic-dim}, on voit que
$\mathcal{O}_{Y, y}/\mathfrak m_x \mathcal{O}_{Y, y}$ a une
dimension $0$ et est réduit. C'est donc un corps.
\end{proof}





\section{Schémas géométriquement réduits}
\label{section-geometrically-reduced}

\noindent
Si $X$ est un schéma réduit sur un corps, il peut arriver que
$X$ devienne non réduit après l'extension du corps de base.
Cela ne se produit pas pour les schémas géométriquement réduits.

\begin{definition}
\label{definition-geometrically-reduced}
Soit $k$ un corps.
Soit $X$ un schéma sur $k$.
\begin{enumerate}
\item Soit $x \in X$. On dit que $X$ est
{\it géométriquement réduit en $x$} si, pour toute
extension de corps $k'/k$ et tout point $x' \in X_{k'}$
au-dessus de $x$, l'anneau local $\mathcal{O}_{X_{k'}, x'}$
est réduit.
\item On dit que $X$ est {\it géométriquement réduit} sur $k$
si $X$ est géométriquement réduit en tout point de $X$.
\end{enumerate}
\end{definition}

\noindent
Cette définition peut paraître quelque peu mystérieuse, mais elle
coïncide en réalité avec la notion étudiée dans le chapitre sur
l'algèbre. Les résultats élémentaires qui suivent précisent ce lien.
Voir aussi le lemme
\ref{lemma-geometrically-reduced-dense-smooth-open}, d'après lequel
«~géométriquement réduit~» équivaut à «~génériquement lisse~» pour les variétés sur un corps.

\begin{lemma}
\label{lemma-geometrically-reduced-at-point}
\begin{slogan}
La propriété d'être géométriquement réduit se vérifie sur les anneaux locaux.
\end{slogan}
Soit $k$ un corps.
Supposons que $X$ soit un schéma
sur $k$. Soit $x \in X$. Les
énoncés suivants sont équivalents
\begin{enumerate}
\item $X$ est géométriquement réduit en $x$, et
\item l’anneau $\mathcal{O}_{X, x}$ est
géométriquement réduit sur $k$ (voir Algèbre,
définition \ref{algebra-definition-geometrically-reduced}).
\end{enumerate}
\end{lemma}

\begin{proof}
Supposons (1). Cela implique en particulier que
$\mathcal{O}_{X, x}$ est réduit. Supposons que $k'/k$ soit une
extension de corps finie purement inséparable. Considérons
l'anneau $\mathcal{O}_{X, x} \otimes_k k'$. D'après l'algèbre,
le lemme \ref{algebra-lemma-p-ring-map}, son spectre est le
même que le spectre de $\mathcal{O}_{X, x}$. Par conséquent,
c'est également un anneau local (Algèbre, lemme
\ref{algebra-lemma-characterize-local-ring}). Donc, il existe un
point unique $x' \in X_{k'}$ au-dessus de $x$ et de
$\mathcal{O}_{X_{k'}, x'} \cong \mathcal{O}_{X, x} \otimes_k k'$, voir
le lemme \ref{lemma-change-fields-flat}. Par hypothèse, il s'agit d'un
anneau réduit. Ainsi, nous déduisons (2) d'après Algèbre, lemme
\ref{algebra-lemma-geometrically-reduced-finite-purely-inseparable-extension}.

\medskip\noindent
Supposons (2). Supposons que $k'/k$ soit une extension
de corps. Puisque $\Spec(k') \to \Spec(k)$ est surjectif,
$X_{k'} \to X$ est aussi surjectif (Morphismes, lemme
\ref{morphisms-lemma-base-change-surjective}). Supposons
que $x' \in X_{k'}$ soit un point quelconque situé
au-dessus de $x$. L'anneau local $\mathcal{O}_{X_{k'}, x'}$
est une localisation de l'anneau $\mathcal{O}_{X, x}
\otimes_k k'$ d'après le lemme \ref{lemma-change-fields-flat}.
Par conséquent, il est réduit par hypothèse et (1) est prouvé.
\end{proof}

\noindent
La notion n'est pas intéressante en caractéristique zéro.

\begin{lemma}
\label{lemma-perfect-reduced}
Supposons que $X$ soit un schéma sur un corps parfait $k$ (par
exemple \ $k$ a une caractéristique nulle). Soit $x \in X$. Si
$\mathcal{O}_{X, x}$ est réduit, alors $X$ est géométriquement réduit
en $x$. Si $X$ est réduit, alors $X$ est géométriquement réduit sur $k$.
\end{lemma}

\begin{proof}
La première affirmation découle du lemme
\ref{lemma-geometrically-reduced-at-point} et de
l'Algèbre, du lemme
\ref{algebra-lemma-separable-extension-preserves-reducedness} et de
la définition d'un corps parfait (Algèbre, définition
\ref{algebra-definition-perfect}). La seconde affirmation découle de la première.
\end{proof}

\begin{lemma}
\label{lemma-geometrically-reduced}
Supposons que $k$ soit un corps de caractéristique $p > 0$. Supposons
que $X$ soit un schéma sur $k$. Les énoncés suivants sont équivalents
\begin{enumerate}
\item $X$ est géométriquement réduit, \item
$X_{k'}$ est réduit pour toute extension de
corps $k'/k$, \item $X_{k'}$ est réduit pour
toute extension de corps finie purement
inséparable $k'/k$, \item $X_{k^{1/p}}$ est
réduit, \item $X_{k^{perf}}$ est réduit, \item
$X_{\bar k}$ est réduit, \item pour tout ouvert
affine $U \subset X$ l'anneau $\mathcal{O}_X(U)$
est géométriquement réduit (voir Algèbre,
définition \ref{algebra-definition-geometrically-reduced}).
\end{enumerate}
\end{lemma}

\begin{proof}
Supposons (1). Alors pour toute extension de corps $k'/k$ et
chaque point $x' \in X_{k'}$, l'anneau local de $X_{k'}$ en $x'$
est réduit. En d'autres termes, $X_{k'}$ est réduit. D'où (2).

\medskip\noindent
Supposons (2). Supposons que $U \subset X$ soit un ouvert affine.
Alors, pour toute extension de corps $k'/k$, le schéma $X_{k'}$
est réduit, donc $U_{k'} = \Spec(\mathcal{O}(U)\otimes_k k')$ est
réduit, donc $\mathcal{O}(U)\otimes_k k'$ est réduit (voir
Propriétés, section \ref{properties-section-integral}). En d'autres
termes, $\mathcal{O}(U)$ est géométriquement réduit, donc (7) est vérifié.

\medskip\noindent
Supposons (7). Pour toute extension de corps $k'/k$, le
changement de base $X_{k'}$ s'obtient par recollement des spectres
des anneaux $\mathcal{O}_X(U) \otimes_k k'$ où $U$ est un
ouvert affine dans $X$ (voir Schémas, section
\ref{schemes-section-fibre-products}). Donc $X_{k'}$ est réduit. Ainsi (1) est vrai.

\medskip\noindent
Cela prouve que (1), (2) et (7) sont équivalents. Ceux-ci sont
équivalents à (3), (4), (5) et (6) car nous pouvons appliquer
Algèbre, lemme
\ref{algebra-lemma-geometrically-reduced-finite-purely-inseparable-extension} à $\mathcal{O}_X(U)$ pour $U \subset X$ ouvert affine.
\end{proof}

\begin{lemma}
\label{lemma-check-only-finite-inseparable-extensions}
Supposons que $k$ soit un corps de caractéristique $p > 0$. Supposons que $X$
soit un schéma sur $k$. Soit $x \in X$. Les énoncés suivants sont équivalents
\begin{enumerate}
\item $X$ est géométriquement réduit en $x$, \item
$\mathcal{O}_{X_{k'}, x'}$ est réduit pour chaque
extension de corps finie purement inséparable $k'$ de
$k$ et $x' \in X_{k'}$ le point unique au-dessus de $x$,
\item $\mathcal{O}_{X_{k^{1/p}}, x'}$ est réduit pour
$x' \in X_{k^{1/p}}$ le point unique au-dessus de $x$,
et \item $\mathcal{O}_{X_{k^{perf}}, x'}$ est réduit pour
$x' \in X_{k^{perf}}$ le point unique au-dessus de $x$.
\end{enumerate}
\end{lemma}

\begin{proof}
Notons que si $k'/k$ est purement inséparable, alors $X_{k'}
\to X$ induit un homéomorphisme sur les espaces
topologiques sous-jacents, voir Algèbre, lemme
\ref{algebra-lemma-p-ring-map}. D'où l'unicité de $x'$
au-dessus de $x$ mentionnée dans l'énoncé. De plus, dans ce cas
$\mathcal{O}_{X_{k'}, x'} = \mathcal{O}_{X, x} \otimes_k k'$.
Par conséquent, le lemme découle du lemme
\ref{lemma-geometrically-reduced-at-point} ci-dessus et du lemme
\ref{algebra-lemma-geometrically-reduced-finite-purely-inseparable-extension} d'Algèbre.
\end{proof}

\begin{lemma}
\label{lemma-geometrically-reduced-upstairs}
Soit $k$ un corps. Supposons que $X$
soit un schéma sur $k$. Supposons que $k'/k$ soit
une extension de corps. Supposons que $x \in X$ soit
un point, et soit $x' \in X_{k'}$ un point situé
au-dessus de $x$. Les énoncés suivants sont équivalents
\begin{enumerate}
\item $X$ est géométriquement réduit en $x$,
\item $X_{k'}$ est géométriquement réduit en $x'$.
\end{enumerate}
En particulier, $X$ est géométriquement réduit sur $k$ si et
seulement si $X_{k'}$ est géométriquement réduit sur $k'$.
\end{lemma}

\begin{proof}
Il est clair que (1) implique (2). Supposons (2). Supposons que
$k''/k$ soit une extension de corps finie purement inséparable et
soit $x'' \in X_{k''}$ un point au-dessus de $x$ (en fait, il
est unique). Choisissons une extension de corps commune $k'''/k$ de
$k''/k$ et $k'/k$, voir corps, lemme
\ref{fields-lemma-common-extension-field}. Choisissons un point $x''' \in X_{k'''}$
au-dessus à la fois de $x'$ et $x''$. Considérons le morphisme d'anneaux locaux
$$
\mathcal{O}_{X_{k''}, x''} \longrightarrow \mathcal{O}_{X_{k'''}, x'''}.
$$
Ceci est un homomorphisme local plat d'anneaux et donc fidèlement
plat. Par (2), on voit que l'anneau local à droite est réduit.
Ainsi, d'après L'Algèbre, lemme
\ref{algebra-lemma-descent-reduced}, nous concluons que
$\mathcal{O}_{X_{k''}, x''}$ est réduit. Ainsi, d'après le lemme
\ref{lemma-check-only-finite-inseparable-extensions}, nous concluons que $X$ est géométriquement réduit en $x$.
\end{proof}

\begin{lemma}
\label{lemma-geometrically-reduced-any-base-change}
Soit $k$ un corps. Supposons que $X$, $Y$ soient des schémas sur $k$.
\begin{enumerate}
\item Si $X$ est géométriquement réduit en $x$, et
$Y$ est réduit, alors $X \times_k Y$ est réduit
en tout point au-dessus de $x$. \item Si $X$ est
géométriquement réduit sur $k$ et $Y$ est réduit,
alors $X \times_k Y$ est réduit. \item Si $X$ et
$Y$ sont géométriquement réduits sur $k$, alors $X
\times_k Y$ est géométriquement réduit. \item Si
$k$ est parfait et $X$ et $Y$ sont réduits, alors
$X \times_k Y$ est réduit. \item Ajouter plus ici.
\end{enumerate}
\end{lemma}

\begin{proof}
Pour démontrer (1), on combine le lemme
\ref{lemma-geometrically-reduced-at-point} avec le lemme
d'Algèbre
\ref{algebra-lemma-geometrically-reduced-any-reduced-base-change}.
Pour démontrer (2), on combine le lemme
\ref{lemma-geometrically-reduced} avec le lemme d'Algèbre
\ref{algebra-lemma-geometrically-reduced-any-reduced-base-change}.
Pour démontrer (3), notons que $(X \times_k Y)_{\overline{k}} =
X_{\overline{k}} \times_{\overline{k}} Y_{\overline{k}}$ et utilisez
(2) ainsi que le lemme \ref{lemma-geometrically-reduced}. Pour démontrer
(4), utilisez (3) combiné avec le lemme \ref{lemma-perfect-reduced}.
\end{proof}

\begin{lemma}
\label{lemma-generic-points-geometrically-reduced}
Soit $k$ un corps.
Soit $X$ un schéma sur $k$.
\begin{enumerate}
\item Si $x' \leadsto x$ est une spécialisation et $X$
est géométriquement réduit en $x$, alors $X$ est
géométriquement réduit en $x'$. \item Si $x \in X$ tel
que (a) $\mathcal{O}_{X, x}$ est réduit, et (b) pour
chaque spécialisation $x' \leadsto x$ où $x'$ est un
point générique d'une composante irréductible de $X$, le
schéma $X$ est géométriquement réduit en $x'$, alors $X$
est géométriquement réduit en $x$. \item Si $X$ est
réduit et géométriquement réduit en tous les points
génériques des composantes irréductibles de $X$, alors $X$
est géométriquement réduit. \item Si $X$ est une variété sur
$k$, alors $X$ est géométriquement réduit si et seulement si
son corps des fonctions est une extension séparable de $k$.
\end{enumerate}
\end{lemma}

\begin{proof}
La partie (1) découle du lemme
\ref{lemma-geometrically-reduced-at-point} et du fait que
si $A$ est une $k$-algèbre géométriquement réduite, alors
$S^{-1}A$ est une $k$-algèbre géométriquement réduite pour
tout sous-ensemble multiplicatif $S$ de $A$, voir Algèbre,
lemme \ref{algebra-lemma-geometrically-reduced-permanence}.

\medskip\noindent
Soit $A = \mathcal{O}_{X, x}$. Les hypothèses (a) et (b) de (2)
impliquent que $A$ est réduit, et que $A_{\mathfrak q}$ est
géométriquement réduit sur $k$ pour chaque idéal premier minimal
$\mathfrak q$ de $A$. Par conséquent, $A$ est géométriquement
réduit sur $k$, voir Algèbre, lemme
\ref{algebra-lemma-generic-points-geometrically-reduced}. Ainsi, $X$ est
géométriquement réduit en $x$, voir lemme \ref{lemma-geometrically-reduced-at-point}.

\medskip\noindent
La partie (3) découle trivialement de la partie (2).
Enfin, (4) découle de (3) par Algèbre, lemme
\ref{algebra-lemma-characterize-separable-field-extensions}.
\end{proof}

\begin{lemma}
\label{lemma-Noetherian-geometrically-reduced-at-point}
Soit $k$ un corps. Supposons
que $X$ soit un schéma sur $k$. Soit $x \in
X$. Supposons que $X$ soit localement
noethérien et géométriquement réduit en $x$.
Alors il existe un voisinage ouvert $U \subset
X$ de $x$ qui est géométriquement réduit sur $k$.
\end{lemma}

\begin{proof}
Supposons que $X$ soit localement noethérien et
géométriquement réduit en $x$. D'après le lemme des
propriétés
\ref{properties-lemma-ring-affine-open-injective-local-ring}, nous
pouvons trouver un voisinage affine ouvert $U \subset X$ de $x$
tel que $R = \mathcal{O}_X(U) \to \mathcal{O}_{X, x}$ soit
injectif. D'après le lemme \ref{lemma-geometrically-reduced-at-point},
l'hypothèse signifie que $\mathcal{O}_{X, x}$ est géométriquement
réduit sur $k$. D'après Algèbre, lemme
\ref{algebra-lemma-subalgebra-separable}, cela implique que $R$ est
géométriquement réduit sur $k$, ce qui implique à son tour que $U$ est géométriquement réduit.
\end{proof}

\begin{example}
\label{example-not-geometrically-reduced}
Soit $k = \mathbf{F}_p(s, t)$, c'est-à-dire une
extension purement transcendantale du corps premier.
Considérons la variété $X = \Spec(k[x, y]/(1 + sx^p +
ty^p))$. Soit $k'/k$ une extension quelconque telle que $s$
et $t$ possèdent toutes deux une racine $p$-ième dans $k'$.
Alors le changement de base $X_{k'}$ n'est pas réduit. En
effet, l'anneau $k'[x, y]/(1 + s x^p + ty^p)$ contient
l'élément $1 + s^{1/p}x + t^{1/p}y$, dont la puissance $p$-ième est
nulle alors que lui-même ne l'est pas (l'idéal $(1 + sx^p + ty^p)$
ne contient certainement aucun élément non nul de degré $< p$).
\end{example}

\begin{lemma}
\label{lemma-finite-extension-geometrically-reduced}
Soit $k$ un corps. Soit $X \to \Spec(k)$ un morphisme
localement de type fini. Supposons que $X$ n'ait il'un nombre
fini de composantes irréductibles. Il existe alors une
extension finie purement inséparable $k'/k$ telle que
$(X_{k'})_{red}$ soit géométriquement réduit sur $k'$.
\end{lemma}

\begin{proof}
On peut remplacer $X$ par sa réduction $X_{red}$, et donc
supposer $X$ réduit et localement de type fini sur $k$. Soient
$x_1, \ldots, x_n \in X$ les points génériques de ses composantes
irréductibles. Pour toute extension algébrique purement
inséparable $k'/k$, le morphisme
$(X_{k'})_{red} \to X$ est un
homéomorphisme ; voir Algèbre, lemme
\ref{algebra-lemma-p-ring-map}. Les points $x'_1, \ldots, x'_n$
au-dessus de $x_1, \ldots, x_n$ sont donc les points génériques
des composantes irréductibles de $(X_{k'})_{red}$. Comme $X$ est
réduit, les anneaux locaux $K_i = \mathcal{O}_{X, x_i}$ sont des
corps ; voir Algèbre, lemme
\ref{algebra-lemma-minimal-prime-reduced-ring}. Puisque $X$ est
localement de type fini sur $k$, les extensions de corps $K_i/k$
sont de type fini. Enfin, les anneaux locaux
$\mathcal{O}_{(X_{k'})_{red}, x'_i}$ sont les corps
$(K_i \otimes_k k')_{red}$. D'après Algèbre, lemme
\ref{algebra-lemma-make-separable}, il existe une extension finie
purement inséparable $k'/k$ telle que les corps
$(K_i \otimes_k k')_{red}$ soient séparables sur $k'$. Chacun est
donc géométriquement réduit sur $k'$, d'après
Algèbre, lemme
\ref{algebra-lemma-characterize-separable-field-extensions}.
Le lemme \ref{lemma-generic-points-geometrically-reduced}, partie (3),
montre alors que $(X_{k'})_{red}$ est géométriquement réduit.
\end{proof}







\section{Schémas géométriquement connexes}
\label{section-geometrically-connected}

\noindent
Si $X$ est un schéma connexe sur un corps, il peut arriver
que $X$ devienne non connexe après l'extension du corps de
base. Cela n'arrive pas pour les schémas géométriquement connexes.

\begin{definition}
\label{definition-geometrically-connected}
Soit $X$ un schéma sur le corps $k$. On dit que
$X$ est {\it géométriquement connexe} sur $k$ si le schéma
$X_{k'}$ est connexe pour toute extension de corps $k'$ de $k$.
\end{definition}

\noindent
Par convention, un espace topologique connexe est non vide ; par
conséquent a fortiori, les schémas géométriquement connexes sont non
vides. Voici un exemple d'une variété qui n'est pas géométriquement connexe.

\begin{example}
\label{example-not-geometrically-irreducible}
Soit $k = \mathbf{Q}$. Le schéma $X =
\Spec(\mathbf{Q}(i))$ est une variété sur
$\Spec(\mathbf{Q})$. Mais le changement de base
$X_{\mathbf{C}}$ est le spectre de $\mathbf{C}
\otimes_{\mathbf{Q}} \mathbf{Q}(i) \cong \mathbf{C} \times \mathbf{C}$
qui est l'union disjointe de deux copies de $\Spec(\mathbf{C})$. Donc
en fait, c'est un exemple d'une variété non géométriquement connexe.
\end{example}

\begin{lemma}
\label{lemma-geometrically-connected-check-after-extension}
Supposons que $X$ soit un schéma sur le corps $k$.
Supposons que $k'/k$ soit une extension de corps.
Alors $X$ est géométriquement connexe sur $k$ si et
seulement si $X_{k'}$ est géométriquement connexe sur $k'$.
\end{lemma}

\begin{proof}
Si $X$ est géométriquement connexe sur $k$, alors il est clair
que $X_{k'}$ est géométriquement connexe sur $k'$. Pour la
réciproque, pour toute extension de corps $k''/k$, il existe
une extension de corps commune $k'''/k'$ et $k'''/k''$, voir
corps, lemme \ref{fields-lemma-common-extension-field}. Comme
le morphisme $X_{k'''} \to X_{k''}$ est surjectif (en tant que
changement de base d'un morphisme surjectif entre spectres de
corps), on voit que la connexité de $X_{k'''}$ implique la
connexité de $X_{k''}$. Ainsi, si $X_{k'}$ est géométriquement
connexe sur $k'$, alors $X$ est géométriquement connexe sur $k$.
\end{proof}

\begin{lemma}
\label{lemma-bijection-connected-components}
Soit $k$ un corps. Supposons
que $X$, $Y$ soient des schémas sur $k$.
Supposons que $X$ soit géométriquement
connexe sur $k$. Alors le morphisme de projection
$$
p : X \times_k Y \longrightarrow Y
$$
induit une bijection entre les composantes connexes.
\end{lemma}

\begin{proof}
Les fibres schématiques de $p$ sont connexes,
puisque ce sont des changements de base du
schéma géométriquement connexe $X$ par
extensions de corps. De plus, les fibres
schématiques sont homéomorphes aux fibres
ensemblistes, voir Schémas, lemme
\ref{schemes-lemma-fibre-topological}. Par morphismes,
lemme
\ref{morphisms-lemma-scheme-over-field-universally-open},
l'application $p$ est ouverte. Ainsi, nous pouvons appliquer Topologie,
lemme \ref{topology-lemma-connected-fibres-connected-components} pour conclure.
\end{proof}

\begin{lemma}
\label{lemma-affine-geometrically-connected}
Soit $k$ un corps. Supposons que $A$
soit une $k$-algèbre. Alors $X = \Spec(A)$ est
géométriquement connexe sur $k$ si et seulement si $A$
est géométriquement connexe sur $k$ (voir Algèbre,
définition \ref{algebra-definition-geometrically-connected} ).
\end{lemma}

\begin{proof}
Immédiat d'après les définitions.
\end{proof}

\begin{lemma}
\label{lemma-separably-closed-field-connected-components}
Supposons que $k'/k$ soit une extension de corps.
Soit $X$ un schéma sur $k$. Supposons
que $k$ soit séparablement algébriquement clos. Alors
le morphisme $X_{k'} \to X$ induit une bijection des
composantes connexes. En particulier, $X$ est
géométriquement connexe sur $k$ si et seulement si $X$ est connexe.
\end{lemma}

\begin{proof}
Puisque $k$ est séparablement algébriquement clos, nous
voyons que $k'$ est géométriquement connexe sur $k$, voir
Algèbre, lemme
\ref{algebra-lemma-separably-closed-connected-implies-geometric}.
Par conséquent, $Z = \Spec(k')$ est géométriquement connexe sur
$k$ d'après le lemme \ref{lemma-affine-geometrically-connected}
ci-dessus. Puisque $X_{k'} = Z \times_k X$, le résultat est un cas
particulier du lemme \ref{lemma-bijection-connected-components}.
\end{proof}

\begin{lemma}
\label{lemma-characterize-geometrically-connected}
Soit $k$ un corps. Supposons que $X$
soit un schéma sur $k$. Supposons que $\overline{k}$
soit une clôture algébrique séparable de $k$.
Alors $X$ est géométriquement connexe si et seulement
si le changement de base $X_{\overline{k}}$ est connexe.
\end{lemma}

\begin{proof}
Supposons que $X_{\overline{k}}$ soit connexe. Supposons
que $k'/k$ soit une extension de corps. Il existe une
extension de corps $\overline{k}'/\overline{k}$ telle que
$k'$ se plonge dans $\overline{k}'$ comme extension de
$k$. D'après le lemme
\ref{lemma-separably-closed-field-connected-components}, on voit que
$X_{\overline{k}'}$ est connexe. Puisque $X_{\overline{k}'} \to X_{k'}$
est surjectif, nous concluons que $X_{k'}$ est connexe comme voulu.
\end{proof}

\begin{lemma}
\label{lemma-descend-open}
Soit $k$ un corps. Supposons
que $X$ soit un schéma sur $k$. Supposons
que $A$ soit une $k$-algèbre. Supposons que
$V \subset X_A$ soit un ouvert
quasi-compact. Alors il existe une
$k$-subalgèbre de type fini $A' \subset A$ et un ouvert
quasi-compact $V' \subset X_{A'}$ tels que $V = V'_A$.
\end{lemma}

\begin{proof}
Nous remarquons que si $X$ est également quasi-séparé, cela
découle de Limites, lemme \ref{limits-lemma-descend-opens}.
Choisissons un nombre fini d'ouverts affines $U_1, \ldots, U_n$
de $X$ tels que $V \subset \bigcup U_{i,
A}$. Disons $U_i = \Spec(R_i)$. Puisque $V$ est quasi-compact,
nous pouvons trouver un nombre fini d'éléments $f_{ij} \in R_i
\otimes_k A$, $j = 1, \ldots, n_i$ tels que $V = \bigcup_i
\bigcup_{j = 1, \ldots, n_i} D(f_{ij})$ où $D(f_{ij}) \subset U_{i,
A}$ est l'ouvert principal correspondant. (Nous ne prétendons pas
que $V \cap U_{i, A}$ est l'union des $D(f_{ij})$, $j = 1, \ldots,
n_i$.) Il est clair que nous pouvons trouver une $k$-sous-algèbre
de type fini $A' \subset A$ telle que $f_{ij}$ soit l'image de
certains $f_{ij}' \in R_i \otimes_k A'$. Posons $V' = \bigcup
D(f_{ij}')$ qui est un ouvert quasi-compact de $X_{A'}$. Notons $\pi
: X_A \to X_{A'}$ le morphisme canonique. Nous avons $\pi(V) \subset
V'$ comme $\pi(D(f_{ij})) \subset D(f_{ij}')$. Si $x \in X_A$ avec
$\pi(x) \in V'$, alors $\pi(x) \in D(f_{ij}')$ pour certains $i, j$
et on voit que $x \in D(f_{ij})$ en tant que $f_{ij}'$ s'envoie
vers $f_{ij}$. Ainsi, on voit que $V = \pi^{-1}(V')$ comme voulu.
\end{proof}

\noindent
Soit $k$ un corps.
Supposons que $\overline{k}/k$ soit
une extension galoisienne
(éventuellement infinie). Par
exemple, $\overline{k}$ pourrait
être la clôture algébrique séparable
de $k$. Pour tout $\sigma \in
\text{Gal}(\overline{k}/k)$, nous
obtenons un automorphisme
correspondant $ \Spec(\sigma) :
\Spec(\overline{k}) \longrightarrow
\Spec(\overline{k}) $. Notons que
$\Spec(\sigma) \circ \Spec(\tau) = \Spec(\tau \circ
\sigma)$. Par conséquent, nous obtenons une action
$$
\text{Gal}(\overline{k}/k)^{opp} \times \Spec(\overline{k})
\longrightarrow
\Spec(\overline{k})
$$
du groupe opposé sur le schéma
$\Spec(\overline{k})$. Supposons que $X$ soit un
schéma sur $k$. Puisque $X_{\overline{k}} =
\Spec(\overline{k}) \times_{\Spec(k)} X$ par définition,
on voit que l'action ci-dessus induit une action canonique
\begin{equation}
\label{equation-galois-action-base-change-kbar}
\text{Gal}(\overline{k}/k)^{opp} \times X_{\overline{k}}
\longrightarrow
X_{\overline{k}}.
\end{equation}

\begin{lemma}
\label{lemma-Galois-action-quasi-compact-open}
Soit $k$ un corps. Supposons que $X$ soit un
schéma sur $k$. Supposons que $\overline{k}$ soit une
extension de Galois (éventuellement infinie) de $k$. Supposons que
$V \subset X_{\overline{k}}$ soit un ouvert quasi-compact. Alors
\begin{enumerate}
\item il existe une sous-extension finie
$\overline{k}/k'/k$ et un ouvert quasi-compact $V'
\subset X_{k'}$ tels que $V = (V')_{\overline{k}}$, \item il
existe un sous-groupe ouvert $H \subset
\text{Gal}(\overline{k}/k)$ tels que $\sigma(V) = V$ pour tout $\sigma \in H$.
\end{enumerate}
\end{lemma}

\begin{proof}
D'après le lemme \ref{lemma-descend-open}, il existe une sous-extension
finie $k'/k \subset \overline{k}$ et un ouvert $V' \subset X_{k'}$ dont
l'image réciproque est $V$. Cela prouve (1). Puisque $\text{Gal}(\overline{k}/k')$ est
ouvert dans $\text{Gal}(\overline{k}/k)$, la partie (2) est également claire.
\end{proof}

\begin{lemma}
\label{lemma-closed-fixed-by-Galois}
Soit $k$ un corps. Supposons que $\overline{k}/k$ soit une
extension de Galois (éventuellement infinie). Supposons que $X$ soit un schéma
sur $k$. Soit $\overline{T} \subset X_{\overline{k}}$ vérifiant les propriétés suivantes
\begin{enumerate}
\item $\overline{T}$ est un sous-ensemble fermé de
$X_{\overline{k}}$, \item pour chaque $\sigma \in
\text{Gal}(\overline{k}/k)$ nous avons $\sigma(\overline{T}) = \overline{T}$.
\end{enumerate}
Alors il existe un sous-ensemble fermé $T \subset X$ dont
l'image réciproque dans $X_{\overline{k}}$ est $\overline{T}$.
\end{lemma}

\begin{proof}
Ce lemme se réduit immédiatement au cas où $X = \Spec(A)$
est affine. Dans ce cas, soit $\overline{I} \subset A
\otimes_k \overline{k}$ l'idéal radical correspondant à
$\overline{T}$. L'hypothèse (2) implique que
$\sigma(\overline{I}) = \overline{I}$ pour tout $\sigma \in
\text{Gal}(\overline{k}/k)$. Choisissons $x \in \overline{I}$. Il
existe une extension galoisienne finie $k'/k$ contenue dans
$\overline{k}$ telle que $x \in A \otimes_k k'$. Posons $G = \text{Gal}(k'/k)$. Posons
$$
P(T) = \prod\nolimits_{\sigma \in G} (T - \sigma(x)) \in (A \otimes_k k')[T]
$$
Il est clair que $P(T)$ est unitaire et est en fait un élément
de $(A \otimes_k k')^G[T] = A[T]$ (selon la théorie de Galois de
base). De plus, si nous écrivons $P(T) = T^d + a_1T^{d - 1} +
\ldots + a_d$ alors on voit que $a_i \in I := A \cap
\overline{I}$. En combinant $P(x) = 0$ et $a_i \in I$ nous
trouvons $x^d = - a_1 x^{d - 1} - \ldots - a_d \in I(A \otimes_k
\overline{k})$. Ainsi, $x$ est contenu dans le radical de $I(A
\otimes_k \overline{k})$. Par conséquent, $\overline{I}$ est le radical
de $I(A \otimes_k \overline{k})$; le choix $T = V(I)$ convient.
\end{proof}

\begin{lemma}
\label{lemma-characterize-geometrically-disconnected}
Soit $k$ un corps.
Supposons que $X$ soit un schéma sur
$k$. Les conditions suivantes sont équivalentes
\begin{enumerate}
\item $X$ est géométriquement connexe,
\item pour toute extension de corps finie
séparable $k'/k$ le schéma $X_{k'}$ est connexe.
\end{enumerate}
\end{lemma}

\begin{proof}
Il découle immédiatement de la définition que (1)
implique (2). Supposons que $X$ ne soit pas
géométriquement connexe. Supposons que $k \subset
\overline{k}$ soit une clôture algébrique séparable de $k$.
D'après le lemme
\ref{lemma-characterize-geometrically-connected}, il s'ensuit que
$X_{\overline{k}}$ est non connexe. Disons $X_{\overline{k}} = \overline{U} \amalg
\overline{V}$ avec $\overline{U}$ et $\overline{V}$ ouverts, fermés et non vides.

\medskip\noindent
Supposons que $W \subset X$ soit un ouvert quasi-compact
quelconque. Alors $W_{\overline{k}} \cap \overline{U}$ et
$W_{\overline{k}} \cap \overline{V}$ sont ouverts et
fermés dans $W_{\overline{k}}$. En particulier,
$W_{\overline{k}} \cap \overline{U}$ et $W_{\overline{k}}
\cap \overline{V}$ sont quasi-compacts, et d'après le lemme
\ref{lemma-Galois-action-quasi-compact-open},
$W_{\overline{k}} \cap \overline{U}$ et $W_{\overline{k}} \cap
\overline{V}$ sont définis sur une sous-extension finie et
invariants sous un sous-groupe ouvert de
$\text{Gal}(\overline{k}/k)$. Nous utiliserons cela sans autre mention dans ce qui suit.

\medskip\noindent
Choisissons $W_0 \subset X$ quasi-compact ouvert de
telle sorte que $W_{0, \overline{k}} \cap
\overline{U}$ et $W_{0, \overline{k}} \cap
\overline{V}$ soient tous deux non vides. Choisissons
une sous-extension finie $\overline{k}/k'/k$ et une
décomposition $W_{0, k'} = U_0' \amalg V_0'$ en
parties ouvertes et fermées de sorte que $W_{0,
\overline{k}} \cap \overline{U} =
(U'_0)_{\overline{k}}$ et $W_{0, \overline{k}} \cap
\overline{V} = (V'_0)_{\overline{k}}$. Soit $H =
\text{Gal}(\overline{k}/k') \subset \text{Gal}(\overline{k}/k)$.
En particulier $\sigma(W_{0, \overline{k}} \cap \overline{U}) =
W_{0, \overline{k}} \cap \overline{U}$ et de même pour $\overline{V}$.

\medskip\noindent
Ayant choisi $W_0$, $k'$ comme ci-dessus, pour
chaque ouvert quasi-compact $W \subset X$ nous posons
$$
U_W =
\bigcap\nolimits_{\sigma \in H} \sigma(W_{\overline{k}} \cap \overline{U}),
\quad
V_W =
\bigcup\nolimits_{\sigma \in H} \sigma(W_{\overline{k}} \cap \overline{V}).
$$
Maintenant, puisque $W_{\overline{k}} \cap \overline{U}$ et
$W_{\overline{k}} \cap \overline{V}$ sont fixés par un sous-groupe ouvert
de $\text{Gal}(\overline{k}/k)$, on voit que l'union et
l'intersection ci-dessus sont finies. Par conséquent, $U_W$ et $V_W$ sont à la
fois ouverts et fermés. De plus, par construction $W_{\bar k} = U_W \amalg V_W$.

\medskip\noindent
Nous affirmons que si $W \subset W' \subset X$ sont
quasi-compacts ouverts, alors $W_{\overline{k}} \cap U_{W'}
= U_W$ et $W_{\overline{k}} \cap V_{W'} = V_W$. Vérification
omise. Ainsi, nous voyons qu'en définissant $U = \bigcup_{W
\subset X} U_W$ et $V = \bigcup_{W \subset X} V_W$ nous
obtenons $X_{\overline{k}} = U \amalg V$ est une union
disjointe de ouverts et fermés. Il est évident que $V$ est
non vide car il est construit en prenant des unions
(localement). D'autre part, $U$ est non vide puisqu'il contient
$W_0 \cap \overline{U}$ par construction. Enfin, $U, V \subset
X_{\bar k}$ sont fermés et $H$-invariants par construction.
Ainsi, par le lemme \ref{lemma-closed-fixed-by-Galois} nous
avons $U = (U')_{\bar k}$, et $V = (V')_{\bar k}$ pour un certain
$U', V' \subset X_{k'}$ fermé. Il est clair que $X_{k'} = U'
\amalg V'$ et on voit que $X_{k'}$ est non connexe comme voulu.
\end{proof}

\begin{lemma}
\label{lemma-tricky}
Soit $k$ un corps. Supposons que $\overline{k}/k$ soit une
extension de Galois (éventuellement infinie). Supposons que $f : T \to X$
soit un morphisme de schémas sur $k$. Supposons que $T_{\overline{k}}$
soit connexe et $X_{\overline{k}}$ soit non connexe. Alors $X$ est non connexe.
\end{lemma}

\begin{proof}
Écrivons $X_{\overline{k}} = \overline{U} \amalg \overline{V}$
avec $\overline{U}$ et $\overline{V}$ ouverts et fermés. Notons
$\overline{f} : T_{\overline{k}} \to X_{\overline{k}}$ le
changement de base de $f$. Puisque $T_{\overline{k}}$ est connexe,
on voit que $T_{\overline{k}}$ est contenu soit dans
$\overline{f}^{-1}(\overline{U})$, soit dans
$\overline{f}^{-1}(\overline{V})$. Disons que $T_{\overline{k}} \subset \overline{f}^{-1}(\overline{U})$.

\medskip\noindent
Fixons un ouvert quasi-compact $W \subset X$. Il
existe une sous-extension galoisienne finie
$\overline{k}/k'/k$ telle que $\overline{U} \cap
W_{\overline{k}}$ et $\overline{V} \cap
W_{\overline{k}}$ proviennent d'ouverts quasi-compacts $U', V'
\subset W_{k'}$. Alors aussi $W_{k'} = U' \amalg V'$. Considérons
$$
U'' = \bigcap\nolimits_{\sigma \in \text{Gal}(k'/k)} \sigma(U'),
\quad
V'' = \bigcup\nolimits_{\sigma \in \text{Gal}(k'/k)} \sigma(V').
$$
Ces parties sont invariantes par Galois, ouvertes et
fermés, et $W_{k'} = U'' \amalg V''$.
D'après le lemme
\ref{lemma-closed-fixed-by-Galois}, nous obtenons des
ouverts et fermés $U_W, V_W \subset W$ tels que $U'' =
(U_W)_{k'}$, $V'' = (V_W)_{k'}$ et $W = U_W \amalg V_W$.

\medskip\noindent
Nous affirmons que si $W \subset W' \subset X$ sont
quasi-compacts ouverts, alors $W \cap U_{W'} = U_W$ et $W
\cap V_{W'} = V_W$. Nous omettons la vérification. Par conséquent,
nous voyons qu'en définissant $U = \bigcup_{W \subset X}
U_W$ et $V = \bigcup_{W \subset X} V_W$, nous obtenons $X =
U \amalg V$. Il est clair que $V$ est non vide car il est
construit en prenant des unions (localement). D'autre part,
$U$ est non vide puisqu'il contient $f(T)$ par construction.
\end{proof}

\begin{lemma}
\label{lemma-geometrically-connected-criterion}
\begin{reference}
\cite[IV Corollary 4.5.13.1(i)]{EGA}
\end{reference}
Soit $k$ un corps. Soit $T \to X$ un
morphisme de schémas sur $k$. Supposons $T$
géométriquement connexe et $X$ connexe. Alors $X$ est géométriquement connexe.
\end{lemma}

\begin{proof}
Ceci est une reformulation
du lemme \ref{lemma-tricky}.
\end{proof}

\begin{lemma}
\label{lemma-geometrically-connected-if-connected-and-point}
Soit $k$ un corps. Supposons que $X$ soit
un schéma sur $k$. Supposons que $X$ soit connexe et
possède un point $x$ tel que $k$ soit algébriquement clos
dans $\kappa(x)$. Alors $X$ est géométriquement connexe.
En particulier, si $X$ possède un point $k$-rationnel et
que $X$ est connexe, alors $X$ est géométriquement connexe.
\end{lemma}

\begin{proof}
Posons $T = \Spec(\kappa(x))$. Supposons que
$\overline{k}$ soit une clôture algébrique séparable
de $k$. L'hypothèse sur $\kappa(x)/k$ implique que
$T_{\overline{k}}$ est irréductible, voir Algèbre, lemme
\ref{algebra-lemma-field-extension-geometrically-irreducible}.
Donc, d'après le lemme
\ref{lemma-geometrically-connected-criterion}, on voit que
$X_{\overline{k}}$ est connexe. D'après le lemme
\ref{lemma-characterize-geometrically-connected}, nous concluons que $X$ est géométriquement connexe.
\end{proof}

\begin{lemma}
\label{lemma-inverse-image-connected-component}
Soit $K/k$ une extension de corps.
Soit $X$ un schéma sur $k$. Pour chaque
composante connexe $T$ de $X$, l'image réciproque $T_K
\subset X_K$ est une union de composantes connexes de $X_K$.
\end{lemma}

\begin{proof}
C'est un énoncé purement topologique. Notons $p
: X_K \to X$ le morphisme de projection. Soit $T \subset
X$ une composante connexe de $X$. Soit $t \in T_K =
p^{-1}(T)$. Soit $C \subset X_K$ une composante connexe
contenant $t$. Alors $p(C)$ est un sous-ensemble connexe de $X$ qui
rencontre $T$, donc $p(C) \subset T$. Par conséquent $C \subset T_K$.
\end{proof}

\noindent
Le lemme suivant sera supplanté par le lemme plus
fort \ref{lemma-image-connected-component} ci-dessous.

\begin{lemma}
\label{lemma-image-connected-component-finite-extension}
Supposons que $K/k$ soit une extension finie de corps et
soit $X$ un schéma sur $k$. Notons $p : X_K \to X$ le
morphisme de projection. Pour chaque composante connexe $T$
de $X_K$, l'image $p(T)$ est une composante connexe de $X$.
\end{lemma}

\begin{proof}
L'image $p(T)$ est contenue dans une certaine composante
connexe $X'$ de $X$. Considérons $X'$ comme un
sous-schéma fermé de $X$ de quelque manière que ce soit.
Alors $T$ est également une composante connexe de $X'_K =
p^{-1}(X')$ et nous pouvons donc supposer que $X$ est
connexe. Le morphisme $p$ est ouvert (Morphismes, lemme
\ref{morphisms-lemma-scheme-over-field-universally-open}),
fermé (Morphismes, lemme
\ref{morphisms-lemma-integral-universally-closed}) et les fibres de
$p$ sont des ensembles finis (Morphismes, lemme
\ref{morphisms-lemma-finite-quasi-finite}). Ainsi, nous pouvons appliquer
Topologie, lemme \ref{topology-lemma-finite-fibre-connected-components} pour conclure.
\end{proof}

\begin{lemma}[Gabber]
\label{lemma-image-connected-component}
\begin{reference}
Email de Ofer Gabber daté du 4 juin 2016
\end{reference}
Soit $K/k$ une extension de corps. Supposons que $X$
soit un schéma sur $k$. Notons $p : X_K \to X$ le morphisme de
projection. Supposons que $\overline{T} \subset X_K$ soit une
composante connexe. Alors $p(\overline{T})$ est une composante connexe de $X$.
\end{lemma}

\begin{proof}
Lorsque $K/k$ est fini, c'est le lemme
\ref{lemma-image-connected-component-finite-extension}.
En général, la preuve est plus difficile.

\medskip\noindent
Soit $T \subset X$ soit la composante connexe de $X$
contenant l'image de $\overline{T}$. Nous pouvons remplacer $X$ par
$T$ (avec la structure de sous-schéma réduit induite). Ainsi, nous
pouvons supposer que $X$ est connexe. Supposons que $A = H^0(X,
\mathcal{O}_X)$. Soit $L \subset A$ la $k$-algèbre maximale faiblement
étale, voir Compléments d'algèbre, lemme
\ref{more-algebra-lemma-max-weakly-etale-subalgebra}. Comme $A$ n'a pas
d'idempotents non triviaux, on voit que $L$ est un corps et une
extension algébrique séparable de $k$ d'après Compléments d'algèbre, lemme
\ref{more-algebra-lemma-class-weakly-etale-over-field}. Observons que $L$
est également l'algèbre $L$ maximale faiblement étale de $A$ (car toute
algèbre faiblement étale $L$ est faiblement étale sur $k$ d'après
Compléments d'algèbre, lemme \ref{more-algebra-lemma-composition-weakly-etale}).
D'après Schémas, lemme \ref{schemes-lemma-morphism-into-affine}, nous
obtenons une factorisation $X \to \Spec(L) \to \Spec(k)$ du morphisme de structure.

\medskip\noindent
Supposons que $L'/L$ soit une extension
finie séparable. D'après Cohomologie des
Schémas, lemme
\ref{coherent-lemma-finite-locally-free-base-change-cohomology} nous avons
$$
A \otimes_L L' =
H^0(X \times_{\Spec(L)} \Spec(L'), \mathcal{O}_{X \times_{\Spec(L)} \Spec(L')})
$$
La $L'$-sous-algèbre faiblement étale maximale de $A \otimes_L
L'$ est $L \otimes_L L' = L'$, d'après Compléments d'algèbre, lemme
\ref{more-algebra-lemma-change-fields-max-weakly-etale-subalgebra}.
En particulier, $A \otimes_L L'$ n'a pas d'idempotents non
triviaux (un tel idempotent engendrerait une sous-algèbre faiblement
étale) et nous concluons que $X \times_{\Spec(L)} \Spec(L')$ est
connexe. Par le lemme
\ref{lemma-characterize-geometrically-disconnected}, nous concluons que $X$ est géométriquement connexe sur $L$.

\medskip\noindent
Donnons à $\overline{T}$ la structure de
schéma induit réduite et considérons la composition
$$
\overline{T} \xrightarrow{i} X_K = X \times_{\Spec(k)} \Spec(K)
\xrightarrow{\pi}
\Spec(L \otimes_k K)
$$
L'image est contenue dans une composante connexe de $\Spec(L
\otimes_k K)$. Puisque $K \to L \otimes_k K$ est intégral,
on voit que les composantes connexes de $\Spec(L \otimes_k
K)$ sont des points et tous les points sont fermés, voir
Algèbre, lemme \ref{algebra-lemma-integral-over-field}. Ainsi,
nous obtenons un corps quotient $L \otimes_k K \to E$ tel que
$\overline{T}$ s'envoie dans $\Spec(E) \subset \Spec(L \otimes_k
K)$. Par conséquent $i(\overline{T}) \subset \pi^{-1}(\Spec(E))$. Mais
$$
\pi^{-1}(\Spec(E)) =
(X \times_{\Spec(k)} \Spec(K)) \times_{\Spec(L \otimes_k K)} \Spec(E) =
X \times_{\Spec(L)} \Spec(E)
$$
qui est connexe parce que $X$ est géométriquement connexe sur
$L$. Ensuite, nous obtenons l'égalité $\overline{T} = X
\times_{\Spec(L)} \Spec(E)$ (au sens des ensembles) et nous
concluons que $\overline{T} \to X$ est surjective comme voulu.
\end{proof}

\noindent
Soit $X$ un schéma. On note
$\pi_0(X)$ l'ensemble des composantes connexes de $X$.

\begin{lemma}
\label{lemma-galois-action-connected-components}
Soit $k$ un corps, muni d'une clôture
algébrique séparable $\overline{k}$. Supposons
que $X$ soit un schéma sur $k$. Il existe une action
$$
\text{Gal}(\overline{k}/k)^{opp} \times \pi_0(X_{\overline{k}})
\longrightarrow
\pi_0(X_{\overline{k}})
$$
avec les propriétés suivantes :
\begin{enumerate}
\item Un élément $\overline{T} \in
\pi_0(X_{\overline{k}})$ est fixé par l'action si et
seulement s'il existe une composante connexe $T \subset X$, qui
est géométriquement connexe sur $k$, telle que
$T_{\overline{k}} = \overline{T}$. \item Pour toute extension de corps
$k'/k$ avec clôture algébrique séparable $\overline{k}'$ le diagramme
$$
\xymatrix{
\text{Gal}(\overline{k}'/k') \times \pi_0(X_{\overline{k}'})
\ar[r] \ar[d] &
\pi_0(X_{\overline{k}'}) \ar[d] \\
\text{Gal}(\overline{k}/k) \times \pi_0(X_{\overline{k}})
\ar[r] &
\pi_0(X_{\overline{k}})
}
$$
est commutatif (où la flèche verticale droite est une bijection
d'après le lemme \ref{lemma-separably-closed-field-connected-components}).
\end{enumerate}
\end{lemma}

\begin{proof}
L'action (\ref{equation-galois-action-base-change-kbar})
de $\text{Gal}(\overline{k}/k)$ sur $X_{\overline{k}}$
induit une action sur ses composantes connexes. Les
composantes connexes sont toujours fermées (Topologie,
lemme \ref{topology-lemma-connected-components}). Donc si
$\overline{T}$ est comme dans (1), alors par le lemme
\ref{lemma-closed-fixed-by-Galois} il existe un sous-ensemble
fermé $T \subset X$ tel que $\overline{T} =
T_{\overline{k}}$. Notons que $T$ est géométriquement connexe
sur $k$, voir le lemme
\ref{lemma-characterize-geometrically-connected}. Pour voir que $T$ est
une composante connexe de $X$, supposons que $T \subset T'$, $T \not =
T'$ où $T'$ est une composante connexe de $X$. Dans ce cas, $T'_{k'}$
contient strictement $\overline{T}$ et est donc non connexe. Par le lemme
\ref{lemma-tricky} cela signifie que $T'$ est non connexe ! Contradiction.

\medskip\noindent
Nous omettons la preuve de la fonctorialité en (2).
\end{proof}

\begin{lemma}
\label{lemma-galois-action-connected-components-continuous}
Soit $k$ un corps, muni d'une clôture
algébrique séparable $\overline{k}$.
Soit $X$ un schéma sur $k$. Supposons
\begin{enumerate}
\item $X$ est quasi-compact, et \item les
composantes connexes de $X_{\overline{k}}$ sont ouvertes.
\end{enumerate}
Alors
\begin{enumerate}
\item [(a)] $\pi_0(X_{\overline{k}})$ est fini,
et \item [(b)] l'action de
$\text{Gal}(\overline{k}/k)$ sur $\pi_0(X_{\overline{k}})$ est continue.
\end{enumerate}
De plus, les hypothèses (1) et (2) sont
satisfaites lorsque $X$ est de type fini sur $k$.
\end{lemma}

\begin{proof}
Comme les composantes connexes, qui sont ouvertes, recouvrent
$X_{\overline{k}}$ (Topologie, lemme
\ref{topology-lemma-connected-components}) et
$X_{\overline{k}}$ est quasi-compact, nous concluons qu'il
n'y en a qu'un nombre fini. Ainsi, (a) est vérifié. D'après
le lemme \ref{lemma-descend-open}, ces composantes connexes
sont chacune définies sur une sous-extension finie de
$\overline{k}/k$ et nous obtenons (b). Si $X$ est de type
fini sur $k$, alors $X_{\overline{k}}$ est de type fini sur
$\overline{k}$ (Morphismes, lemme
\ref{morphisms-lemma-base-change-finite-type}). Par conséquent,
$X_{\overline{k}}$ est un schéma noethérien (Morphismes, lemme
\ref{morphisms-lemma-finite-type-noetherian}). Ainsi,
$X_{\overline{k}}$ a un nombre fini de composantes irréductibles
(Propriétés, lemme \ref{properties-lemma-Noetherian-irreducible-components})
et a fortiori un nombre fini de composantes connexes (qui sont donc ouvertes).
\end{proof}









\section{Schémas géométriquement irréductibles}
\label{section-geometrically-irreducible}

\noindent
Si $X$ est un schéma irréductible sur un corps, il peut arriver
que $X$ devienne réductible après l'extension du corps de base.
Cela n'arrive pas pour les schémas géométriquement irréductibles.

\begin{definition}
\label{definition-geometrically-irreducible}
Soit $X$ un schéma sur le corps $k$. On dit que
$X$ est {\it géométriquement irréductible} sur $k$ si le
schéma $X_{k'}$ est irréductible \footnote { Un espace
irréductible est non vide. } pour toute extension de corps $k'$ de $k$.
\end{definition}

\begin{lemma}
\label{lemma-geometrically-irreducible-check-after-extension}
Supposons que $X$ soit un schéma sur le corps $k$.
Supposons que $k'/k$ soit une extension de corps. Alors
$X$ est géométriquement irréductible sur $k$ si et
seulement si $X_{k'}$ est géométriquement irréductible sur $k'$.
\end{lemma}

\begin{proof}
Si $X$ est géométriquement irréductible sur $k$, alors il est
clair que $X_{k'}$ est géométriquement irréductible sur $k'$.
Pour la réciproque, pour toute extension de corps $k''/k$, il
existe une extension de corps commune $k'''/k'$ et $k'''/k''$,
voir corps, lemme \ref{fields-lemma-common-extension-field}. Comme
le morphisme $X_{k'''} \to X_{k''}$ est surjectif (en tant que
changement de base d'un morphisme surjectif entre spectres de
corps), on voit que l'irréductibilité de $X_{k'''}$ implique
l'irréductibilité de $X_{k''}$. Ainsi, si $X_{k'}$ est géométriquement
irréductible sur $k'$, alors $X$ est géométriquement irréductible sur $k$.
\end{proof}

\begin{lemma}
\label{lemma-separably-closed-irreducible}
Supposons que $X$ soit un schéma sur un corps
séparablement clos $k$. Si $X$ est irréductible, alors
$X_K$ est irréductible pour toute extension de corps
$K/k$. C.-à-d., $X$ est géométriquement irréductible sur $k$.
\end{lemma}

\begin{proof}
Cela résulte de Propriétés, lemme
\ref{properties-lemma-characterize-irreducible} et Algèbre, lemme \ref{algebra-lemma-separably-closed-irreducible}.
\end{proof}

\begin{lemma}
\label{lemma-bijection-irreducible-components}
Soit $k$ un corps. Supposons
que $X$, $Y$ soient des schémas sur $k$.
Supposons que $X$ soit géométriquement
irréductible sur $k$. Ensuite, le morphisme de projection
$$
p : X \times_k Y \longrightarrow Y
$$
induit une bijection entre les composantes irréductibles.
\end{lemma}

\begin{proof}
Tout d'abord, notons que les fibres schématiques de $p$
sont irréductibles, puisqu'elles sont des changements de
base du schéma géométriquement irréductible $X$ par des
extensions de corps. De plus, les fibres schématiques
sont homéomorphes aux fibres ensemblistes, voir Schémas,
lemme \ref{schemes-lemma-fibre-topological}. Selon
Morphismes, lemme
\ref{morphisms-lemma-scheme-over-field-universally-open}, l'application
$p$ est ouverte. Ainsi, nous pouvons appliquer Topologie, lemme
\ref{topology-lemma-irreducible-fibres-irreducible-components} pour conclure.
\end{proof}

\begin{lemma}
\label{lemma-geometrically-irreducible-local}
\begin{slogan}
L'irréductibilité géométrique est Zariski locale modulo la connexité.
\end{slogan}
Soit $k$ un corps. Supposons que $X$ soit
un schéma sur $k$. Les conditions suivantes sont équivalentes
\begin{enumerate}
\item $X$ est géométriquement irréductible sur $k$, \item
pour chaque ouvert affine non vide $U$ la $k$-algèbre
$\mathcal{O}_X(U)$ est géométriquement irréductible sur $k$
(voir Algèbre, définition
\ref{algebra-definition-geometrically-irreducible}), \item $X$
est irréductible et il existe un recouvrement ouvert affine $X =
\bigcup U_i$ tel que chaque $k$-algèbre $\mathcal{O}_X(U_i)$ est
géométriquement irréductible, et \item il existe un recouvrement
ouvert $X = \bigcup_{i \in I} X_i$ avec $I \not = \emptyset$
tel que $X_i$ est géométriquement irréductible pour chaque $i$ et
tel que $X_i \cap X_j \not = \emptyset$ pour tout $i, j \in I$.
\end{enumerate}
De plus, si $X$ est géométriquement irréductible, il en
est de même pour tout sous-schéma ouvert non vide de $X$.
\end{lemma}

\begin{proof}
Un schéma affine $\Spec(A)$ sur $k$ est géométriquement
irréductible si et seulement si $A$ est géométriquement
irréductible sur $k$ ; cela découle immédiatement des
définitions. Rappelons que si un schéma est irréductible, alors
tout sous-schéma ouvert non vide de $X$ l'est aussi, et que deux
ouverts non vides ont une intersection non vide. De plus, si tout
ouvert affine est irréductible, alors le schéma est irréductible,
voir Propriétés, lemme
\ref{properties-lemma-characterize-irreducible}. Ainsi, l'énoncé final du
lemme est clair, ainsi que les implications (1) $\Rightarrow$ (2), (2)
$\Rightarrow$ (3), et (3) $\Rightarrow$ (4). Si (4) est vrai, alors pour
toute extension de corps $k'/k$, le schéma $X_{k'}$ possède un recouvrement
par des ouverts irréductibles dont les intersections deux à deux sont non
vides. Par conséquent, $X_{k'}$ est irréductible. Ainsi, (4) implique (1).
\end{proof}

\begin{lemma}
\label{lemma-geometrically-irreducible-function-field}
Supposons que $X$ soit un schéma irréductible sur le corps $k$. Supposons que
$\xi \in X$ soit son point générique. Les énoncés suivants sont équivalents :
\begin{enumerate}
\item $X$ est géométriquement irréductible sur $k$ ; \item
$\kappa(\xi)$ est géométriquement irréductible sur $k$ ; \item la
clôture algébrique séparable de $k$ dans $\kappa(\xi)$ est égale à $k$.
\end{enumerate}
\end{lemma}

\begin{proof}
D'après Algèbre, lemme
\ref{algebra-lemma-geometrically-irreducible-separable-elements},
(2) et (3) sont équivalents.

\medskip\noindent
Supposons (1). Rappelons que $\mathcal{O}_{X, \xi}$ est
la colimite filtrante de $\mathcal{O}_X(U)$ où $U$
parcourt les sous-schémas affines ouverts non vides de
$X$. En combinant le lemme
\ref{lemma-geometrically-irreducible-local} et le lemme
d'Algèbre
\ref{algebra-lemma-subalgebra-geometrically-irreducible}, nous voyons
que $\mathcal{O}_{X, \xi}$ est géométriquement irréductible sur $k$.
Puisque $\mathcal{O}_{X, \xi} \to \kappa(\xi)$ est une surjection avec
noyau localement nilpotent (voir le lemme d'Algèbre
\ref{algebra-lemma-minimal-prime-reduced-ring}), il s'ensuit que $\kappa(\xi)$ est
géométriquement irréductible, voir le lemme d'Algèbre \ref{algebra-lemma-p-ring-map}.

\medskip\noindent
Supposons (2). On peut supposer que $X$ est réduit. Supposons que
$U \subset X$ est un ouvert affine non vide. Alors $U = \Spec(A)$ où $A$
est un anneau intègre avec corps des fractions $\kappa(\xi)$. Ainsi $A$ est
une $k$-sous-algèbre d'une $k$-algèbre géométriquement irréductible. Par
conséquent, d'après le lemme d'Algèbre
\ref{algebra-lemma-subalgebra-geometrically-irreducible}, on voit que $A$ est
géométriquement irréductible sur $k$. D'après le lemme
\ref{lemma-geometrically-irreducible-local}, nous en concluons que $X$ est géométriquement irréductible sur $k$.
\end{proof}

\begin{lemma}
\label{lemma-geometrically-irreducible-self-product}
Un schéma $X$ sur un corps $k$ est géométriquement
irréductible sur $k$ si et seulement si $X\times_k X$ est irréductible.
\end{lemma}

\begin{proof}
Si $X$ est géométriquement irréductible sur $k$,
alors $X\times_k X$ est irréductible selon le
lemme \ref{lemma-bijection-irreducible-components}.

\medskip\noindent
Inversement, supposons que $X \times_k X$ soit
irréductible. En remplaçant $X$ par le schéma réduit
sous-jacent, on peut supposer que $X$ est réduit sans
changer le problème. Puisque la première projection $X
\times_k X \to X$ est surjective, $X$ est également
irréductible, donc intégral. Supposons qu'il existe un élément
$\alpha$ dans le corps des fonctions $k(X)$ qui soit
algébrique séparable sur $k$ mais pas dans $k$. Soit $E = k(\alpha)$.

\medskip\noindent
Puis il existe un sous-schéma affine ouvert non vide $U$ de $X$
tel que $k \subset E \subset O(U)$, et donc $E \otimes_k E
\subset O(U)\otimes_k O(U) = O(U\times_k U)$. Ainsi nous avons
un morphisme dominant $U \times_k U \to \Spec(E\otimes_k E)$.
Puisque $X\times_k X$ est irréductible, $U \times_k U$ l'est
aussi, et donc $\Spec(E \otimes_k E)$ est connexe. Si $E$ est
strictement plus grand que $k$, alors $E \otimes_k E$ est un
produit de corps incluant $E$ comme un facteur. Il a une
dimension $(\dim_k(E))^2>\dim_k(E)$ comme espace vectoriel sur
$k$, et donc $E\otimes_k E$ est un produit d'au moins deux corps,
ce qui contredit que $\Spec(E \otimes_k E)$ est connexe. Donc en
fait $E$ est égal à $k$. D'après le lemme
\ref{lemma-geometrically-irreducible-function-field}, $X$ est géométriquement irréductible sur $k$.
\end{proof}

\begin{lemma}
\label{lemma-separably-closed-field-irreducible-components}
Supposons que $k'/k$ soit une extension de corps. Supposons que
$X$ est un schéma sur $k$. Posons $X' = X_{k'}$. Supposons que
$k$ est séparablement algébriquement clos. Alors le morphisme
$X' \to X$ induit une bijection des composantes irréductibles.
\end{lemma}

\begin{proof}
Puisque $k$ est séparablement algébriquement clos, on voit que
$k'$ est géométriquement irréductible sur $k$, voir Algèbre, lemme
\ref{algebra-lemma-separably-closed-irreducible-implies-geometric}.
Par conséquent, $Z = \Spec(k')$ est géométriquement irréductible
sur $k$, d'après le lemme
\ref{lemma-geometrically-irreducible-local} ci-dessus. Puisque $X' = Z \times_k X$, le
résultat est un cas particulier du lemme \ref{lemma-bijection-irreducible-components}.
\end{proof}

\begin{lemma}
\label{lemma-characterize-geometrically-irreducible}
\begin{slogan}
L'irréductibilité géométrique peut être testée sur
une clôture algébrique séparable du corps de base.
\end{slogan}
Soit $k$ un corps. Supposons que $X$
soit un schéma sur $k$. Les conditions suivantes sont équivalentes :
\begin{enumerate}
\item $X$ est géométriquement irréductible sur $k$,
\item pour toute extension de corps finie séparable
$k'/k$ le schéma $X_{k'}$ est irréductible, et \item
$X_{\overline{k}}$ est irréductible, où $k \subset
\overline{k}$ est une clôture algébrique séparable de $k$.
\end{enumerate}
\end{lemma}

\begin{proof}
Supposons que $X_{\overline{k}}$ soit irréductible, c'est-à-dire,
supposons (3). Supposons que $k'/k$ soit une extension de corps.
Il existe une extension de corps $\overline{k}'/\overline{k}$
telle que $k'$ se plonge dans $\overline{k}'$ comme une extension
de $k$. D'après le lemme
\ref{lemma-separably-closed-field-irreducible-components}, on voit que
$X_{\overline{k}'}$ est irréductible. Puisque $X_{\overline{k}'} \to X_{k'}$ est
surjective, nous concluons que $X_{k'}$ est irréductible. Par conséquent, (1) est vrai.

\medskip\noindent
Supposons que $k \subset \overline{k}$ soit une clôture
algébrique séparable de $k$. Supposons que non (3),
c'est-à-dire supposons que $X_{\overline{k}}$ soit
réductible. Notre but est de montrer que $X_{k'}$ est
également réductible pour une certaine sous-extension
finie $\overline{k}/k'/k$. Supposons que $X = \bigcup_{i
\in I} U_i$ soit un recouvrement affine ouvert avec $U_i$
non vide. Si pour un certain $i$ le schéma $U_i$ est
réductible, ou si pour une certaine paire $i \not = j$
l'intersection $U_i \cap U_j$ est vide, alors $X$ est
réductible (Propriétés, lemme
\ref{properties-lemma-characterize-irreducible}) et nous
avons terminé. En particulier, on peut supposer que
$U_{i, \overline{k}} \cap U_{j, \overline{k}}$ pour tout $i, j
\in I$ est non vide et nous en concluons que $U_{i,
\overline{k}}$ doit être réductible pour un certain $i$. Selon
Algèbre, lemme \ref{algebra-lemma-geometrically-irreducible},
cela signifie que $U_{i, k'}$ est réductible pour une certaine
extension de corps finie séparable $k'/k$. Par conséquent, $X_{k'}$
est également réductible. Ainsi, on voit que (2) implique (3).

\medskip\noindent
L'implication (1) $\Rightarrow$ (2)
est immédiate. Cela prouve le lemme.
\end{proof}

\begin{lemma}
\label{lemma-inverse-image-irreducible}
Soit $K/k$ une extension de corps.
Soit $X$ un schéma sur $k$. Pour chaque
composante irréductible $T$ de $X$, l'image réciproque $T_K
\subset X_K$ est une union de composantes irréductibles de $X_K$.
\end{lemma}

\begin{proof}
Soit $T \subset X$ une composante irréductible de
$X$. Le morphisme $T_K \to T$ est plat, donc les généralisations
se relèvent le long de $T_K \to T$. Ainsi, chaque $\xi \in T_K$
qui est un point générique d'une composante irréductible de
$T_K$ s'envoie sur le point générique $\eta$ de $T$. Si $\xi'
\leadsto \xi$ est une spécialisation dans $X_K$, alors $\xi'$ se
mappe à $\eta$ puisqu'il n'y a pas de points se spécialisant en
$\eta$ dans $X$. Ainsi $\xi' \in T_K$ et nous concluons que $\xi =
\xi'$. En d'autres termes, $\xi$ est le point générique d'une
composante irréductible de $X_K$. Cela signifie que les composantes
irréductibles de $T_K$ sont toutes des composantes irréductibles de $X_K$.
\end{proof}

\noindent
Pour un schéma $X$, nous notons
$\text{IrredComp}(X)$ l'ensemble des composantes irréductibles de $X$.

\begin{lemma}
\label{lemma-image-irreducible}
Soit $K/k$ une extension de corps.
Soit $X$ un schéma sur $k$. Pour chaque
composante irréductible $\overline{T} \subset X_K$,
l'image de $\overline{T}$ dans $X$ est une composante
irréductible dans $X$. Cela définit une application canonique
$$
\text{IrredComp}(X_K)
\longrightarrow
\text{IrredComp}(X)
$$
qui est surjective.
\end{lemma}

\begin{proof}
Considérons le diagramme
$$
\xymatrix{
X_K \ar[d] & X_{\overline{K}} \ar[d] \ar[l] \\
X & X_{\overline{k}} \ar[l]
}
$$
où $\overline{K}$ est la clôture algébrique séparable de $K$,
et où $\overline{k}$ est la clôture algébrique séparable de
$k$. D'après le lemme
\ref{lemma-separably-closed-field-irreducible-components}, le
morphisme $X_{\overline{K}} \to X_{\overline{k}}$ induit une bijection
entre les composantes irréductibles. Il suffit donc de démontrer le
lemme pour les morphismes $X_{\overline{k}} \to X$ et $X_{\overline{K}}
\to X_K$. En d'autres termes, on peut supposer que $K = \overline{k}$.

\medskip\noindent
Le morphisme $p : X_{\overline{k}} \to X$ est intégral, plat et
surjectif. La platitude implique que les généralisations se relèvent
le long de $p$, voir Morphismes, lemme
\ref{morphisms-lemma-generalizations-lift-flat}. Par conséquent, les
points génériques des composantes irréductibles de $X_{\overline{k}}$ se
projettent sur des points génériques des composantes irréductibles de
$X$. L'intégralité implique que $p$ est universellement fermé, voir
Morphismes, lemme \ref{morphisms-lemma-integral-universally-closed}.
Ainsi, nous concluons que l'image $p(\overline{T})$ d'une composante
irréductible est un sous-ensemble fermé irréductible qui contient un point
générique d'une composante irréductible de $X$, donc $p(\overline{T})$ est
une composante irréductible de $X$. Cela prouve la première assertion. Si
$T \subset X$ est une composante irréductible, alors $p^{-1}(T) =T_K$ est
une union non vide de composantes irréductibles, voir lemme
\ref{lemma-inverse-image-irreducible}. Chacune d'elles se projette nécessairement
sur $T$ selon la première partie. Par conséquent, l'application est surjective.
\end{proof}

\begin{lemma}
\label{lemma-irreducible-dense-rational-points-geometrically-irreducible}
Soit $k$ un corps. Supposons que $X$ soit un schéma
sur $k$. Si $X$ est irréductible et possède un ensemble dense de
points $k$-rationnels, alors $X$ est géométriquement irréductible.
\end{lemma}

\begin{proof}
Supposons que $k'/k$ soit une extension finie de corps et soit
$Z, Z' \subset X_{k'}$ des composantes irréductibles. Il
suffit de montrer $Z = Z'$, voir le lemme
\ref{lemma-characterize-geometrically-irreducible}. Par le lemme
\ref{lemma-image-irreducible}, nous avons $p(Z) = p(Z') = X$ où
$p : X_{k'} \to X$ est la projection. Si $Z \not = Z'$ alors $Z
\cap Z'$ est nulle part dense dans $X_{k'}$ et donc $p(Z \cap
Z')$ n'est pas dense d'après Morphismes, lemme
\ref{morphisms-lemma-image-nowhere-dense-finite} ; ici nous
utilisons aussi que $p$ est un morphisme fini comme le changement
de base du morphisme fini $\Spec(k') \to \Spec(k)$, voir
Morphismes, lemme \ref{morphisms-lemma-base-change-finite}. Ainsi
nous pouvons choisir un point $k$-rationnel $x \in X$ avec $x \not \in
p(Z \cap Z')$. Puisque le corps résiduel de $x$ est $k$, nous voyons
que $p^{-1}(\{x\}) = \{x'\}$ où $x' \in X_{k'}$ est un point dont le
corps résiduel est $k'$. Puisque $x \in p(Z) = p(Z')$, nous concluons
que $x' \in Z \cap Z'$, ce qui est la contradiction que nous cherchions.
\end{proof}

\begin{lemma}
\label{lemma-galois-action-irreducible-components}
Soit $k$ un corps, muni d'une clôture
algébrique séparable $\overline{k}$. Supposons
que $X$ soit un schéma sur $k$. Il existe une action
$$
\text{Gal}(\overline{k}/k)^{opp} \times \text{IrredComp}(X_{\overline{k}})
\longrightarrow
\text{IrredComp}(X_{\overline{k}})
$$
avec les propriétés suivantes :
\begin{enumerate}
\item Un élément $\overline{T} \in
\text{IrredComp}(X_{\overline{k}})$ est fixé par l'action si et
seulement s'il existe une composante irréductible $T \subset X$,
qui est géométriquement irréductible sur $k$, telle que
$T_{\overline{k}} = \overline{T}$. \item Pour toute extension de corps
$k'/k$ avec clôture algébrique séparable $\overline{k}'$, le diagramme
$$
\xymatrix{
\text{Gal}(\overline{k}'/k') \times \text{IrredComp}(X_{\overline{k}'})
\ar[r] \ar[d] &
\text{IrredComp}(X_{\overline{k}'}) \ar[d] \\
\text{Gal}(\overline{k}/k) \times \text{IrredComp}(X_{\overline{k}})
\ar[r] &
\text{IrredComp}(X_{\overline{k}})
}
$$
est commutatif (où la flèche verticale droite est une bijection selon
le lemme \ref{lemma-separably-closed-field-irreducible-components}).
\end{enumerate}
\end{lemma}

\begin{proof}
L'action ( \ref{equation-galois-action-base-change-kbar} ) de
$\text{Gal}(\overline{k}/k)$ sur $X_{\overline{k}}$ induit
une action sur ses composantes irréductibles. Les composantes
irréductibles sont toujours fermées (Topologie, lemme
\ref{topology-lemma-connected-components}). Par conséquent, si
$\overline{T}$ est comme dans (1), alors par le lemme
\ref{lemma-closed-fixed-by-Galois} il existe un sous-ensemble
fermé $T \subset X$ tel que $\overline{T} = T_{\overline{k}}$.
Notons que $T$ est géométriquement irréductible sur $k$, voir le
lemme \ref{lemma-characterize-geometrically-irreducible}. Pour
voir que $T$ est une composante irréductible de $X$, supposons
que $T \subset T'$, $T \not = T'$ où $T'$ est une composante
irréductible de $X$. Supposons que $\overline{\eta}$ soit le
point générique de $\overline{T}$. Il s'envoie sur le point
générique $\eta$ de $T$. Ensuite, le point générique $\xi \in
T'$ se spécialise en $\eta$. Comme $X_{\overline{k}} \to X$ est
plat, il existe un point $\overline{\xi} \in X_{\overline{k}}$
qui s'envoie sur $\xi$ et se spécialise en $\overline{\eta}$. Il
s'ensuit que l'adhérence du singleton $\{\overline{\xi}\}$ est un
sous-ensemble fermé irréductible de $X_{\overline{\xi}}$ qui
contient strictement $\overline{T}$. C'est la contradiction souhaitée.

\medskip\noindent
Nous omettons la preuve de la fonctorialité en (2).
\end{proof}

\begin{lemma}
\label{lemma-orbit-irreducible-components}
Soit $k$ un corps, muni d'une clôture
algébrique séparable $\overline{k}$. Supposons que
$X$ soit un schéma sur $k$. Les fibres de l'application
$$
\text{IrredComp}(X_{\overline{k}})
\longrightarrow
\text{IrredComp}(X)
$$
du lemme \ref{lemma-image-irreducible}
sont exactement les orbites de
$\text{Gal}(\overline{k}/k)$ sous l'action du lemme
\ref{lemma-galois-action-irreducible-components}.
\end{lemma}

\begin{proof}
Soit $T \subset X$ une composante
irréductible de $X$. Supposons que $\eta \in T$ soit
son point générique. D'après les lemmes
\ref{lemma-inverse-image-irreducible} et
\ref{lemma-image-irreducible}, les points génériques des
composantes irréductibles de $\overline{T}$ qui se
projettent sur $T$ se projettent sur $\eta$. D'après le lemme
d'algèbre \ref{algebra-lemma-Galois-orbit}, le groupe de
Galois agit transitivement sur tous les points de
$X_{\overline{k}}$ se projetant sur $\eta$. D'où le lemme en découle.
\end{proof}

\begin{lemma}
\label{lemma-galois-action-irreducible-components-locally-finite-type}
Soit $k$ un corps.
Supposons que $X \to \Spec(k)$ soit
localement de type fini. Dans ce cas
\begin{enumerate}
\item l'action
$$
\text{Gal}(\overline{k}/k)^{opp} \times \text{IrredComp}(X_{\overline{k}})
\longrightarrow
\text{IrredComp}(X_{\overline{k}})
$$
est continue si nous donnons à
$\text{IrredComp}(X_{\overline{k}})$ la topologie
discrète, \item chaque composante irréductible de
$X_{\overline{k}}$ peut être définie sur une extension finie
de $k$, et \item étant donnée une composante irréductible
$T \subset X$, le schéma $T_{\overline{k}}$ est une union
finie de composantes irréductibles de $X_{\overline{k}}$ qui
sont toutes dans le même $\text{Gal}(\overline{k}/k)$-orbite.
\end{enumerate}
\end{lemma}

\begin{proof}
Supposons que $\overline{T}$ soit une composante irréductible
de $X_{\overline{k}}$. Nous pouvons choisir un ouvert affine $U
\subset X$ tel que $\overline{T} \cap U_{\overline{k}}$ ne soit
pas vide. Écrivons $U = \Spec(A)$, ainsi $A$ est une
$k$-algèbre de type fini, voir Morphismes, lemme
\ref{morphisms-lemma-locally-finite-type-characterize}. Par
conséquent, $A_{\overline{k}}$ est une $\overline{k}$-algèbre de
type fini, et en particulier noethérienne. Supposons que $\mathfrak
p = (f_1, \ldots, f_n)$ soit l'idéal premier correspondant à
$\overline{T} \cap U_{\overline{k}}$. Comme $A_{\overline{k}} = A
\otimes_k \overline{k}$, nous voyons qu'il existe une sous-extension
finie $\overline{k}/k'/k$ telle que chaque $f_i \in A_{k'}$. Il est clair
que $\text{Gal}(\overline{k}/k')$ fixe $\overline{T}$, ce qui prouve (1).

\medskip\noindent
La partie (2) suit en appliquant le lemme
\ref{lemma-galois-action-irreducible-components} (1) à
la situation sur $k'$, ce qui implique que la composante
irréductible $\overline{T}$ est de la forme
$T'_{\overline{k}}$ pour un certain irréductible $T' \subset X_{k'}$.

\medskip\noindent
Pour démontrer (3), soit $T \subset X$ une composante
irréductible. Choisissons une composante irréductible
$\overline{T} \subset X_{\overline{k}}$ qui s'envoie sur
$T$, voir le lemme \ref{lemma-image-irreducible}. D'après ce
qui précède, l'orbite de $\overline{T}$ est finie, disons
qu'elle est $\overline{T}_1, \ldots, \overline{T}_n$. Alors
$\overline{T}_1 \cup \ldots \cup \overline{T}_n$ est un
sous-ensemble fermé invariant par $\text{Gal}(\overline{k}/k)$
de $X_{\overline{k}}$ et donc de la forme $W_{\overline{k}}$
pour un certain $W \subset X$ fermé d'après le lemme
\ref{lemma-closed-fixed-by-Galois}. Clairement $W = T$ et cela conclut.
\end{proof}

\begin{lemma}
\label{lemma-finite-extension-geometrically-irreducible-components}
Soit $k$ un corps. Supposons que $X
\to \Spec(k)$ soit localement de type fini.
Supposons que $X$ ait un nombre fini de composantes
irréductibles. Alors il existe une extension finie
séparable $k'/k$ telle que chaque composante irréductible
de $X_{k'}$ soit géométriquement irréductible sur $k'$.
\end{lemma}

\begin{proof}
Supposons que $\overline{k}$ soit une clôture algébrique
séparable de $k$. L'hypothèse que $X$ a un nombre fini de
composantes irréductibles combinée avec le lemme
\ref{lemma-galois-action-irreducible-components-locally-finite-type}
(3) montre que $X_{\overline{k}}$ a un nombre fini de
composantes irréductibles $\overline{T}_1, \ldots,
\overline{T}_n$. Par le lemme
\ref{lemma-galois-action-irreducible-components-locally-finite-type} (2), il existe
une extension finie $\overline{k}/k'/k$ et des composantes irréductibles $T_i
\subset X_{k'}$ telles que $\overline{T}_i = T_{i, \overline{k}}$ et nous avons terminé.
\end{proof}

\begin{lemma}
\label{lemma-irreducible-components-geometrically-irreducible}
Supposons que $X$ soit un schéma sur le corps
$k$. Supposons que $X$ ait un nombre fini de
composantes irréductibles, toutes géométriquement
irréductibles. Alors $X$ a un nombre fini de
composantes connexes, chacune étant géométriquement connexe.
\end{lemma}

\begin{proof}
Ceci est clair parce qu'une composante connexe est
une union de composantes irréductibles. Nous omettons les détails.
\end{proof}







\section{Schémas géométriquement intègres}
\label{section-geometrically-integral}

\noindent
Si $X$ est un schéma intègre sur un corps, il peut arriver que $X$
devienne soit non réduit soit réductible après extension du corps de
base. Cela n'arrive pas pour les schémas géométriquement intègres.

\begin{definition}
\label{definition-geometrically-integral}
Soit $X$ un schéma sur le corps $k$.
\begin{enumerate}
\item Soit $x \in X$. Nous disons que $X$ est {\it
géométriquement ponctuellement intègre en $x$ } si pour toute
extension de corps $k'/k$ et tout $x' \in X_{k'}$ au-dessus de $x$,
l'anneau local $\mathcal{O}_{X_{k'}, x'}$ est intègre. \item Nous
disons que $X$ est {\it géométriquement ponctuellement intègre} si
$X$ est géométriquement ponctuellement intègre en tout point. \item
Nous disons que $X$ est {\it géométriquement intègre} sur $k$ si le
schéma $X_{k'}$ est intègre pour toute extension de corps $k'$ de $k$.
\end{enumerate}
\end{definition}

\noindent
La distinction entre les notions (2) et (3) est nécessaire.
Par exemple, si $k = \mathbf{R}$ et $X =
\Spec(\mathbf{C}[x])$, alors $X$ est géométriquement ponctuellement
intégral sur $\mathbf{R}$ mais bien sûr pas géométriquement intègre.

\begin{lemma}
\label{lemma-geometrically-integral}
Soit $k$ un corps. Supposons
que $X$ soit un schéma sur $k$. Alors $X$ est
géométriquement intègre sur $k$ si et
seulement si $X$ est à la fois géométriquement
réduit et géométriquement irréductible sur $k$.
\end{lemma}

\begin{proof}
Voir Propriétés, lemme \ref{properties-lemma-characterize-integral}.
\end{proof}

\begin{lemma}
\label{lemma-proper-geometrically-reduced-global-sections}
Soit $k$ un corps. Supposons que $X$ soit un schéma propre sur $k$.
\begin{enumerate}
\item $A = H^0(X, \mathcal{O}_X)$ est une $k$-algèbre
de dimension finie, \item $A = \prod_{i = 1, \ldots,
n} A_i$ est un produit d'$k$-algèbres locales
artiniennes, un facteur pour chaque composante connexe
de $X$, \item si $X$ est réduit, alors $A = \prod_{i =
1, \ldots, n} k_i$ est un produit de corps, chacun
étant une extension finie de $k$, \item si $X$ est
géométriquement réduit, alors $k_i$ est finie et séparable
sur $k$, \item si $X$ est géométriquement connexe, alors
$A$ est géométriquement irréductible sur $k$, \item si
$X$ est géométriquement irréductible, alors $A$ est
géométriquement irréductible sur $k$, \item si $X$ est
géométriquement réduit et géométriquement connexe, alors $A =
k$, et \item si $X$ est géométriquement intègre, alors $A = k$.
\end{enumerate}
\end{lemma}

\begin{proof}
D'après Cohomologie des Schémas, lemme
\ref{coherent-lemma-proper-over-affine-cohomology-finite},
on voit que $A = H^0(X, \mathcal{O}_X)$ est
une $k$-algèbre de dimension finie. Cela prouve (1).

\medskip\noindent
Alors $A$ est un produit de $k$-algèbres artiniennes locales par
Algèbre, lemme
\ref{algebra-lemma-finite-dimensional-algebra} et la proposition
\ref{algebra-proposition-dimension-zero-ring}. Si $X = Y \amalg Z$
avec $Y$ et $Z$ ouverts dans $X$, alors nous obtenons un idempotent $e
\in A$ en prenant la section de $\mathcal{O}_X$ qui est $1$ sur $Y$ et
$0$ sur $Z$. Réciproquement, si $e \in A$ est un idempotent, alors nous
obtenons une décomposition correspondante de $X$. Enfin, comme $X$ a un
espace topologique sous-jacent noethérien, ses composantes connexes
sont ouvertes. Par conséquent, les composantes connexes de $X$
correspondent de $1$ à $1$ avec des idempotents primitifs de $A$. Cela prouve (2).

\medskip\noindent
Si $X$ est réduit, alors $A$ est réduit. Par
conséquent, les anneaux locaux $A_i = k_i$ sont réduits
et donc corps (par exemple d'après Algèbre, lemme
\ref{algebra-lemma-minimal-prime-reduced-ring}). Cela prouve (3).

\medskip\noindent
Si $X$ est géométriquement réduit, alors $A
\otimes_k \overline{k} = H^0(X_{\overline{k}},
\mathcal{O}_{X_{\overline{k}}})$ (égalité par la
cohomologie des Schémas, lemme
\ref{coherent-lemma-flat-base-change-cohomology})
est réduit. Cela implique que $k_i \otimes_k
\overline{k}$ est un produit de corps et donc $k_i/k$
est séparable, par exemple d'après Algèbre, lemmes
\ref{algebra-lemma-characterize-separable-field-extensions}
et
\ref{algebra-lemma-geometrically-reduced-finite-purely-inseparable-extension}. Cela prouve (4).

\medskip\noindent
Si $X$ est géométriquement connexe, alors $A \otimes_k
\overline{k} = H^0(X_{\overline{k}},
\mathcal{O}_{X_{\overline{k}}})$ est un anneau local de
dimension zéro d'après la partie (2) et donc son spectre a un
point, en particulier il est irréductible. Ainsi $A$ est
géométriquement irréductible. Cela prouve (5). Bien sûr, (5) implique (6).

\medskip\noindent
Si $X$ est géométriquement réduit et géométriquement connexe, alors
$A = k_1$ est un corps et l'extension $k_1/k$ est finie séparable et
géométriquement irréductible. Cependant, alors $k_1 \otimes_k
\overline{k}$ est un produit de $[k_1 : k]$ copies de $\overline{k}$ et
nous concluons que $k_1 = k$. Cela prouve (7). Bien sûr, (7) implique (8).
\end{proof}

\noindent
Voici une version simplifiée de la factorisation de Stein ; la
factorisation de Stein proprement dite sera abordée dans Plus sur les
Morphismes, section \ref{more-morphisms-section-stein-factorization}.

\begin{lemma}
\label{lemma-baby-stein}
Supposons que $X$ soit un schéma propre sur un corps $k$.
Posons $A = H^0(X, \mathcal{O}_X)$. Les fibres du morphisme
canonique $X \to \Spec(A)$ sont géométriquement connexes.
\end{lemma}

\begin{proof}
Posons $S = \Spec(A)$. Le morphisme canonique $X \to S$
est le morphisme correspondant à $\Gamma(S, \mathcal{O}_S)
= A = \Gamma(X, \mathcal{O}_X)$ d'après Schémas, lemme
\ref{schemes-lemma-morphism-into-affine}. La $k$-algèbre
$A$ est un produit fini $A = \prod A_i$ de $k$-algèbres
locales artiniennes finies sur $k$, voir le lemme
\ref{lemma-proper-geometrically-reduced-global-sections}.
Désignons par $s_i \in S$ le point correspondant à l'idéal
maximal de $A_i$. Choisissons une clôture algébrique
$\overline{k}$ de $k$ et posons $\overline{A} = A \otimes_k
\overline{k}$. Choisissons un plongement $\kappa(s_i) \to
\overline{k}$ sur $k$ ; cela détermine un morphisme d'algèbres $\overline{k}$.
$$
\sigma_i : \overline{A} = A \otimes_k \overline{k} \to
\kappa(s_i) \otimes_k \overline{k} \to \overline{k}
$$
Considérons le changement de base
$$
\xymatrix{
\overline{X} \ar[r] \ar[d] & X \ar[d] \\
\overline{S} \ar[r] & S
}
$$
de $X$ à $\overline{S} = \Spec(\overline{A})$. Par la
cohomologie des Schémas, lemme
\ref{coherent-lemma-flat-base-change-cohomology} nous avons
$\Gamma(\overline{X}, \mathcal{O}_{\overline{X}}) =
\overline{A}$. Si $\overline{s}_i \in \Spec(\overline{A})$
désigne le point $\overline{k}$-rationnel correspondant à
$\sigma_i$, alors on voit que $\overline{s}_i$ s'envoie sur $s_i
\in S$ et $\overline{X}_{\overline{s}_i}$ est le changement de
base de $X_{s_i}$ par $\Spec(\sigma_i)$. Ainsi, nous voyons qu'il
suffit de prouver le lemme dans le cas où $k$ est algébriquement clos.

\medskip\noindent
Supposons que $k$ soit algébriquement clos. Dans ce cas,
$\kappa(s_i)$ est algébriquement clos et nous devons montrer
que $X_{s_i}$ est connexe. La décomposition en produit $A =
\prod A_i$ correspond à une décomposition en union disjointe
$\Spec(A) = \coprod \Spec(A_i)$, voir Algèbre, lemme
\ref{algebra-lemma-spec-product}. Notons $X_i$ l'image
inverse de $\Spec(A_i)$. Il découle du lemme
\ref{lemma-proper-geometrically-reduced-global-sections},
partie (2), que $A_i = \Gamma(X_i, \mathcal{O}_{X_i})$. Il faut
observer que $X_{s_i} \to X_i$ est une immersion fermée
induisant un isomorphisme sur les espaces topologiques
sous-jacents (parce que $\Spec(A_i)$ est un singleton). Ainsi, si
$X_{s_i}$ n'est pas connexe, alors $X_i$ ne l'est pas non plus.
Donc soit $X_i$ est vide et $A_i = 0$ ou $X_i$ peut être écrit
comme $U \amalg V$ avec $U$ et $V$ ouverts et non vides, ce qui
impliquerait que $A_i$ possède un idempotent non trivial. Puisque
$A_i$ est local, c'est une contradiction et la démonstration est complète.
\end{proof}

\begin{lemma}
\label{lemma-geometrically-reduced-stein}
Soit $k$ un corps. Supposons que $X$ soit un schéma propre
géométriquement réduit sur $k$. Les conditions suivantes sont équivalentes
\begin{enumerate}
\item $H^0(X, \mathcal{O}_X) = k$, et
\item $X$ est géométriquement connexe.
\end{enumerate}
\end{lemma}

\begin{proof}
D'après le lemme \ref{lemma-baby-stein}, nous avons (1)
$\Rightarrow$ (2). D'après le lemme
\ref{lemma-proper-geometrically-reduced-global-sections}, nous avons (2) $\Rightarrow$ (1).
\end{proof}




\section{Schémas géométriquement normaux}
\label{section-geometrically-normal}

\noindent
Dans Propriétés, définition
\ref{properties-definition-normal} nous avons défini la
notion de schéma normal. Cette notion est définie même pour
des schémas non noethériens. Ainsi, contrairement à notre
discussion des schémas «~géométriquement réguliers~», nous
considérons toutes les extensions de corps du corps de base.

\begin{definition}
\label{definition-geometrically-normal}
Soit $X$ un schéma sur le corps $k$.
\begin{enumerate}
\item Soit $x \in X$. On dit que $X$ est {\it
géométriquement normal en $x$ } si pour toute extension
de corps $k'/k$ et tout $x' \in X_{k'}$ au-dessus de $x$
l'anneau local $\mathcal{O}_{X_{k'}, x'}$ est normal.
\item On dit que $X$ est {\it géométriquement normal} sur
$k$ si $X$ est géométriquement normal en tout $x \in X$.
\end{enumerate}
\end{definition}

\begin{lemma}
\label{lemma-geometrically-normal-at-point}
Soit $k$ un corps.
Supposons que $X$ soit un schéma
sur $k$. Soit $x \in X$. Les
énoncés suivants sont équivalents
\begin{enumerate}
\item $X$ est géométriquement normal en $x$, \item
pour toute extension de corps finie purement
inséparable $k'$ de $k$ et $x' \in X_{k'}$ au-dessus
de $x$, l'anneau local $\mathcal{O}_{X_{k'}, x'}$
est normal, et \item l'anneau $\mathcal{O}_{X, x}$
est géométriquement normal sur $k$ (voir Algèbre,
définition \ref{algebra-definition-geometrically-normal}).
\end{enumerate}
\end{lemma}

\begin{proof}
Il est clair que (1) implique (2). Supposons (2).
Supposons que $k'/k$ soit une extension de corps finie
purement inséparable (par exemple $k = k'$). Considérons
l'anneau $\mathcal{O}_{X, x} \otimes_k k'$. D'après
Algèbre, lemme \ref{algebra-lemma-p-ring-map}, son spectre
est le même que le spectre de $\mathcal{O}_{X, x}$. Par
conséquent, c'est aussi un anneau local (Algèbre, lemme
\ref{algebra-lemma-characterize-local-ring}). Il existe donc
un point unique $x' \in X_{k'}$ au-dessus de $x$ et
$\mathcal{O}_{X_{k'}, x'} \cong \mathcal{O}_{X, x} \otimes_k k'$.
Par hypothèse, c'est un anneau normal. Nous en déduisons donc (3)
d'après Algèbre, lemme \ref{algebra-lemma-geometrically-normal}.

\medskip\noindent
Supposons (3). Supposons que $k'/k$ soit une extension de
corps. Comme $\Spec(k') \to \Spec(k)$ est surjectif,
$X_{k'} \to X$ est aussi surjectif (Morphismes, lemme
\ref{morphisms-lemma-base-change-surjective}). Soit
$x' \in X_{k'}$ un point quelconque au-dessus
de $x$. L'anneau local $\mathcal{O}_{X_{k'}, x'}$ est une
localisation de l'anneau $\mathcal{O}_{X, x} \otimes_k k'$.
Par conséquent, il est normal par hypothèse et (1) est prouvé.
\end{proof}

\begin{lemma}
\label{lemma-geometrically-normal}
Soit $k$ un corps.
Supposons que $X$ soit un schéma sur
$k$. Les conditions suivantes sont équivalentes
\begin{enumerate}
\item $X$ est géométriquement normal, \item $X_{k'}$
est un schéma normal pour toute extension de corps
$k'/k$, \item $X_{k'}$ est un schéma normal pour
chaque extension de corps de type fini $k'/k$, \item
$X_{k'}$ est un schéma normal pour chaque extension
de corps finie purement inséparable $k'/k$, \item
pour chaque ouvert affine $U \subset X$ l'anneau
$\mathcal{O}_X(U)$ est géométriquement normal (voir
Algèbre, définition
\ref{algebra-definition-geometrically-normal}), et \item $X_{k^{perf}}$ est un schéma normal.
\end{enumerate}
\end{lemma}

\begin{proof}
Supposons (1). Alors pour toute extension de corps
$k'/k$ et tout point $x' \in X_{k'}$, l'anneau
local de $X_{k'}$ en $x'$ est normal. Par
définition, cela signifie que $X_{k'}$ est normal. D'où (2).

\medskip\noindent
Il est clair que (2) implique (3) implique (4).

\medskip\noindent
Supposons (4) et soit $U \subset X$ un sous-schéma affine
ouvert. Alors $U_{k'}$ est un schéma normal pour toute
extension finie purement inséparable $k'/k$ (y compris $k =
k'$). Cela signifie que $k' \otimes_k \mathcal{O}(U)$ est un
anneau normal pour toutes les extensions finies purement
inséparables $k'/k$. Par conséquent, $\mathcal{O}(U)$ est une algèbre
géométriquement normale $k$ par définition. Ainsi, (4) implique (5).

\medskip\noindent
Supposons (5). Pour toute extension de corps $k'/k$, le
changement de base $X_{k'}$ s'obtient par recollement des spectres
des anneaux $\mathcal{O}_X(U) \otimes_k k'$ où $U$ est un
ouvert affine dans $X$ (voir Schémas, section
\ref{schemes-section-fibre-products}). Donc $X_{k'}$ est normal. Ainsi (1) est vrai.

\medskip\noindent
L'équivalence de (5) et (6) découle de la
définition des algèbres géométriquement normales
et de l'équivalence que l'on vient de démontrer de (3) et (4).
\end{proof}

\begin{lemma}
\label{lemma-geometrically-normal-upstairs}
Soit $k$ un corps. Supposons que $X$
soit un schéma sur $k$. Supposons que $k'/k$ soit
une extension de corps. Supposons que $x \in X$ soit
un point, et soit $x' \in X_{k'}$ un point situé
au-dessus de $x$. Les énoncés suivants sont équivalents
\begin{enumerate}
\item $X$ est géométriquement normal en $x$,
\item $X_{k'}$ est géométriquement normal en $x'$.
\end{enumerate}
En particulier, $X$ est géométriquement normal sur $k$ si et
seulement si $X_{k'}$ est géométriquement normal sur $k'$.
\end{lemma}

\begin{proof}
Il est clair que (1) implique (2). Supposons (2). Supposons que
$k''/k$ soit une extension de corps finie purement inséparable et
soit $x'' \in X_{k''}$ un point au-dessus de $x$ (en fait, il
est unique). Choisissons une extension de corps commune $k'''/k$ de
$k'/k$ et $k''/k$, voir corps, lemme
\ref{fields-lemma-common-extension-field}. Choisissons un point $x''' \in X_{k'''}$
au-dessus à la fois de $x'$ et $x''$. Considérons le morphisme d'anneaux locaux
$$
\mathcal{O}_{X_{k''}, x''} \longrightarrow \mathcal{O}_{X_{k'''}, x''''}.
$$
Ceci est un homomorphisme local plat d'anneaux et donc
fidèlement plat. D'après (2), on voit que l'anneau local à
droite est normal. Ainsi, d'après L'Algèbre, lemme
\ref{algebra-lemma-descent-normal}, nous concluons que
$\mathcal{O}_{X_{k''}, x''}$ est normal. D'après le lemme
\ref{lemma-geometrically-normal-at-point}, on voit que $X$ est géométriquement normal en $x$.
\end{proof}

\begin{lemma}
\label{lemma-fibre-product-normal}
Soit $k$ un corps. Supposons que $X$ soit un
schéma géométriquement normal sur $k$ et soit $Y$ un schéma
normal sur $k$. Alors $X \times_k Y$ est un schéma normal.
\end{lemma}

\begin{proof}
Cela se ramène à l'algèbre, lemme
\ref{algebra-lemma-geometrically-normal-tensor-normal}
par le lemme
\ref{lemma-geometrically-normal}.
\end{proof}

\begin{lemma}
\label{lemma-base-change-normal-by-separable}
Soit $k$ un corps. Supposons que $X$ soit un schéma normal sur $k$.
Supposons que $K/k$ soit une extension de corps séparable. Alors $X_K$ est un schéma normal.
\end{lemma}

\begin{proof}
Cela résulte du lemme
\ref{lemma-fibre-product-normal} et du lemme d'Algèbre
\ref{algebra-lemma-separable-field-extension-geometrically-normal}.
\end{proof}

\begin{lemma}
\label{lemma-geometrically-normal-stein}
Soit $k$ un corps. Supposons que $X$ soit un schéma propre
géométriquement normal sur $k$. Les conditions suivantes sont équivalentes
\begin{enumerate}
\item $H^0(X, \mathcal{O}_X) = k$, \item
$X$ est géométriquement connexe, \item
$X$ est géométriquement irréductible, et
\item $X$ est géométriquement intègre.
\end{enumerate}
\end{lemma}

\begin{proof}
D'après le lemme
\ref{lemma-geometrically-reduced-stein}, nous avons
l'équivalence de (1) et (2). Un schéma normal localement
noethérien (tel que $X_{\overline{k}}$) est une réunion
disjointe de ses composantes irréductibles (Propriétés, lemme
\ref{properties-lemma-normal-Noetherian}). Ainsi, on voit que
(2) et (3) sont équivalents. Puisque $X_{\overline{k}}$ est
supposé réduit, on voit que (3) et (4) sont également équivalents.
\end{proof}









\section{Changement de corps et schémas localement noethériens}
\label{section-locally-Noetherian}

\noindent
Soit $X$ un schéma localement noethérien sur un corps $k$. Il n'est pas
toujours vrai que $X_{k'}$ soit aussi localement noethérien. Par
exemple, si $X = \Spec(\overline{\mathbf{Q}})$ et $k = \mathbf{Q}$,
alors $X_{\overline{\mathbf{Q}}}$ est le spectre de
$\overline{\mathbf{Q}} \otimes_{\mathbf{Q}} \overline{\mathbf{Q}}$ qui n'est
pas noethérien. (Indice : Il a trop d'idempotents). Mais si nous ne faisons
que des changements de base utilisant des extensions de corps de type fini,
alors la propriété noethérienne est préservée. (Ou si $X$ est localement de
type fini sur $k$, puisque cette propriété est préservée par changement de base.)

\begin{lemma}
\label{lemma-locally-Noetherian-base-change}
Soit $k$ un corps. Supposons
que $X$ soit un schéma sur $k$. Supposons que
$k'/k$ soit une extension de corps de type
fini. Alors $X$ est localement noethérien si et
seulement si $X_{k'}$ est localement noethérien.
\end{lemma}

\begin{proof}
D'après Propriétés, lemme
\ref{properties-lemma-locally-Noetherian} nous réduisons au cas où
$X$ est affine, disons $X = \Spec(A)$. Dans ce cas, nous devons
prouver que $A$ est noethérien si et seulement si $A_{k'}$ est
noethérien. Puisque $A \to A_{k'} = k' \otimes_k A$ est fidèlement
plat, on voit que si $A_{k'}$ est noethérien, alors $A$ l'est
aussi, d'après Algèbre, lemme
\ref{algebra-lemma-descent-Noetherian}. Réciproquement, si $A$ est noethérien, alors
$A_{k'}$ est noethérien d'après Algèbre, lemme \ref{algebra-lemma-Noetherian-field-extension}.
\end{proof}







\section{Schémas géométriquement réguliers}
\label{section-geometrically-regular}

\noindent
Un schéma géométriquement régulier sur un corps $k$ est un schéma
localement noethérien sur $k$ qui reste régulier après des changements
appropriés de corps de base. Un schéma de type fini sur $k$ est
géométriquement régulier si et seulement s'il est lisse sur $k$ (voir le lemme
\ref{lemma-geometrically-regular-smooth}). La notion de régularité géométrique
est particulièrement intéressante dans les situations où la lissité ne peut pas
être utilisée, comme pour les fibres formelles (insérer référence future ici).

\medskip\noindent
Dans la définition suivante, nous nous limitons aux schémas
localement noethériens, puisque la propriété d'être un anneau
local régulier n'est définie que pour les anneaux locaux
noethériens. D'après le lemme
\ref{lemma-locally-Noetherian-base-change} ci-dessus, si nous nous limitons
aux extensions de corps de type fini, alors cette propriété est préservée par
changement de corps de base. Ce commentaire sera utilisé sans référence
supplémentaire dans cette section. En particulier, la définition suivante a un sens.

\begin{definition}
\label{definition-geometrically-regular}
Soit $k$ un corps. Supposons que $X$ soit un schéma localement noethérien sur $k$.
\begin{enumerate}
\item Soit $x \in X$. On dit que $X$ est {\it
géométriquement régulier en $x$} sur $k$ si pour toute
extension de corps de type fini $k'/k$ et tout $x' \in X_{k'}$
au-dessus de $x$, l'anneau local $\mathcal{O}_{X_{k'}, x'}$ est
régulier. \item On dit que $X$ est {\it géométriquement régulier
sur $k$} si $X$ est géométriquement régulier en tout point.
\end{enumerate}
\end{definition}

\noindent
Une définition analogue permet de définir les schémas géométriquement
de Cohen--Macaulay, $(R_k)$ et $(S_k)$ sur un corps. Nous ajouterons
des sections séparées lorsque cela sera nécessaire.

\begin{lemma}
\label{lemma-geometrically-regular-at-point}
Soit $k$ un corps.
Supposons que $X$ soit un schéma
localement noethérien sur $k$. Soit $x \in
X$. Les assertions suivantes sont équivalentes
\begin{enumerate}
\item $X$ est géométriquement régulier en $x$, \item
pour toute extension de corps finie purement
inséparable $k'$ de $k$ et $x' \in X_{k'}$ située
au-dessus de $x$ l'anneau local $\mathcal{O}_{X_{k'},
x'}$ est régulier, et \item l'anneau $\mathcal{O}_{X,
x}$ est géométriquement régulier sur $k$ (voir Algèbre,
définition \ref{algebra-definition-geometrically-regular}).
\end{enumerate}
\end{lemma}

\begin{proof}
Il est clair que (1) implique (2). Supposons (2). Cela implique
en particulier que $\mathcal{O}_{X, x}$ est un anneau local
régulier. Supposons que $k'/k$ soit une extension de corps finie
purement inséparable. Considérons l'anneau $\mathcal{O}_{X, x}
\otimes_k k'$. D'après Algèbre, lemme
\ref{algebra-lemma-p-ring-map}, son spectre est le même que le
spectre de $\mathcal{O}_{X, x}$. Par conséquent, c'est aussi un
anneau local (Algèbre, lemme
\ref{algebra-lemma-characterize-local-ring}). Donc, il y a un point unique $x'
\in X_{k'}$ au-dessus de $x$ et $\mathcal{O}_{X_{k'}, x'} \cong
\mathcal{O}_{X, x} \otimes_k k'$. Par hypothèse, il s'agit d'un anneau régulier.
Ainsi, nous déduisons (3) de la définition d'un anneau géométriquement régulier.

\medskip\noindent
Supposons (3). Supposons que $k'/k$ soit une extension de
corps. Comme $\Spec(k') \to \Spec(k)$ est surjectif,
$X_{k'} \to X$ est aussi surjectif (Morphismes, lemme
\ref{morphisms-lemma-base-change-surjective}). Supposons que
$x' \in X_{k'}$ est un point quelconque situé au-dessus de
$x$. L'anneau local $\mathcal{O}_{X_{k'}, x'}$ est une
localisation de l'anneau $\mathcal{O}_{X, x} \otimes_k k'$.
Par conséquent, il est régulier par hypothèse et (1) est prouvé.
\end{proof}

\begin{lemma}
\label{lemma-geometrically-regular}
Soit $k$ un corps. Supposons que
$X$ soit un schéma localement noethérien sur
$k$. Les affirmations suivantes sont équivalentes
\begin{enumerate}
\item $X$ est géométriquement régulier, \item $X_{k'}$
est un schéma régulier pour toute extension de corps
de type fini $k'/k$, \item $X_{k'}$ est un schéma
régulier pour toute extension de corps finie purement
inséparable $k'/k$, \item pour tout ouvert affine $U
\subset X$ l'anneau $\mathcal{O}_X(U)$ est
géométriquement régulier (voir Algèbre, définition
\ref{algebra-definition-geometrically-regular}), et \item il
existe un recouvrement ouvert affine $X = \bigcup U_i$ tel que
chaque $\mathcal{O}_X(U_i)$ est géométriquement régulier sur $k$.
\end{enumerate}
\end{lemma}

\begin{proof}
Supposons (1). Alors pour toute extension de corps de
type fini $k'/k$ et pour tout point $x' \in X_{k'}$,
l'anneau local de $X_{k'}$ en $x'$ est régulier. D'après
les Propriétés, lemme
\ref{properties-lemma-characterize-regular}, cela signifie que $X_{k'}$ est régulier. D'où (2).

\medskip\noindent
Il est clair que (2) implique (3).

\medskip\noindent
Supposons (3) et soit $U \subset X$ un sous-schéma affine
ouvert. Alors $U_{k'}$ est un schéma régulier pour toute
extension finie purement inséparable $k'/k$ (y compris $k =
k'$). Cela signifie que $k' \otimes_k \mathcal{O}(U)$ est un
anneau régulier pour toutes les extensions finies purement
inséparables $k'/k$. Par conséquent, $\mathcal{O}(U)$ est une
algèbre $k$ géométriquement régulière et on voit que (4) est vrai.

\medskip\noindent
Il est clair que (4) implique (5). Supposons que $X = \bigcup
U_i$ soit un recouvrement affine ouvert comme dans (5). Pour toute
extension de corps $k'/k$, le changement de base $X_{k'}$ est
obtenu en recollant les spectres des anneaux $\mathcal{O}_X(U_i)
\otimes_k k'$ (voir Schémas, section
\ref{schemes-section-fibre-products}). Par conséquent, $X_{k'}$ est régulier. Donc (1) est vérifié.
\end{proof}

\begin{lemma}
\label{lemma-geometrically-regular-upstairs}
Soit $k$ un corps. Supposons que $X$ soit
un schéma sur $k$. Supposons que $k'/k$ est une
extension de corps de type fini. Supposons que $x \in
X$ est un point, et soit $x' \in X_{k'}$ un point situé
au-dessus de $x$. Les énoncés suivants sont équivalents
\begin{enumerate}
\item $X$ est géométriquement régulier en $x$,
\item $X_{k'}$ est géométriquement régulier en $x'$.
\end{enumerate}
En particulier, $X$ est géométriquement régulier sur $k$ si et
seulement si $X_{k'}$ est géométriquement régulier sur $k'$.
\end{lemma}

\begin{proof}
Il est clair que (1) implique (2). Supposons (2). Supposons que
$k''/k$ soit une extension de corps finie purement inséparable
et soit $x'' \in X_{k''}$ un point au-dessus de $x$ (en fait il
est unique). Nous pouvons trouver une extension de corps
commune, de type fini, $k'''/k$ de $k'/k$ et $k''/k$, voir
corps, lemme \ref{fields-lemma-common-extension-field} et sa
démonstration, et un point $x''' \in X_{k'''}$ se trouvant au-dessus
à la fois de $x'$ et $x''$. Considérons le morphisme d’anneaux locaux
$$
\mathcal{O}_{X_{k''}, x''} \longrightarrow \mathcal{O}_{X_{k'''}, x''''}.
$$
Ceci est un homomorphisme plat local d'anneaux locaux
noethériens et donc fidèlement plat. D'après (2), nous
voyons que l'anneau local à droite est régulier. Ainsi,
d'après Algèbre, lemme
\ref{algebra-lemma-flat-under-regular}, nous concluons que
$\mathcal{O}_{X_{k''}, x''}$ est régulier. D'après le lemme
\ref{lemma-geometrically-regular-at-point}, on voit que $X$ est géométriquement régulier en $x$.
\end{proof}

\noindent
Le lemme suivant est une variante géométrique de
l'Algèbre, lemme \ref{algebra-lemma-geometrically-regular-descent}.

\begin{lemma}
\label{lemma-flat-under-geometrically-regular}
Soit $k$ un corps. Supposons que $f : X \to Y$
soit un morphisme de schémas localement noethériens sur $k$.
Supposons que $x \in X$ soit un point et posons $y = f(x)$. Si
$X$ est géométriquement régulier en $x$ et que $f$ est plat en
$x$ alors $Y$ est géométriquement régulier en $y$. En
particulier, si $X$ est géométriquement régulier sur $k$ et que $f$
est plat et surjectif, alors $Y$ est géométriquement régulier sur $k$.
\end{lemma}

\begin{proof}
Supposons que $k'$ soit une extension finie purement
inséparable de $k$. Supposons que $f' : X_{k'} \to
Y_{k'}$ soit le changement de base de $f$. Supposons que
$x' \in X_{k'}$ soit le point unique situé au-dessus de
$x$. Si nous montrons que $Y_{k'}$ est régulier en $y' =
f'(x')$, alors $Y$ est géométriquement régulier sur $k$
en $y'$, voir le lemme
\ref{lemma-geometrically-regular}. D'après Morphismes, lemme
\ref{morphisms-lemma-base-change-module-flat} le morphisme $X_{k'}
\to Y_{k'}$ est plat en $x'$. Par conséquent, l'homomorphisme d'anneaux
$$
\mathcal{O}_{Y_{k'}, y'}
\longrightarrow
\mathcal{O}_{X_{k'}, x'}
$$
est un homomorphisme local plat d'anneaux locaux
noethériens avec le côté droit régulier par hypothèse. Par
conséquent, le côté gauche est un anneau local régulier
d'après Algèbre, lemme \ref{algebra-lemma-flat-under-regular}.
\end{proof}

\begin{lemma}
\label{lemma-geometrically-regular-smooth}
Soit $k$ un corps. Supposons que $X$
soit un schéma localement de type fini sur $k$. Soit $x
\in X$. Alors $X$ est géométriquement régulier en $x$
si et seulement si $X \to \Spec(k)$ est lisse en $x$
(Morphismes, définition \ref{morphisms-definition-smooth}).
\end{lemma}

\begin{proof}
La question est locale autour de $x$, donc
on peut supposer que $X = \Spec(A)$
pour une certaine $k$-algèbre de type fini.
Soit $\mathfrak p$ l'idéal premier correspondant à $x$.

\medskip\noindent
Si $A$ est lisse sur $k$ à $\mathfrak p$, alors nous pouvons
localiser $A$ et supposer que $A$ est lisse sur $k$. Dans ce cas,
$k' \otimes_k A$ est lisse sur $k'$ pour tout corps d'extension
$k'/k$, et chacun de ces anneaux noethériens est régulier selon
l'Algèbre, lemme \ref{algebra-lemma-characterize-smooth-over-field}.

\medskip\noindent
Supposons que $X$ soit géométriquement régulier en $x$.
Considérons le corps résiduel $K := \kappa(x) =
\kappa(\mathfrak p)$ de $x$. C'est une extension de type
fini de $k$. D'après Algèbre, lemme
\ref{algebra-lemma-make-separable}, il existe une
extension finie purement inséparable $k'/k$ telle que le
compositum $k'K$ soit une extension de corps séparable de
$k'$. Supposons que $\mathfrak p' \subset A' = k'
\otimes_k A$ soit un idéal premier au-dessus de $\mathfrak
p$. C'est le seul premier au-dessus de $\mathfrak p$, voir
Algèbre, lemme \ref{algebra-lemma-p-ring-map}. Ainsi, le
corps résiduel $K' := \kappa(\mathfrak p')$ est le
compositum $k'K$. Par hypothèse, l'anneau local
$(A')_{\mathfrak p'}$ est régulier. Ainsi, d'après l'Algèbre,
lemme \ref{algebra-lemma-separable-smooth}, on voit que $k'
\to A'$ est lisse en $\mathfrak p'$. Cela implique à son tour
que $k \to A$ est lisse en $\mathfrak p$ d'après Algèbre, lemme
\ref{algebra-lemma-smooth-field-change-local}. Le lemme est prouvé.
\end{proof}


\begin{example}
\label{example-geometrically-reduced-not-normal}
Soit $k =\mathbf{F}_p(t)$. Il est assez facile de donner un
exemple d'une variété régulière $V$ sur $k$ qui n'est pas
géométriquement réduite. Par exemple, nous pouvons prendre
$\Spec(k[x]/(x^p - t))$. En fait, il existe un exemple d'une
variété régulière $V$ qui est géométriquement réduite, mais pas même
géométriquement normale. À savoir, prenons pour $p > 2$ le schéma $V
= \Spec(k[x, y]/(y^2 - x^p + t))$. C'est une variété puisque le
polynôme $y^2 - x^p + t \in k[x, y]$ est irréductible. Le morphisme
$V \to \Spec(k)$ est lisse en tous les points sauf au point $v_0 \in
V$ correspondant à l'idéal maximal $(y, x^p - t)$ (car $2y$ est
inversible). En particulier, on voit que $V$ est (géométriquement)
régulière en tous les points, sauf éventuellement $v_0$. L'anneau local
$$
\mathcal{O}_{V, v_0} = \left(k[x, y]/(y^2 - x^p + t)\right)_{(y, x^p - t)}
$$
est un anneau intègre de dimension $1$. Son idéal maximal est engendré
par un seul élément, à savoir $y$. Par conséquent, c'est un anneau
à valuation discrète et régulier. Soit $k' = k[t^{1/p}]$. On note
$t' = t^{1/p} \in k'$, $V' = V_{k'}$, $v'_0 \in V'$ le point
unique au-dessus de $v_0$. Au-dessus de $k'$, nous pouvons écrire
$x^p - t = (x - t')^p$, mais le polynôme $y^2 - (x - t')^p$ est
toujours irréductible et $V'$ est toujours une variété. Mais l'élément
$$
\frac{y}{x - t'} \in (\text{fraction field of }\mathcal{O}_{V', v'_0})
$$
est intégral sur $\mathcal{O}_{V', v'_0}$ (il suffit
de calculer son carré) et n'y est pas contenu, donc
$V'$ n'est pas normal en $v'_0$. Cela conclut l'exemple.
\end{example}







\section{Changement de corps et propriété de Cohen--Macaulay}
\label{section-CM}

\noindent
Le lemme suivant dit qu'il n'a pas de sens de définir
des schémas géométriquement de Cohen--Macaulay, puisque
ceux-ci seraient les mêmes que les schémas de Cohen-Macaulay.

\begin{lemma}
\label{lemma-CM-base-change}
Supposons que $X$ soit un schéma localement noethérien sur le
corps $k$. Supposons que $k'/k$ soit une extension de corps de
type fini. Supposons que $x \in X$ soit un point, et soit $x'
\in X_{k'}$ un point au-dessus de $x$. Alors nous avons
$$
\mathcal{O}_{X, x}\text{ is Cohen-Macaulay}
\Leftrightarrow
\mathcal{O}_{X_{k'}, x'}\text{ is Cohen-Macaulay}
$$
Si $X$ est localement de type fini sur $k$, il en
va de même pour toute extension de corps $k'/k$.
\end{lemma}

\begin{proof}
Le premier cas du lemme découle de l'Algèbre,
lemme \ref{algebra-lemma-CM-geometrically-CM}. Le
deuxième cas du lemme est équivalent à l'Algèbre,
lemme \ref{algebra-lemma-extend-field-CM-locus}.
\end{proof}







\section{Changement de corps et propriété de Jacobson}
\label{section-overfield}

\noindent
Un schéma localement de type fini sur un corps a beaucoup
de points fermés, à savoir qu'il est de Jacobson. De plus, les
corps résiduels sont des extensions finies du corps de base.

\begin{lemma}
\label{lemma-locally-finite-type-Jacobson}
Supposons que $X$ soit un schéma qui est
localement de type fini sur $k$. Alors
\begin{enumerate}
\item pour tout point fermé $x \in X$
l'extension $\kappa(x)/k$ est algébrique, et
\item $X$ est un schéma de Jacobson (Propriétés,
définition \ref{properties-definition-jacobson} ).
\end{enumerate}
\end{lemma}

\begin{proof}
Un schéma est de Jacobson si et seulement s'il possède un
recouvrement par des ouverts affines de schémas de Jacobson,
voir Propriétés, lemme
\ref{properties-lemma-locally-jacobson}. La propriété sur les corps
résiduels aux points fermés est également locale sur $X$. Par
conséquent, on peut supposer que $X$ est affine. Dans ce cas,
le résultat est une conséquence du Nullstellensatz de Hilbert, voir
Algèbre, théorème \ref{algebra-theorem-nullstellensatz}. Il découle
également d'une combinaison de Morphismes, lemmes
\ref{morphisms-lemma-jacobson-finite-type-points},
\ref{morphisms-lemma-Jacobson-universally-Jacobson} et \ref{morphisms-lemma-ubiquity-Jacobson-schemes}.
\end{proof}

\noindent
Il s'avère que si $X$ n'est pas localement de type fini, alors nous pouvons obtenir
le même résultat après avoir fait une extension du corps de base suffisamment grande.

\begin{lemma}
\label{lemma-make-Jacobson}
Supposons que $X$ soit un schéma sur un corps
$k$. Pour toute extension de corps $K/k$ dont la
cardinalité est suffisamment grande, nous avons
\begin{enumerate}
\item pour tout point fermé $x \in X_K$
l'extension $\kappa(x)/K$ est algébrique, et
\item $X_K$ est un schéma de Jacobson (Propriétés,
définition \ref{properties-definition-jacobson} ).
\end{enumerate}
\end{lemma}

\begin{proof}
Choisissons un recouvrement ouvert affine $X =
\bigcup U_i$. D'après Algèbre, lemme
\ref{algebra-lemma-base-change-Jacobson} et les
Propriétés, lemme
\ref{properties-lemma-affine-jacobson}, il existe des
cardinaux $\kappa_i$ tels que $U_{i, K}$ possède les
propriétés souhaitées sur $K$ si $\#(K) \geq \kappa_i$.
Posons $\kappa = \max\{\kappa_i\}$. Alors, si la cardinalité
de $K$ est supérieure à $\kappa$, on voit que chaque
$U_{i, K}$ satisfait les conclusions du lemme. Par
conséquent, $X_K$ est Jacobson d'après Propriétés, lemme
\ref{properties-lemma-locally-jacobson}. L'énoncé sur les corps
résiduels aux points fermés de $X_K$ découle des énoncés
correspondants pour les corps résiduels des points fermés de $U_{i, K}$.
\end{proof}





\section{Changement de corps et faisceaux inversibles amples}
\label{section-change-fields-ample}

\noindent
Le résultat suivant est typique des résultats de cette section.

\begin{lemma}
\label{lemma-ample-after-field-extension}
Soit $k$ un corps. Supposons que $X$ soit
un schéma sur $k$. S'il existe un faisceau inversible
ample sur $X_K$ pour une certaine extension de corps
$K/k$, alors $X$ possède un faisceau inversible ample.
\end{lemma}

\begin{proof}
Supposons que $K/k$ soit une extension de corps telle que
$X_K$ possède un faisceau inversible ample $\mathcal{L}$. Le
morphisme $X_K \to X$ est surjectif. Par conséquent, $X$ est
quasi-compact en tant qu'image d'un schéma quasi-compact
(Propriétés, définition \ref{properties-definition-ample}).
Puisque $X_K$ est quasi-séparé (d'après Propriétés, lemme
\ref{properties-lemma-affine-s-opens-cover-quasi-separated}), nous
voyons que $X$ est quasi-séparé : si $U, V \subset X$ sont ouverts
affines, alors $(U \cap V)_K = U_K \cap V_K$ est quasi-compact et
$(U \cap V)_K \to U \cap V$ est surjectif. Ainsi, le lemme
\ref{schemes-lemma-characterize-quasi-separated} sur les schémas s'applique.

\medskip\noindent
Écrivons $K = \colim A_i$ comme colimite des sous-algèbres de
$K$ qui sont de type fini sur $k$. Notons $X_i = X
\times_{\Spec(k)} \Spec(A_i)$. Puisque $X_K = \lim X_i$, nous
trouvons un $i$ et un faisceau inversible $\mathcal{L}_i$ sur
$X_i$ dont le image réciproque sur $X_K$ est $\mathcal{L}$
(Limites, lemme \ref{limits-lemma-descend-invertible-modules} ;
ici et ci-dessous, nous utilisons que $X$ est quasi-compact et
quasi-séparé comme montré précédemment). D'après Limites, lemme
\ref{limits-lemma-limit-ample}, on peut supposer que
$\mathcal{L}_i$ est ample après avoir éventuellement augmenté
$i$. Fixons un tel $i$ et soit $\mathfrak m \subset A_i$ un idéal
maximal. Par le Nullstellensatz de Hilbert (Algèbre, théorème
\ref{algebra-theorem-nullstellensatz}), le corps résiduel $k' =
A_i/\mathfrak m$ est une extension finie de $k$. Ainsi $X_{k'}
\subset X_i$ est un sous-schéma fermé et possède donc un faisceau
inversible ample (Propriétés, lemme
\ref{properties-lemma-ample-on-closed}). Puisque $X_{k'} \to X$ est localement
libre de rang fini, nous en concluons que $X$ possède un faisceau inversible
ample d’après Diviseurs, proposition \ref{divisors-proposition-push-down-ample}.
\end{proof}

\begin{lemma}
\label{lemma-quasi-affine-after-field-extension}
Soit $k$ un corps. Soit $X$ un schéma sur $k$. Si $X_K$ est
quasi-affine pour une certaine extension de corps $K/k$, alors $X$ est quasi-affine.
\end{lemma}

\begin{proof}
Supposons que $K/k$ soit une extension de corps telle que $X_K$ soit
quasi-affine. Le morphisme $X_K \to X$ est surjectif. Ainsi $X$ est
quasi-compact comme l'image d'un schéma quasi-compact (Propriétés,
définition \ref{properties-definition-quasi-affine}). Puisque $X_K$
est quasi-séparé (en tant que sous-schéma ouvert d'un schéma affine),
on voit que $X$ est quasi-séparé : si $U, V \subset X$ sont
ouverts affines, alors $(U \cap V)_K = U_K \cap V_K$ est quasi-compact
et $(U \cap V)_K \to U \cap V$ est surjectif. Ainsi, le lemme sur les
schémas, \ref{schemes-lemma-characterize-quasi-separated}, s'applique.

\medskip\noindent
Écrivons $K = \colim A_i$ comme colimite des
sous-algèbres de $K$ qui sont de type fini sur $k$. Noter
$X_i = X \times_{\Spec(k)} \Spec(A_i)$. Puisque $X_K = \lim
X_i$, nous trouvons un $i$ tel que $X_i$ soit quasi-affine
(Limites, lemme \ref{limits-lemma-limit-quasi-affine} ; ici
nous utilisons que $X$ est quasi-compact et quasi-séparé
comme montré précédemment). Par le Nullstellensatz de
Hilbert (Algèbre, théorème
\ref{algebra-theorem-nullstellensatz}), le corps résiduel $k' =
A_i/\mathfrak m$ est une extension finie de $k$. Ainsi $X_{k'}
\subset X_i$ est un sous-schéma fermé donc est quasi-affine
(Propriétés, lemme
\ref{properties-lemma-quasi-affine-locally-closed}). Puisque $X_{k'} \to X$ est localement libre
de type fini, nous concluons par Diviseurs, lemme \ref{divisors-lemma-push-down-quasi-affine}.
\end{proof}

\begin{lemma}
\label{lemma-quasi-projective-after-field-extension}
Soit $k$ un corps. Supposons que $X$ soit un schéma
sur $k$. Si $X_K$ est quasi-projectif sur $K$ pour une certaine
extension de corps $K/k$, alors $X$ est quasi-projectif sur $k$.
\end{lemma}

\begin{proof}
Par définition, un morphisme de schémas $g : Y \to T$ est
quasi-projectif s'il est localement de type fini, quasi-compact,
et s'il existe un faisceau inversible $g$-ample sur $Y$.
Supposons que $K/k$ soit une extension de corps telle que $X_K$
soit quasi-projectif sur $K$. Supposons que $\Spec(A) \subset X$
soit un ouvert affine. Alors $U_K$ est un sous-schéma ouvert
affine de $X_K$, donc $A_K$ est une $K$-algèbre de type fini.
Alors $A$ est une $k$-algèbre de type fini d'après le lemme
d'algèbre \ref{algebra-lemma-finite-type-descends}. Par conséquent,
$X \to \Spec(k)$ est localement de type fini. Puisque $X_K \to
\Spec(K)$ est quasi-compact, on voit que $X_K$ est
quasi-compact, donc $X$ est quasi-compact, donc $X \to \Spec(k)$ est
de type fini. D'après le lemme sur les morphismes
\ref{morphisms-lemma-finite-type-over-affine-ample-very-ample}, nous
voyons que $X_K$ possède un faisceau inversible ample. Ensuite, $X$ possède
un faisceau inversible ample d'après le lemme
\ref{lemma-ample-after-field-extension}. Par conséquent, $X \to \Spec(k)$ est
quasi-projectif d'après Morphismes, lemme \ref{morphisms-lemma-finite-type-over-affine-ample-very-ample}.
\end{proof}

\noindent
Le lemme suivant est un cas particulier de la Descente,
lemme \ref{descent-lemma-descending-property-proper}.

\begin{lemma}
\label{lemma-proper-after-field-extension}
Soit $k$ un corps. Supposons que $X$ soit un
schéma sur $k$. Si $X_K$ est propre sur $K$ pour une
certaine extension de corps $K/k$, alors $X$ est propre sur $k$.
\end{lemma}

\begin{proof}
Supposons que $K/k$ soit une extension de corps telle que $X_K$ soit
propre sur $K$. Rappelons que cela implique que $X_K$ est séparé et
quasi-compact (Morphismes, définition
\ref{morphisms-definition-proper}). Le morphisme $X_K \to X$ est
surjectif. Ainsi, $X$ est quasi-compact comme l'image d'un schéma
quasi-compact (Propriétés, définition \ref{properties-definition-ample}).
Puisque $X_K$ est séparé, on voit que $X$ est quasi-séparé : si $U, V
\subset X$ sont ouverts affines, alors $(U \cap V)_K = U_K \cap V_K$ est
quasi-compact et $(U \cap V)_K \to U \cap V$ est surjectif. Ainsi, le
lemme des schémas \ref{schemes-lemma-characterize-quasi-separated} s'applique.

\medskip\noindent
Écrivons $K = \colim A_i$ comme colimite des
sous-algèbres de $K$ qui sont de type fini sur $k$.
Notons $X_i = X \times_{\Spec(k)} \Spec(A_i)$. D'après
Limites, lemme \ref{limits-lemma-eventually-proper}, il
existe un $i$ tel que $X_i \to \Spec(A_i)$ soit propre.
Ici, nous utilisons que $X$ est quasi-compact et
quasi-séparé comme montré précédemment. Choisissons un
idéal maximal $\mathfrak m \subset A_i$. D'après le
Nullstellensatz de Hilbert (Algèbre, théorème
\ref{algebra-theorem-nullstellensatz}), le corps résiduel
$k' = A_i/\mathfrak m$ est une extension finie de $k$. Le
changement de base $X_{k'} \to \Spec(k')$ est propre
(Morphismes, lemme
\ref{morphisms-lemma-base-change-proper}). Puisque $k'/k$ est
fini, à la fois $X_{k'} \to X$ et la composition $X_{k'} \to
\Spec(k)$ sont également propres (Morphismes, lemmes
\ref{morphisms-lemma-finite-proper},
\ref{morphisms-lemma-base-change-proper} et
\ref{morphisms-lemma-composition-proper}). Le premier implique que $X$ est
séparé sur $k$ puisque $X_{k'}$ est séparé (Morphismes, lemme
\ref{morphisms-lemma-image-universally-closed-separated}). Le second implique que $X \to
\Spec(k)$ est propre d'après Morphismes, lemme \ref{morphisms-lemma-image-proper-is-proper}.
\end{proof}

\begin{lemma}
\label{lemma-projective-after-field-extension}
Soit $k$ un corps. Supposons que $X$ soit un
schéma sur $k$. Si $X_K$ est projectif sur $K$ pour une
certaine extension de corps $K/k$, alors $X$ est projectif sur $k$.
\end{lemma}

\begin{proof}
Un schéma sur $k$ est projectif sur $k$ si et
seulement s'il est quasi-projectif et propre sur $k$.
Voir Morphismes, lemme
\ref{morphisms-lemma-projective-is-quasi-projective-proper}.
Ainsi, le lemme découle des lemmes
\ref{lemma-quasi-projective-after-field-extension} et \ref{lemma-proper-after-field-extension}.
\end{proof}




\section{Espaces tangents}
\label{section-tangent-spaces}

\noindent
Dans cette section, nous définissons l'espace tangent d'un morphisme de schémas
en un point de la source à l'aide de points à valeurs dans les nombres duaux.

\begin{definition}
\label{definition-dual-numbers}
Pour tout anneau $R$, l'{\it algèbre des nombres duaux} sur $R$
est la $R$-algèbre notée $R[\epsilon]$. Comme $R$-module,
elle est libre de base $1$, $\epsilon$; la structure de
$R$-algèbre est définie par la relation $\epsilon^2 = 0$.
\end{definition}

\noindent
Soit $f : X \to S$ un morphisme de schémas.
Soit $x \in X$ un point d'image $s =
f(x)$ dans $S$. Considérons le diagramme commutatif formé de flèches pleines
\begin{equation}
\label{equation-tangent-space}
\vcenter{
\xymatrix{
\Spec(\kappa(x)) \ar[r] \ar[dr] \ar@/^1pc/[rr] &
\Spec(\kappa(x)[\epsilon]) \ar@{.>}[r] \ar[d]&
X \ar[d] \\
&
\Spec(\kappa(s)) \ar[r] &
S
}
}
\end{equation}
où la flèche courbe est le morphisme
canonique de $\Spec(\kappa(x))$ vers $X$.

\begin{lemma}
\label{lemma-tangent-space}
L'ensemble des flèches pointillées rendant commutatif le diagramme
(\ref{equation-tangent-space}) possède une structure canonique de $\kappa(x)$-espace vectoriel.
\end{lemma}

\begin{proof}
Posons $\kappa = \kappa(x)$. Observons que la somme
amalgamée suivante existe dans la catégorie des schémas :
$$
\Spec(\kappa[\epsilon]) \amalg_{\Spec(\kappa)} \Spec(\kappa[\epsilon])
= \Spec(\kappa[\epsilon_1, \epsilon_2])
$$
où $\kappa[\epsilon_1, \epsilon_2]$ est la $\kappa$-algèbre
libre comme module, de base $1, \epsilon_1, \epsilon_2$, avec
$\epsilon_1^2 = \epsilon_1\epsilon_2 =
\epsilon_2^2 = 0$. Cela découle immédiatement du
résultat correspondant pour les anneaux et de la
description des morphismes du spectre d'un anneau
local vers un schéma; voir Schémas, lemme
\ref{schemes-lemma-morphism-from-spec-local-ring}.
Étant donné deux flèches $\theta_1, \theta_2 :
\Spec(\kappa[\epsilon]) \to X$, nous pouvons considérer le morphisme
$$
\theta_1 + \theta_2 :
\Spec(\kappa[\epsilon]) \to
\Spec(\kappa[\epsilon_1, \epsilon_2])
\xrightarrow{\theta_1, \theta_2} X
$$
où la première flèche est donnée par $\epsilon_i \mapsto \epsilon$.
D'autre part, pour $\lambda \in \kappa$, il existe un endomorphisme de
$\Spec(\kappa[\epsilon])$ correspondant à l'endomorphisme
de $\kappa$-algèbres de $\kappa[\epsilon]$ qui envoie $\epsilon$ sur
$\lambda \epsilon$. La précomposition de $\theta : \Spec(\kappa[\epsilon]) \to X$ par
cet endomorphisme donne $\lambda \theta$. On vérifie
les axiomes d'un espace vectoriel au moyen de diagrammes
commutatifs convenables de schémas. Nous omettons les détails. (Une
autre preuve consisterait à tout exprimer en termes d'anneaux locaux, puis à
vérifier les axiomes de l'espace vectoriel au niveau des homomorphismes d'anneaux.)
\end{proof}

\begin{definition}
\label{definition-tangent-space}
Soit $f : X \to S$ un morphisme de schémas. Soit $x \in X$.
L'ensemble des flèches pointillées rendant commutatif le diagramme \ref{equation-tangent-space},
muni de sa structure canonique d'espace vectoriel sur $\kappa(x)$, est
appelé le {\it espace tangent de $X$ sur $S$ en $x$} et nous le notons $T_{X/S,
x}$. Un élément de cet espace est appelé un {\it vecteur tangent} de $X/S$ en $x$.
\end{definition}

\noindent
Puisque les vecteurs tangents en $x \in X$ sont contenus dans la fibre schématique $X_s$ de
$f : X \to S$ au-dessus de $s = f(x)$, nous obtenons une identification canonique
\begin{equation}
\label{equation-tangent-space-fibre}
T_{X/S, x} = T_{X_s/s, x}
\end{equation}
Cette agréable définition faisant intervenir le
foncteur des points admet la description algébrique
suivante, qui suggère de définir {\it l'espace cotangent de
$X$ sur $S$ en $x$} comme l'espace vectoriel sur $\kappa(x)$
$$
T^*_{X/S, x} = \Omega_{X/S, x} \otimes_{\mathcal{O}_{X, x}} \kappa(x)
$$
puisqu'il est canoniquement le dual, sur
$\kappa(x)$, de l'espace tangent de $X$ sur $S$ en $x$.

\begin{lemma}
\label{lemma-tangent-space-cotangent-space}
Soit $f : X \to S$ un morphisme de
schémas. Soit $x \in X$. Il existe un isomorphisme canonique
$$
T_{X/S, x} = \Hom_{\mathcal{O}_{X, x}}(\Omega_{X/S, x}, \kappa(x))
$$
d'espaces vectoriels sur $\kappa(x)$.
\end{lemma}

\begin{proof}
Posons $\kappa = \kappa(x)$. Étant donné
$\theta \in T_{X/S, x}$, nous obtenons un homomorphisme
$$
\theta^*\Omega_{X/S} \to
\Omega_{\Spec(\kappa[\epsilon])/\Spec(\kappa(s))} \to
\Omega_{\Spec(\kappa[\epsilon])/\Spec(\kappa)}
$$
En prenant les sections, nous obtenons un morphisme
$\mathcal{O}_{X, x}$-linéaire $\xi_\theta : \Omega_{X/S, x}
\to \kappa \text{d}\epsilon$, c'est-à-dire un élément du
membre de droite de la formule du lemme. Pour montrer que
$\theta \mapsto \xi_\theta$ est un isomorphisme, nous
pouvons remplacer $S$ par $s$ et $X$ par la fibre schématique
$X_s$. En effet, les deux membres de la formule ne dépendent
que de la fibre schématique ; cela est clair pour $T_{X/S,
x}$ et, pour le membre de droite, voir Morphismes, lemme
\ref{morphisms-lemma-base-change-differentials}. Nous pouvons
aussi remplacer $X$ par le spectre de $\mathcal{O}_{X, x}$,
car cela ne change ni $T_{X/S, x}$ (Schémas, lemme
\ref{schemes-lemma-morphism-from-spec-local-ring}) ni $\Omega_{X/S,
x}$ (Modules, lemme \ref{modules-lemma-stalk-module-differentials}).

\medskip\noindent
Soit $(A, \mathfrak m, \kappa)$ un anneau local sur un
corps $k$. Il faut montrer que tout
homomorphisme $A$-linéaire $\xi : \Omega_{A/k} \to \kappa$ provient
d'un unique morphisme de $k$-algèbres $\varphi : A \to
\kappa[\epsilon]$ coïncidant avec le morphisme canonique $c : A \to
\kappa$ modulo $\epsilon$. Écrivons $\varphi(a) = c(a) + D(a) \epsilon$
: on constate que $a \mapsto D(a)$ est une $k$-dérivation. En
utilisant la propriété universelle de $\Omega_{A/k}$, on voit que
chaque $D$ correspond à un $\xi$ unique et vice versa. Cela termine la preuve.
\end{proof}

\begin{lemma}
\label{lemma-tangent-space-rational-point}
Soit $f : X \to S$ un morphisme de schémas. Soit
$x \in X$ un point et posons $s = f(x) \in S$. Supposons que
$\kappa(x) = \kappa(s)$. Alors il existe des isomorphismes canoniques
$$
\mathfrak m_x/(\mathfrak m_x^2 + \mathfrak m_s\mathcal{O}_{X, x})
=
\Omega_{X/S, x} \otimes_{\mathcal{O}_{X, x}} \kappa(x)
$$
et
$$
T_{X/S, x} =
\Hom_{\kappa(x)}(
\mathfrak m_x/(\mathfrak m_x^2 + \mathfrak m_s\mathcal{O}_{X, x}),
\kappa(x))
$$
Cela fonctionne de manière plus générale si
$\kappa(x)/\kappa(s)$ est une extension algébrique séparable.
\end{lemma}

\begin{proof}
Le deuxième isomorphisme découle du premier par le lemme
\ref{lemma-tangent-space-cotangent-space}. Pour le premier, nous
pouvons remplacer $S$ par $s$ et $X$ par $X_s$, voir Morphismes,
lemme \ref{morphisms-lemma-base-change-differentials}. Nous
pouvons également remplacer $X$ par le spectre de $\mathcal{O}_{X,
x}$, voir Modules, lemme
\ref{modules-lemma-stalk-module-differentials}. Il suffit donc de démontrer le fait
algébrique suivant : soit $(A, \mathfrak m, \kappa)$ un anneau local sur un corps
$k$ tel que $\kappa/k$ soit algébrique séparable. Alors l'application canonique
$$
\mathfrak m/\mathfrak m^2
\longrightarrow
\Omega_{A/k} \otimes \kappa
$$
est un isomorphisme. Observons que $\mathfrak
m/\mathfrak m^2 = H_1(\NL_{\kappa/A})$. D'après
Algèbre, lemme \ref{algebra-lemma-exact-sequence-NL},
il suffit de montrer que $\Omega_{\kappa/k} = 0$ et
$H_1(\NL_{\kappa/k}) = 0$. Comme $\kappa$ est la réunion des
extensions finies séparables de $k$ qu'il contient, il suffit de
prouver cela lorsque $\kappa$ est une extension séparable
finie de $k$ (Algèbre, lemme
\ref{algebra-lemma-colimits-NL}). Dans ce cas, l'homomorphisme
d'anneaux $k \to \kappa$ est étale et donc $\NL_{\kappa/k} = 0$ (à peu
près par définition, voir Algèbre, section \ref{algebra-section-etale}).
\end{proof}

\begin{lemma}
\label{lemma-map-tangent-spaces}
Soit $f : X \to Y$ un morphisme de schémas sur un schéma de
base $S$. Soit $x \in X$ un point. Posons $y = f(x)$. Si
$\kappa(y) = \kappa(x)$, alors $f$ induit une application linéaire naturelle
$$
\text{d}f : T_{X/S, x} \longrightarrow T_{Y/S, y}
$$
qui est duale de l'application linéaire $\Omega_{Y/S, y} \otimes
\kappa(y) \to \Omega_{X/S, x} \otimes \kappa(x)$ via les
identifications du lemme \ref{lemma-tangent-space-cotangent-space}.
\end{lemma}

\begin{proof}
Démonstration omise.
\end{proof}

\begin{lemma}
\label{lemma-tangent-space-product}
Soient $X$, $Y$ des schémas sur une base $S$. Soient $x \in
X$ et $y \in Y$ avec le même point image $s \in S$ tel que $\kappa(s) =
\kappa(x)$ et $\kappa(s) = \kappa(y)$. Il existe un isomorphisme canonique
$$
T_{X \times_S Y/S, (x, y)} = T_{X/S, x} \oplus T_{Y/S, y}
$$
L'application de gauche à droite est induite par les
applications sur les espaces tangents provenant des
projections $X \times_S Y \to X$ et $X \times_S Y \to Y$. La
application de droite à gauche est induite par les applications $1
\times y : X_s \to X_s \times_s Y_s$ et $x \times 1 : Y_s \to X_s
\times_s Y_s$ via l'identification
(\ref{equation-tangent-space-fibre}) des espaces tangents avec les espaces tangents des fibres.
\end{lemma}

\begin{proof}
La décomposition en somme directe découle de Morphismes,
lemme \ref{morphisms-lemma-differential-product} via le lemme
\ref{lemma-tangent-space-rational-point}. La compatibilité avec
les applications provient du lemme \ref{lemma-map-tangent-spaces}.
\end{proof}

\begin{lemma}
\label{lemma-injective-tangent-spaces-unramified}
Supposons que $f : X \to Y$ soit un morphisme de schémas localement de type fini sur
un schéma de base $S$. Supposons que $x \in X$ soit un point. Posons $y = f(x)$ et
supposons que $\kappa(y) = \kappa(x)$. Alors les assertions suivantes sont équivalentes
\begin{enumerate}
\item $\text{d}f : T_{X/S, x} \longrightarrow T_{Y/S, y}$ est injective, et
\item $f$ est non ramifié en $x$.
\end{enumerate}
\end{lemma}

\begin{proof}
Le morphisme $f$ est localement de type fini d'après
Morphismes, lemme
\ref{morphisms-lemma-permanence-finite-type}. L'application
$\text{d}f$ est injective si et seulement si $\Omega_{Y/S, y}
\otimes \kappa(y) \to \Omega_{X/S, x} \otimes \kappa(x)$ est
surjective (lemme \ref{lemma-map-tangent-spaces}). La suite
exacte $f^*\Omega_{Y/S} \to \Omega_{X/S} \to \Omega_{X/Y} \to
0$ (Morphismes, lemme
\ref{morphisms-lemma-triangle-differentials}) montre alors que cela se
produit si et seulement si $\Omega_{X/Y, x} \otimes \kappa(x) = 0$. Le
résultat découle donc de Morphismes, lemme \ref{morphisms-lemma-unramified-at-point}.
\end{proof}







\section{Morphismes génériquement finis}
\label{section-generically-finite}

\noindent
Dans cette section, nous revenons sur la notion de
morphisme génériquement fini de schémas telle qu'elle est étudiée dans
Morphismes, section \ref{morphisms-section-generically-finite}.

\begin{lemma}
\label{lemma-quasi-finite-in-codim-1}
Soit $f : X \to Y$ un morphisme localement de type fini. Soit $y
\in Y$ un point tel que $\mathcal{O}_{Y, y}$ soit noethérien de dimension
$\leq 1$. Supposons en outre que l'une des conditions suivantes soit satisfaite :
\begin{enumerate}
\item pour tout point générique $\eta$ d'une composante irréductible
de $X$, l'extension de corps $\kappa(\eta)/\kappa(f(\eta))$
est finie (ou algébrique),
\item pour tout point générique $\eta$ d'une composante irréductible
de $X$ tel que $f(\eta) \leadsto y$, l'extension de corps
$\kappa(\eta)/\kappa(f(\eta))$ est finie (ou algébrique),
\item $f$ est quasi-fini en tout point générique d'une
composante irréductible de $X$,
\item $Y$ est localement noethérien et $f$
est quasi-fini sur une partie dense de $X$,
\item ajouter d'autres conditions ici.
\end{enumerate}
Alors $f$ est quasi-fini en tout point de $X$ au-dessus de $y$.
\end{lemma}

\begin{proof}
La condition (4) implique que $X$ est localement noethérien
(Morphismes, lemme \ref{morphisms-lemma-finite-type-noetherian}).
L'ensemble des points en lesquels le morphisme est quasi-fini est
ouvert (Morphismes, lemme
\ref{morphisms-lemma-quasi-finite-points-open}). Un ouvert dense d'un
schéma localement noethérien contient tous les points génériques des
composantes irréductibles, donc (4) implique (3). La condition (3)
implique la condition (1) d'après Morphismes, lemme
\ref{morphisms-lemma-residue-field-quasi-finite}. La condition (1) implique la condition
(2). Ainsi, il suffit de prouver le lemme dans le cas où la condition (2) est vérifiée.

\medskip\noindent
Supposons (2). Rappelons que
$\Spec(\mathcal{O}_{Y, y})$ est l'ensemble des points de $Y$
se spécialisant en $y$, voir Schémas, lemme
\ref{schemes-lemma-specialize-points}. En combinant avec Morphismes,
lemme \ref{morphisms-lemma-base-change-quasi-finite}, on voit
que nous pouvons remplacer $Y$ par $\Spec(\mathcal{O}_{Y, y})$.
Ainsi nous pouvons supposer $Y = \Spec(B)$ où $B$ est un anneau
local noethérien de dimension $\leq 1$ et $y$ est le point fermé.

\medskip\noindent
Soit $X = \bigcup X_i$ la décomposition en composantes irréductibles
de $X$, munies de leur structure réduite. Si nous pouvons
montrer que chaque fibre $X_{i, y}$ est un espace discret, alors
$X_y = \bigcup X_{i, y}$ est également discret et nous en concluons
que $X \to Y$ est quasi-fini en tout point de $X_y$ d'après
Morphismes, lemme
\ref{morphisms-lemma-quasi-finite-at-point-characterize}. Ainsi on peut supposer que $X$ est un schéma intègre.

\medskip\noindent
Si $X \to Y$ envoie le point générique $\eta$ de $X$ à
$y$, alors $X$ est le spectre d'une extension finie de
$\kappa(y)$ et le résultat est vrai. Supposons que $X$
envoie $\eta$ à un point correspondant à un idéal premier
minimal $\mathfrak q$ de $B$ différent de $\mathfrak m_B$.
Nous obtenons une factorisation $X \to \Spec(B/\mathfrak
q) \to \Spec(B)$. Supposons que $x \in X$ soit un point
au-dessus de $y$. Par la formule de la dimension
(Morphismes, lemme \ref{morphisms-lemma-dimension-formula}) nous avons
$$
\dim(\mathcal{O}_{X, x}) \leq \dim(B/\mathfrak q) +
\text{trdeg}_{\kappa(\mathfrak q)}(R(X)) - \text{trdeg}_{\kappa(y)} \kappa(x)
$$
Nous savons que $\dim(B/\mathfrak q) = 1$, que le point
générique de $X$ n'est pas égal à $x$ et se spécialise en $x$ et que
$R(X)$ est algébrique sur $\kappa(\mathfrak q)$. Ainsi nous obtenons
$$
1 \leq 1 - \text{trdeg}_{\kappa(y)} \kappa(x)
$$
Ainsi, chaque point $x$ de $X_y$ est fermé dans $X_y$ par
Morphismes, lemme
\ref{morphisms-lemma-algebraic-residue-field-extension-closed-point-fibre}
et donc $X \to Y$ est quasi-fini en tout point $x$ de $X_y$ par
Morphismes, lemme
\ref{morphisms-lemma-quasi-finite-at-point-characterize} (ce qui implique également que $X_y$ est un espace topologique discret).
\end{proof}

\begin{lemma}
\label{lemma-finite-in-codim-1}
Soit $f : X \to Y$ un morphisme propre. Soit $y \in Y$
un point tel que $\mathcal{O}_{Y, y}$ soit noethérien de dimension
$\leq 1$. Supposons en outre que l'une des conditions suivantes soit satisfaite :
\begin{enumerate}
\item pour tout point générique $\eta$ d'une composante irréductible
de $X$, l'extension de corps $\kappa(\eta)/\kappa(f(\eta))$
est finie (ou algébrique),
\item pour tout point générique $\eta$ d'une composante irréductible
de $X$ tel que $f(\eta) \leadsto y$, l'extension de corps
$\kappa(\eta)/\kappa(f(\eta))$ est finie (ou algébrique),
\item $f$ est quasi-fini en tout point générique de $X$,
\item $Y$ est localement noethérien et $f$ est quasi-fini
sur une partie dense de $X$,
\item ajouter d'autres conditions ici.
\end{enumerate}
Alors il existe un voisinage ouvert $V \subset
Y$ de $y$ tel que $f^{-1}(V) \to V$ soit fini.
\end{lemma}

\begin{proof}
D'après le lemme \ref{lemma-quasi-finite-in-codim-1}, le
morphisme $f$ est quasi-fini en tout point de la fibre
$X_y$. Par conséquent, $X_y$ est un espace topologique
discret (Morphismes, lemme
\ref{morphisms-lemma-quasi-finite-at-point-characterize}). Comme $f$
est propre, la fibre $X_y$ est quasi-compacte, c'est-à-dire finie.
Ainsi, nous pouvons appliquer Cohomologie des schémas, lemme
\ref{coherent-lemma-proper-finite-fibre-finite-in-neighbourhood} pour conclure.
\end{proof}

\begin{lemma}
\label{lemma-modification-normal-iso-over-codimension-1}
Soit $X$ un schéma noethérien. Soit $f : Y \to X$ un
morphisme propre birationnel, avec $Y$ réduit. Soit $U \subset
X$ le plus grand ouvert sur lequel $f$ est un isomorphisme. Alors $U$ contient
\begin{enumerate}
\item tout point de codimension $0$ dans $X$,
\item tout $x \in X$ de codimension $1$ dans $X$ tel que
$\mathcal{O}_{X, x}$ soit un anneau de valuation discrète,
\item tout $x \in X$ tel que la fibre de $Y \to X$ au-dessus
de $x$ soit finie et que $\mathcal{O}_{X, x}$ soit normal, et
\item tout $x \in X$ tel que $f$ soit quasi-fini en un
$y \in Y$ au-dessus de $x$ et que $\mathcal{O}_{X, x}$ soit normal.
\end{enumerate}
\end{lemma}

\begin{proof}
La partie (1) découle de Morphismes, lemme
\ref{morphisms-lemma-birational-isomorphism-over-dense-open}. La
partie (2) découle de la partie (3) et du lemme
\ref{lemma-finite-in-codim-1} (et du fait que les morphismes finis ont des fibres finies).

\medskip\noindent
La partie (3) découle de la partie (4) et de Morphismes,
lemme \ref{morphisms-lemma-finite-fibre}, mais nous allons
également donner une preuve directe. Soit $x \in
X$ comme dans (3). D'après Cohomologie des schémas,
lemme
\ref{coherent-lemma-proper-finite-fibre-finite-in-neighbourhood},
on peut supposer que $f$ est fini. Nous pouvons supposer
que $X$ est affine. Cela nous ramène au cas d'un morphisme
birationnel fini de schémas affines noethériens $Y \to X$ et $x
\in X$ tels que $\mathcal{O}_{X, x}$ est un anneau intègre normal.
Comme $\mathcal{O}_{X, x}$ est un anneau intègre et $X$ est
noethérien, nous pouvons remplacer $X$ par un voisinage ouvert affine de
$x$ qui est intègre. Ensuite, puisque $Y \to X$ est birationnel
et $Y$ est réduit, on voit que $Y$ est intègre. Écrivons
$X = \Spec(A)$ et $Y = \Spec(B)$, on voit que $A \subset B$
est une inclusion finie d'anneaux intègres ayant le même corps de
fractions. Si $\mathfrak p \subset A$ est l'idéal premier correspondant
à $x$, alors le fait que $A_\mathfrak p$ soit normal implique que
$A_\mathfrak p \subset B_\mathfrak p$ est une égalité. Puisque $B$
est un $A$-module fini, nous voyons qu'il existe un $a \in A$, $a
\not \in \mathfrak p$ tel que $A_a \to B_a$ est un isomorphisme.

\medskip\noindent
Soient $x \in X$ et $y \in Y$ comme dans (4). Après
avoir remplacé $X$ par un voisinage affine ouvert, nous pouvons
supposer $X = \Spec(A)$ et $A \subset \mathcal{O}_{X, x}$, voir
Propriétés, lemme
\ref{properties-lemma-ring-affine-open-injective-local-ring}.
Ensuite, $A$ est un anneau intègre et donc $X$ est intègre. Puisque $f$
est birationnel et $Y$ est réduit, il s'ensuit que $Y$ est
également intègre. Considérons l'homomorphisme d'anneaux
$\mathcal{O}_{X, x} \to \mathcal{O}_{Y, y}$. Il s'agit d'un
homomorphisme d'anneaux qui est essentiellement de type fini,
l'extension des corps résiduels est finie, et $\dim(\mathcal{O}_{Y,
y}/\mathfrak m_x\mathcal{O}_{Y, y}) = 0$ (pour le voir, remontons
aux définitions des morphismes quasi-finis dans
Morphismes, définition \ref{morphisms-definition-quasi-finite} et
Algèbre, définition \ref{algebra-definition-quasi-finite}). D'après
Algèbre, lemme
\ref{algebra-lemma-essentially-finite-type-fibre-dim-zero},
$\mathcal{O}_{Y, y}$ est la localisation d'une
$\mathcal{O}_{X, x}$-algèbre finie $B$. Bien sûr, nous pouvons remplacer $B$ par
l'image de $B$ dans $\mathcal{O}_{Y, y}$ et supposer que $B$ est un
anneau intègre ayant le même corps des fractions que $\mathcal{O}_{Y, y}$. Comme
Les anneaux de l'inclusion $\mathcal{O}_{X, x} \subset B$ ont le même corps de fractions puisque $f$ est birationnel.
Comme $\mathcal{O}_{X, x}$ est normal, nous concluons que
$\mathcal{O}_{X, x} = B$ (car fini implique intégral), en particulier, nous
voyons que $\mathcal{O}_{X, x} = \mathcal{O}_{Y, y}$. D'après Morphismes,
lemme \ref{morphisms-lemma-morphism-defined-local-ring}, après avoir rétréci
$X$, nous pouvons supposer qu'il existe une section $X \to Y$ de $f$
envoyant $x$ vers $y$ et induisant l'isomorphisme donné sur les anneaux
locaux. Comme $X \to Y$ est une immersion fermée (d'après Schémas, lemme
\ref{schemes-lemma-section-immersion}) et envoie nécessairement le point
générique de $X$ sur le point générique de $Y$, il s'ensuit que l'image
de $X \to Y$ est $Y$. Ainsi $Y = X$, ce qu'il fallait démontrer.
\end{proof}






\section{Variantes de la normalisation de Noether}
\label{section-noether-normalization}

\noindent
La normalisation de Noether affirme que si $k$ est un corps
et $A$ une $k$-algèbre de type fini de dimension $d$, alors
il existe un homomorphisme injectif et fini de $k$-algèbres
$k[x_1,
\ldots, x_d] \to A$. Voir Algèbre, lemme
\ref{algebra-lemma-Noether-normalization}.
Géométriquement, cela signifie qu'il existe un morphisme
surjectif fini $\Spec(A) \to \mathbf{A}^d_k$ au-dessus de $\Spec(k)$.

\begin{lemma}
\label{lemma-noether-normalization}
Soit $f : X \to S$ un morphisme de schémas.
Soit $x \in X$, d'image $s \in S$. Soit $V \subset S$
un voisinage ouvert affine de $s$. Si $f$ est localement de type
fini et si $\dim_x(X_s) = d$, alors il existe un ouvert affine $U
\subset X$, avec $x \in U$ et $f(U) \subset V$, et une factorisation
$$
U \xrightarrow{\pi} \mathbf{A}^d_V \to V
$$
de $f|_U : U \to V$ telle que $\pi$ soit quasi-fini.
\end{lemma}

\begin{proof}
Cela résulte d'Algèbre, lemme
\ref{algebra-lemma-quasi-finite-over-polynomial-algebra}.
\end{proof}

\begin{lemma}
\label{lemma-noether-normalization-affine}
Soit $f : X \to S$ un morphisme de type fini de schémas
affines. Soit $s \in S$. Si $\dim(X_s) = d$, alors il existe une factorisation
$$
X \xrightarrow{\pi} \mathbf{A}^d_S \to S
$$
de $f$ de telle sorte que le morphisme $\pi_s : X_s \to
\mathbf{A}^d_{\kappa(s)}$ des fibres au-dessus de $s$ soit fini.
\end{lemma}

\begin{proof}
Écrivons $S = \Spec(A)$ et $X = \Spec(B)$, et notons $A \to B$
l'homomorphisme d'anneaux correspondant à $f$. Soit
$\mathfrak p \subset A$ l'idéal premier correspondant à $s$.
Nous pouvons choisir une surjection $A[x_1, \ldots, x_r] \to
B$. Par Algèbre, lemme
\ref{algebra-lemma-Noether-normalization}, il existe des éléments
$y_1, \ldots, y_d \in A$ appartenant à la $\mathbf{Z}$-sous-algèbre de $A$
engendrée par $x_1, \ldots, x_r$ tels que l'homomorphisme de $A$-algèbres
$A[t_1, \ldots, t_d] \to B$ envoyant $t_i$ sur $y_i$ induise un
homomorphisme fini de $\kappa(\mathfrak p)$-algèbres $\kappa(\mathfrak
p)[t_1, \ldots, t_d] \to B \otimes_A \kappa(\mathfrak p)$. Ceci démontre le lemme.
\end{proof}

\begin{lemma}
\label{lemma-geometric-structure-unramified}
Soit $f : X \to S$ un morphisme de schémas.
Soit $x \in X$. Soit $V = \Spec(A)$ un voisinage
ouvert affine de $f(x)$ dans $S$. Si $f$ est non ramifié en
$x$, alors il existe un ouvert affine $U \subset X$ avec $x \in
U$ et $f(U) \subset V$ tel que nous ayons un diagramme commutatif
$$
\xymatrix{
X \ar[d] & U \ar[l] \ar[rd] \ar[r]^-j &
\Spec(A[t]_{g'}/(g)) \ar[d] \ar[r] &
\Spec(A[t]) = \mathbf{A}^1_V \ar[ld] \\
Y & & V \ar[ll]
}
$$
où $j$ est une immersion, $g \in A[t]$ est un polynôme unitaire, et $g'$ est
la dérivée de $g$ par rapport à $t$. Si $f$ est étale en $x$, alors nous
pouvons choisir le diagramme de manière à ce que $j$ soit une immersion ouverte.
\end{lemma}

\begin{proof}
Le cas non ramifié est la traduction géométrique d'Algèbre,
proposition
\ref{algebra-proposition-unramified-locally-standard}. Dans le cas
étale, on traduit de même Algèbre, proposition
\ref{algebra-proposition-etale-locally-standard}; de façon
équivalente, le résultat découle de Morphismes, lemme
\ref{morphisms-lemma-etale-locally-standard-etale} bien que les énoncés diffèrent légèrement.
\end{proof}

\begin{lemma}
\label{lemma-unramfied-over-affine}
Soit $f : X \to S$ un morphisme de type fini
de schémas affines. Soit $x \in X$, d'image $s \in S$. Posons
$$
r =
\dim_{\kappa(x)} \Omega_{X/S, x} \otimes_{\mathcal{O}_{X, x}} \kappa(x) =
\dim_{\kappa(x)} \Omega_{X_s/s, x} \otimes_{\mathcal{O}_{X_s, x}} \kappa(x) =
\dim_{\kappa(x)} T_{X/S, x}
$$
Alors il existe une factorisation
$$
X \xrightarrow{\pi} \mathbf{A}^r_S \to S
$$
de $f$ de telle sorte que $\pi$ soit non ramifié en $x$.
\end{lemma}

\begin{proof}
D'après Morphismes, lemme
\ref{morphisms-lemma-finite-type-differentials} la
première dimension est finie. La première égalité résulte
du fait que la restriction de $\Omega_{X/S}$ à la fibre
est le module des différentielles, d'après Morphismes, lemme
\ref{morphisms-lemma-base-change-differentials}. La
dernière égalité découle du lemme
\ref{lemma-tangent-space-cotangent-space}. Ainsi, on voit que l'énoncé a du sens.

\medskip\noindent
Pour démontrer le lemme, écrivons $S = \Spec(A)$ et $X = \Spec(B)$ et
notons $A \to B$ l'homomorphisme d'anneaux correspondant à $f$.
Soit $\mathfrak q \subset B$ l'idéal premier
correspondant à $x$. Choisissons une surjection de $A$-algèbres
$A[x_1, \ldots, x_t] \to B$. Puisque $\Omega_{B/A}$ est engendré
par $\text{d}x_1, \ldots, \text{d}x_t$, on voit que leurs images
dans $\Omega_{X/S, x} \otimes_{\mathcal{O}_{X, x}} \kappa(x)$
engendrent cet espace comme espace vectoriel sur $\kappa(x)$. Après
renumérotation, on peut supposer que $\text{d}x_1, \ldots,
\text{d}x_r$ s'envoient sur une base de $\Omega_{X/S, x}
\otimes_{\mathcal{O}_{X, x}} \kappa(x)$. Nous affirmons que $P = A[x_1,
\ldots, x_r] \to B$ est non ramifié en $\mathfrak q$. Pour le voir, il
suffit de montrer que $\Omega_{B/P, \mathfrak q} = 0$ (Algèbre, lemme
\ref{algebra-lemma-unramified}). Notons que $\Omega_{B/P}$ est le
quotient de $\Omega_{B/A}$ par le sous-module généré par $\text{d}x_1,
\ldots, \text{d}x_r$. Par conséquent, $\Omega_{B/P, \mathfrak q}
\otimes_{B_\mathfrak q} \kappa(\mathfrak q) = 0$ par notre choix de $x_1,
\ldots, x_r$. Par le lemme de Nakayama, plus précisément Algèbre, lemme
\ref{algebra-lemma-NAK} partie (2) qui s'applique car $\Omega_{B/P}$ est fini
(voir référence ci-dessus), nous concluons que $\Omega_{B/P, \mathfrak q} = 0$.
\end{proof}

\begin{lemma}
\label{lemma-immersion-into-affine}
Soit $f : X \to S$ un morphisme de
schémas. Soit $x \in X$, d'image $s
\in S$. Soit $V \subset S$ un voisinage ouvert affine
ouvert de $s$. Si $f$ est localement de type fini et
$$
r =
\dim_{\kappa(x)} \Omega_{X/S, x} \otimes_{\mathcal{O}_{X, x}} \kappa(x) =
\dim_{\kappa(x)} \Omega_{X_s/s, x} \otimes_{\mathcal{O}_{X_s, x}} \kappa(x) =
\dim_{\kappa(x)} T_{X/S, x}
$$
alors il existe
\begin{enumerate}
\item un ouvert affine $U \subset X$ tel que $x \in U$ et $f(U) \subset V$, et une
factorisation
$$
U \xrightarrow{j} \mathbf{A}^{r + 1}_V \to V
$$
de $f|_U$ tel que $j$ soit une immersion, ou
\item un ouvert affine $U \subset X$ tel que $x \in U$ et $f(U) \subset V$, et une
factorisation
$$
U \xrightarrow{j} D \to V
$$
de $f|_U$ telle que $j$ soit une immersion fermée et $D \to V$
soit lisse de dimension relative $r$.
\end{enumerate}
\end{lemma}

\begin{proof}
Choisissons un ouvert affine quelconque $U \subset X$ avec $x \in U$ et
$f(U) \subset V$. Appliquons le lemme
\ref{lemma-unramfied-over-affine} à $U \to V$ pour obtenir $U \to
\mathbf{A}^r_V \to V$ comme dans l'énoncé de ce lemme. D'après le lemme
\ref{lemma-geometric-structure-unramified}, nous obtenons une factorisation
$$
U \xrightarrow{j} D \xrightarrow{j'} \mathbf{A}^{r + 1}_V
\xrightarrow{p} \mathbf{A}^r_V \to V
$$
où $j$ et $j'$ sont des immersions, $p$ est la
projection, et $p \circ j'$ est étale standard.
Ainsi, les assertions (1) et (2) sont vérifiées.
\end{proof}





\section{Dimension des fibres}
\label{section-dimension-fibres}

\noindent
Nous avons déjà vu que la dimension des fibres des morphismes de type fini peut
augmenter. Dans cette section, nous observons qu'en codimension $1$,
ce phénomène ne se produit pas. Plus généralement, nous étudions l'amplitude
possible de ce saut de dimension. Voici une liste de résultats connexes :
\begin{enumerate}
\item Pour un morphisme de type fini $X \to S$, la partie formée des
$x \in X$ tels que $\dim_x(X_{f(x)}) \leq d$ est ouverte; voir
Algèbre, lemme \ref{algebra-lemma-dimension-fibres-bounded-open-upstairs}
et
Morphismes, lemme \ref{morphisms-lemma-openness-bounded-dimension-fibres}.
\item On dispose de la formule de dimension; voir
Algèbre, lemme \ref{algebra-lemma-dimension-formula} et
Morphismes, lemme \ref{morphisms-lemma-dimension-formula}.
\item Les fibres d'un schéma intègre de type fini dominant
le spectre d'un anneau de valuation sont de dimension constante; voir Algèbre, lemme
\ref{algebra-lemma-finite-type-domain-over-valuation-ring-dim-fibres}.
\item Si $X \to S$ est de type fini et quasi-fini en tout
point générique de $X$, alors $X \to S$ est quasi-fini en codimension $1$; voir
Algèbre, lemme \ref{algebra-lemma-finite-in-codim-1} et
lemme \ref{lemma-quasi-finite-in-codim-1}.
\end{enumerate}
Le dernier résultat mentionné ci-dessus se généralise comme suit.

\begin{lemma}
\label{lemma-dimension-fibre-in-codim-1}
Soit $f : X \to Y$ un morphisme localement de type fini. Soit $x \in
X$ un point d'image $y \in Y$ tel que $\mathcal{O}_{Y, y}$ soit
noethérien de dimension $\leq 1$. Soit $d \geq 0$ un entier tel
que pour tout point générique $\eta$ d'une composante irréductible de $X$
contenant $x$, nous ayons $\dim_\eta(X_{f(\eta)}) = d$. Alors $\dim_x(X_y) = d$.
\end{lemma}

\begin{proof}
Rappelons que $\Spec(\mathcal{O}_{Y, y})$ est l'ensemble
des points de $Y$ se spécialisant en $y$, voir Schémas,
lemme \ref{schemes-lemma-specialize-points}. Ainsi, nous
pouvons remplacer $Y$ par $\Spec(\mathcal{O}_{Y, y})$ et
supposer $Y = \Spec(B)$ où $B$ est un anneau local
noethérien de dimension $\leq 1$ et $y$ est le point fermé. Nous
pouvons également remplacer $X$ par un voisinage affine de $x$.

\medskip\noindent
Soit $X = \bigcup X_i$ la décomposition de $X$ en composantes irréductibles,
munies de leur structure réduite. S'il est possible de montrer que
chaque fibre $X_{i, y}$ a dimension $d$, alors $X_y = \bigcup X_{i, y}$ a
également dimension $d$. Ainsi, on peut supposer que $X$ est un schéma intègre.

\medskip\noindent
Si $X \to Y$ envoie le point générique $\eta$ de $X$ sur
$y$, alors $X$ est un schéma sur $\kappa(y)$ et le
résultat est vrai par hypothèse. Supposons que $X$ envoie
$\eta$ sur un point $\xi \in Y$ correspondant à un idéal
premier minimal $\mathfrak q$ de $B$ différent de $\mathfrak
m_B$. Nous obtenons une factorisation $X \to
\Spec(B/\mathfrak q) \to \Spec(B)$. Par la formule de dimension
(Morphismes, lemme \ref{morphisms-lemma-dimension-formula}) nous avons
$$
\dim(\mathcal{O}_{X, x}) + \text{trdeg}_{\kappa(y)} \kappa(x) \leq
\dim(B/\mathfrak q) + \text{trdeg}_{\kappa(\mathfrak q)}(R(X))
$$
Nous avons $\dim(B/\mathfrak q) = 1$. Nous avons
$\text{trdeg}_{\kappa(\mathfrak q)}(R(X)) = d$ par
hypothèse, puisque $\dim_\eta(X_\xi) = d$; voir
Morphismes, lemme
\ref{morphisms-lemma-dimension-fibre-at-a-point}. Puisque
$\mathcal{O}_{X, x} \to \mathcal{O}_{X_s, x}$ a un noyau
(comme $\eta \mapsto \xi \not = y$) et puisque
$\mathcal{O}_{X, x}$ est un anneau intègre noethérien, on voit que
$\dim(\mathcal{O}_{X, x}) > \dim(\mathcal{O}_{X_y, x})$. Nous concluons que
$$
\dim_x(X_s) =
\dim(\mathcal{O}_{X_s, x}) + \text{trdeg}_{\kappa(y)} \kappa(x) \leq d
$$
(Morphismes, lemme
\ref{morphisms-lemma-dimension-fibre-at-a-point}). D'autre part,
nous avons $\dim_x(X_s) \geq \dim_\eta(X_{f(\eta)}) = d$ d'après
Morphismes, lemme \ref{morphisms-lemma-openness-bounded-dimension-fibres}.
\end{proof}

\begin{lemma}
\label{lemma-dominate-valuation-ring-dimension-fibres}
Soit $f : X \to \Spec(R)$ un morphisme d'un schéma
irréductible vers le spectre d'un anneau de valuation. Si $f$ est
localement de type fini et surjectif, alors la fibre spéciale est
équidimensionnelle de dimension égale à la dimension de la fibre générique.
\end{lemma}

\begin{proof}
Nous pouvons remplacer $X$ par sa réduction car cela ne change
pas la dimension de $X$ ni celle de la fibre spéciale. Alors
$X$ est intègre et le lemme résulte d'Algèbre, lemme
\ref{algebra-lemma-finite-type-domain-over-valuation-ring-dim-fibres}.
\end{proof}

\noindent
Le lemme suivant généralise le lemme \ref{lemma-dimension-fibre-in-codim-1}.

\begin{lemma}
\label{lemma-dimension-fibre-in-higher-codimension}
Soit $f : X \to Y$ un morphisme localement de type fini. Soit
$x \in X$ un point d'image $y \in Y$ tel que $\mathcal{O}_{Y,
y}$ soit noethérien. Soit $d \geq 0$ un entier tel que
pour tout point générique $\eta$ d'une composante irréductible de $X$
qui contient $x$, nous ayons $f(\eta) \not = y$ et
$\dim_\eta(X_{f(\eta)}) = d$. Alors $\dim_x(X_y) \leq d + \dim(\mathcal{O}_{Y, y}) - 1$.
\end{lemma}

\begin{proof}
Comme dans la preuve du lemme
\ref{lemma-dimension-fibre-in-codim-1}, nous réduisons au cas $X =
\Spec(A)$ avec $A$ un anneau intègre et $Y = \Spec(B)$ où $B$ est un anneau
local noethérien dont l'idéal maximal correspond à $y$. Après avoir
remplacé $B$ par $B/\Ker(B \to A)$, on peut supposer que $B$ est un anneau
intègre et que $B \subset A$. Ensuite, nous utilisons la formule de
dimension (Morphismes, lemme \ref{morphisms-lemma-dimension-formula}) pour obtenir
$$
\dim(\mathcal{O}_{X, x}) + \text{trdeg}_{\kappa(y)} \kappa(x) \leq
\dim(B) + \text{trdeg}_B(A)
$$
Nous avons $\text{trdeg}_B(A) = d$ par hypothèse,
puisque $\dim_\eta(X_\xi) = d$; voir
Morphismes, lemme
\ref{morphisms-lemma-dimension-fibre-at-a-point}. Puisque
$\mathcal{O}_{X, x} \to \mathcal{O}_{X_y, x}$ a un noyau
(comme $f(\eta) \not = y$) et puisque $\mathcal{O}_{X, x}$
est un anneau intègre noethérien, on voit que
$\dim(\mathcal{O}_{X, x}) > \dim(\mathcal{O}_{X_y, x})$. Nous concluons que
$$
\dim_x(X_y) =
\dim(\mathcal{O}_{X_y, x}) + \text{trdeg}_{\kappa(y)} \kappa(x)
< \dim(B) + d
$$
(l'égalité vient de Morphismes, lemme
\ref{morphisms-lemma-dimension-fibre-at-a-point}), ce qui donne le résultat.
\end{proof}




\section{Schémas algébriques}
\label{section-algebraic-schemes}

\noindent
La définition suivante est tirée de
\cite[I, définition 6.4.1]{EGA}.

\begin{definition}
\label{definition-algebraic-scheme}
Soit $k$ un corps. Un {\it $k$-schéma
algébrique} est un schéma $X$ sur $k$ tel que le morphisme de
structure $X \to \Spec(k)$ soit de type fini. Un {\it $k$-schéma
localement algébrique} est un schéma $X$ sur $k$ tel que le
morphisme de structure $X \to \Spec(k)$ soit localement de type fini.
\end{definition}

\noindent
Notons que tout $k$-schéma (localement) algébrique est (localement)
noethérien, voir Morphismes, lemme
\ref{morphisms-lemma-finite-type-noetherian}. La catégorie des $k$-schémas
algébriques possède tous les produits et produits fibrés (contrairement à la catégorie
des variétés sur $k$). De même pour la catégorie des $k$-schémas localement algébriques.

\begin{lemma}
\label{lemma-algebraic-scheme-dim-0}
Soit $k$ un corps. Soit $X$ un $k$-schéma
localement algébrique de dimension $0$. Alors $X$ est une réunion
disjointe de spectres de $k$-algèbres locales artiniennes $A$ avec
$\dim_k(A) < \infty$. Si $X$ est un $k$-schéma algébrique de dimension
$0$, alors en plus $X$ est affine et le morphisme $X \to \Spec(k)$ est fini.
\end{lemma}

\begin{proof}
Soit $X$ un $k$-schéma localement algébrique
de dimension $0$. Soit $U = \Spec(A) \subset X$
un sous-schéma ouvert affine. Comme $\dim(X) = 0$, on
a $\dim(A) = 0$. Par normalisation de Noether,
voir Algèbre, lemme
\ref{algebra-lemma-Noether-normalization}, nous voyons qu'il
existe une injection finie $k \to A$, c'est-à-dire $\dim_k(A) <
\infty$. Par conséquent, $A$ est artinien, voir Algèbre, lemme
\ref{algebra-lemma-finite-dimensional-algebra}. Cela implique que
$A = A_1 \times \ldots \times A_r$ est un produit d'un nombre
fini d'anneaux locaux artiniens, voir Algèbre, lemme
\ref{algebra-lemma-artinian-finite-length}. Bien sûr, $\dim_k(A_i) <
\infty$ pour chaque $i$ puisque la somme de ces dimensions égale $\dim_k(A)$.

\medskip\noindent
Les arguments ci-dessus montrent que $X$ possède un recouvrement ouvert
dont les membres sont des espaces topologiques discrets finis. Par
conséquent, $X$ est un espace topologique discret. Il s'ensuit que $X$ est
isomorphe à la réunion disjointe de ses composantes connexes, chacune étant un
espace à un point. Puisqu'un tel schéma est affine, nous concluons (d'après
les résultats du paragraphe précédent) que chacune de ces composantes est le
spectre d'une $k$-algèbre locale artinienne $A$ telle que $\dim_k(A) < \infty$.

\medskip\noindent
Enfin, si $X$ est un $k$-schéma algébrique de
dimension $0$, alors $X$ est quasi-compact donc
est une réunion disjointe finie $X = \Spec(A_1)
\amalg \ldots \amalg \Spec(A_r)$ donc affine (voir
Schémas, lemme
\ref{schemes-lemma-disjoint-union-affines}); nous avons vu la
finitude de $X \to \Spec(k)$ dans le premier paragraphe de la démonstration.
\end{proof}

\noindent
Le lemme suivant rassemble quelques énoncés sur la théorie
de la dimension pour les schémas localement algébriques.

\begin{lemma}
\label{lemma-dimension-locally-algebraic}
Soit $k$ un corps. Soit $X$ un $k$-schéma localement algébrique.
\begin{enumerate}
\item
\label{item-catenary}
L'espace topologique sous-jacent à $X$ est caténaire
(Topologie, définition \ref{topology-definition-catenary}).
\item
\label{item-dimension-at-closed-point}
Pour tout point fermé $x \in X$, on a $\dim_x(X) = \dim(\mathcal{O}_{X, x})$.
\item
\label{item-dimension-irreducible}
Si $X$ est irréductible, on a $\dim(X) = \dim(U)$ pour tout
ouvert non vide $U \subset X$, et $\dim(X) = \dim_x(X)$
pour tout $x \in X$.
\item
\label{item-irreducible-maximal-chains}
Si $X$ est irréductible, toute chaîne de parties fermées irréductibles peut être
prolongée en une chaîne maximale, et toutes les chaînes maximales de parties
fermées irréductibles sont de longueur $\dim(X)$.
\item
\label{item-dimension-irreducibles-passing-through}
Pour $x \in X$, on a
$\dim_x(X) = \max \dim(Z) = \min \dim(\mathcal{O}_{X, x'})$,
où le maximum porte sur les composantes irréductibles
$Z \subset X$ contenant $x$, et le minimum sur les
spécialisations $x \leadsto x'$ telles que $x'$ soit fermé dans $X$.
\item
\label{item-dimension-irreducible-trdeg}
Si $X$ est irréductible de point générique $x$, alors
$\dim(X) = \text{trdeg}_k(\kappa(x))$.
\item
\label{item-immediate-specialization}
Si $x \leadsto x'$ est une spécialisation immédiate
de points de $X$, alors
$\text{trdeg}_k(\kappa(x)) = \text{trdeg}_k(\kappa(x')) + 1$.
\item
\label{item-dimension-sup-trdeg}
La dimension de $X$ est le suprémum des nombres
$\text{trdeg}_k(\kappa(x))$, où $x$ parcourt les
points génériques des composantes irréductibles de $X$.
\item
\label{item-specialization}
Si $x \leadsto x'$ est une spécialisation non triviale de points de $X$, alors
\begin{enumerate}
\item $\dim_x(X) \leq \dim_{x'}(X)$,
\item $\dim(\mathcal{O}_{X, x}) < \dim(\mathcal{O}_{X, x'})$,
\item $\text{trdeg}_k(\kappa(x)) > \text{trdeg}_k(\kappa(x'))$, et
\item toute chaîne maximale de spécialisations non triviales
$x = x_0 \leadsto x_1 \leadsto \ldots \leadsto x_n = x'$ est de
longueur $n = \text{trdeg}_k(\kappa(x)) - \text{trdeg}_k(\kappa(x'))$.
\end{enumerate}
\item
\label{item-dimension-formula}
Pour $x \in X$, on a
$\dim_x(X) = \text{trdeg}_k(\kappa(x)) + \dim(\mathcal{O}_{X, x})$.
\item
\label{item-immediate-specialization-local-ring}
Si $x \leadsto x'$ est une spécialisation immédiate
de points de $X$ et si $X$ est irréductible ou équidimensionnel, alors
$\dim(\mathcal{O}_{X, x'}) = \dim(\mathcal{O}_{X, x}) + 1$.
\end{enumerate}
\end{lemma}

\begin{proof}
Au lieu de nous appuyer sur les résultats plus généraux
démontrés précédemment, nous allons ramener ces énoncés aux
énoncés correspondants sur les $k$-algèbres de type fini et
citer des résultats du chapitre d'algèbre commutative.

\medskip\noindent
Preuve de (\ref{item-catenary}). Cette propriété est locale sur $X$ d'après
Topologie, lemme \ref{topology-lemma-catenary}. Ainsi, nous
pouvons supposer $X = \Spec(A)$ où $A$ est une $k$-algèbre de
type fini. Nous devons montrer que $A$ est caténaire (Algèbre,
lemme \ref{algebra-lemma-catenary}). Nous pouvons réduire à
$k[x_1, \ldots, x_n]$ en utilisant Algèbre, lemme
\ref{algebra-lemma-quotient-catenary} et ensuite appliquer
Algèbre, lemme
\ref{algebra-lemma-dimension-height-polynomial-ring}. On peut aussi conclure
du fait que $k$ est de Cohen--Macaulay (trivialement) et que les anneaux de Cohen--Macaulay
sont universellement caténaires (Algèbre, lemme \ref{algebra-lemma-CM-ring-catenary}).

\medskip\noindent
Preuve de (\ref{item-dimension-at-closed-point}). Choisissons un
voisinage affine $U = \Spec(A)$ de $x$. Alors $\dim_x(X) =
\dim_x(U)$. Par conséquent, nous nous réduisons au cas affine, qui est
Algèbre, lemme \ref{algebra-lemma-dimension-closed-point-finite-type-field}.

\medskip\noindent
Preuve de (\ref{item-dimension-irreducible}). Il suffit de
montrer que deux ouverts affines non vides $U, U' \subset X$
quelconques ont la même dimension (toute chaîne finie de
sous-ensembles irréductibles rencontre un ouvert affine).
Choisissons un point fermé $x$ de $X$ avec $x \in U \cap U'$. Cela
est possible car $X$ est irréductible, donc $U \cap U'$ est non
vide, donc il existe un tel point fermé car $X$ est Jacobson
d'après le lemme \ref{lemma-locally-finite-type-Jacobson}. Ensuite
$\dim(U) = \dim(\mathcal{O}_{X, x}) = \dim(U')$ d'après Algèbre,
lemme \ref{algebra-lemma-dimension-spell-it-out} (strictement
parlant, il faut remplacer $X$ par sa réduction avant d'appliquer le lemme).

\medskip\noindent
Preuve de (\ref{item-irreducible-maximal-chains}). Étant
donné une chaîne de sous-ensembles fermés irréductibles, nous
pouvons trouver un ouvert affine $U \subset X$ qui rencontre
le plus petit. Ainsi, l'énoncé résulte d'Algèbre, lemme
\ref{algebra-lemma-dimension-spell-it-out} et $\dim(U) =
\dim(X)$, comme on l'a vu dans (\ref{item-dimension-irreducible}).

\medskip\noindent
Preuve de
(\ref{item-dimension-irreducibles-passing-through}).
Choisissons un voisinage affine $U = \Spec(A)$ de $x$. Alors
$\dim_x(X) = \dim_x(U)$. La règle $Z \mapsto Z \cap U$ est une
bijection entre les composantes irréductibles de $X$ passant par
$x$ et les composantes irréductibles de $U$ passant par $x$. De
plus, $\dim(Z \cap U) = \dim(Z)$ pour un tel $Z$ d'après
(\ref{item-dimension-irreducible}). Ainsi, l'énoncé découle d'Algèbre,
lemme \ref{algebra-lemma-dimension-at-a-point-finite-type-over-field}.

\medskip\noindent
Preuve de (\ref{item-dimension-irreducible-trdeg}).
D'après (\ref{item-dimension-irreducible}), cela se
réduit au cas où $X = \Spec(A)$ est affine. Dans ce
cas, cela résulte d'Algèbre, lemme
\ref{algebra-lemma-dimension-prime-polynomial-ring} appliqué à $A_{red}$.

\medskip\noindent
Preuve de (\ref{item-immediate-specialization}).
Soit $Z = \overline{\{x\}} \supset Z' =
\overline{\{x'\}}$. Il résulte de
(\ref{item-irreducible-maximal-chains}) que $Z \supset Z'$
est le début d'une chaîne maximale de parties fermées
irréductibles dans $Z$ et par conséquent $\dim(Z) = \dim(Z') +
1$. On conclut par (\ref{item-dimension-irreducible-trdeg}).

\medskip\noindent
Preuve de (\ref{item-dimension-sup-trdeg}). Un argument
topologique simple montre que $\dim(X) = \sup \dim(Z)$ où le
suprémum porte sur les composantes irréductibles de $X$ (indication :
utiliser Topologie, lemme \ref{topology-lemma-irreducible}).
Le résultat découle donc de (\ref{item-dimension-irreducible-trdeg}).

\medskip\noindent
Preuve de (\ref{item-specialization}). La partie (a) découle du fait
que tout ouvert $U \subset X$ contenant $x'$ contient également $x$. La
partie (b) découle du fait que $\mathcal{O}_{X, x}$ est une localisation
de $\mathcal{O}_{X, x'}$; toute chaîne d'idéaux premiers de
$\mathcal{O}_{X, x}$ correspond donc à une chaîne d'idéaux premiers de $\mathcal{O}_{X,
x'}$, qui peut être prolongée en ajoutant $\mathfrak m_{x'}$ à la fin. Les
parties (c) et (d) découlent formellement de (\ref{item-immediate-specialization}).

\medskip\noindent
Preuve de (\ref{item-dimension-formula}). Choisissons un
voisinage affine $U = \Spec(A)$ de $x$. Alors $\dim_x(X) =
\dim_x(U)$. Par conséquent, nous nous réduisons au cas affine, qui est
Algèbre, lemme \ref{algebra-lemma-dimension-at-a-point-finite-type-field}.

\medskip\noindent
Preuve de
(\ref{item-immediate-specialization-local-ring}). Si $X$ est
équidimensionnel (Topologie, définition
\ref{topology-definition-equidimensional}) alors $\dim(X)$ est égal à la
dimension de chaque composante irréductible de $X$, d'où $\dim_x(X) =
\dim(X) = \dim_{x'}(X)$ d'après
(\ref{item-dimension-irreducibles-passing-through}). Cela découle donc de (\ref{item-immediate-specialization}).
\end{proof}

\begin{lemma}
\label{lemma-dimension-fibres-locally-algebraic}
Soit $k$ un corps. Soit $f : X \to Y$ un morphisme de
$k$-schémas localement algébriques.
\begin{enumerate}
\item Pour $y \in Y$, la fibre $X_y$ est un schéma localement
algébrique sur $\kappa(y)$; tous les résultats du
lemme \ref{lemma-dimension-locally-algebraic} s'appliquent donc.
\item Supposons $X$ irréductible. Posons $Z = \overline{f(X)}$ et
$d = \dim(X) - \dim(Z)$. Alors
\begin{enumerate}
\item $\dim_x(X_{f(x)}) \geq d$ pour tout $x \in X$,
\item la partie des $x \in X$ tels que $\dim_x(X_{f(x)}) = d$ est ouverte dense,
\item si $\dim(\mathcal{O}_{Z, f(x)}) \geq 1$, alors
$\dim_x(X_{f(x)}) \leq d + \dim(\mathcal{O}_{Z, f(x)}) - 1$,
\item si $\dim(\mathcal{O}_{Z, f(x)}) = 1$, alors $\dim_x(X_{f(x)}) = d$,
\end{enumerate}
\item Pour $x \in X$ avec $y = f(x)$ nous
avons $\dim_x(X_y) \geq \dim_x(X) - \dim_y(Y)$.
\end{enumerate}
\end{lemma}

\begin{proof}
Le morphisme $f$ est localement de type fini d'après
Morphismes, lemme
\ref{morphisms-lemma-permanence-finite-type}. Par conséquent,
le changement de base $X_y \to \Spec(\kappa(y))$ est localement
de type fini. Cela prouve (1). Dans le reste de la
démonstration, nous utiliserons librement les résultats du lemme
\ref{lemma-dimension-locally-algebraic} pour $X$, $Y$, et les fibres de $f$.

\medskip\noindent
Preuve de (2). Soit $\eta \in X$ le point générique et
posons $\xi = f(\eta)$. Alors $Z = \overline{\{\xi\}}$. Par conséquent
$$
d = \dim(X) - \dim(Z) =
\text{trdeg}_k \kappa(\eta) - \text{trdeg}_k \kappa(\xi) =
\text{trdeg}_{\kappa(\xi)} \kappa(\eta) = \dim_\eta(X_\xi)
$$
Ainsi, les parties (2)(a) et (2)(b) découlent des
Morphismes, lemme
\ref{morphisms-lemma-openness-bounded-dimension-fibres}. Les
parties (2)(c) et (2)(d) découlent des lemmes
\ref{lemma-dimension-fibre-in-higher-codimension} et \ref{lemma-dimension-fibre-in-codim-1}.

\medskip\noindent
Preuve de (3). Soit $x \in X$. Soit $X' \subset X$ une
composante irréductible de $X$ passant par $x$ de dimension
$\dim_x(X)$. Alors (2) implique que $\dim_x(X_y) \geq \dim(X') -
\dim(Z')$ où $Z' \subset Y$ est l'adhérence de l'image de $X'$. Cela prouve (3).
\end{proof}

\begin{lemma}
\label{lemma-dimension-product-locally-algebraic}
\begin{slogan}
La dimension du produit est la somme des dimensions.
\end{slogan}
Soit $k$ un corps. Soient $X$, $Y$ des $k$-schémas localement algébriques.
\begin{enumerate}
\item Pour $z \in X \times Y$ au-dessus de $(x, y)$,
nous avons $\dim_z(X \times Y) = \dim_x(X) + \dim_y(Y)$.
\item Nous avons $\dim(X \times Y) = \dim(X) + \dim(Y)$.
\end{enumerate}
\end{lemma}

\begin{proof}
Preuve de (1). Considérons la factorisation
$$
X \times Y \longrightarrow Y \longrightarrow \Spec(k)
$$
du morphisme de structure. Le premier morphisme $p : X \times Y \to
Y$ est plat par changement de base du morphisme plat $X \to
\Spec(k)$ d'après Morphismes, lemme
\ref{morphisms-lemma-base-change-flat}. De plus, nous avons
$\dim_z(p^{-1}(y)) = \dim_x(X)$ d'après Morphismes, lemme
\ref{morphisms-lemma-dimension-fibre-after-base-change}. Par conséquent $\dim_z(X
\times Y) = \dim_x(X) + \dim_y(Y)$ d'après Morphismes, lemme
\ref{morphisms-lemma-dimension-fibre-at-a-point-additive}. La partie (2) est une conséquence directe de (1).
\end{proof}






\section{Anneaux locaux complets}
\label{section-complete-local-rings}

\noindent
Quelques résultats sur les anneaux locaux complets des schémas sur un corps.

\begin{lemma}
\label{lemma-complete-local-ring-structure-as-algebra}
Soit $k$ un corps. Soit $X$ un
schéma localement noethérien sur $k$. Soit $x \in X$
un point de corps résiduel $\kappa$. Il existe un isomorphisme
\begin{equation}
\label{equation-complete-local-ring}
\kappa[[x_1, \ldots, x_n]]/I \longrightarrow \mathcal{O}_{X, x}^\wedge
\end{equation}
induisant l'identité sur les corps résiduels. En général, on ne
peut choisir (\ref{equation-complete-local-ring}) comme
isomorphisme de $k$-algèbres. Toutefois, si l'extension
$\kappa/k$ est séparable, on peut choisir
(\ref{equation-complete-local-ring}) comme isomorphisme de $k$-algèbres.
\end{lemma}

\begin{proof}
L'existence de l'isomorphisme est une conséquence immédiate du
théorème de structure de Cohen\footnote{Notons que si $\kappa$ a
pour caractéristique $p$, alors le théorème dit simplement que nous
obtenons une surjection $\Lambda[[x_1, \ldots, x_n]] \to
\mathcal{O}_{X, x}^\wedge$ où $\Lambda$ est un anneau de Cohen pour
$\kappa$. Mais bien sûr, dans ce cas, l'application se factorise par
$\Lambda/p\Lambda[[x_1, \ldots, x_n]]$ et $\Lambda/p\Lambda = \kappa$.
} (Algèbre, théorème \ref{algebra-theorem-cohen-structure-theorem}).

\medskip\noindent
Soit $p$ un nombre premier impair, posons $k =
\mathbf{F}_p(t)$ et $A = k[x, y]/(y^2 + x^p - t)$. Alors le
complété $A^\wedge$ de $A$ pour l'idéal maximal $\mathfrak m
= (y)$ est isomorphe à $k(t^{1/p})[[z]]$ comme anneau,
mais non comme $k$-algèbre. En effet, $A^\wedge$
ne contient pas d'élément dont la $p$-ième puissance est $t$
(comme le lecteur peut le voir en calculant modulo $y^2$). Cela
montre également qu'aucun isomorphisme
(\ref{equation-complete-local-ring}) ne peut être un isomorphisme de $k$-algèbres.

\medskip\noindent
Si $\kappa/k$ est séparable, alors il existe un
homomorphisme de $k$-algèbres $\kappa \to \mathcal{O}_{X,
x}^\wedge$ induisant l'identité sur les corps résiduels
d'après Compléments d'algèbre, lemme
\ref{more-algebra-lemma-lift-residue-field}. Soient $f_1,
\ldots, f_n \in \mathfrak m_x$ des générateurs. Considérons l'homomorphisme
$$
\kappa[[x_1, \ldots, x_n]] \longrightarrow \mathcal{O}_{X, x}^\wedge,\quad
x_i \longmapsto f_i
$$
Puisque les deux membres sont complets pour la topologie
$(x_1, \ldots, x_n)$-adique (le membre de droite d'après
Algèbre, lemme
\ref{algebra-lemma-hathat-finitely-generated}), cette
application est surjective d'après Algèbre, lemme
\ref{algebra-lemma-completion-generalities} car elle est
surjective modulo $(x_1, \ldots, x_n)$ par construction.
\end{proof}

\begin{lemma}
\label{lemma-base-change-complete-local-ring}
Soit $K/k$ une extension de corps. Soit
$X$ un $k$-schéma localement algébrique. Posons $Y = X_K$.
Soit $y \in Y$ un point d'image $x \in X$.
Supposons que $\dim(\mathcal{O}_{X, x}) = \dim(\mathcal{O}_{Y,
y})$ et que $\kappa(x)/k$ soit séparable. Choisissons un isomorphisme
$$
\kappa(x)[[x_1, \ldots, x_n]]/(g_1, \ldots, g_m) \longrightarrow
\mathcal{O}_{X, x}^\wedge
$$
de $k$-algèbres comme dans
(\ref{equation-complete-local-ring}). On a alors un isomorphisme
$$
\kappa(y)[[x_1, \ldots, x_n]]/(g_1, \ldots, g_m) \longrightarrow
\mathcal{O}_{Y, y}^\wedge
$$
de $K$-algèbres comme dans
(\ref{equation-complete-local-ring}). On utilise ici $\kappa(x)
\to \kappa(y)$ pour considérer $g_j$ comme une série formelle sur $\kappa(y)$.
\end{lemma}

\begin{proof}
L'homomorphisme local $\mathcal{O}_{X, x}
\to \mathcal{O}_{Y, y}$ induit un homomorphisme
local $\mathcal{O}_{X, x}^\wedge \to
\mathcal{O}_{Y, y}^\wedge$. L'homomorphisme induit
$$
\kappa(x) \to \kappa(x)[[x_1, \ldots, x_n]]/(g_1, \ldots, g_m)
\to \mathcal{O}_{X, x}^\wedge \to \mathcal{O}_{Y, y}^\wedge
$$
composée avec la projection vers $\kappa(y)$ est
l'homomorphisme canonique $\kappa(x) \to \kappa(y)$. D'après
le lemme \ref{lemma-change-fields-flat}, le corps résiduel
$\kappa(y)$ est une localisation de $\kappa(x) \otimes_k K$
au noyau $\mathfrak p_0$ de $\kappa(x) \otimes_k K \to
\kappa(y)$. D'autre part, d'après le lemme
\ref{lemma-change-fields-algebraic-unramified}, l'anneau local $(\kappa(x)
\otimes_k K)_{\mathfrak p_0}$ est égal à $\kappa(y)$. Ainsi, l'application
$$
\kappa(x) \otimes_k K \to \mathcal{O}_{Y, y}^\wedge
$$
se factorise canoniquement par
$\kappa(y)$. Nous obtenons un diagramme commutatif
$$
\xymatrix{
\kappa(y) \ar[rr] & & \mathcal{O}_{Y, y}^\wedge \\
\kappa(x) \ar[r] \ar[u] &
\kappa(x)[[x_1, \ldots, x_n]]/(g_1, \ldots, g_m) \ar[r] &
\mathcal{O}_{X, x}^\wedge \ar[u]
}
$$
Soit $f_i \in \mathfrak m_x^\wedge \subset
\mathcal{O}_{X, x}^\wedge$ soit l'image de $x_i$.
Observons que $\mathfrak m_x^\wedge = (f_1, \ldots, f_n)$ puisque
l'homomorphisme est surjectif. Considérons l'homomorphisme
$$
\kappa(y)[[x_1, \ldots, x_n]] \longrightarrow \mathcal{O}_{Y, y}^\wedge,\quad
x_i \longmapsto f_i
$$
où ici $f_i$ signifie vraiment l'image de $f_i$ dans
$\mathfrak m_y^\wedge$. Puisque $\mathfrak m_x
\mathcal{O}_{Y, y} = \mathfrak m_y$ d'après le lemme
\ref{lemma-change-fields-algebraic-unramified}, nous voyons
que le membre de droite est complet pour la topologie $(x_1, \ldots,
x_n)$-adique (on utilise Algèbre, lemme
\ref{algebra-lemma-hathat-finitely-generated} pour voir que le membre
de droite est un anneau local complet). Puisque les deux membres sont complets pour
la topologie $(x_1, \ldots, x_n)$-adique, notre application est
surjective par Algèbre, lemme
\ref{algebra-lemma-completion-generalities} car elle est surjective modulo
$(x_1, \ldots, x_n)$. Bien sûr, les séries formelles $g_1, \ldots, g_m$ sont
envoyées sur zéro par cet homomorphisme, puisqu'elles le sont déjà
dans $\mathcal{O}_{X, x}^\wedge$. On obtient donc le diagramme commutatif
$$
\xymatrix{
\kappa(y)[[x_1, \ldots, x_n]]/(g_1, \ldots, g_m) \ar[r] &
\mathcal{O}_{Y, y}^\wedge \\
\kappa(x)[[x_1, \ldots, x_n]]/(g_1, \ldots, g_m) \ar[r] \ar[u] &
\mathcal{O}_{X, x}^\wedge \ar[u]
}
$$
Nous devons encore montrer que la flèche horizontale supérieure
est un isomorphisme. Nous savons déjà qu'elle est surjective.
Nous savons que $\mathcal{O}_{X, x} \to \mathcal{O}_{Y, y}$ est
plat (lemme \ref{lemma-change-fields-flat}), ce qui implique
que $\mathcal{O}_{X, x}^\wedge \to \mathcal{O}_{Y, y}^\wedge$
est plat (Compléments d'algèbre, lemme
\ref{more-algebra-lemma-flat-completion}). Ainsi, nous pouvons
appliquer Algèbre, lemme \ref{algebra-lemma-mod-injective} avec $R =
\kappa(x)[[x_1, \ldots, x_n]]/(g_1, \ldots, g_m)$, avec $S =
\kappa(y)[[x_1, \ldots, x_n]]/(g_1, \ldots, g_m)$, avec $M =
\mathcal{O}_{Y, y}^\wedge$ et avec $N = S$ pour conclure que l'application est injective.
\end{proof}


\section{Engendrement par les sections globales}
\label{section-global-generation}

\noindent
Quelques lemmes sur l'engendrement des modules quasi-cohérents par leurs sections globales.

\begin{lemma}
\label{lemma-globally-generated-base-change}
Soit $X \to \Spec(A)$ un morphisme de schémas. Soit $A \subset
A'$ une inclusion d'anneaux fidèlement plate. Soit $\mathcal{F}$
un $\mathcal{O}_X$-module quasi-cohérent. Alors $\mathcal{F}$ est engendré par ses sections globales
si et seulement si $\mathcal{F}_{A'}$ est engendré par ses sections globales.
\end{lemma}

\begin{proof}
Plus précisément, posons $X_{A'} = X \times_{\Spec(A)}
\Spec(A')$. Soit $\mathcal{F}_{A'} = p^*\mathcal{F}$ où $p :
X_{A'} \to X$ est la projection. D'après Cohomologie des
schémas, lemme
\ref{coherent-lemma-flat-base-change-cohomology} nous avons
$H^0(X_{k'}, \mathcal{F}_{A'}) = H^0(X, \mathcal{F}) \otimes_A
A'$. Ainsi, si $s_i$, $i \in I$ sont des générateurs de $H^0(X,
\mathcal{F})$ comme module $A$, alors leurs images dans
$H^0(X_{A'}, \mathcal{F}_{A'})$ sont des générateurs de
$H^0(X_{A'}, \mathcal{F}_{A'})$ comme module $A'$. Ainsi, nous
devons montrer que l'application $\alpha : \bigoplus_{i \in I}
\mathcal{O}_X \to \mathcal{F}$, $(f_i) \mapsto \sum f_i s_i$ est
surjective si et seulement si $p^*\alpha$ est surjective. Cela,
nous pouvons le vérifier sur un ouvert affine $U = \Spec(B)$ de
$X$. Alors $\mathcal{F}|_U$ correspond à un $B$-module $M$ et
$s_i|_U$ aux éléments $x_i \in M$. Ainsi, nous devons montrer que
$\bigoplus_{i \in I} B \to M$ est surjective si et seulement si le
changement de base $\bigoplus_{i \in I} B \otimes_A A' \to M \otimes_A
A'$ est surjective. C'est vrai parce que $A \to A'$ est fidèlement plat.
\end{proof}

\begin{lemma}
\label{lemma-very-ample-vanish-at-point}
Soit $k$ un corps infini. Soit $X$
un schéma de type fini sur $k$. Soit $\mathcal{L}$
un faisceau inversible très ample sur $X$. Soient
$n \geq 0$ et $x, x_1, \ldots, x_n \in X$, avec
$x$ un point $k$-rationnel, c'est-à-dire $\kappa(x) = k$,
et $x \not = x_i$ pour $i = 1, \ldots, n$. Alors il existe un $s
\in H^0(X, \mathcal{L})$ qui s'annule en $x$ mais en aucun des $x_i$.
\end{lemma}

\begin{proof}
Si $n = 0$ le résultat est trivial, nous supposons donc $n >
0$. Par définition d'un faisceau inversible très ample, le
lemme se réduit immédiatement au cas où $X = \mathbf{P}^r_k$
pour un certain $r > 0$ et $\mathcal{L} = \mathcal{O}_X(1)$.
Écrivons $\mathbf{P}^r_k = \text{Proj}(k[T_0, \ldots, T_r])$.
Posons $V = H^0(X, \mathcal{L}) = kT_0 \oplus \ldots \oplus
kT_r$. Puisque $x$ est un point $k$-rationnel, on voit que
l'ensemble des $s \in V$ qui s'annulent en $x$ est un sous-espace
de codimension $1$, à savoir $W \subset V$, et que $W$ engendre l'idéal premier
homogène correspondant à $x$. Puisque $x_i \not = x$, l'idéal premier
homogène correspondant $\mathfrak p_i \subset k[T_0, \ldots, T_r]$
ne contient pas $W$. Comme $k$ est infini, nous voyons alors que $W
\not = \bigcup W \cap \mathfrak q_i$ et la démonstration est complète.
\end{proof}

\begin{lemma}
\label{lemma-generated-by-dim-plus-1-sections}
Soit $k$ un corps infini. Soit $X$ un $k$-schéma
algébrique. Soit $\mathcal{L}$ un $\mathcal{O}_X$-module
inversible. Soit $V \to \Gamma(X, \mathcal{L})$ une application
linéaire de $k$-espaces vectoriels dont l'image génère $\mathcal{L}$. Alors il existe
un sous-espace $W \subset V$ avec $\dim_k(W) \leq \dim(X) + 1$ qui génère $\mathcal{L}$.
\end{lemma}

\begin{proof}
Tout au long de la démonstration, nous utiliserons
que pour tout $x \in X$, l'application linéaire
$$
\psi_x : V \to \Gamma(X, \mathcal{L}) \to \mathcal{L}_x \to
\mathcal{L}_x \otimes_{\mathcal{O}_{X, x}} \kappa(x)
$$
n'est pas nulle. La démonstration se fait par récurrence sur $\dim(X)$.

\medskip\noindent
Le cas de base est $\dim(X) = 0$. Dans ce cas, $X$ a un
nombre fini de points $X = \{x_1, \ldots, x_n\}$ (voir par
exemple le lemme \ref{lemma-algebraic-scheme-dim-0}). Puisque
$k$ est infini, il existe un vecteur $v \in V$ tel que
$\psi_{x_i}(v) \not = 0$ pour tout $i$. Alors $W = k\cdot v$ convient.

\medskip\noindent
Supposons $\dim(X) > 0$. Soient $X_i \subset X$ les composantes
irréductibles de dimension égale à $\dim(X)$. Comme $X$
est noethérien, elles sont en nombre fini. Pour tout $i$,
choisissons un point $x_i \in X_i$. Comme ci-dessus, choisissons $v
\in V$ tel que $\psi_{x_i}(v) \not = 0$ pour tout $i$. Soit
$Z \subset X$ le schéma des zéros de l'image de $v$ dans $\Gamma(X,
\mathcal{L})$, voir Diviseurs, définition
\ref{divisors-definition-zero-scheme-s}. Par construction $\dim(Z) < \dim(X)$.
Par induction, nous pouvons trouver $W \subset V$ avec $\dim(W) \leq \dim(X)$
tel que $W$ génère $\mathcal{L}|_Z$. Alors $W + k\cdot v$ génère $\mathcal{L}$.
\end{proof}






\section{Séparation des points et des vecteurs tangents}
\label{section-separating-points-tangent-vectors}

\noindent
C'est simplement le résultat suivant.

\begin{lemma}
\label{lemma-separate-points-tangent-vectors}
Soit $k$ un corps algébriquement clos. Soit
$X$ un $k$-schéma propre. Soit $\mathcal{L}$
un $\mathcal{O}_X$-module inversible. Soit $V
\subset H^0(X, \mathcal{L})$ soit un $k$-sous-espace vectoriel. Si
\begin{enumerate}
\item pour toute paire de points fermés distincts $x, y \in X$, il
existe une section $s \in V$ qui s'annule en $x$ mais non en $y$, et
\item pour tout point fermé $x \in X$ et tout vecteur tangent non
nul $\theta \in T_{X/k, x}$, il existe une section $s \in V$ qui
s'annule en $x$ mais dont l'image réciproque par $\theta$ est non nulle,
\end{enumerate}
alors $\mathcal{L}$ est très ample et le
morphisme canonique $\varphi_{\mathcal{L}, V} :
X \to \mathbf{P}(V)$ est une immersion fermée.
\end{lemma}

\begin{proof}
La condition (1) implique en particulier que les éléments de $V$ engendrent
$\mathcal{L}$ sur $X$. On obtient donc un morphisme canonique
$$
\varphi = \varphi_{\mathcal{L}, V} :
X
\longrightarrow
\mathbf{P}(V)
$$
par Constructions, exemple
\ref{constructions-example-projective-space}. Le morphisme
$\varphi$ est propre d'après Morphismes, lemme
\ref{morphisms-lemma-image-proper-scheme-closed}. D'après (1),
le morphisme $\varphi$ est injectif sur les points fermés (calcul
omis). En particulier, la fibre au-dessus de tout point fermé de
$\mathbf{P}(V)$ est réduite à un point (petit détail omis). Ainsi, nous voyons
que $\varphi$ est fini, par exemple en utilisant Cohomologie des
schémas, lemme
\ref{coherent-lemma-proper-finite-fibre-finite-in-neighbourhood}. Pour terminer la démonstration, il suffit de montrer que l'application
$$
\varphi^\sharp :
\mathcal{O}_{\mathbf{P}(V)}
\longrightarrow
\varphi_*\mathcal{O}_X
$$
est surjective. Nous pouvons le vérifier sur les germes aux
points fermés. Soit $x \in X$ un point fermé d'image le point
fermé $p = \varphi(x) \in \mathbf{P}(V)$. Puisque
$\varphi^{-1}(\{p\}) = \{x\}$ par (1) et que $\varphi$ est
propre (donc fermé), on voit que $\varphi^{-1}(U)$ décrit un
système fondamental de voisinages ouverts de $x$ lorsque $U$
décrit un système fondamental de voisinages ouverts de $p$. Nous
concluons qu'en passant aux germes en $p$, nous obtenons l'homomorphisme
$$
\varphi^\sharp_x :
\mathcal{O}_{\mathbf{P}(V), p}
\longrightarrow
\mathcal{O}_{X, x}
$$
En particulier, $\mathcal{O}_{X, x}$ est un
$\mathcal{O}_{\mathbf{P}(V), p}$-module fini. De plus, les
corps résiduels de $x$ et $p$ sont égaux à $k$ (comme $k$
est algébriquement clos --- on utilise le Nullstellensatz de
Hilbert). Enfin, la condition (2) implique que l'application
$$
T_{X/k, x} \longrightarrow T_{\mathbf{P}(V)/k, p}
$$
est injective puisque tout $\theta$ non nul dans le noyau de cette
application ne saurait satisfaire la conclusion de (2). En
termes du morphisme d'anneaux locaux ci-dessus, cela signifie que
$$
\mathfrak m_p/\mathfrak m_p^2 \longrightarrow
\mathfrak m_x/\mathfrak m_x^2
$$
est surjective, voir le lemme
\ref{lemma-tangent-space-rational-point}. La preuve s'achève
alors par Algèbre, lemme \ref{algebra-lemma-when-surjective-local}.
\end{proof}

\begin{lemma}
\label{lemma-variant-separate-points-tangent-vectors}
Soit $k$ un corps algébriquement clos.
Soit $X$ un $k$-schéma propre. Soit
$\mathcal{L}$ un $\mathcal{O}_X$-module
inversible. Supposons que, pour tout sous-schéma fermé $Z
\subset X$ de dimension $0$ et de degré $2$ sur $k$, l'application
$$
H^0(X, \mathcal{L}) \longrightarrow H^0(Z, \mathcal{L}|_Z)
$$
est surjective. Alors $\mathcal{L}$ est très ample sur $X$ relativement à $k$.
\end{lemma}

\begin{proof}
C'est une reformulation du lemme
\ref{lemma-separate-points-tangent-vectors}. En effet,
étant donnés deux points fermés distincts $x, y
\in X$, le choix $Z = x \cup y$ (considéré comme
sous-schéma fermé) donne la condition (1) du
lemme. Étant donné un vecteur tangent non nul
$\theta \in T_{X/k, x}$ le morphisme $\theta :
\Spec(k[\epsilon]) \to X$ est une immersion fermée. En
posant $Z = \Im(\theta)$ nous obtenons la condition (2) du lemme.
\end{proof}









\section{Adhérences de produits}
\label{section-closure-of-products}

\noindent
Quelques résultats sur les rapports entre l'adhérence et les produits.

\begin{lemma}
\label{lemma-closure-of-product}
Soit $k$ un corps.
Soient $X$, $Y$ des schémas sur $k$, et soient
$A \subset X$, $B \subset Y$ des
parties. Posons
$$
AB =
\{z \in X \times_k Y \mid \text{pr}_X(z) \in A, \ \text{pr}_Y(z) \in B\}
\subset X \times_k Y
$$
Alors, ensemblistement, nous avons
$$
\overline{A} \times_k \overline{B} = \overline{AB}
$$
\end{lemma}

\begin{proof}
L'inclusion $\overline{AB} \subset \overline{A}
\times_k \overline{B}$ est immédiate. Nous pouvons
remplacer $X$ et $Y$ par les sous-schémas fermés réduits
$\overline{A}$ et $\overline{B}$. Soit $W
\subset X \times_k Y$ soit un ouvert non vide. Par
Morphismes, lemme
\ref{morphisms-lemma-scheme-over-field-universally-open}, le
sous-ensemble $U = \text{pr}_X(W)$ est ouvert non vide dans
$X$. Ainsi $A \cap U$ n'est pas vide. Choisissons $a \in A \cap
U$. Notons $Y_a = \{a\} \times_k Y \cong
Y_{\kappa(a)}$ la fibre de $\text{pr}_X : X \times_k Y \to X$
au-dessus de $a$. D'après Morphismes, lemme
\ref{morphisms-lemma-scheme-over-field-universally-open}, le morphisme
$Y_a \to Y$ est encore ouvert puisque
$\Spec(\kappa(a)) \to \Spec(k)$ est universellement ouvert. Ainsi
l'ouvert non vide $W_a = W \times_{X \times_k Y} Y_a$ se projette
sur un ouvert non vide de $Y$. Nous en concluons qu'il existe un $b
\in B$ dans l'image. Ainsi $AB \cap W \not = \emptyset$ comme voulu.
\end{proof}

\begin{lemma}
\label{lemma-closure-image-product-map}
Soit $k$ un corps. Soient $f : A
\to X$, $g : B \to Y$ des morphismes de schémas
sur $k$. Alors, ensemblistement, on a
$$
\overline{f(A)} \times_k \overline{g(B)} =
\overline{(f \times g)(A \times_k B)}
$$
\end{lemma}

\begin{proof}
Cela découle du lemme
\ref{lemma-closure-of-product}
puisque l'image de $f \times g$ est
$f(A)g(B)$ dans la notation de ce lemme.
\end{proof}

\begin{lemma}
\label{lemma-scheme-theoretic-image-product-map}
Soit $k$ un corps. Soient $f : A \to X$, $g
: B \to Y$ des morphismes quasi-compacts de schémas sur
$k$. Soit $Z \subset X$ l'image schématique de $f$;
voir Morphismes, définition
\ref{morphisms-definition-scheme-theoretic-image}. De même, soit $Z' \subset Y$
l'image schématique de $g$. Alors $Z \times_k Z'$ est l'image schématique de $f \times g$.
\end{lemma}

\begin{proof}
Rappelons que $Z$ est le plus petit sous-schéma fermé de $X$ par
lequel $f$ se factorise. De même pour $Z'$. Soit $W \subset X
\times_k Y$ l'image schématique de $f \times g$. Comme $f \times g$ se
factorise par $Z \times_k Z'$, on voit que $W \subset Z \times_k Z'$.

\medskip\noindent
Pour démontrer l'autre inclusion, soient $U \subset X$ et $V
\subset Y$ des ouverts affines. D'après Morphismes, lemme
\ref{morphisms-lemma-quasi-compact-scheme-theoretic-image},
le schéma $Z \cap U$ est l'image schématique de
$f|_{f^{-1}(U)} : f^{-1}(U) \to U$, et de même pour $Z'
\cap V$ et $W \cap U \times_k V$. On peut donc
supposer que $X$ et $Y$ sont affines. Comme $f$ et
$g$ sont quasi-compacts, cela implique que $A = \bigcup
U_i$ est une union finie d'affines et $B = \bigcup V_j$ est
une union finie d'affines. On peut alors remplacer
$A$ par $\coprod U_i$ et $B$ par $\coprod V_j$,
c'est-à-dire qu'on peut supposer que $A$ et $B$ sont
également affines. Dans ce cas, $Z$ est défini par
$\Ker(\Gamma(X, \mathcal{O}_X) \to \Gamma(A, \mathcal{O}_A))$ et de
même pour $Z'$ et $W$. Par conséquent, le résultat découle de l'égalité
$$
\Gamma(A \times_k B, \mathcal{O}_{A \times_k B})
=
\Gamma(A, \mathcal{O}_A) \otimes_k \Gamma(B, \mathcal{O}_B)
$$
qui est vraie puisque $A$ et $B$ sont affines. Nous omettons les détails.
\end{proof}





\section{Schémas lisses sur des corps} 
\label{section-smooth}

\noindent
 Voici deux lemmes caractérisant les schémas lisses sur les corps.

\begin{lemma}
\label{lemma-char-zero-differentials-free-smooth}
 Soit $k$ un corps. Soit $X$ un schéma sur $k$.
Supposons 
\begin{enumerate}
\item $X$ est localement de type fini sur $k$, 
\item $\Omega_{X/k}$ est localement libre, et 
\item $k$ a caractéristique zéro. 
\end{enumerate}
 Alors le morphisme de structure $X \to \Spec(k)$ est lisse. 
\end{lemma}

\begin{proof}
 Cela résulte
d'Algèbre, lemme \ref{algebra-lemma-characteristic-zero-local-smooth}. 
\end{proof}

\noindent
 En caractéristique positive, il existe des schémas de type fini
non réduits dont le faisceau des différentielles est libre, par exemple 
$\Spec(\mathbf{F}_p[t]/(t^p))$ sur $\Spec(\mathbf{F}_p)$.
Si le corps de base $k$ est non parfait de caractéristique $p$,
il existe des schémas réduits $X/k$ avec $\Omega_{X/k}$ libre qui
ne sont pas lisses, par exemple $\Spec(k[t]/(t^p-a)$ où $a \in k$
 n'est pas une puissance $p$.

\begin{lemma}
\label{lemma-char-p-differentials-free-smooth}
 Soit $k$ un corps. Soit $X$ un schéma sur $k$.
Supposons 
\begin{enumerate}
\item $X$ est localement de type fini sur $k$, 
\item $\Omega_{X/k}$ est localement libre, 
\item $X$ est réduit, et 
\item $k$ est parfait. 
\end{enumerate}
 Alors le morphisme de structure $X \to \Spec(k)$ est lisse. 
\end{lemma}

\begin{proof}
 Soit $x \in X$ un point. Comme $X$ est localement noethérien (voir
Morphismes, lemme \ref{morphisms-lemma-finite-type-noetherian}),
il y a un nombre fini de composantes irréductibles 
$X_1, \ldots, X_n$ passant par $x$ (voir
Propriétés, lemme \ref{properties-lemma-Noetherian-topology} et
Topologie, lemme \ref{topology-lemma-Noetherian}).
Soit $\eta_i \in X_i$ le point générique. Comme $X$ est réduit, on a 
$\mathcal{O}_{X, \eta_i} = \kappa(\eta_i)$,
voir Algèbre, lemme \ref{algebra-lemma-minimal-prime-reduced-ring}.
De plus, $\kappa(\eta_i)$ est une extension de corps de type
fini du corps parfait $k$, donc séparablement engendrée sur $k$ (voir
Algèbre, section \ref{algebra-section-separability}).
Il s'ensuit que $\Omega_{X/k, \eta_i} = \Omega_{\kappa(\eta_i)/k}$
 est libre de rang égal au degré de transcendance de $\kappa(\eta_i)$ sur $k$.
D'après Morphismes, lemme \ref{morphisms-lemma-dimension-fibre-at-a-point},
nous concluons que $\dim_{\eta_i}(X_i) = \text{rank}_{\eta_i}(\Omega_{X/k})$.
Puisque $x \in X_1 \cap \ldots \cap X_n$, on voit que 
$$
\text{rank}_x(\Omega_{X/k}) = \text{rank}_{\eta_i}(\Omega_{X/k}) = \dim(X_i).
$$
 Par conséquent $\dim_x(X) = \text{rank}_x(\Omega_{X/k})$,
voir Algèbre, lemme \ref{algebra-lemma-dimension-at-a-point-finite-type-over-field}.
Il s'ensuit que $X \to \Spec(k)$ est lisse en $x$ par exemple
d'après Algèbre, lemme \ref{algebra-lemma-characterize-smooth-over-field}. 
\end{proof}

\begin{lemma}
\label{lemma-smooth-regular}
\begin{slogan}
 Être lisse sur un corps implique être régulier 
\end{slogan}
 Soit $X \to \Spec(k)$ un morphisme lisse, où $k$ est un corps.
Alors $X$ est un schéma régulier. 
\end{lemma}

\begin{proof}
 (Voir aussi
le lemme \ref{lemma-geometrically-regular-smooth}.)

D'après Algèbre, lemme \ref{algebra-lemma-characterize-smooth-over-field},
tout anneau local $\mathcal{O}_{X, x}$ est régulier.
Comme $X$ est localement de type fini sur $k$, il est
localement noethérien. Par conséquent, $X$ est régulier
d'après Propriétés, lemme \ref{properties-lemma-characterize-regular}. 
\end{proof}

\begin{lemma}
\label{lemma-smooth-geometrically-normal}
 Soit $X \to \Spec(k)$ un morphisme lisse, où $k$ est un corps.
Alors $X$ est géométriquement régulier, géométriquement normal et
géométriquement réduit sur $k$. 
\end{lemma}

\begin{proof}
 (Voir aussi
le lemme \ref{lemma-geometrically-regular-smooth}.)
Soit $k'$ une extension finie purement inséparable de $k$.
Il suffit de montrer que $X_{k'}$ est régulier, normal et réduit ; voir les
lemmes \ref{lemma-geometrically-regular}, 
\ref{lemma-geometrically-normal}, et 
\ref{lemma-check-only-finite-inseparable-extensions}.

D'après Morphismes, lemme \ref{morphisms-lemma-base-change-smooth},
le morphisme $X_{k'} \to \Spec(k')$ est également lisse.
Par conséquent, il suffit de montrer qu'un schéma $X$ lisse sur un corps est régulier,
normal et réduit. On voit que $X$ est régulier d'après
le lemme \ref{lemma-smooth-regular}.

Par conséquent, Propriétés, lemme \ref{properties-lemma-regular-normal}
 garantit que $X$ est normal. 
\end{proof}

\begin{lemma}
\label{lemma-affine-space-over-field}
 Soit $k$ un corps, soit $d \geq 0$ et soit $W \subset \mathbf{A}^d_k$
 ouvert non vide. Alors il existe un point fermé $w \in W$ tel que 
$k \subset \kappa(w)$ soit finie et séparable. 
\end{lemma}

\begin{proof}
 Après avoir éventuellement rétréci $W$, on peut supposer que 
$W = \mathbf{A}^d_k \setminus V(f)$ pour certains $f \in k[x_1, \ldots, x_d]$.
Si le lemme est faux, alors $f(a_1, \ldots, a_d) = 0$ pour tous les 
$(a_1, \ldots, a_d) \in (k^{sep})^d$. C'est absurde car $k^{sep}$
 est un corps infini. 
\end{proof}

\begin{lemma}
\label{lemma-smooth-separable-closed-points-dense}
 Soit $k$ un corps. Si $X$ est lisse sur $\Spec(k)$ alors
l'ensemble 
$$
\{x \in X\text{ fermé tel que }k \subset \kappa(x)
\text{ est finie séparable}\}
$$
 est dense en $X$. 
\end{lemma}

\begin{proof}
 Il suffit de montrer que, pour tout schéma lisse non vide $X$ sur $k$,
il existe au moins un point fermé dont le corps résiduel
est fini et séparable sur $k$. Pour le voir, choisissons un diagramme 
$$
\xymatrix{
X & U \ar[l] \ar[r]^-\pi & \mathbf{A}^d_k
}
$$
 où $\pi$ est étale, voir
Morphismes, lemme \ref{morphisms-lemma-smooth-etale-over-affine-space}.
Le morphisme $\pi : U \to \mathbf{A}^d_k$ est ouvert,
voir Morphismes, lemme \ref{morphisms-lemma-etale-open}.
D'après
le lemme \ref{lemma-affine-space-over-field},
nous pouvons choisir un point fermé $w \in \pi(U)$ dont le corps résiduel
est fini et séparable sur $k$. Choisissons un $x \in U$ avec $\pi(x) = w$.
D'après Morphismes, lemme \ref{morphisms-lemma-etale-over-field},
l'extension de corps $\kappa(x)/\kappa(w)$ est finie et séparable. Par
conséquent, $\kappa(x)/k$ est finie et séparable. Le point $x$ est
un point fermé de $X$
 d'après Morphismes, lemme 
\ref{morphisms-lemma-algebraic-residue-field-extension-closed-point-fibre}. 
\end{proof}

\begin{lemma}
\label{lemma-geometrically-reduced-dense-smooth-open}
 Soit $X$ un schéma réduit localement de type fini sur un corps $k$.
Alors $X$ est géométriquement réduit sur $k$ si et seulement si 
$X$ contient un ouvert dense qui est lisse sur $k$. 
\end{lemma}

\begin{proof}
 La question étant locale sur $X$, on peut supposer que $X$ est quasi-compact. Soit 
$X = X_1 \cup \ldots \cup X_n$ la réunion des composantes irréductibles de $X$.
Supposons d'abord que $X$ contient un ouvert dense qui est lisse sur $k$.
Alors $X$ est géométriquement réduit au point générique de chaque $X_i$
 (par exemple d'après le lemme \ref{lemma-smooth-geometrically-normal}).
D'après le lemme \ref{lemma-generic-points-geometrically-reduced},
il s'ensuit que $X$ est géométriquement réduit.

\medskip\noindent
 Réciproquement, supposons que $X$ soit géométriquement réduit.
Comme $Z = \bigcup_{i \not = j} X_i \cap X_j$ est nulle part dense dans $X$,
nous pouvons remplacer $X$ par $X \setminus Z$. Comme $X \setminus Z$ est
une réunion disjointe de schémas irréductibles, cela nous ramène au cas
où $X$ est irréductible. Comme $X$ est irréductible et réduit, il
est intègre,
voir Propriétés, lemme \ref{properties-lemma-characterize-integral}.
Soit $\eta \in X$ son point générique.
Alors le corps des fonctions $K = k(X) = \kappa(\eta) = \mathcal{O}_{X, \eta}$
 est géométriquement réduit sur $k$, donc séparable sur $k$,
voir Algèbre, lemme \ref{algebra-lemma-characterize-separable-field-extensions}.
Soit $U = \Spec(A) \subset X$ un ouvert affine non vide quelconque, de
sorte que $K = A_{(0)}$ soit le corps des fractions de $A$.
Appliquons Algèbre, lemme \ref{algebra-lemma-separable-smooth}
 pour en déduire que $A$ est lisse en $(0)$ sur $k$.
Par définition, cela signifie qu'une certaine localisation principale de 
$A$ est lisse sur $k$, ce qui achève la démonstration. 
\end{proof}

\begin{lemma}
\label{lemma-dense-smooth-open-variety-over-perfect-field}
 Soit $k$ un corps parfait. Soit $X$
un $k$-schéma réduit localement algébrique, par exemple une variété sur $k$. Alors nous avons 
$$
\{x \in X \mid X \to \Spec(k)\text{ is smooth at }x\} =
\{x \in X \mid \mathcal{O}_{X, x}\text{ is regular}\}
$$
 et c'est un ouvert dense de $X$. 
\end{lemma}

\begin{proof}
 L'égalité des deux ensembles résulte immédiatement
d'Algèbre, lemme \ref{algebra-lemma-separable-smooth} et des définitions
(voir Algèbre, définition \ref{algebra-definition-perfect} pour la définition
d'un corps parfait). L'ensemble est ouvert car l'ensemble des points où
un morphisme de schémas est lisse est
ouvert, voir Morphismes, définition \ref{morphisms-definition-smooth}.
Enfin, nous donnons deux arguments pour voir qu'il est dense :
(1) Les points génériques de $X$ sont dans l'ensemble car les anneaux locaux
aux points génériques sont des corps (Algèbre, lemme 
\ref{algebra-lemma-minimal-prime-reduced-ring}), donc réguliers. (2)
Nous utilisons que $X$ est géométriquement réduit d'après le
lemme \ref{lemma-perfect-reduced} et donc
le lemme \ref{lemma-geometrically-reduced-dense-smooth-open} s'applique. 
\end{proof}

\begin{lemma}
\label{lemma-flat-under-smooth}
 Soit $k$ un corps. Soit $f : X \to Y$ un morphisme de schémas localement
de type fini sur $k$. Soit $x \in X$ un point et posons $y = f(x)$.
Si $X \to \Spec(k)$ est lisse en $x$ et que $f$ est plat en $x$
 alors $Y \to \Spec(k)$ est lisse en $y$. En particulier, si $X$
 est lisse sur $k$ et que $f$ est plat et surjectif, alors $Y$ est lisse sur $k$. 
\end{lemma}

\begin{proof}
 Il suffit de montrer que $Y$ est géométriquement régulier en $y$, voir
le lemme \ref{lemma-geometrically-regular-smooth}.
Cela découle du
lemme \ref{lemma-flat-under-geometrically-regular}
 (et
du lemme \ref{lemma-geometrically-regular-smooth}
 appliqué à $(X, x)$). 
\end{proof}

\begin{lemma}
\label{lemma-variety-with-smooth-rational-point}
 Soit $k$ un corps. Soit $X$ une variété sur $k$ possédant un point 
$k$-rationnel $x$ tel que $X$ soit lisse en $x$.
Alors $X$ est géométriquement intègre sur $k$. 
\end{lemma}

\begin{proof}
 Soit $U \subset X$ le lieu lisse de $X$. Par hypothèse, $U$ n'est pas
vide et est donc dense et schématiquement dense. Alors 
$U_{\overline{k}} \subset X_{\overline{k}}$ est également
dense et schématiquement dense (certains détails omis). Par conséquent,
il suffit de montrer que $U$ est géométriquement intègre.
Comme $U$ possède un point $k$-rationnel, il est géométriquement connexe d'après le
lemme \ref{lemma-geometrically-connected-if-connected-and-point}.
D'autre part, $U_{\overline{k}}$ est réduit et normal
(lemme \ref{lemma-smooth-geometrically-normal}).
Puisqu'un schéma noethérien connexe et
normal est intègre (Propriétés, lemme \ref{properties-lemma-normal-Noetherian}),
la démonstration est complète. 
\end{proof}

\begin{lemma}
\label{lemma-smooth-stratification}
 Soit $X$ un schéma de type fini sur un corps $k$.
Il existe une extension finie purement inséparable $k'/k$,
un entier $t \geq 0$, et des sous-schémas fermés 
$$
X_{k'} \supset Z_0 \supset Z_1 \supset \ldots \supset Z_t = \emptyset
$$
 de sorte que $Z_0 = (X_{k'})_{red}$ et $Z_i \setminus Z_{i + 1}$
 soient lisses sur $k'$ pour tout $i$. 
\end{lemma}

\begin{proof}
 Procédons par récurrence sur $\dim(X)$. D'après
le lemme \ref{lemma-finite-extension-geometrically-reduced},
nous pouvons trouver une extension finie purement inséparable $k'/k$
 telle que $(X_{k'})_{red}$ soit géométriquement réduit sur $k'$.
D'après le lemme \ref{lemma-geometrically-reduced-dense-smooth-open},
il existe un sous-schéma fermé nulle part dense $X' \subset (X_{k'})_{red}$
 tel que $(X_{k'})_{red} \setminus X'$ soit lisse sur $k'$.
Ensuite $\dim(X') < \dim(X)$. Par hypothèse d'induction, il existe
une extension finie purement inséparable $k''/k'$,
un entier $t' \geq 0$, et des sous-schémas fermés 
$$
X'_{k''} \supset Y_0 \supset Y_1 \supset \ldots \supset Y_{t'} = \emptyset
$$
 tels que $Y_0 = (X'_{k''})_{red}$ et $Y_i \setminus Y_{i + 1}$
 soit lisse sur $k''$ pour tout $i$. Posons alors $t = t' + 1$
 et nous considérons 
$$
X_{k''} \supset Z_0 \supset Z_1 \supset \ldots \supset Z_t = \emptyset
$$
 où $Z_0 = (X_{k''})_{red}$ et $Z_i = Y_{i - 1}$ pour $i > 0$;
cela a du sens car $X'_{k''}$ est un sous-schéma fermé de $X_{k''}$.
Nous omettons de vérifier toutes les propriétés énoncées. 
\end{proof}




\section{Types de variétés} 
\label{section-types}

\noindent
 Cette courte section traite de quelques propriétés globales élémentaires des variétés.

\begin{definition}
\label{definition-variety-type}
 Soit $k$ un corps. Soit $X$ une variété sur $k$. 
\begin{enumerate}
\item On dit que $X$ est une {\it variété affine} si $X$ est un schéma
affine. Cela équivaut à exiger que $X$ soit isomorphe à un sous-schéma
fermé de $\mathbf{A}^n_k$ pour un certain $n$. 
\item On dit que $X$ est une {\it variété projective} si le
morphisme de structure $X \to \Spec(k)$ est projectif.
D'après Morphismes, lemme \ref{morphisms-lemma-characterize-locally-projective},
cela est vrai si et seulement si $X$ est isomorphe à un sous-schéma
fermé de $\mathbf{P}^n_k$ pour un certain $n$. 
\item On dit que $X$ est une {\it variété quasi-projective} si le
morphisme de structure $X \to \Spec(k)$ est quasi-projectif. D'après Morphismes, lemme
\ref{morphisms-lemma-characterize-locally-quasi-projective},
cela est vrai si et seulement si $X$ est isomorphe à
un sous-schéma localement fermé de $\mathbf{P}^n_k$ pour un certain $n$. 
\item Une {\it variété propre} est une variété telle que
le morphisme $X \to \Spec(k)$ est propre. 
\item Une {\it variété lisse} est une variété telle que
le morphisme $X \to \Spec(k)$ est lisse. 
\end{enumerate}
\end{definition}

\noindent
 Notons qu'une variété projective est une variété propre,
voir Morphismes, lemme \ref{morphisms-lemma-locally-projective-proper}.
De plus, une variété affine est quasi-projective car $\mathbf{A}^n_k$
 est isomorphe à un sous-schéma ouvert de $\mathbf{P}^n_k$,
voir Constructions,
lemme \ref{constructions-lemma-standard-covering-projective-space}.

\begin{lemma}
\label{lemma-regular-functions-proper-variety}
 Soit $X$ une variété propre sur $k$. Alors 
\begin{enumerate}
\item $K = H^0(X, \mathcal{O}_X)$ est un corps qui
est une extension finie du corps $k$, 
\item si $X$ est géométriquement réduit, alors $K/k$ est séparable, 
\item si $X$ est géométriquement irréductible, alors $K/k$
 est purement inséparable, 
\item si $X$ est géométriquement intègre, alors $K = k$. 
\end{enumerate}
\end{lemma}

\begin{proof}
 Ceci est un cas particulier du
lemme \ref{lemma-proper-geometrically-reduced-global-sections}. 
\end{proof}




\section{Normalisation} 
\label{section-normalization}

\noindent
 Quelques questions relatives à la normalisation.

\begin{lemma}
\label{lemma-normalization-locally-algebraic}
 Soit $k$ un corps. Soit $X$ un schéma localement algébrique sur $k$.
Soit $\nu : X^\nu \to X$ le morphisme de normalisation,
voir Morphismes, définition \ref{morphisms-definition-normalization}.
Alors 
\begin{enumerate}
\item $\nu$ est fini, dominant, et $X^\nu$ est une
réunion disjointe de schémas irréductibles normaux localement algébriques sur $k$, 
\item $\nu$ se factorise en $X^\nu \to X_{red} \to X$ et le
premier morphisme est le morphisme de normalisation de $X_{red}$, 
\item si $X$ est un schéma algébrique réduit, alors $\nu$
 est birationnel, 
\item si $X$ est une variété, alors $X^\nu$ est une variété et 
$\nu$ est un morphisme birationnel fini de variétés. 
\end{enumerate}
\end{lemma}

\begin{proof}
 Puisque $X$ est localement de type fini sur un corps, on voit que 
$X$ est localement noethérien
(Morphismes, lemme \ref{morphisms-lemma-finite-type-noetherian}),
donc chaque ouvert quasi-compact a un nombre fini
de composantes irréductibles (Propriétés, lemme 
\ref{properties-lemma-Noetherian-irreducible-components}).
Ainsi, Morphismes, définition \ref{morphisms-definition-normalization}, s'applique.
La normalisation $X^\nu$ est toujours une réunion disjointe de schémas
intègres normaux et le morphisme de normalisation $\nu$ est toujours dominant,
voir Morphismes, lemme \ref{morphisms-lemma-normalization-normal}.
Comme $X$ est universellement Nagata
(Morphismes, lemme \ref{morphisms-lemma-ubiquity-nagata}),
on voit que $\nu$ est
fini (Morphismes, lemme \ref{morphisms-lemma-nagata-normalization}).
Ainsi $X^\nu$ est également localement algébrique. À
ce stade, nous avons prouvé (1).

\medskip\noindent
 La partie (2) est donnée par Morphismes, lemme \ref{morphisms-lemma-normalization-reduced}.

\medskip\noindent
 La partie (3) est donnée par Morphismes, lemme \ref{morphisms-lemma-normalization-birational}.

\medskip\noindent
 La partie (4) découle de (1), (2), (3), et du fait que $X^\nu$
 est séparé en tant que schéma fini sur un schéma séparé. 
\end{proof}

\begin{lemma}
\label{lemma-normalize-projective}
 Soit $k$ un corps. Soit $X$ un schéma propre sur $k$.
Soit $\nu : X^\nu \to X$ le morphisme de normalisation,
voir Morphismes, définition \ref{morphisms-definition-normalization}.
Alors $X^\nu$ est propre sur $k$. Si $X$ est projectif sur $k$,
alors $X^\nu$ est projectif sur $k$. 
\end{lemma}

\begin{proof}
 D'après le lemme \ref{lemma-normalization-locally-algebraic} le morphisme $\nu$
 est fini. Par conséquent, $X^\nu$ est propre sur $k$
 d'après Morphismes, lemmes \ref{morphisms-lemma-finite-proper} et 
\ref{morphisms-lemma-composition-proper}.
Le morphisme $\nu$ est projectif d'après
Morphismes, lemme \ref{morphisms-lemma-finite-projective}
 et donc si $X$ est projectif sur $k$, alors $X^\nu$ est projectif sur $k$
 d'après Morphismes, le lemme \ref{morphisms-lemma-composition-projective}. 
\end{proof}

\begin{lemma}
\label{lemma-relative-normalization-finite}
 Soit $k$ un corps. Soit $f : Y \to X$ un morphisme
quasi-compact de schémas localement algébriques sur $k$. Soit $X'$
 la normalisation de $X$ dans $Y$. Si $Y$ est réduit, alors 
$X' \to X$ est fini. 
\end{lemma}

\begin{proof}
 Puisque $Y$ est quasi-séparé (par
Propriétés, lemme \ref{properties-lemma-locally-Noetherian-quasi-separated}
 et Morphismes, lemme \ref{morphisms-lemma-finite-type-noetherian}),
le morphisme $f$ est quasi-séparé
(Schémas, lemme \ref{schemes-lemma-compose-after-separated}).
Ainsi, Morphismes, définition \ref{morphisms-definition-normalization-X-in-Y}
 s'applique. Le résultat découle de Morphismes, lemme 
\ref{morphisms-lemma-nagata-normalization-finite-general}.
Cela utilise que les schémas localement algébriques sont localement noethériens
(et ont donc localement un nombre fini de composantes irréductibles)
et que les schémas localement algébriques sont de
Nagata (Morphismes, lemme \ref{morphisms-lemma-ubiquity-nagata}).
Quelques petits détails sont omis. 
\end{proof}

\begin{lemma}
\label{lemma-finite-extension-geometrically-normal}
 Soit $k$ un corps. Soit $X$ un $k$-schéma
algébrique. Alors il existe une extension finie purement inséparable $k'/k$
 telle que la normalisation $Y$ de $X_{k'}$ soit géométriquement normale sur $k'$. 
\end{lemma}

\begin{proof}
 Soit $K = k^{perf}$ la clôture parfaite. Soit $Y_K$
la normalisation de $X_K$, voir le lemme \ref{lemma-normalization-locally-algebraic}.
D'après Limites, lemme \ref{limits-lemma-descend-finite-presentation},
il existe une sous-extension finie $K/k'/k$ et un morphisme 
$\nu : Y \to X_{k'}$ de présentation finie dont le changement de base vers $K$
 est le morphisme de normalisation $\nu_K : Y_K \to X_K$.
On observe que $Y$ est géométriquement normal sur $k'$
 (lemme \ref{lemma-geometrically-normal}).
Après avoir augmenté $k'$, on peut supposer que $Y \to X_{k'}$ est fini
(Limites, lemme \ref{limits-lemma-descend-finite-finite-presentation}).
Comme $\nu_K : Y_K \to X_K$ est le morphisme de normalisation,
il induit un morphisme birationnel $Y_K \to (X_K)_{red}$.
Par conséquent, il existe un ouvert dense $V_K \subset X_K$ tel que 
$\nu_K^{-1}(V_K) \to V_K$ soit une immersion
fermée (induisant un isomorphisme de $\nu_K^{-1}(V_K)$ avec $V_{K, red}$,
voir par exemple Morphismes, lemme 
\ref{morphisms-lemma-birational-isomorphism-over-dense-open}).
Après avoir augmenté $k'$, nous trouvons que $V_K$ est le changement de base d'un ouvert dense 
$V \subset Y$ et que le morphisme $\nu^{-1}(V) \to V$ est une immersion
fermée (Limites, lemmes \ref{limits-lemma-descend-opens} et 
\ref{limits-lemma-descend-closed-immersion-finite-presentation}).
Il s'ensuit facilement de cela que $\nu$ est le morphisme de
normalisation et la preuve est complète. 
\end{proof}

\begin{lemma}
\label{lemma-normalization-and-change-of-fields}
 Soit $k$ un corps. Soit $X$ un $k$-schéma localement algébrique.
Soit $K/k$ une extension de corps. Soit $\nu : X^\nu \to X$
 la normalisation de $X$, et soit $Y^\nu \to X_K$ la
normalisation du changement de base. Alors le morphisme canonique 
$$
Y^\nu \longrightarrow X^\nu \times_{\Spec(k)} \Spec(K)
$$
 est un isomorphisme si $K/k$ est séparable et un homéomorphisme universel
en général. 
\end{lemma}

\begin{proof}
 Posons $Y = X_K$. Soient $X^{(0)}$, resp.\ $Y^{(0)}$, les ensembles des points
génériques des composantes irréductibles de $X$, resp. \ $Y$. Alors le morphisme de projection 
$\pi : Y \to X$ satisfait $\pi(Y^{(0)}) = X^{(0)}$. Cela est vrai car 
$\pi$ est surjectif, ouvert et générisant,
voir Morphismes, lemme \ref{morphisms-lemma-scheme-over-field-universally-open} et 
\ref{morphisms-lemma-open-generizing}.
Si nous considérons $X^{(0)}$, resp. \ $Y^{(0)}$ comme des schémas (réduits), alors 
$X^\nu$, resp. \ $Y^\nu$ est la normalisation de $X$, resp. \ $Y$ dans 
$X^{(0)}$, resp. \ $Y^{(0})$.
Ainsi, Morphismes, lemme \ref{morphisms-lemma-functoriality-normalization}
 donne un morphisme canonique $Y^\nu \to X^\nu$ sur $Y \to X$ qui
à son tour donne le morphisme canonique du lemme par la
propriété universelle du produit fibré.

\medskip\noindent
 Pour montrer que ce morphisme possède les propriétés énoncées dans le lemme, on peut
supposer que $X = \Spec(A)$ est affine. Soit $Q(A_{red})$
l'anneau total des fractions de $A_{red}$. Alors $X^\nu$ est le spectre
de la clôture intégrale $A'$ de $A$ dans $Q(A_{red})$,
voir Morphismes, lemmes \ref{morphisms-lemma-normalization-reduced} et 
\ref{morphisms-lemma-description-normalization}.
De même, $Y^\nu$ est le spectre de la clôture intégrale $B'$ de 
$A \otimes_k K$ dans $Q((A \otimes_k K)_{red})$. Il existe un homomorphisme
canonique $Q(A_{red}) \to Q((A \otimes_k K)_{red})$, un homomorphisme canonique 
$A' \to B'$, et le morphisme du lemme correspond à l'homomorphisme induit. 
$$
A' \otimes_k K \longrightarrow B'
$$
 de $K$-algèbres. Le noyau est constitué d'éléments nilpotents
car le noyau de $Q(A_{red}) \otimes_k K \to Q((A \otimes_k K)_{red})$
 est l'ensemble des éléments nilpotents.

\medskip\noindent
 Si $K/k$ est séparable, alors $A' \otimes_k K$ est normal d'après
le lemme \ref{lemma-base-change-normal-by-separable}. En particulier,
il est réduit, donc $Q((A \otimes_k K)_{red}) = Q(A' \otimes_k K)$
 et $B' = A' \otimes_k K$
 d'après Algèbre, lemme \ref{algebra-lemma-characterize-reduced-ring-normal}.

\medskip\noindent
 Supposons que $K/k$ ne soit pas séparable. Alors la caractéristique de $k$ est $p > 0$.
Nous allons montrer que pour chaque $b \in B'$ il existe une puissance $q$ de $p$
 telle que $b^q$ appartienne à l'image de $A' \otimes_k K$. Cela montrera
que l'application affichée est un homéomorphisme universel d'après le
lemme d'Algèbre \ref{algebra-lemma-p-ring-map}.
Pour $b$ fixé, il existe un sous-corps $F \subset K$
 avec $F/k$ de type fini tel que $b$ soit contenu dans 
$Q((A \otimes_k F)_{red})$ et soit entier sur $A \otimes_k F$.
Choisissons un polynôme unitaire $P = T^d + \alpha_1 T^{d - 1} + \ldots + \alpha _d$
 avec $P(b) = 0$ et $\alpha_i \in A \otimes_k F$.
Choisissons une base de transcendance $t_1, \ldots, t_r$ de $F$ sur $k$.
Soit $F/F'/k(t_1, \ldots, t_r)$ la sous-extension
séparable maximale (Corps, lemme \ref{fields-lemma-separable-first}).
Comme $F/F'$ est finie purement inséparable, il existe un $q$ tel
que $\lambda^q \in F'$ pour tout $\lambda \in F$. Alors $b^q$ est
dans $Q((A \otimes_k F')_{red})$ et satisfait le polynôme 
$T^d + \alpha_1^q T^{d - 1} + \ldots + \alpha _d^q$ avec 
$\alpha_i^q \in A \otimes_k F'$. Par le cas séparable,
nous voyons que $b^q \in A' \otimes_k F'$ et la démonstration est complète. 
\end{proof}

\begin{lemma}
\label{lemma-geometrically-normal-in-codim-1}
 Soit $k$ un corps. Soit $X$ un $k$-schéma localement algébrique. Soit
$\nu : X^\nu \to X$ la normalisation de $X$.
Soit $x \in X$ un point tel que (a) $\mathcal{O}_{X, x}$
 soit réduit, (b) $\dim(\mathcal{O}_{X, x}) = 1$, et
(c) pour tout $x' \in X^\nu$ avec $\nu(x') = x$ l'extension 
$\kappa(x')/k$ soit séparable. Alors $X$ est géométriquement réduit en $x$
 et $X^\nu$ est géométriquement régulier en $x'$ avec $\nu(x') = x$. 
\end{lemma}

\begin{proof}
 Nous utiliserons les résultats du lemme \ref{lemma-normalization-locally-algebraic}
 sans autre mention. Soit $x' \in X^\nu$ un point au-dessus de $x$.
Par la théorie de la dimension (section \ref{section-algebraic-schemes}), nous avons 
$\dim(\mathcal{O}_{X^\nu, x'}) = 1$. Puisque $X^\nu$ est normal, on
voit que $\mathcal{O}_{X^\nu, x'}$ est un anneau de valuation
discrète (Propriétés, lemme \ref{properties-lemma-criterion-normal}).
Ainsi, $\mathcal{O}_{X^\nu, x'}$ est une $k$-algèbre locale
régulière dont le corps résiduel est séparable sur $k$. Par conséquent, 
$k \to \mathcal{O}_{X^\nu, x'}$ est formellement lisse dans la
topologie $\mathfrak m_{x'}$-adique, voir
Compléments d'algèbre, lemme \ref{more-algebra-lemma-regular-implies-fs}.
Alors $\mathcal{O}_{X^\nu, x'}$ est géométriquement régulier
sur $k$
 par Compléments d'algèbre, théorème \ref{more-algebra-theorem-regular-fs}.
Ainsi $X^\nu$ est géométriquement régulier en $x'$ d'après
le lemme \ref{lemma-geometrically-regular-at-point}.

\medskip\noindent
 Puisque $\mathcal{O}_{X, x}$ est réduit, la famille d'homomorphismes 
$\mathcal{O}_{X, x} \to \mathcal{O}_{X^\nu, x'}$ est injective.
Puisque $\mathcal{O}_{X^\nu, x'}$ est une 
$k$-algèbre géométriquement réduite, il s'ensuit immédiatement que $\mathcal{O}_{X, x}$
 est une $k$-algèbre géométriquement réduite. Par conséquent, $X$ est géométriquement
réduit en $x$ d'après le lemme \ref{lemma-geometrically-reduced-at-point}. 
\end{proof}









\section{Groupes de fonctions inversibles} 
\label{section-units}

\noindent
 Il arrive souvent (mais pas toujours) que $\mathcal{O}^*(X)/k^*$
 soit un groupe abélien de type fini si $X$ est une variété sur $k$.
Nous le montrons par une série
de lemmes. Tout repose sur le cas spécial suivant.

\begin{lemma}
\label{lemma-open-in-normal-proper}
 Soit $k$ un corps algébriquement clos. Soit
$\overline{X}$ une variété propre sur $k$.
Soit $X \subset \overline{X}$ un sous-schéma ouvert.
Supposons que $X$ soit normal.
Alors $\mathcal{O}^*(X)/k^*$ est un groupe abélien de type fini. 
\end{lemma}

\begin{proof}
 Puisque l'énoncé ne concerne que $X$, nous pouvons remplacer $\overline{X}$
 par une autre variété propre sur $k$. Soit $\nu : \overline{X}^\nu \to
\overline{X}$ le morphisme de normalisation. D'après le
lemme \ref{lemma-normalization-locally-algebraic},
on sait que $\nu$ est fini et que $\overline{X}^\nu$ est
une variété. Comme $X$ est normale, on voit que $\nu^{-1}(X) \to X$ est
un isomorphisme (petit détail omis). Enfin, on voit que 
$\overline{X}^\nu$ est propre sur $k$, car un morphisme fini
est propre (Morphismes, lemme \ref{morphisms-lemma-finite-proper})
et les compositions de morphismes
propres sont propres (Morphismes, lemme \ref{morphisms-lemma-composition-proper}).
Ainsi, nous pouvons et faisons l'hypothèse que $\overline{X}$ est normal.

\medskip\noindent
 Nous utiliserons sans plus le rappeler que, pour tout ouvert affine $U$ de 
$\overline{X}$, l'anneau $\mathcal{O}(U)$ est une 
$k$-algèbre de type fini, qui est noethérienne, intègre et normale,
voir Algèbre, lemme \ref{algebra-lemma-Noetherian-permanence},
Propriétés, définition \ref{properties-definition-integral},
Propriétés, lemmes \ref{properties-lemma-locally-Noetherian} et 
\ref{properties-lemma-locally-normal},
Morphismes, lemme \ref{morphisms-lemma-locally-finite-type-characterize}.

\medskip\noindent
 Soient $\xi_1, \ldots, \xi_r$ les points génériques du complément de $X$
 dans $\overline{X}$. Ils sont en nombre fini puisque $\overline{X}$ possède
un espace topologique sous-jacent noethérien (voir
Morphismes, lemme \ref{morphisms-lemma-finite-type-noetherian},
Propriétés, lemme \ref{properties-lemma-Noetherian-topology},
et Topologie, lemme \ref{topology-lemma-Noetherian}).
Pour chaque $i$, l'anneau local $\mathcal{O}_i = \mathcal{O}_{X, \xi_i}$
 est un anneau intègre local normal noethérien (comme une localisation d’un
anneau intègre normal noethérien). Soit $J \subset \{1, \ldots, r\}$ l'ensemble
des indices $i$ tels que $\dim(\mathcal{O}_i) = 1$. Pour $j \in J$,
l'anneau local $\mathcal{O}_j$ est un anneau de valuation discrète,
voir Algèbre, lemme \ref{algebra-lemma-characterize-dvr}.
Par conséquent, nous obtenons une valuation. 
$$
v_j : k(\overline{X})^* \longrightarrow \mathbf{Z}
$$
 avec la propriété que $v_j(f) \geq 0 \Leftrightarrow f \in \mathcal{O}_j$.

\medskip\noindent
 Identifions $\mathcal{O}(X)$ à une sous-$k$-algèbre de $k(X) = k(\overline{X})$.
Nous affirmons que le noyau de l'application 
$$
\mathcal{O}(X)^* \longrightarrow
\prod\nolimits_{j \in J} \mathbf{Z},
\quad
f \longmapsto \prod v_j(f)
$$
 est $k^*$. Il est clair que cette affirmation prouve le lemme.
 En effet, soit $f \in \mathcal{O}(X)$ un élément du noyau. Soit
$U = \Spec(B) \subset \overline{X}$ un ouvert affine quelconque.
Alors $B$ est un anneau intègre normal noethérien. Pour
chaque idéal premier de hauteur un $\mathfrak q \subset B$ avec le point
correspondant $\xi \in X$, on a ou bien $\xi = \xi_j$ pour un certain $j \in J$,
ou bien $\xi \in X$. En effet, 
$\text{codim}(\overline{\{\xi\}}, \overline{X}) = 1$ d'après
Propriétés, lemme \ref{properties-lemma-codimension-local-ring},
 et donc si $\xi \in \overline{X} \setminus X$, cela doit
être un point générique de $\overline{X} \setminus X$, donc égal à un certain 
$\xi_j$, $j \in J$.
Nous concluons que $f \in \mathcal{O}_{X, \xi} = B_{\mathfrak q}$
 dans les deux cas puisque $f$ est dans le noyau de l'application. Ainsi 
$f \in \bigcap_{\text{ht}(\mathfrak q) = 1} B_{\mathfrak q} = B$,
voir Algèbre, lemme 
\ref{algebra-lemma-normal-domain-intersection-localizations-height-1}.
En d'autres termes, on voit que 
$f \in \Gamma(\overline{X}, \mathcal{O}_{\overline{X}})$.
Mais puisque $k$ est algébriquement clos, nous concluons que 
$f \in k$ d'après
le lemme \ref{lemma-regular-functions-proper-variety}. 
\end{proof}

\noindent
 Nous généralisons maintenant le cas précédent par quelques arguments élémentaires, en
supposant toujours le corps algébriquement clos.

\begin{lemma}
\label{lemma-units-integral-finite-type-algebraically-closed}
 Soit $k$ un corps algébriquement clos. Soit
$X$ un schéma intègre localement de type fini sur $k$.
Alors $\mathcal{O}^*(X)/k^*$ est un groupe abélien de type fini. 
\end{lemma}

\begin{proof}
 Comme $X$ est intègre, l'homomorphisme de restriction 
$\mathcal{O}(X) \to \mathcal{O}(U)$ est injectif pour tout
sous-schéma ouvert non vide $U \subset X$. Par conséquent, on peut supposer
que $X$ est affine. Choisissons une immersion fermée 
$X \to \mathbf{A}^n_k$
 et notons $\overline{X}$ l'adhérence de $X$ dans $\mathbf{P}^n_k$
 via l'immersion usuelle $\mathbf{A}^n_k \to \mathbf{P}^n_k$.
Ainsi, on peut supposer que $X$ est un ouvert affine d'une variété
projective $\overline{X}$.

\medskip\noindent
 Soit $\nu : \overline{X}^\nu \to \overline{X}$ le morphisme
de normalisation,
voir Morphismes, définition \ref{morphisms-definition-normalization}.
Nous savons que $\nu$ est fini, dominant, et que $\overline{X}^\nu$
 est un schéma irréductible normal,
voir Morphismes, lemmes \ref{morphisms-lemma-normalization-normal}, 
\ref{morphisms-lemma-Japanese-normalization}, et 
\ref{morphisms-lemma-ubiquity-nagata}.
Il s'ensuit que $\overline{X}^\nu$ est une variété propre,
parce que $\overline{X} \to \Spec(k)$ est propre comme composition d'un
morphisme fini et d'un morphisme propre (voir les
résultats dans Morphismes, sections \ref{morphisms-section-proper} et 
\ref{morphisms-section-integral}).
Il s'ensuit également que $\nu$ est un morphisme surjectif, parce que
l'image de $\nu$ est fermée et contient le point générique de $\overline{X}$.
Ainsi, en posant $X^\nu = \nu^{-1}(X)$, nous voyons qu'il suffit de prouver le
résultat pour $X^\nu$. En d'autres termes, on peut supposer que $X$ est un ouvert
non vide d'une variété propre normale $\overline{X}$. Ce cas est traité par
le lemme \ref{lemma-open-in-normal-proper}. 
\end{proof}

\noindent
 Le lemme précédent implique la légère généralisation suivante.

\begin{lemma}
\label{lemma-units-general-algebraically-closed}
 Soit $k$ un corps algébriquement clos. Soit
$X$ un schéma réduit connexe localement de type fini
sur $k$ avec un nombre fini de composantes irréductibles.
Alors $\mathcal{O}^*(X)/k^*$ est un groupe abélien de type fini. 
\end{lemma}

\begin{proof}
 Écrivons $X = \bigcup X_i$ pour la décomposition en composantes irréductibles. Par le
lemme \ref{lemma-units-integral-finite-type-algebraically-closed},
nous voyons que $\mathcal{O}(X_i)^*/k^*$ est un groupe abélien de
type fini. Soit $f \in \mathcal{O}(X)^*$ dans le noyau
de l'application 
$$
\mathcal{O}(X)^* \longrightarrow \prod \mathcal{O}(X_i)^*/k^*.
$$
 Alors pour chaque $i$ il existe un élément $\lambda_i \in k$
 tel que $f|_{X_i} = \lambda_i$.
En se restreignant à $X_i \cap X_j$, nous concluons que 
$\lambda_i = \lambda_j$ si $X_i \cap X_j \not = \emptyset$.
Puisque $X$ est connexe, nous concluons que tous les $\lambda_i$ sont
égaux et donc que $f \in k^*$. Cela prouve que 
$$
\mathcal{O}(X)^*/k^* \subset \prod \mathcal{O}(X_i)^*/k^*
$$
 et le lemme en découle, car le membre de droite est un produit d'un
nombre fini de groupes abéliens de type fini. 
\end{proof}

\begin{lemma}
\label{lemma-integral-closure-ground-field}
 Soit $k$ un corps. Soit
$X$ un schéma sur $k$, connexe et réduit.
Alors la clôture intégrale de $k$ dans $\Gamma(X, \mathcal{O}_X)$
 est un corps. 
\end{lemma}

\begin{proof}
 Soit $k' \subset \Gamma(X, \mathcal{O}_X)$ la clôture intégrale de 
$k$. Alors $X \to \Spec(k)$ se factorise à travers $\Spec(k')$,
voir Schémas, lemme \ref{schemes-lemma-morphism-into-affine}.
Comme $X$ est réduit, on voit que $k'$ n'a pas d'éléments nilpotents non nuls. Comme 
$k \to k'$ est intègre, on voit que tout idéal premier de $k'$ est
à la fois un idéal maximal et un premier minimal,
et $\Spec(k')$ est totalement discontinu,
voir Algèbre, lemmes \ref{algebra-lemma-integral-no-inclusion} et 
\ref{algebra-lemma-ring-with-only-minimal-primes}.
Comme $X$ est connexe, le morphisme $X \to \Spec(k')$
 est constant, d'image le point correspondant à 
$\mathfrak p \subset k'$. Alors tout $f \in k'$, $f \not \in \mathfrak p$
 s'envoie sur un élément inversible de $\mathcal{O}_X$. Par définition
de $k'$, il s'ensuit que $f$ est une unité de $k'$. Ainsi, on voit
que $k'$ est local avec idéal maximal $\mathfrak p$,
voir Algèbre, lemme \ref{algebra-lemma-characterize-local-ring}.
Puisque nous avons déjà vu que $k'$ est réduit, cela implique que 
$k'$ est un corps,
voir Algèbre, lemme \ref{algebra-lemma-minimal-prime-reduced-ring}. 
\end{proof}

\begin{proposition}
\label{proposition-units-general}
 Soit $k$ un corps. Soit $X$ un schéma sur $k$. Supposons que $X$
 soit localement de type fini sur $k$, connexe et réduit, et possède un nombre
fini de composantes irréductibles. Alors $\mathcal{O}(X)^*/k^*$ est un groupe abélien
de type fini si, en plus des conditions ci-dessus, au
moins l'une des conditions suivantes est satisfaite : 
\begin{enumerate}
\item la clôture intégrale de $k$ dans $\Gamma(X, \mathcal{O}_X)$ est $k$, 
\item $X$ a un point $k$-rationnel, ou 
\item $X$ est géométriquement intègre. 
\end{enumerate}
\end{proposition}

\begin{proof}
 Soit $\overline{k}$ une clôture algébrique de $k$.
Soit $Y$ une composante connexe de $(X_{\overline{k}})_{red}$.
Notons que le morphisme canonique $p : Y \to X$ est ouvert
(d'après Morphismes, lemme \ref{morphisms-lemma-scheme-over-field-universally-open})
et fermé
(d'après Morphismes, lemme \ref{morphisms-lemma-integral-universally-closed}).
Ainsi, $p(Y) = X$, puisque $X$ est connexe. En particulier, comme 
$X$ est réduit, cela implique $\mathcal{O}(X) \subset \mathcal{O}(Y)$. D'après
le lemme \ref{lemma-galois-action-irreducible-components-locally-finite-type},
nous voyons que $Y$ possède un nombre fini de composantes irréductibles (le corps 
$\overline{k}$ dans le lemme est la clôture algébrique séparable, mais cela
n'a pas d'importance car les deux changements de base $X_{\overline{k}}$
 ont les mêmes composantes irréductibles d'après le
lemme \ref{lemma-separably-closed-field-irreducible-components}). Ainsi,
le lemme \ref{lemma-units-general-algebraically-closed}
 s'applique à $Y$. Cela implique que si $\mathcal{O}(X)^*/k^*$
 n'est pas un groupe abélien de type fini, alors il existe un élément 
$f \in \mathcal{O}(X)$, $f \not \in k$, qui s'envoie sur un élément de 
$\overline{k}$ via l'application $\mathcal{O}(X) \to \mathcal{O}(Y)$.
Dans ce cas, $f$ est algébrique sur $k$,
donc intégral sur $k$. Ainsi, si la condition (1) est vérifiée, cela ne peut pas se
produire. Pour terminer la démonstration, nous montrons que les conditions (2) et (3) impliquent (1).

\medskip\noindent
 Soit $k \subset k' \subset \Gamma(X, \mathcal{O}_X)$ la clôture
intégrale de $k$ dans $\Gamma(X, \mathcal{O}_X)$. D'après
le lemme \ref{lemma-integral-closure-ground-field},
on voit que $k'$ est un corps.
Si $e : \Spec(k) \to X$ est un point $k$-rationnel, alors 
$e^\sharp : \Gamma(X, \mathcal{O}_X) \to k$ est une section
de l'application d'inclusion $k \to \Gamma(X, \mathcal{O}_X)$. En particulier,
la restriction de $e^\sharp$ à $k'$ est un homomorphisme de corps $k' \to k$ sur $k$,
ce qui montre clairement que (2) implique (1).

\medskip\noindent
 Si la clôture intégrale $k'$ de $k$ dans $\Gamma(X, \mathcal{O}_X)$
 n'est pas triviale, alors on voit que $X$ n'est pas géométriquement connexe 
(si $k'/k$ n'est pas purement inséparable), ou bien que $X$ n'est pas
géométriquement réduit (si $k'/k$ est purement inséparable non triviale). Nous omettons
les détails. Par conséquent, (3) implique (1). 
\end{proof}

\begin{lemma}
\label{lemma-units-variety}
 Soit $k$ un corps. Soit
$X$ une variété sur $k$.
Le groupe $\mathcal{O}(X)^*/k^*$ est un groupe abélien de type fini
à condition qu'au moins une des conditions suivantes soit remplie : 
\begin{enumerate}
\item $k$ est intégralement clos dans $\Gamma(X, \mathcal{O}_X)$, 
\item $k$ est algébriquement clos dans $k(X)$, 
\item $X$ est géométriquement intègre sur $k$, ou 
\item $k$ est l'intersection `` '' des extensions de corps 
$\kappa(x)/k$ où $x$ parcourt les points fermés de $x$. 
\end{enumerate}
\end{lemma}

\begin{proof}
 On voit que (1) suffit
d'après la proposition \ref{proposition-units-general}.
Nous omettons la vérification que chacun de (2), (3), (4) implique (1). 
\end{proof}







\section{Formule de K\"unneth, I} 
\label{section-kunneth}

\noindent
 Dans cette section, nous démontrons la formule de K\"unneth lorsque
la base est un corps et que nous considérons la cohomologie
des modules quasi-cohérents. Pour une version plus générale, voir
Catégories dérivées de schémas, section \ref{perfect-section-kunneth}.

\begin{lemma}
\label{lemma-kunneth}
 Soit $k$ un corps. Soient $X$ et $Y$ des schémas sur $k$, et
soient $\mathcal{F}$, resp.\ $\mathcal{G}$, un module quasi-cohérent sur 
$\mathcal{O}_X$, resp.\ sur $\mathcal{O}_Y$.
Alors nous avons un isomorphisme canonique 
$$
H^n(X \times_{\Spec(k)} Y, \text{pr}_1^*\mathcal{F}
\otimes_{\mathcal{O}_{X \times_{\Spec(k)} Y}} \text{pr}_2^*\mathcal{G}) =
\bigoplus\nolimits_{p + q = n}
H^p(X, \mathcal{F}) \otimes_k H^q(Y, \mathcal{G})
$$
 si $X$ et $Y$ sont quasi-compacts et ont une
diagonale affine\footnote{Le cas où $X$ et $Y$ sont quasi-séparés sera
traité au lemme \ref{lemma-kunneth-general} ci-dessous.}
(par exemple si $X$ et $Y$ sont séparés). 
\end{lemma}

\begin{proof}
 Dans cette démonstration, les produits non décorés et les produits tensoriels sont pris sur $k$. Pour construire les applications 
$$
H^p(X, \mathcal{F}) \otimes H^q(Y, \mathcal{G})
\longrightarrow
H^n(X \times Y, \text{pr}_1^*\mathcal{F}
\otimes_{\mathcal{O}_{X \times Y}} \text{pr}_2^*\mathcal{G})
$$
 on utilise la fonctorialité de la cohomologie pour obtenir les homomorphismes 
$H^p(X, \mathcal{F}) \to H^p(X \times Y, \text{pr}_1^*\mathcal{F})$ et 
$H^p(Y, \mathcal{G}) \to H^p(X \times Y, \text{pr}_2^*\mathcal{G})$,
puis le cup-produit 
$$
\cup :
H^p(X \times Y, \text{pr}_1^*\mathcal{F})
\otimes H^q(X \times Y, \text{pr}_2^*\mathcal{G})
\longrightarrow
H^n(X \times Y, \text{pr}_1^*\mathcal{F}
\otimes_{\mathcal{O}_{X \times Y}} \text{pr}_2^*\mathcal{G})
$$
 Le résultat est vrai lorsque $X$ et $Y$ sont affines en raison de la
nullité des groupes de cohomologie supérieure sur les affines (Cohomologie des Schémas, lemme 
\ref{coherent-lemma-quasi-coherent-affine-cohomology-zero})
et des définitions (des images réciproques de modules quasi-cohérents et
des produits tensoriels de modules quasi-cohérents).

\medskip\noindent
 Choisissons des recouvrements affines ouverts finis 
$\mathcal{U} : X = \bigcup_{i \in I} U_i$ et 
$\mathcal{V} : Y = \bigcup_{j \in J} V_j$.
Cela détermine un recouvrement affine ouvert 
$\mathcal{W} : X \times Y = \bigcup_{(i, j) \in I \times J} U_i \times V_j$.
Notons que $\mathcal{W}$ est un raffinement de 
$\text{pr}_1^{-1}\mathcal{U}$ et de $\text{pr}_2^{-1}\mathcal{V}$.
 Ainsi, d'après Cohomologie, lemme \ref{cohomology-lemma-functoriality-cech},
on obtient des homomorphismes 
$$
\check{\mathcal{C}}^\bullet(\mathcal{U}, \mathcal{F}) \to
\check{\mathcal{C}}^\bullet(\mathcal{W}, \text{pr}_1^*\mathcal{F})
\quad\text{et}\quad
\check{\mathcal{C}}^\bullet(\mathcal{V}, \mathcal{G}) \to
\check{\mathcal{C}}^\bullet(\mathcal{W}, \text{pr}_2^*\mathcal{G})
$$
 compatibles avec les homomorphismes d'image réciproque en cohomologie. Dans Cohomologie, équation ( 
\ref{cohomology-equation-needs-signs})
nous avons construit une application de complexes 
$$
\text{Tot}(
\check{\mathcal{C}}^\bullet(\mathcal{W}, \text{pr}_1^*\mathcal{F})
\otimes
\check{\mathcal{C}}^\bullet(\mathcal{W}, \text{pr}_2^*\mathcal{G}))
\longrightarrow
\check{\mathcal{C}}^\bullet(\mathcal{W},
\text{pr}_1^*\mathcal{F} \otimes_{\mathcal{O}_{X \times Y}}
\text{pr}_2^*\mathcal{G})
$$
 qui définit le cup-produit en cohomologie. En combinant ce qui précède, nous
obtenons une application de complexes 
\begin{equation}
\label{equation-kunneth-on-cech}
\text{Tot}(
\check{\mathcal{C}}^\bullet(\mathcal{U}, \mathcal{F})
\otimes
\check{\mathcal{C}}^\bullet(\mathcal{V}, \mathcal{G}))
\longrightarrow
\check{\mathcal{C}}^\bullet(\mathcal{W},
\text{pr}_1^*\mathcal{F} \otimes_{\mathcal{O}_{X \times Y}}
\text{pr}_2^*\mathcal{G})
\end{equation}
 Cette application n'est pas
un isomorphisme de complexes. Rappelons que nous pouvons calculer les
cohomologies de nos faisceaux quasi-cohérents en utilisant nos
recouvrements (Cohomologie des Schémas, lemmes 
\ref{coherent-lemma-affine-diagonal} et 
\ref{coherent-lemma-cech-cohomology-quasi-coherent}).
Ainsi, en cohomologie, (\ref{equation-kunneth-on-cech}) reproduit l'application
du lemme.

\medskip\noindent
 Considérons une suite exacte courte
$0 \to \mathcal{F} \to \mathcal{F}' \to \mathcal{F}'' \to 0$
 de modules quasi-cohérents. Puisque la construction de ( 
\ref{equation-kunneth-on-cech} ) est fonctorielle en $\mathcal{F}$ et puisque
la formation des complexes de {\v C}ech concernés est exacte
dans la variable $\mathcal{F}$ (car nous prenons des sections sur des
ouverts affines), on obtient un morphisme entre deux suites exactes courtes
de complexes. 
$$
\xymatrix{
\text{Tot}(
\check{\mathcal{C}}^\bullet(\mathcal{U}, \mathcal{F})
\otimes
\check{\mathcal{C}}^\bullet(\mathcal{V}, \mathcal{G})) \ar[r] \ar[d] &
\text{Tot}(
\check{\mathcal{C}}^\bullet(\mathcal{U}, \mathcal{F}')
\otimes
\check{\mathcal{C}}^\bullet(\mathcal{V}, \mathcal{G})) \ar[r] \ar[d] &
\text{Tot}(
\check{\mathcal{C}}^\bullet(\mathcal{U}, \mathcal{F}'')
\otimes
\check{\mathcal{C}}^\bullet(\mathcal{V}, \mathcal{G})) \ar[d] \\
\check{\mathcal{C}}^\bullet(\mathcal{W},
\text{pr}_1^*\mathcal{F} \otimes_{\mathcal{O}_{X \times Y}}
\text{pr}_2^*\mathcal{G}) \ar[r] &
\check{\mathcal{C}}^\bullet(\mathcal{W},
\text{pr}_1^*\mathcal{F}' \otimes_{\mathcal{O}_{X \times Y}}
\text{pr}_2^*\mathcal{G}) \ar[r] &
\check{\mathcal{C}}^\bullet(\mathcal{W},
\text{pr}_1^*\mathcal{F}'' \otimes_{\mathcal{O}_{X \times Y}}
\text{pr}_2^*\mathcal{G})
}
$$
 (les zéros extérieurs ont été omis). En
considérant les suites exactes longues de cohomologie, on constate que, si le résultat
du lemme vaut pour $2$ des $3$ faisceaux 
$\mathcal{F}, \mathcal{F}', \mathcal{F}''$, alors il vaut aussi
pour le troisième.

\medskip\noindent
 Notons que $X$ est de dimension cohomologique finie pour
les modules quasi-cohérents, voir Cohomologie des Schémas, lemme 
\ref{coherent-lemma-vanishing-nr-affines}.
En utilisant l'induction sur 
$d(\mathcal{F}) = \max \{d \mid H^d(X, \mathcal{F}) \not = 0\}$,
nous allons réduire au cas $d(\mathcal{F}) = 0$.
Supposons $d(\mathcal{F}) > 0$.
 D'après Cohomologie des Schémas, lemme 
\ref{coherent-lemma-affine-diagonal-universal-delta-functor},
nous avons
vu qu'il existe un plongement $\mathcal{F} \to \mathcal{F}'$
 tel que $H^p(X, \mathcal{F}') = 0$ pour tout $p \geq 1$.
En posant $\mathcal{F}'' = \Coker(\mathcal{F} \to \mathcal{F}')$,
on voit que $d(\mathcal{F}'') < d(\mathcal{F})$.
Ensuite, nous pouvons appliquer le résultat du paragraphe précédent pour
voir qu'il suffit de prouver le lemme pour $\mathcal{F}'$
 et $\mathcal{F}''$, ce qui établit l'étape d'induction.

\medskip\noindent
 En raisonnant de même pour $\mathcal{G}$, on constate
qu'on peut supposer que $\mathcal{F}$ et $\mathcal{G}$
 n'ont une cohomologie non nulle qu'au degré $0$.
 Soit $V \subset Y$ un ouvert affine. Considérons le recouvrement affine ouvert 
$\mathcal{U}_V : X \times V = \bigcup_{i \in I} U_i \times V$.
Il est immédiat que 
$$
\check{\mathcal{C}}^\bullet(\mathcal{U}, \mathcal{F}) \otimes \mathcal{G}(V) =
\check{\mathcal{C}}^\bullet(\mathcal{U}_V,
\text{pr}_1^*\mathcal{F} \otimes_{\mathcal{O}_{X \times Y}}
\text{pr}_2^*\mathcal{G})
$$
 (égalité des complexes). Nous concluons que 
$$
R\text{pr}_{2, *}(\text{pr}_1^*\mathcal{F} \otimes_{\mathcal{O}_{X \times Y}}
\text{pr}_2^*\mathcal{G})
\cong
\Gamma(X, \mathcal{F}) \otimes_k \mathcal{G} \cong
\bigoplus\nolimits_{\alpha \in A} \mathcal{G}
$$
 sur $Y$. Ici $A$ est une base du $k$-espace vectoriel $\Gamma(X, \mathcal{F})$.
La cohomologie sur $Y$ commute avec les sommes
directes (Cohomologie, lemme \ref{cohomology-lemma-quasi-separated-cohomology-colimit}).
En utilisant la suite spectrale de Leray pour $\text{pr}_2$
 (via Cohomologie, lemme \ref{cohomology-lemma-apply-Leray}),
nous concluons que 
$H^n(X \times Y, \text{pr}_1^*\mathcal{F} \otimes_{\mathcal{O}_{X \times Y}}
\text{pr}_2^*\mathcal{G})$
 est nul pour $n > 0$ et isomorphe à 
$H^0(X, \mathcal{F}) \otimes H^0(Y, \mathcal{G})$ pour $n = 0$.
 Cela achève la démonstration (sauf qu'il faudrait vérifier que
l’isomorphisme est bien donné par le cup-produit en degré $0$ ;
nous omettons la vérification). 
\end{proof}

\begin{lemma}
\label{lemma-kunneth-general}
 Soit $k$ un corps. Soient $X$ et $Y$ des schémas sur $k$, et
soient $\mathcal{F}$, resp.\ $\mathcal{G}$, un module quasi-cohérent sur 
$\mathcal{O}_X$, resp.\ sur $\mathcal{O}_Y$.
Alors nous avons un isomorphisme canonique 
$$
H^n(X \times_{\Spec(k)} Y, \text{pr}_1^*\mathcal{F}
\otimes_{\mathcal{O}_{X \times_{\Spec(k)} Y}} \text{pr}_2^*\mathcal{G}) =
\bigoplus\nolimits_{p + q = n}
H^p(X, \mathcal{F}) \otimes_k H^q(Y, \mathcal{G})
$$
 à condition que $X$ et $Y$ soient quasi-compacts et quasi-séparés. 
\end{lemma}

\begin{proof}
 Si $X$ et $Y$ sont séparés ou, plus généralement, ont une diagonale affine, voir
le lemme \ref{lemma-kunneth} pour une « meilleure » démonstration
(dont l'avantage sur celle-ci est d'identifier les
applications comme des images réciproques suivies de cup-produits). Soient
$X'$, resp.\ $Y'$, les épaississements infinitésimaux de $X$, resp.\ $Y$,
de faisceaux de structure 
$\mathcal{O}_{X'} = \mathcal{O}_X \oplus \mathcal{F}$,
respectivement \ $\mathcal{O}_{Y'} = \mathcal{O}_Y \oplus \mathcal{G}$, où 
$\mathcal{F}$, resp.\ $\mathcal{G}$, est un idéal de carré nul.
Alors 
$$
\mathcal{O}_{X' \times Y'} =
\mathcal{O}_{X \times Y}
\oplus \text{pr}_1^*\mathcal{F}
\oplus \text{pr}_2^*\mathcal{G}
\oplus \text{pr}_1^*\mathcal{F}
\otimes_{\mathcal{O}_{X \times Y}} \text{pr}_2^*\mathcal{G}
$$
 en tant que faisceaux sur $X \times Y$. On voit ainsi qu'il suffit de
prouver que 
$$
H^n(X \times Y, \mathcal{O}_{X \times Y}) =
\bigoplus\nolimits_{p + q = n}
H^p(X, \mathcal{O}_X) \otimes_k H^q(Y, \mathcal{O}_Y)
$$
 pour toute paire de schémas quasi-compacts et quasi-séparés sur $k$.
Certains détails sont omis.

\medskip\noindent
 Pour démontrer cette assertion, nous utilisons la cohomologie et le changement de base sous
la forme de Cohomologie des Schémas, lemme \ref{coherent-lemma-hypercoverings}.
Ce lemme fournit un complexe de $k$-espaces vectoriels borné inférieurement, c'est-à-dire
un complexe $\mathcal{K}^\bullet$ de modules quasi-cohérents sur $\Spec(k)$,
qui calcule universellement la cohomologie de $Y$ sur $\Spec(k)$.
En particulier, on voit que 
$$
R\text{pr}_{1, *}(\mathcal{O}_{X \times Y}) \cong
(X \to \Spec(k))^*\mathcal{K}^\bullet
$$
 dans $D(\mathcal{O}_X)$. À homotopie près, le complexe $\mathcal{K}^\bullet$
 est isomorphe à $\bigoplus_{q \geq 0} H^q(Y, \mathcal{O}_Y)[-q]$ car
cela est vrai pour tout complexe d'espaces vectoriels sur un corps.
Nous concluons que 
$$
R\text{pr}_{1, *}(\mathcal{O}_{X \times Y}) \cong
\bigoplus\nolimits_{q \geq 0}
H^q(Y, \mathcal{O}_Y)[-q] \otimes_k \mathcal{O}_X
$$
 dans $D(\mathcal{O}_X)$. Puis nous avons 
\begin{align*}
R\Gamma(X \times Y, \mathcal{O}_{X \times Y})
& =
R\Gamma(X, R\text{pr}_{1, *}(\mathcal{O}_{X \times Y})) \\
& =
R\Gamma(X, \bigoplus\nolimits_{q \geq 0}
H^q(Y, \mathcal{O}_Y)[-q] \otimes_k \mathcal{O}_X) \\
& =
\bigoplus\nolimits_{q \geq 0}
R\Gamma(X,  H^q(Y, \mathcal{O}_Y) \otimes \mathcal{O}_X)[-q] \\
& =
\bigoplus\nolimits_{q \geq 0}
R\Gamma(X, \mathcal{O}_X) \otimes_k H^q(Y, \mathcal{O}_Y)[-q] \\
& =
\bigoplus\nolimits_{p, q \geq 0}
H^p(X, \mathcal{O}_X)[-p] \otimes_k H^q(Y, \mathcal{O}_Y)[-q]
\end{align*}
 comme voulu. La première égalité résulte de Leray pour $\text{pr}_1$
 (Cohomologie, lemme \ref{cohomology-lemma-before-Leray}).
La deuxième résulte de la décomposition de l'image directe totale donnée ci-dessus. La
troisième vient de ce que la cohomologie commute toujours avec les sommes directes finies (et
la cohomologie de $Y$ s'annule en degré suffisamment grand d'après
Cohomologie des Schémas, lemme 
\ref{coherent-lemma-vanishing-nr-affines-quasi-separated} ).
La quatrième vient de ce que la cohomologie sur $X$ commute avec les sommes
directes infinies d'après Cohomologie, lemme 
\ref{cohomology-lemma-quasi-separated-cohomology-colimit}.
Enfin, la dernière égalité résulte de la remarque faite plus haut sur la catégorie dérivée
d'un corps. 
\end{proof}








\section{Groupes de Picard des variétés} 
\label{section-picard-groups}

\noindent
 Dans cette section, nous recueillons quelques résultats élémentaires sur les groupes de
Picard des variétés algébriques.

\begin{lemma}
\label{lemma-change-rings-pic-pre}
 Soit $A \to B$ un homomorphisme d'anneaux fidèlement plat. Soit $X$ un schéma quasi-compact
et quasi-séparé sur $A$. Soit $\mathcal{L}$ un 
$\mathcal{O}_X$-module inversible dont l'image réciproque sur $X_B$ est triviale. Alors 
$H^0(X, \mathcal{L})$ et $H^0(X, \mathcal{L}^{\otimes -1})$ sont des 
$H^0(X, \mathcal{O}_X)$-modules inversibles,
et l'application de multiplication induit un isomorphisme 
$$
H^0(X, \mathcal{L}) \otimes_{H^0(X, \mathcal{O}_X)}
H^0(X, \mathcal{L}^{\otimes -1}) \longrightarrow
H^0(X, \mathcal{O}_X)
$$
\end{lemma}

\begin{proof}
 Notons $\mathcal{L}_B$ l'image réciproque de $\mathcal{L}$ sur $X_B$.
Choisissons un isomorphisme $\mathcal{L}_B \to \mathcal{O}_{X_B}$.
Posons $R = H^0(X, \mathcal{O}_X)$, $M = H^0(X, \mathcal{L})$ et considérons 
$M$ comme un $R$-module. Pour tout $\mathcal{O}_X$-module quasi-cohérent 
$\mathcal{F}$, dont $\mathcal{F}_B$ désigne le tiré en arrière sur $X_B$, il
existe un isomorphisme canonique 
$H^0(X_B, \mathcal{F}_B) = H^0(X, \mathcal{F}) \otimes_A B$,
voir Cohomologie des Schémas, lemme \ref{coherent-lemma-flat-base-change-cohomology}.
Ainsi, nous avons 
$$
M \otimes_R (R \otimes_A B) =
M \otimes_A B = H^0(X_B, \mathcal{L}_B) \cong
H^0(X_B, \mathcal{O}_{X_B}) = R \otimes_A B
$$
 Puisque $R \to R \otimes_A B$ est fidèlement plat (en tant que changement de
base de l'application fidèlement plate $A \to B$), nous concluons
que $M$ est un $R$-module inversible
d'après Algèbre, proposition \ref{algebra-proposition-ffdescent-finite-projectivity}.
De même, $N = H^0(X, \mathcal{L}^{\otimes -1})$ est un $R$-module inversible.
 Pour l'assertion sur les produits tensoriels, on utilise qu'elle devient vraie après
passage à $X_B$, ainsi que la fidèle platitude de $R \to R \otimes_A B$.
Quelques détails sont omis. 
\end{proof}

\begin{lemma}
\label{lemma-change-rings-pic}
 Soit $A \to B$ un homomorphisme d'anneaux fidèlement plat. Soit
$X$ un schéma sur $A$ tel que 
\begin{enumerate}
\item $X$ est quasi-compact et quasi-séparé, et 
\item $R = H^0(X, \mathcal{O}_X)$ est un anneau semi-local. 
\end{enumerate}
 Alors l'homomorphisme induit par image réciproque $\Pic(X) \to \Pic(X_B)$ est injectif. 
\end{lemma}

\begin{proof}
 Soit $\mathcal{L}$ un $\mathcal{O}_X$-module inversible
dont l'image réciproque $\mathcal{L}'$ sur $X_B$ est triviale.
Posons $M = H^0(X, \mathcal{L})$ et $N = H^0(X, \mathcal{L}^{\otimes - 1})$.
D'après le lemme \ref{lemma-change-rings-pic-pre}, les $R$-modules $M$ et $N$
 sont inversibles. Puisque $R$ est semi-local, $M \cong R$ et $N \cong R$,
voir Algèbre, lemme \ref{algebra-lemma-locally-free-semi-local-free}.
Choisissons des générateurs $s \in M$ et $t \in N$. Alors 
$st \in R = H^0(X, \mathcal{O}_X)$ est une unité d'après la dernière assertion
du lemme \ref{lemma-change-rings-pic-pre}.
Nous concluons que $s$ et $t$ définissent des trivialisations de $\mathcal{L}$ et 
$\mathcal{L}^{\otimes -1}$ sur $X$. 
\end{proof}

\begin{lemma}
\label{lemma-change-fields-pic}
 Soit $k'/k$ une extension de corps. Soit
$X$ un schéma sur $k$ tel que 
\begin{enumerate}
\item $X$ est quasi-compact et quasi-séparé, et 
\item $R = H^0(X, \mathcal{O}_X)$ est semi-local, par exemple, si $\dim_k R < \infty$. 
\end{enumerate}
 Alors l'homomorphisme induit par image réciproque $\Pic(X) \to \Pic(X_{k'})$ est injectif. 
\end{lemma}

\begin{proof}
 Cas particulier du lemme \ref{lemma-change-rings-pic}.
Si $\dim_k R < \infty$, alors 
$R$ est artinien et donc semi-local (Algèbre, lemmes 
\ref{algebra-lemma-finite-dimensional-algebra} et 
\ref{algebra-lemma-artinian-finite-nr-max}). 
\end{proof}

\begin{example}
\label{example-need-condition}
 Le lemme \ref{lemma-change-fields-pic} devient faux sans hypothèse supplémentaire
sur le schéma $X$ sur le corps $k$. Voici un exemple. Soit
$k$ un corps. Soit $t \in \mathbf{P}^1_k$ un point fermé.
Posons $X = \mathbf{P}^1 \setminus \{t\}$. Alors nous avons une surjection 
$$
\mathbf{Z} = \Pic(\mathbf{P}^1_k) \longrightarrow \Pic(X)
$$
 La première égalité résulte du lemme
des Diviseurs \ref{divisors-lemma-Pic-projective-space-UFD},
 et la surjectivité résulte du lemme
l'homomorphisme résulte de Diviseurs, lemme \ref{divisors-lemma-extend-invertible-module}
 (puisque $\mathbf{P}^1_k$ est lisse de dimension $1$ sur $k$ et
donc tous ses anneaux locaux sont des anneaux de valuation discrète). Si 
$\mathcal{L}$ est dans le noyau de l'application affichée,
alors $\mathcal{L} \cong \mathcal{O}_{\mathbf{P}^1_k}(nt)$
 pour un certain $n \in \mathbf{Z}$. Nous laissons au lecteur le soin
de montrer que 
$\mathcal{O}_{\mathbf{P}^1_k}(t) \cong \mathcal{O}_{\mathbf{P}^1_k}(d)$
 où $d = [\kappa(t) : k]$. Par conséquent 
$$
\Pic(X) = \mathbf{Z}/d\mathbf{Z}
$$
 Ainsi, si $t$ n'est pas un point $k$-rationnel, alors $d > 1$ et ce
groupe de Picard est non
nul. En revanche, si l'on étend le corps de base $k$ à une extension
de corps $k'$ telle qu'il existe un $k$-plongement 
$\kappa(t) \to k'$, alors $\mathbf{P}^1_{k'} \setminus X_{k'}$
 possède un point $k'$-rationnel $t'$. Par conséquent, 
$\mathcal{O}_{\mathbf{P}^1_{k'}}(1) = \mathcal{O}_{\mathbf{P}^1_{k'}}(t')$
 sera dans le noyau de l'application $\mathbf{Z} \to \Pic(X_{k'})$
 et il s'ensuivra de la même manière que ci-dessus que 
$\Pic(X_{k'}) = 0$. 
\end{example}

\noindent
 Le lemme suivant affirme que les «
modules inversibles rationnellement équivalents » sont isomorphes sur les variétés normales.

\begin{lemma}
\label{lemma-rational-equivalence-for-Pic}
 Soit $k$ un corps. Soit $X$ une variété normale sur $k$.
Soit $U \subset \mathbf{A}^n_k$ un sous-schéma ouvert possédant des points 
$k$-rationnels $p, q \in U(k)$. Pour chaque module
inversible $\mathcal{L}$ sur $X \times_{\Spec(k)} U$, les restrictions 
$\mathcal{L}|_{X \times p}$ et $\mathcal{L}|_{X \times q}$
 sont isomorphes. 
\end{lemma}

\begin{proof}
 Les fibres de $X \times_{\Spec(k)} U \to X$ sont des sous-schémas ouverts de
l'espace affine de dimension $n$ sur un corps. Par conséquent, ces fibres ont
des groupes de Picard triviaux
d'après Diviseurs, lemme \ref{divisors-lemma-open-subscheme-UFD}.
En appliquant Diviseurs, lemme \ref{divisors-lemma-in-image-pullback},
on voit que $\mathcal{L}$ est l'image réciproque d'un module
inversible $\mathcal{N}$ sur $X$. 
\end{proof}










\section{Unicité du corps de base} 
\label{section-base-field}

\noindent
 La phrase « soit $X$ un schéma sur $k$ » signifie que $X$ est un
schéma muni d'un morphisme $X \to \Spec(k)$. On peut alors
se demander si le corps $k$ est déterminé de manière unique par le schéma $X$.
Bien sûr, ce n'est pas le cas, car par exemple 
$\mathbf{A}^1_{\mathbf{C}}$ que nous considérons habituellement comme un schéma
sur le corps $\mathbf{C}$ des nombres complexes, pourrait également être considéré comme
un schéma sur $\mathbf{Q}$. Mais que se passe-t-il si nous demandons que le morphisme 
$X \to \Spec(k)$ ne se factorise pas sous la forme 
$X \to \Spec(k') \to \Spec(k)$ pour une extension de corps
non triviale $k'/k$ ? En d'autres termes, nous demandons que $k$ soit
d'une certaine manière maximal pour que $X$ soit défini sur $k$.

\medskip\noindent
 Un exemple montrant que cela ne garantit toujours pas l'unicité de $k$
 est le schéma 
$$
X =
\Spec\left(
\mathbf{Q}(x)[y]\left[\frac{1}{P(y)}, P \in \mathbf{Q}[y], P \not = 0\right]
\right)
$$
 À première vue, cela semble être un schéma sur $\mathbf{Q}(x)$, mais
à y regarder de plus près, il est clair que c'est aussi un schéma sur $\mathbf{Q}(y)$.
De plus, les corps $\mathbf{Q}(x)$ et $\mathbf{Q}(y)$ sont des sous-corps
de $R = \Gamma(X, \mathcal{O}_X)$ qui sont maximaux parmi les sous-corps de 
$R$ (détails omis). En particulier, $\mathbf{Q}(x)$ et 
$\mathbf{Q}(y)$ sont tous deux maximaux au sens ci-dessus. Remarquons que les deux morphismes 
$X \to \Spec(\mathbf{Q}(x))$
 et $X \to \Spec(\mathbf{Q}(y))$ sont « essentiellement de type fini
» (c'est-à-dire que l'homomorphisme d'anneaux correspondant est essentiellement de type fini).
Ainsi $X$ est un schéma noethérien de dimension finie, c'est-à-dire
qu'il n'est pas complètement pathologique.

\medskip\noindent
 Un autre problème qui peut empêcher l'unicité est que le schéma $X$ peut
ne pas être réduit. Dans ce cas, il peut y avoir de nombreux morphismes différents
de $X$ vers le spectre d'un corps donné. À titre d'exemple explicite, considérons
l'algèbre des nombres duaux 
$D = \mathbf{C}[y]/(y^2) = \mathbf{C} \oplus \epsilon \mathbf{C}$.
Étant donnée toute dérivation $\theta : \mathbf{C} \to \mathbf{C}$ sur $\mathbf{Q}$,
nous obtenons un homomorphisme d'anneaux 
$$
\mathbf{C} \longrightarrow D, \quad
c \longmapsto c + \epsilon \theta(c).
$$
 Le sous-corps de $\mathbf{C}$ sur lequel tous ces homomorphismes sont les
mêmes est la clôture algébrique de $\mathbf{Q}$. Cela signifie que prendre
l'intersection de tous les corps sur lesquels $X$ peut être défini peut
être un très petit corps si $X$ n'est pas réduit.

\medskip\noindent
Une observation à ce sujet est la suivante : étant donné un corps
$k$ et deux sous-corps $k_1, k_2$ de $k$ tels que $k$ soit fini
sur $k_1$ et sur $k_2$, il n'est alors en général {\it pas} vrai
que $k$ soit fini sur $k_1 \cap k_2$. Un exemple est le corps
$k = \mathbf{Q}(t)$ et ses sous-corps $k_1 = \mathbf{Q}(t^2)$ et
$\mathbf{Q}((t + 1)^2)$. En effet, nous avons $k_1 \cap k_2 = \mathbf{Q}$
dans ce cas. Ainsi, dans la suite, il faut être prudent lorsqu'on prend
des intersections de corps.

\medskip\noindent
 Cela dit, nous montrons maintenant que si $X$ est localement de type
fini sur un corps, alors une certaine unicité s'applique. Voici le
résultat précis.

\begin{proposition}
\label{proposition-unique-base-field}
 Soit $X$ un schéma. Soient $a : X \to \Spec(k_1)$ et 
$b : X \to \Spec(k_2)$ des morphismes de $X$ vers des spectres de corps.
Supposons que $a, b$ soient localement de type fini et que 
$X$ soit réduit et connexe. Alors nous avons 
$k_1' = k_2'$, où $k_i' \subset \Gamma(X, \mathcal{O}_X)$
 est la clôture intégrale de $k_i$ dans $\Gamma(X, \mathcal{O}_X)$. 
\end{proposition}

\begin{proof}
 Admettons d'abord le résultat lorsque $X$ est quasi-compact (nous établirons
ce cas ci-dessous). Comme 
$X$ est localement de type fini sur un corps, il est localement noethérien,
voir Morphismes, lemme \ref{morphisms-lemma-finite-type-noetherian}.
En particulier, cela signifie qu'il est localement
connexe, les composantes connexes des ouverts sont
ouvertes, et les intersections d'ouverts quasi-compacts sont quasi-compactes,
voir Propriétés, lemme \ref{properties-lemma-Noetherian-topology},
Topologie, lemme \ref{topology-lemma-locally-connected},
Topologie, section \ref{topology-section-noetherian},
et Topologie, lemme \ref{topology-lemma-constructible-Noetherian-space}.
Choisissons un recouvrement ouvert $X = \bigcup_{i \in I} U_i$
 tel que chaque $U_i$ soit quasi-compact et connexe. Pour
chaque $i$, soit $K_i \subset \mathcal{O}_X(U_i)$ la clôture
intégrale de $k_1$ et de $k_2$.
Pour chaque paire $i, j \in I$ nous décomposons 
$$
U_i \cap U_j = \coprod U_{i, j, l}
$$
 en ses composantes connexes, qui sont en nombre fini. Écrivons 
$K_{i, j, l} \subset \mathcal{O}(U_{i, j, l})$
 pour la clôture intégrale de $k_1$ et de $k_2$. D'après
le lemme \ref{lemma-integral-closure-ground-field},
les anneaux $K_i$ et $K_{i, j, l}$ sont des corps.
Nous affirmons maintenant que $k_1'$ et $k_2'$ sont tous deux égaux au noyau de l'application 
$$
\prod K_i \longrightarrow \prod K_{i, j, l}, \quad
(x_i)_i \longmapsto x_i|_{U_{i, j, l}} - x_j|_{U_{i, j, l}}
$$
ce qui donne le résultat. En effet, il est clair que
$k_1'$ est contenu dans ce noyau. D'autre part, supposons que
$(x_i)_i$ soit dans le noyau. Par la condition de faisceau,
$(x_i)_i$ correspond à $f \in \mathcal{O}(X)$.
Choisissons $i_0 \in I$ et un polynôme unitaire $P(T) \in k_1[T]$
tel que $P(x_{i_0}) = 0$. Nous affirmons alors que $P(f) =
0$, ce qui prouve que $f \in k_1$. Il suffit de
montrer que $P(x_i) = 0$ pour tout $i$. Choisissons $i \in I$.
Comme $X$ est connexe, il existe une suite $i_0, i_1, \ldots, i_n =
i \in I$ telle que $U_{i_t} \cap U_{i_{t + 1}} \not = \emptyset$.
Cela signifie que, pour chaque $t$, il existe un $l_t$
tel que $x_{i_t}$ et $x_{i_{t + 1}}$ se projettent sur le même
élément du corps $K_{i_t, i_{t + 1}, l_t}$. Par conséquent, si
$P(x_{i_t}) = 0$, alors $P(x_{i_{t + 1}}) = 0$. Par induction, en
partant de $P(x_{i_0}) = 0$, nous obtenons $P(x_i) = 0$, comme annoncé.

\medskip\noindent
Pour terminer la démonstration du lemme, nous prouvons
le lemme sous l’hypothèse supplémentaire que $X$ est
quasi-compact. D’après le lemme
\ref{lemma-integral-closure-ground-field}, après avoir
remplacé $k_i$ par $k_i'$, on peut supposer que $k_i$
est intégralement clos dans $\Gamma(X, \mathcal{O}_X)$.
Cela implique que $\mathcal{O}(X)^*/k_i^*$ est un groupe
abélien de type fini, voir la proposition
\ref{proposition-units-general}. Soit $k_{12} = k_1 \cap k_2$ un
sous-anneau de $\mathcal{O}(X)$. Notons que $k_{12}$ est un corps. Puisque
$$
k_1^*/k_{12}^* \longrightarrow \mathcal{O}(X)^*/k_2^*
$$
on voit que $k_1^*/k_{12}^*$ est également un groupe abélien
de type fini. Par conséquent, il existe $\alpha_1, \ldots,
\alpha_n \in k_1^*$ tels que chaque élément $\lambda \in k_1$ ait la forme
$$
\lambda = c \alpha_1^{e_1} \ldots \alpha_n^{e_n}
$$
pour certains $e_i \in \mathbf{Z}$ et $c \in
k_{12}$. En particulier, l'homomorphisme d'anneaux
$$
k_{12}[x_1, \ldots, x_n, \frac{1}{x_1 \ldots x_n}] \longrightarrow k_1, \quad
x_i \longmapsto \alpha_i
$$
est surjectif. D'après le Nullstellensatz de Hilbert, Algèbre,
théorème \ref{algebra-theorem-nullstellensatz}, nous concluons
que $k_1$ est une extension finie de $k_{12}$. De la même
manière, nous concluons que $k_2$ est une extension finie de
$k_{12}$. En particulier, à la fois $k_1$ et $k_2$ sont contenus
dans la clôture intégrale $k_{12}'$ de $k_{12}$ dans $\Gamma(X,
\mathcal{O}_X)$. Mais puisque $k_{12}'$ est un corps selon le
lemme \ref{lemma-integral-closure-ground-field} et que nous avons
choisi $k_i$ pour être intégralement clos dans $\Gamma(X,
\mathcal{O}_X)$, nous concluons que $k_1 = k_{12} = k_2$, comme annoncé.
\end{proof}





\section{Automorphismes}
\label{section-automorphisms}

\noindent
Une section sur les automorphismes des schémas sur les
corps. Pour certaines informations sur les automorphismes
(infinitésimaux) des courbes, voir Courbes Algébriques,
section \ref{curves-section-vector-fields} et Modules de
Courbes, section \ref{moduli-curves-section-finite-aut}.

\begin{lemma}
\label{lemma-automorphism}
Supposons que $X$ soit un schéma de type fini réduit sur un corps $k$.
Supposons que $f : X \to X$ soit un automorphisme sur $k$ qui induit
l'identité sur l'espace topologique sous-jacent de $X$. Alors
\begin{enumerate}
\item $f^*\mathcal{F} \cong \mathcal{F}$ pour
chaque $\mathcal{O}_X$-module cohérent, et
\item si $\dim(Z) > 0$ pour chaque composante
irréductible $Z \subset X$, alors $f$ est l'identité.
\end{enumerate}
\end{lemma}

\begin{proof}
La partie (1) découle de la partie (2) et du fait que les
composantes connexes de $X$ de dimension $0$ sont des spectres de corps.

\medskip\noindent
Supposons que $Z \subset X$ soit une composante irréductible considérée
comme un sous-schéma fermé intègre. Évidemment, $f(Z) \subset Z$
et $f|_Z : Z \to Z$ sont un automorphisme sur $k$ qui induit
l'identité sur l'espace topologique sous-jacent de $Z$.
Puisque $X$ est réduit, il suffit de montrer que les flèches $f|_Z : Z \to Z$
sont l'identité. Cela nous ramène au cas discuté dans le paragraphe suivant.

\medskip\noindent
Supposons que $X$ soit irréductible de dimension $> 0$.
Choisissons un ouvert affine non vide $U \subset X$. Puisque
$f(U) \subset U$ et puisque $U \subset X$ est schématiquement
dense, il suffit de prouver que $f|_U : U \to U$ est l'identité.

\medskip\noindent
Supposons que $X = \Spec(A)$ soit affine, irréductible, de dimension $> 0$ et que
$k$ soit un corps infini. Soit $g \in A$ non constant. L'ensemble
$$
S = \bigcup\nolimits_{\lambda \in k} V(g - \lambda)
$$
est dense dans $X$ parce que c'est l'image réciproque du
sous-ensemble dense $\mathbf{A}^1_k(k)$ par le morphisme
non constant $g : X \to \mathbf{A}^1_k$. Si $x \in S$,
alors l'image $g(x)$ de $g$ dans $\kappa(x)$ est dans
l'image de $k \to \kappa(x)$. Par conséquent, $f^\sharp :
\kappa(x) \to \kappa(x)$ fixe $g(x)$. Ainsi, l'image de
$f^\sharp(g)$ dans $\kappa(x)$ est égale à $g(x)$. Nous en concluons que
$$
S \subset V(g - f^\sharp(g))
$$
et puisque $X$ est réduit et $S$ est dense, nous concluons
$g=f^\sharp(g)$. Cela prouve $f^\sharp = \text{id}_A$ car $A$ est
engendré comme une $k$-algèbre par les éléments $g$ comme ci-dessus (détails
omis ; indication : l'ensemble des fonctions constantes est un $k$-sous-espace
vectoriel de dimension finie de $A$). Nous concluons que $f = \text{id}_X$.

\medskip\noindent
Supposons que $X = \Spec(A)$ soit affine, irréductible, de
dimension $> 0$ et que $k$ soit un corps fini. Si pour tout
sous-schéma fermé intègre de dimension $1$, $C \subset X$, la
restriction $f|_C : C \to C$ est l'identité, alors $f$ est
l'identité. Cela nous ramène au cas où $X$ est une courbe. Une courbe
sur un corps fini possède un groupe d'automorphismes fini (détails
omis). Ainsi, $f$ a un ordre fini, disons $n$. Ensuite, nous choisissons
$g : X \to \mathbf{A}^1_k$ non constant comme ci-dessus et considérons
$$
S = \{x \in X\text{ fermé tel que }[\kappa(g(x)) : k]
\text{ est premier à }n\}
$$
En argumentant comme précédemment, nous trouvons que
$S$ est dense dans $X$. Puisque pour $x \in X$ fermé,
l'application $f^\sharp : \kappa(x) \to \kappa(x)$
est un automorphisme d'ordre divisant $n$, nous
voyons que pour $x \in S$, cet automorphisme agit
trivialement sur le sous-corps engendré par l'image de
$g$ dans $\kappa(x)$. Ainsi nous concluons que $S \subset
V(g - f^\sharp(g))$ et on conclut comme précédemment.
\end{proof}






\section{Caractéristiques d'Euler--Poincaré}
\label{section-euler}

\noindent
Dans cette section, nous prouvons quelques propriétés élémentaires des
caractéristiques d'Euler--Poincaré des faisceaux cohérents sur des schémas propres sur un corps.

\begin{definition}
\label{definition-euler-characteristic}
Soit $k$ un corps. Soit $X$ un schéma propre sur $k$. Soit
$\mathcal{F}$ un $\mathcal{O}_X$-module cohérent. Dans cette
situation, la {\it caractéristique d'Euler--Poincaré de $\mathcal{F}$ } est l'entier
$$
\chi(X, \mathcal{F}) = \sum\nolimits_i (-1)^i \dim_k H^i(X, \mathcal{F}).
$$
Pour la justification de la formule, voir ci-dessous.
\end{definition}

\noindent
Dans le cas de la définition, seul un nombre fini des
espaces vectoriels $H^i(X, \mathcal{F})$ sont non nuls
(Cohomologie des Schémas, lemme
\ref{coherent-lemma-quasi-coherence-higher-direct-images})
et chacun de ces espaces est de dimension finie (Cohomologie
des Schémas, lemme
\ref{coherent-lemma-proper-over-affine-cohomology-finite}). Ainsi
$\chi(X, \mathcal{F}) \in \mathbf{Z}$ est bien défini. Observons que cette
définition dépend du corps $k$ et pas seulement de la paire $(X, \mathcal{F})$.

\begin{lemma}
\label{lemma-euler-characteristic-additive}
Soit $k$ un corps. Supposons que $X$ soit un schéma propre
sur $k$. Supposons que $0 \to \mathcal{F}_1 \to \mathcal{F}_2 \to
\mathcal{F}_3 \to 0$ soit une suite exacte courte de modules cohérents sur $X$. Alors
$$
\chi(X, \mathcal{F}_2) = \chi(X, \mathcal{F}_1) + \chi(X, \mathcal{F}_3)
$$
\end{lemma}

\begin{proof}
Considérons la longue suite exacte de cohomologie
$$
0 \to H^0(X, \mathcal{F}_1) \to H^0(X, \mathcal{F}_2) \to
H^0(X, \mathcal{F}_3) \to H^1(X, \mathcal{F}_1) \to \ldots
$$
associée à la suite exacte courte du lemme. Le théorème
du rang en algèbre linéaire montre que
$$
0 = \dim H^0(X, \mathcal{F}_1) - \dim H^0(X, \mathcal{F}_2)
+ \dim H^0(X, \mathcal{F}_3) - \dim H^1(X, \mathcal{F}_1) + \ldots
$$
Cela implique immédiatement le lemme.
\end{proof}

\begin{lemma}
\label{lemma-chi-tensor-finite}
Soit $k$ un corps. Supposons que $X$ soit un
schéma propre sur $k$. Supposons que $\mathcal{F}$ soit un
faisceau cohérent avec $\dim(\text{Supp}(\mathcal{F})) \leq 0$. Alors
\begin{enumerate}
\item $\mathcal{F}$ est engendré par ses sections
globales, \item $H^0(X, \mathcal{F}) =
\bigoplus_{x \in \text{Supp}(\mathcal{F})}
\mathcal{F}_x$, \item $H^i(X, \mathcal{F}) = 0$ pour
$i > 0$, \item $\chi(X, \mathcal{F}) = \dim_k
H^0(X, \mathcal{F})$, et \item $\chi(X, \mathcal{F}
\otimes \mathcal{E}) = n\chi(X, \mathcal{F})$ pour
chaque module localement libre $\mathcal{E}$ de rang $n$.
\end{enumerate}
\end{lemma}

\begin{proof}
Par Cohomologie des Schémas, lemme
\ref{coherent-lemma-coherent-support-closed}, on voit que
$\mathcal{F} = i_*\mathcal{G}$ où $i : Z \to X$ est l'inclusion
du support schématique de $\mathcal{F}$ et où $\mathcal{G}$ est
un $\mathcal{O}_Z$-module cohérent. Par définition du support
schématique, l'espace topologique sous-jacent de $Z$ est
$\text{Supp}(\mathcal{F})$. Puisque la dimension de $Z$ est
$0$, on voit que $Z$ est affine (Propriétés, lemme
\ref{properties-lemma-locally-Noetherian-dimension-0}). Ainsi
$\mathcal{G}$ est engendré par ses sections globales et les groupes de
cohomologie supérieurs de $\mathcal{G}$ sont nuls (Cohomologie
des Schémas, lemme
\ref{coherent-lemma-quasi-coherent-affine-cohomology-zero}). En
fait, par le lemme \ref{lemma-algebraic-scheme-dim-0}, le schéma $Z$
est une réunion finie disjointe de spectres d'anneaux locaux
artiniens. Par conséquent, $H^0(Z, \mathcal{G}) =
\bigoplus_{z \in Z} \mathcal{G}_z$. Les cohomologies de $\mathcal{F}$
et $\mathcal{G}$ coïncident d'après Cohomologie des Schémas, lemme
\ref{coherent-lemma-relative-affine-cohomology}. Ainsi $H^i(X,
\mathcal{F}) = 0$ pour $i > 0$ et $H^0(X, \mathcal{F}) = H^0(Z,
\mathcal{G})$. En particulier, (3) est vérifiée. Pour $z \in Z$
correspondant à $x \in \text{Supp}(\mathcal{F})$ nous avons
$\mathcal{G}_z = (i_*\mathcal{G})_x = \mathcal{F}_x$. Nous concluons que (2)
est vérifiée. Évidemment, (2) implique (1). Nous obtenons (4) par définition de la
caractéristique d'Euler--Poincaré $\chi(X, \mathcal{F})$ et (3). Par la formule de
projection (Cohomologie, lemme \ref{cohomology-lemma-projection-formula}) nous avons
$$
i_*(\mathcal{G} \otimes i^*\mathcal{E}) = \mathcal{F} \otimes \mathcal{E}.
$$
Puisque $Z$ a pour dimension $0$, le faisceau localement
libre $i^*\mathcal{E}$ est isomorphe à $\mathcal{O}_Z^{\oplus
n}$ et, en raisonnant comme ci-dessus, on voit que (5) est vérifiée.
\end{proof}

\begin{lemma}
\label{lemma-euler-characteristic-extend-base-field}
Supposons que $k'/k$ soit une extension de corps. Supposons que $X$ soit un
schéma propre sur $k$. Supposons que $\mathcal{F}$ soit un faisceau
cohérent sur $X$. Supposons que $\mathcal{F}'$ soit l'image réciproque de
$\mathcal{F}$ sur $X_{k'}$. Alors $\chi(X, \mathcal{F}) = \chi(X_{k'}, \mathcal{F}')$.
\end{lemma}

\begin{proof}
En effet,
$$
H^i(X_{k'}, \mathcal{F}') = H^i(X, \mathcal{F}) \otimes_k k'
$$
par changement de base plat, voir Cohomologie des
Schémas, lemme \ref{coherent-lemma-flat-base-change-cohomology}.
\end{proof}

\begin{lemma}
\label{lemma-euler-characteristic-morphism}
Soit $k$ un corps. Supposons que $f : Y \to X$ soit un morphisme de schémas
propres sur $k$. Supposons que $\mathcal{G}$ soit un $\mathcal{O}_Y$-module cohérent. Alors
$$
\chi(Y, \mathcal{G}) = \sum (-1)^i \chi(X, R^if_*\mathcal{G})
$$
\end{lemma}

\begin{proof}
La formule a du sens : les faisceaux
$R^if_*\mathcal{G}$ sont cohérents et seul un nombre
fini d'entre eux est non nul, voir Cohomologie des
schémas, proposition
\ref{coherent-proposition-proper-pushforward-coherent} et lemme
\ref{coherent-lemma-quasi-coherence-higher-direct-images}. D'après
Cohomologie, lemme \ref{cohomology-lemma-Leray}, il existe une suite spectrale avec
$$
E_2^{p, q} = H^p(X, R^qf_*\mathcal{G})
$$
convergeant vers $H^{p + q}(Y, \mathcal{G})$. Par la finitude de la
cohomologie sur $X$, les termes $E_2^{p, q}$ non nuls sont en nombre
fini et chaque $E_2^{p, q}$ est un espace vectoriel de dimension finie.
Il en va de même des termes $E_r^{p, q}$ pour $r \geq 2$, et
$$
\sum (-1)^{p + q} \dim_k E_r^{p, q}
$$
est indépendant de $r$. Puisque pour $r$ suffisamment grand
nous avons $E_r^{p, q} = E_\infty^{p, q}$ et puisque la
convergence signifie qu'il existe une filtration sur $H^n(Y,
\mathcal{G})$ dont les morceaux gradués sont $E_\infty^{p, q}$
avec $p + q = n$ (c'est le sens de la convergence de la suite
spectrale), nous concluons. Voir aussi le lemme plus général du chapitre
Homologie générale, lemme \ref{homology-lemma-biregular-ss-relation-in-K0}.
\end{proof}








\section{Espace projectif}
\label{section-projective-space}

\noindent
Quelques résultats sur l'espace projectif sur un corps.

\begin{lemma}
\label{lemma-projective-space-smooth}
\begin{slogan}
L'espace projectif est lisse.
\end{slogan}
Supposons que $k$ soit un corps et $n \geq 0$. Alors
$\mathbf{P}^n_k$ est une variété projective lisse de dimension $n$ sur $k$.
\end{lemma}

\begin{proof}
Démonstration omise.
\end{proof}

\begin{lemma}
\label{lemma-intersection-in-affine-space}
Supposons que $k$ soit un corps et $n \geq 0$. Supposons que $X, Y
\subset \mathbf{A}^n_k$ soit des sous-ensembles fermés. Supposons que $X$
et $Y$ soient équidimensionnels, $\dim(X) = r$ et $\dim(Y) = s$. Alors
chaque composante irréductible de $X \cap Y$ a une dimension $\geq r + s - n$.
\end{lemma}

\begin{proof}
Considérons le sous-schéma fermé $X \times Y \subset
\mathbf{A}^{2n}_k$ où nous utilisons les coordonnées
$x_1, \ldots, x_n, y_1, \ldots, y_n$. Alors $X \cap Y = X
\times Y \cap V(x_1 - y_1, \ldots, x_n - y_n)$. Supposons
que $t \in X \cap Y \subset X \times Y$ soit un point
fermé. Par le lemme
\ref{lemma-dimension-product-locally-algebraic}, nous avons $\dim_t(X
\times Y) = \dim(X) + \dim(Y)$. Ainsi $\dim(\mathcal{O}_{X \times Y,
t}) = r + s$ par le lemme \ref{lemma-dimension-locally-algebraic}.
Par Algèbre, lemme \ref{algebra-lemma-one-equation}, nous concluons que
$$
\dim(\mathcal{O}_{X \cap Y, t}) =
\dim(\mathcal{O}_{X \times Y, t}/(x_1 - y_1, \ldots, x_n - y_n)) \geq
r + s - n
$$
Cela implique le résultat par le lemme \ref{lemma-dimension-locally-algebraic}.
\end{proof}

\begin{lemma}
\label{lemma-intersection-in-projective-space}
Supposons que $k$ soit un corps et $n \geq 0$. Supposons que
$X, Y \subset \mathbf{P}^n_k$ soit des sous-ensembles fermés
non vides. Si $\dim(X) = r$ et $\dim(Y) = s$ et $r + s \geq n$,
alors $X \cap Y$ est non vide et $\dim(X \cap Y) \geq r + s - n$.
\end{lemma}

\begin{proof}
Écrivons $\mathbf{A}^{n + 1} = \Spec(k[x_0, \ldots, x_n])$ et
$\mathbf{P}^n = \text{Proj}(k[T_0, \ldots, T_n])$.
Considérons le morphisme $\pi : \mathbf{A}^{n + 1}
\setminus \{0\} \to \mathbf{P}^n$ qui envoie $(x_0,
\ldots, x_n)$ vers le point $[x_0 : \ldots : x_n]$. Plus
précisément, c'est le morphisme associé à la paire
$(\mathcal{O}_{\mathbf{A}^{n + 1} \setminus \{0\}}, (x_0,
\ldots, x_n))$, voir Constructions, lemme
\ref{constructions-lemma-projective-space}. Sur l'ouvert affine standard
$D_+(T_i)$, nous obtenons le morphisme associé à l'homomorphisme d'anneaux
$$
k\left[\frac{T_0}{T_i}, \ldots, \frac{T_n}{T_i}\right]
\longrightarrow
k\left[T_0, \ldots, T_n, \frac{1}{T_i}\right] \cong
k\left[\frac{T_0}{T_i}, \ldots, \frac{T_n}{T_i}\right]
\left[T_i, \frac{1}{T_i}\right]
$$
qui est surjectif et lisse de dimension relative $1$, à fibres
irréductibles (détails omis). Par conséquent, $\pi^{-1}(X)$
et $\pi^{-1}(Y)$ sont des sous-ensembles fermés non vides de
dimension $r + 1$ et $s + 1$. Choisissons une composante
irréductible $V \subset \pi^{-1}(X)$ de dimension $r + 1$ et une
composante irréductible $W \subset \pi^{-1}(Y)$ de dimension $s + 1$.
Observons que cela implique que $V$ et $W$ contiennent chaque fibre
de $\pi$ qu'elles rencontrent (puisque $\pi$ a des fibres
irréductibles de dimension $1$ et puisque le lemme
\ref{lemma-dimension-fibres-locally-algebraic} dit que les fibres de $V
\to \pi(V)$ et $W \to \pi(W)$ ont pour dimension $\geq 1$). Supposons
que $\overline{V}$ et $\overline{W}$ soient l'adhérence de $V$ et $W$
dans $\mathbf{A}^{n + 1}$. Puisque $0 \in \mathbf{A}^{n + 1}$ est dans
l'adhérence de chaque fibre de $\pi$, on voit que $0 \in
\overline{V} \cap \overline{W}$. Par le lemme
\ref{lemma-intersection-in-affine-space}, nous avons $\dim(\overline{V} \cap
\overline{W}) \geq r + s - n + 1$. En raisonnant comme ci-dessus et en appliquant à
nouveau le lemme \ref{lemma-dimension-fibres-locally-algebraic}, nous concluons que
$\pi(V \cap W) \subset X \cap Y$ est de dimension au moins $r + s - n$, comme annoncé.
\end{proof}

\begin{lemma}
\label{lemma-equation-codim-1-in-projective-space}
Soit $k$ un corps. Supposons que $Z \subset
\mathbf{P}^n_k$ soit un sous-schéma fermé sans points immergés
et dont chaque composante irréductible de $Z$ est de dimension $n - 1$. Alors
l'idéal $I(Z) \subset k[T_0, \ldots, T_n]$ correspondant à $Z$ est principal.
\end{lemma}

\begin{proof}
Ceci est un cas particulier des Diviseurs, lemme
\ref{divisors-lemma-equation-codim-1-in-projective-space}.
\end{proof}




\section{Faisceaux cohérents sur l'espace projectif}
\label{section-coherent-Pn}

\noindent
Dans cette section, nous prouvons certains résultats sur la
cohomologie des faisceaux cohérents sur $\mathbf{P}^n$ sur un
corps, qui peuvent être trouvés dans \cite{Mum}. Ceux-ci seront
utiles plus tard lors de la discussion des schémas de Quot et de Hilbert.

\subsection{Préliminaires}
\label{subsection-preliminaries}

\noindent
Soit $k$ un corps, $n \geq 1$, $d \geq 1$, et soit $s \in
\Gamma(\mathbf{P}_k^n, \mathcal{O}(d))$ une section non
nulle. Dans cette section, nous noterons $\mathcal{O}(d)$ le
$d$-ième tordu du faisceau structural sur l'espace projectif
(Constructions, définitions
\ref{constructions-definition-twist} et
\ref{constructions-definition-projective-space}). Comme $\mathbf{P}^n_k$
est une variété, cette section est régulière ; ainsi $s$ est une section
régulière de $\mathcal{O}(d)$ et définit un diviseur de Cartier effectif
$H = Z(s) \subset \mathbf{P}^n_k$, voir Diviseurs, section
\ref{divisors-section-effective-Cartier-divisors}. Un tel diviseur $H$ est
appelé une {\it hypersurface} et, si $d = 1$, il est appelé un {\it hyperplan}.

\begin{lemma}
\label{lemma-hyperplane}
Soit $k$ un corps. Supposons
que $n \geq 1$. Soit $i : H \to \mathbf{P}^n_k$
un hyperplan. Alors il existe un isomorphisme
$$
\varphi : \mathbf{P}^{n - 1}_k \longrightarrow H
$$
tel que l'image réciproque de $i^*\mathcal{O}(1)$ soit isomorphe à $\mathcal{O}(1)$.
\end{lemma}

\begin{proof}
Nous avons $\mathbf{P}^n_k = \text{Proj}(k[T_0,
\ldots, T_n])$. La section $s$ correspond à une
forme homogène dans $T_0, \ldots, T_n$ de degré
$1$, voir Cohomologie des schémas, section
\ref{coherent-section-cohomology-projective-space}.
Écrivons $s = \sum a_i T_i$. Constructions, le lemme
\ref{constructions-lemma-closed-in-projective-space}
montre que $H = \text{Proj}(k[T_0, \ldots, T_n]/I)$
pour l'idéal gradué $I$ défini en posant $I_d$ égal
au noyau de l'application $\Gamma(\mathbf{P}^n_k,
\mathcal{O}(d)) \to \Gamma(H, i^*\mathcal{O}(d))$.
D'après notre construction de $Z(s)$ dans Diviseurs,
définition \ref{divisors-definition-zero-scheme-s},
on voit que sur $D_{+}(T_j)$ l'idéal de $H$ est
engendré par $\sum a_i T_i/T_j$ dans l'anneau de polynômes
$k[T_0/T_j, \ldots, T_n/T_j]$. Ainsi, il est clair que
$I$ est l'idéal engendré par $\sum a_i T_i$. Remarquez que
$$
k[T_0, \ldots, T_n]/I = k[T_0, \ldots, T_n]/(\sum a_i T_i) \cong
k[S_0, \ldots, S_{n - 1}]
$$
en tant qu'anneaux gradués. Par exemple, si $a_n \not = 0$,
alors envoyer $S_i$ sur la classe de $T_i$ convient.
Nous obtenons l'isomorphisme annoncé par fonctorialité de
$\text{Proj}$. L'égalité des tordus du faisceau structural
résulte par exemple du chapitre Constructions, lemme
\ref{constructions-lemma-surjective-graded-rings-generated-degree-1-map-proj}.
\end{proof}

\begin{lemma}
\label{lemma-exact-sequence-induction}
Supposons que $k$ soit un corps infini.
Supposons que $n \geq 1$. Soit $\mathcal{F}$ un
module cohérent sur $\mathbf{P}^n_k$. Alors il
existe une section non nulle $s \in
\Gamma(\mathbf{P}^n_k, \mathcal{O}(1))$ et une suite exacte courte
$$
0 \to \mathcal{F}(-1) \to \mathcal{F} \to i_*\mathcal{G} \to 0
$$
où $i : H \to \mathbf{P}^n_k$ est l'hyperplan $H$
associé à $s$ et $\mathcal{G} = i^*\mathcal{F}$.
\end{lemma}

\begin{proof}
L'application $\mathcal{F}(-1) \to \mathcal{F}$ provient
de Constructions, Équation (
\ref{constructions-equation-multiply-on-sheaf} ) avec $n =
1$, $m = -1$ et la section $s$ de $\mathcal{O}(1)$. Voyons à
quoi ressemble cette application si nous la restreignons à
$D_{+}(T_0)$. Écrivons $D_{+}(T_0) = \Spec(k[x_1, \ldots,
x_n])$ avec $x_i = T_i/T_0$. Identifions
$\mathcal{O}(1)|_{D_{+}(T_0)}$ avec $\mathcal{O}$ à l'aide de la
section $T_0$. Par conséquent, si $s = \sum a_iT_i$ alors
$s|_{D_{+}(T_0)} = a_0 + \sum a_ix_i$ avec l'identification
choisie ci-dessus. De plus, supposons que
$\mathcal{F}|_{D_{+}(T_0)}$ corresponde au $k[x_1, \ldots, x_n]$-module
fini $M$. Via l'identification $\mathcal{F}(-1) = \mathcal{F} \otimes
\mathcal{O}(-1)$ et notre trivialisation choisie de $\mathcal{O}(1)$,
on voit que $\mathcal{F}(-1)$ correspond également à $M$. Ainsi,
restreindre $\mathcal{F}(-1) \to \mathcal{F}$ à $D_{+}(T_0)$ donne l'application
$$
M \xrightarrow{a_0 + \sum a_ix_i} M
$$
Pour que la flèche soit injective, il suffit de choisir $a_0
+ \sum a_ix_i$ n'appartenant à aucun des idéaux premiers
associés de $M$, voir Algèbre, lemme
\ref{algebra-lemma-ass-zero-divisors}. D'après Algèbre, lemme
\ref{algebra-lemma-finite-ass}, l'ensemble $\text{Ass}(M)$ des
idéaux premiers associés de $M$ est fini. Notons que pour
$\mathfrak p \in \text{Ass}(M)$, l'intersection $\mathfrak p \cap
\{a_0 + \sum a_i x_i\}$ est un sous-espace vectoriel propre de $k$.
Nous en concluons qu'il existe une famille finie de sous-espaces
vectoriels propres $V_1, \ldots, V_m \subset \Gamma(\mathbf{P}^n_k,
\mathcal{O}(1))$ tels que si nous prenons $s$ en dehors de $\bigcup
V_i$, alors la multiplication par $s$ est injective sur
$D_{+}(T_0)$. De même pour la restriction à $D_{+}(T_j)$ pour $j =
1, \ldots, n$. Comme $k$ est infini, une union finie de sous-espaces
vectoriels propres n'est jamais égale à l'espace tout entier ; nous
pouvons donc choisir $s$ de sorte que l'application soit
injective. Le conoyau de $\mathcal{F}(-1) \to \mathcal{F}$ est annulé
par $\Im(s : \mathcal{O}(-1) \to \mathcal{O})$ qui est l'idéal
faisceau de $H$ selon Diviseurs, définition
\ref{divisors-definition-zero-scheme-s}. Ainsi, nous obtenons $\mathcal{G}$ sur $H$ en
utilisant Cohomologie des Schémas, lemme \ref{coherent-lemma-i-star-equivalence}.
\end{proof}

\begin{remark}
\label{remark-exact-sequence-induction}
Supposons que $k$ soit un corps infini. Soit $n \geq 1$. Étant
donné un nombre fini de modules cohérents $\mathcal{F}_i$ sur
$\mathbf{P}^n_k$, nous pouvons choisir un seul $s \in
\Gamma(\mathbf{P}^n_k, \mathcal{O}(1))$ de sorte que l'énoncé du lemme
\ref{lemma-exact-sequence-induction} fonctionne pour chacun d'eux. Pour
le prouver, il suffit d'appliquer le lemme à $\bigoplus \mathcal{F}_i$.
\end{remark}

\begin{remark}
\label{remark-exact-sequence-induction-cohomology}
Dans la situation des lemmes \ref{lemma-hyperplane} et
\ref{lemma-exact-sequence-induction}, nous avons $H \cong
\mathbf{P}^{n - 1}_k$ avec les tordus de Serre
$\mathcal{O}_H(d) = i^*\mathcal{O}_{\mathbf{P}^n_k}(d)$. Pour
chaque $d \in \mathbf{Z}$, nous avons une suite exacte courte
$$
0 \to \mathcal{F}(d - 1) \to \mathcal{F}(d) \to i_*(\mathcal{G}(d)) \to 0
$$
En effet, le produit tensoriel par
$\mathcal{O}_{\mathbf{P}^n_k}(d)$ est un foncteur exact et,
d'après la formule de projection (Cohomologie, lemme
\ref{cohomology-lemma-projection-formula}), nous avons
$i_*(\mathcal{G}(d)) = i_*\mathcal{G} \otimes
\mathcal{O}_{\mathbf{P}^n_k}(d)$. Nous obtenons des suites exactes longues correspondantes
$$
H^i(\mathbf{P}^n_k, \mathcal{F}(d - 1)) \to
H^i(\mathbf{P}^n_k, \mathcal{F}(d)) \to
H^i(H, \mathcal{G}(d)) \to
H^{i + 1}(\mathbf{P}^n_k, \mathcal{F}(d - 1))
$$
Cela résulte de ce qui précède et de l'égalité
$H^i(\mathbf{P}^n_k, i_*\mathcal{G}(d)) = H^i(H,
\mathcal{G}(d))$ d'après Cohomologie des Schémas, lemme
\ref{coherent-lemma-relative-affine-cohomology} (les immersions fermées sont affines).
\end{remark}





\subsection{Régularité}
\label{subsection-regularity}

\noindent
Voici la définition.

\begin{definition}
\label{definition-regularity}
Soit $k$ un corps. Supposons que $n \geq 0$.
Soit $\mathcal{F}$ un faisceau cohérent sur
$\mathbf{P}^n_k$. Nous disons que $\mathcal{F}$ est {\it $m$-régulier} si
$$
H^i(\mathbf{P}^n_k, \mathcal{F}(m - i)) = 0
$$
pour $i = 1, \ldots, n$.
\end{definition}

\noindent
Notons que $\mathcal{F} = \mathcal{O}(d)$ est $m$-régulier si
et seulement si $d \geq -m$. Cela découle du calcul des groupes
de cohomologie dans Cohomologie des schémas, Équation (
\ref{coherent-equation-identify} ). En effet, on voit que
$H^n(\mathbf{P}^n_k, \mathcal{O}(d)) = 0$ si et seulement si $d \geq -n$.

\begin{lemma}
\label{lemma-m-regular-extend-base-field}
Supposons que $k'/k$ soit une extension de corps. Supposons que
$n \geq 0$. Soit $\mathcal{F}$ un faisceau cohérent sur
$\mathbf{P}^n_k$. Supposons que $\mathcal{F}'$ soit l'image réciproque de
$\mathcal{F}$ sur $\mathbf{P}^n_{k'}$. Alors $\mathcal{F}$
est $m$-régulier si et seulement si $\mathcal{F}'$ est $m$-régulier.
\end{lemma}

\begin{proof}
Cela résulte de l'isomorphisme
$$
H^i(\mathbf{P}^n_{k'}, \mathcal{F}') =
H^i(\mathbf{P}^n_k, \mathcal{F}) \otimes_k k'
$$
par changement de base plat, voir Cohomologie des
Schémas, lemme \ref{coherent-lemma-flat-base-change-cohomology}.
\end{proof}

\begin{lemma}
\label{lemma-m-regular}
Dans la situation du lemme
\ref{lemma-exact-sequence-induction}, si $\mathcal{F}$ est $m$-régulier,
alors $\mathcal{G}$ est $m$-régulier sur $H \cong \mathbf{P}^{n - 1}_k$.
\end{lemma}

\begin{proof}
Rappelons que $H^i(\mathbf{P}^n_k, i_*\mathcal{G}) = H^i(H,
\mathcal{G})$ d'après Cohomologie des Schémas, lemme
\ref{coherent-lemma-relative-affine-cohomology}. Ainsi, pour $i \geq 1$, nous obtenons
$$
H^i(\mathbf{P}^n_k, \mathcal{F}(m - i)) \to
H^i(H, \mathcal{G}(m - i)) \to
H^{i + 1}(\mathbf{P}^n_k, \mathcal{F}(m - 1 - i))
$$
par la Remarque
\ref{remark-exact-sequence-induction-cohomology}. Le lemme suit.
\end{proof}

\begin{lemma}
\label{lemma-m-regular-up}
Soit $k$ un corps. Supposons que $n
\geq 0$. Soit $\mathcal{F}$ un faisceau cohérent
sur $\mathbf{P}^n_k$. Si $\mathcal{F}$ est
$m$-régulier, alors $\mathcal{F}$ est $(m + 1)$-régulier.
\end{lemma}

\begin{proof}
Nous procédons par récurrence sur $n$. Si $n = 0$, chaque
faisceau est $m$-régulier pour tout $m$, il n'y a rien à
démontrer. D'après le lemme
\ref{lemma-m-regular-extend-base-field}, nous pouvons remplacer $k$ par
une extension infinie et supposer que $k$ est infini. Nous pouvons alors
appliquer le lemme \ref{lemma-exact-sequence-induction}. D'après le lemme
\ref{lemma-m-regular}, nous savons que $\mathcal{G}$ est $m$-régulier.
Par induction sur $n$, on voit que $\mathcal{G}$ est $(m +
1)$-régulier. En considérant la suite exacte longue de cohomologie associée à la suite
$$
0 \to \mathcal{F}(m - i) \to \mathcal{F}(m + 1 - i)
\to i_*\mathcal{G}(m + 1 - i) \to 0
$$
comme dans la remarque
\ref{remark-exact-sequence-induction-cohomology}, on déduit
pour $i \geq 1$ l'annulation de $H^i(\mathbf{P}^n_k,
\mathcal{F}(m + 1 - i))$ à partir de l'annulation de
$H^i(\mathbf{P}^n_k, \mathcal{F}(m - i))$ et $H^i(\mathbf{P}^n_k, \mathcal{G}(m + 1 - i))$.
\end{proof}

\begin{lemma}
\label{lemma-m-regular-multiply}
Soit $k$ un corps. Supposons que $n \geq 0$. Soit
$\mathcal{F}$ un faisceau cohérent sur $\mathbf{P}^n_k$. Si
$\mathcal{F}$ est $m$-régulier, alors l'application de multiplication
$$
H^0(\mathbf{P}^n_k, \mathcal{F}(m)) \otimes_k
H^0(\mathbf{P}^n_k, \mathcal{O}(1))
\longrightarrow
H^0(\mathbf{P}^n_k, \mathcal{F}(m + 1))
$$
est surjective.
\end{lemma}

\begin{proof}
Supposons que $k'/k$ soit une extension de corps. Supposons que
$\mathcal{F}'$ soit comme dans le lemme
\ref{lemma-m-regular-extend-base-field}. D'après le chapitre Cohomologie des schémas,
lemme \ref{coherent-lemma-flat-base-change-cohomology}, l'application linéaire obtenue
par changement de base vers $k'$ coïncide avec l'application linéaire associée au
faisceau $\mathcal{F}'$. Puisque $k \to k'$ est fidèlement plat, il suffit de
prouver le lemme sur $k'$, c'est-à-dire que l'on peut supposer $k$ infini.

\medskip\noindent
Supposons que $k$ soit infini. Nous prouvons le lemme par
induction sur $n$. Le cas $n = 0$ est trivial car
$\mathcal{O}(1) \cong \mathcal{O}$ est généré par $T_0$. Pour
$n > 0$, appliquons le lemme \ref{lemma-exact-sequence-induction}
et tensorisons la suite par $\mathcal{O}(m + 1)$ ; on obtient
$$
0 \to \mathcal{F}(m) \xrightarrow{s} \mathcal{F}(m + 1) \to
i_*\mathcal{G}(m + 1) \to 0
$$
voir Remarque
\ref{remark-exact-sequence-induction-cohomology}. Soit $t \in
H^0(\mathbf{P}^n_k, \mathcal{F}(m + 1))$. Par induction, l'image
$\overline{t} \in H^0(H, \mathcal{G}(m + 1))$ est l'image de $\sum g_i
\otimes \overline{s}_i$ avec $\overline{s}_i \in \Gamma(H,
\mathcal{O}(1))$ et $g_i \in H^0(H, \mathcal{G}(m))$. Puisque $\mathcal{F}$
est $m$-régulier, nous avons $H^1(\mathbf{P}^n_k, \mathcal{F}(m - 1)) = 0$,
d'où la suite exacte longue de cohomologie associée à la suite exacte courte
$$
0 \to \mathcal{F}(m - 1) \xrightarrow{s} \mathcal{F}(m) \to
i_*\mathcal{G}(m) \to 0
$$
montre que nous pouvons relever $g_i$ en $f_i \in
H^0(\mathbf{P}^n_k, \mathcal{F}(m))$. Nous pouvons également
relever $\overline{s}_i$ en $s_i \in H^0(\mathbf{P}^n_k,
\mathcal{O}(1))$ (voir la preuve du lemme \ref{lemma-hyperplane} par
exemple). Après avoir soustrait l'image de $\sum f_i \otimes s_i$ de
$t$, on voit que nous pouvons supposer $\overline{t} = 0$. Mais
cela signifie exactement que $t$ est l'image de $f \otimes s$ pour un
certain $f \in H^0(\mathbf{P}^n_k, \mathcal{F}(m))$, comme annoncé.
\end{proof}

\begin{lemma}
\label{lemma-m-regular-globally-generated}
Soit $k$ un corps. Supposons que $n
\geq 0$. Soit $\mathcal{F}$ un faisceau cohérent
sur $\mathbf{P}^n_k$. Si $\mathcal{F}$ est
$m$-régulier, alors $\mathcal{F}(m)$ est engendré par ses sections globales.
\end{lemma}

\begin{proof}
Pour tout $d \gg 0$, le faisceau $\mathcal{F}(d)$ est engendré par ses
sections globales. Cela découle par exemple de la première partie de
Cohomologie des Schémas, lemme
\ref{coherent-lemma-coherent-projective}. Choisissons $d \geq m$ de telle sorte
que $\mathcal{F}(d)$ soit engendré par ses sections globales. Choisissons une base $f_1,
\ldots, f_r \in H^0(\mathbf{P}^n_k, \mathcal{F})$. D'après le lemme
\ref{lemma-m-regular-multiply}, chaque élément $f \in H^0(\mathbf{P}^n_k,
\mathcal{F}(d))$ peut être écrit comme $f = \sum P_if_i$ pour un certain $P_i \in
k[T_0, \ldots, T_n]$ homogène de degré $d - m$. Puisque les sections $f$ génèrent
$\mathcal{F}(d)$, il en découle que les sections $f_i$ génèrent $\mathcal{F}(m)$.
\end{proof}

\subsection{Polynômes de Hilbert}
\label{subsection-hilbert}

\noindent
Le lemme suivant sera englobé par le lemme plus général
\ref{lemma-numerical-polynomial-from-euler}.

\begin{lemma}
\label{lemma-hilbert-polynomial}
Soit $k$ un corps. Supposons que $n \geq 0$. Soit
$\mathcal{F}$ un faisceau cohérent sur $\mathbf{P}^n_k$. La fonction
$$
d \longmapsto \chi(\mathbf{P}^n_k, \mathcal{F}(d))
$$
est un polynôme.
\end{lemma}

\begin{proof}
Nous procédons par récurrence sur $n$. Si $n =
0$, alors $\mathbf{P}^n_k = \Spec(k)$ et
$\mathcal{F}(d) = \mathcal{F}$. Par conséquent, dans
ce cas, la fonction est constante, c’est-à-dire un
polynôme de degré $0$. Supposons $n > 0$. Par le lemme
\ref{lemma-euler-characteristic-extend-base-field},
on peut supposer que $k$ est infini. Appliquons
le lemme \ref{lemma-exact-sequence-induction}. En
appliquant le lemme
\ref{lemma-euler-characteristic-additive} aux suites tordues $0 \to
\mathcal{F}(d - 1) \to \mathcal{F}(d) \to i_*\mathcal{G}(d) \to 0$, nous obtenons
$$
\chi(\mathbf{P}^n_k, \mathcal{F}(d)) -
\chi(\mathbf{P}^n_k, \mathcal{F}(d - 1)) =
\chi(H, \mathcal{G}(d))
$$
Voir la Remarque
\ref{remark-exact-sequence-induction-cohomology}. Comme $H
\cong \mathbf{P}^{n - 1}_k$ par induction, le côté droit est un
polynôme. On conclut en observant que toute fonction $f :
\mathbf{Z} \to \mathbf{Z}$ ayant la propriété que l'application
$d \mapsto f(d) - f(d - 1)$ est un polynôme, est elle-même un
polynôme. Nous omettons la preuve de ce fait (indication : comparer
avec Algèbre, lemme \ref{algebra-lemma-numerical-polynomial}).
\end{proof}

\begin{definition}
\label{definition-hilbert-polynomial}
Soit $k$ un corps. Soit $n \geq 0$. Soit $\mathcal{F}$
un faisceau cohérent sur $\mathbf{P}^n_k$. La fonction
$d \mapsto \chi(\mathbf{P}^n_k, \mathcal{F}(d))$ est
appelée le {\it polynôme de Hilbert} de $\mathcal{F}$.
\end{definition}

\noindent
Le polynôme de Hilbert a des coefficients dans
$\mathbf{Q}$ et pas en général dans $\mathbf{Z}$. Par exemple,
le polynôme de Hilbert de $\mathcal{O}_{\mathbf{P}^n_k}$ est
$$
d \longmapsto {d + n \choose n} = \frac{d^n}{n!} + \ldots
$$
Cela découle du lemme suivant et du fait que
$$
H^0(\mathbf{P}^n_k, \mathcal{O}_{\mathbf{P}^n_k}(d)) = k[T_0, \ldots, T_n]_d
$$
(sa partie de degré $d$), dont la dimension sur $k$ est ${d + n \choose n}$.

\begin{lemma}
\label{lemma-hilbert-polynomial-H0}
Soit $k$ un corps. Supposons que $n \geq 0$.
Soit $\mathcal{F}$ un faisceau cohérent sur
$\mathbf{P}^n_k$ avec le polynôme de Hilbert $P \in \mathbf{Q}[t]$. Alors
$$
P(d) = \dim_k H^0(\mathbf{P}^n_k, \mathcal{F}(d))
$$
pour tous les $d \gg 0$.
\end{lemma}

\begin{proof}
Cela résulte de l'annulation de la cohomologie
des tordus suffisamment élevés de $\mathcal{F}$.
Voir Cohomologie des schémas,
lemme \ref{coherent-lemma-coherent-projective}.
\end{proof}



\subsection{Bornitude des quotients}
\label{subsection-boundedness}

\noindent
Dans cette sous-section, nous bornons la
régularité des quotients d'un faisceau cohérent donné
sur $\mathbf{P}^n$ en termes du polynôme de Hilbert.

\begin{lemma}
\label{lemma-bound-quotients-free}
Soit $k$ un corps. Soit $n \geq 0$. Soit $r
\geq 1$. Soit $P \in \mathbf{Q}[t]$. Il existe un entier $m$
dépendant de $n$, $r$ et $P$ ayant la propriété suivante : si
$$
0 \to \mathcal{K} \to \mathcal{O}^{\oplus r} \to \mathcal{F} \to 0
$$
est une suite exacte courte de faisceaux cohérents
sur $\mathbf{P}^n_k$ et $\mathcal{F}$ a pour polynôme
de Hilbert $P$, alors $\mathcal{K}$ est $m$-régulier.
\end{lemma}

\begin{proof}
Nous procédons par récurrence sur $n$. Si $n = 0$, alors
$\mathbf{P}^n_k = \Spec(k)$ et tout module cohérent est
$0$-régulier et tout épimorphisme induit une surjection
sur les sections globales. Supposons $n > 0$. Considérons une
suite exacte comme dans le lemme. Supposons que $P' \in
\mathbf{Q}[t]$ est le polynôme $P'(t) = P(t) - P(t - 1)$.
Supposons que $m'$ soit l'entier qui fonctionne pour $n - 1$,
$r$ et $P'$. D'après les lemmes
\ref{lemma-m-regular-extend-base-field} et
\ref{lemma-euler-characteristic-extend-base-field}, nous pouvons
remplacer $k$ par une extension de corps, donc nous pouvons
supposer que $k$ est infini. Appliquons le lemme
\ref{lemma-exact-sequence-induction} au faisceau cohérent
$\mathcal{F}$. Le polynôme de Hilbert de $\mathcal{F}' =
i^*\mathcal{F}$ est $P'$ (voir la preuve du lemme
\ref{lemma-hilbert-polynomial}). Puisque $i^*$ est exact à droite, nous
voyons que $\mathcal{F}'$ est un quotient de $\mathcal{O}_H^{\oplus r}
= i^*\mathcal{O}^{\oplus r}$. Ainsi, l'hypothèse d'induction s'applique
à $\mathcal{F}'$ sur $H \cong \mathbf{P}^{n - 1}_k$ (lemme
\ref{lemma-hyperplane}). Notons que l'application $\mathcal{K}(-1) \to
\mathcal{K}$ est injective comme $\mathcal{K} \subset \mathcal{O}^{\oplus
r}$ et a un conoyau $i_*\mathcal{H}$ où $\mathcal{H} = i^*\mathcal{K}$.
Par le lemme du serpent (Homologie, lemme \ref{homology-lemma-snake}) nous
obtenons un diagramme commutatif avec des colonnes et des lignes exactes.
$$
\xymatrix{
& 0 \ar[d] & 0 \ar[d] & 0 \ar[d] \\
0 \ar[r] &
\mathcal{K}(-1) \ar[r] \ar[d] &
\mathcal{O}^{\oplus r}(-1) \ar[r] \ar[d] &
\mathcal{F}(-1) \ar[d] \ar[r] & 0 \\
0 \ar[r] &
\mathcal{K} \ar[r] \ar[d] &
\mathcal{O}^{\oplus r} \ar[r] \ar[d] &
\mathcal{F} \ar[d] \ar[r] & 0\\
0 \ar[r] &
i_*\mathcal{H} \ar[r] \ar[d] &
i_*\mathcal{O}_H^{\oplus r} \ar[r] \ar[d] &
i_*\mathcal{F}' \ar[r] \ar[d] & 0 \\
& 0 & 0 & 0
}
$$
Ainsi, l'hypothèse d'induction s'applique à la suite exacte
$0 \to \mathcal{H} \to \mathcal{O}_H^{\oplus r} \to
\mathcal{F}' \to 0$ sur $H \cong \mathbf{P}^{n - 1}_k$ (lemme
\ref{lemma-hyperplane}) et $\mathcal{H}$ est $m'$-régulier.
Rappelons que cela implique que $\mathcal{H}$ est
$d$-régulier pour tous les $d \geq m'$ (lemme \ref{lemma-m-regular-up}).

\medskip\noindent
Soient $i \geq 2$ et $d \geq m'$. Il découle de la
suite exacte longue de cohomologie associée à la
colonne de gauche du diagramme ci-dessus et de la
annulation de $H^{i - 1}(H, \mathcal{H}(d))$ que l'application
$$
H^i(\mathbf{P}^n_k, \mathcal{K}(d - 1))
\longrightarrow
H^i(\mathbf{P}^n_k, \mathcal{K}(d))
$$
est injective. Comme ces groupes sont nuls pour
$d \gg 0$ (Cohomologie des Schémas, lemme
\ref{coherent-lemma-coherent-projective}), nous en
concluons que $H^i(\mathbf{P}^n_k, \mathcal{K}(d))$
sont nuls pour tous les $d \geq m'$ et $i \geq 2$.

\medskip\noindent
Nous devons encore contrôler $H^1$. Tout d'abord, nous observons que toutes les applications
$$
H^1(\mathbf{P}^n_k, \mathcal{K}(m' - 1)) \to
H^1(\mathbf{P}^n_k, \mathcal{K}(m')) \to
H^1(\mathbf{P}^n_k, \mathcal{K}(m' + 1)) \to \ldots
$$
sont surjectives en raison de l'annulation de $H^1(H,
\mathcal{H}(d))$ pour $d \geq m'$. Supposons que $d > m'$ soit tel que
$$
H^1(\mathbf{P}^n_k, \mathcal{K}(d - 1))
\longrightarrow
H^1(\mathbf{P}^n_k, \mathcal{K}(d))
$$
est injective. Ensuite, $H^0(\mathbf{P}^n_k,
\mathcal{K}(d)) \to H^0(H, \mathcal{H}(d))$ est
surjective. Considérons le diagramme commutatif
$$
\xymatrix{
H^0(\mathbf{P}^n_k, \mathcal{K}(d)) \otimes_k
H^0(\mathbf{P}^n_k, \mathcal{O}(1))
\ar[r] \ar[d] &
H^0(\mathbf{P}^n_k, \mathcal{K}(d + 1)) \ar[d] \\
H^0(H, \mathcal{H}(d)) \otimes_k
H^0(H, \mathcal{O}_H(1))
\ar[r] &
H^0(H, \mathcal{H}(d + 1))
}
$$
D'après le lemme \ref{lemma-m-regular-multiply}, nous voyons
que la flèche horizontale du bas est surjective. Par conséquent,
la flèche verticale de droite est surjective. Nous concluons que
$$
H^1(\mathbf{P}^n_k, \mathcal{K}(d))
\longrightarrow
H^1(\mathbf{P}^n_k, \mathcal{K}(d + 1))
$$
est injective. Par induction, on voit que
$$
H^1(\mathbf{P}^n_k, \mathcal{K}(d - 1)) \to
H^1(\mathbf{P}^n_k, \mathcal{K}(d)) \to
H^1(\mathbf{P}^n_k, \mathcal{K}(d + 1)) \to \ldots
$$
sont toutes injectives et nous concluons que
$H^1(\mathbf{P}^n_k, \mathcal{K}(d - 1)) = 0$ en
raison du fait que ces groupes finissent par s'annuler.
Ainsi, les dimensions des groupes $H^1(\mathbf{P}^n_k,
\mathcal{K}(d))$ pour $d \geq m'$ diminuent strictement
jusqu'à devenir nulles. Il s'ensuit que la régularité
de $\mathcal{K}$ est bornée par $m' + \dim_k
H^1(\mathbf{P}^n_k, \mathcal{K}(m'))$. D'autre part, par la
annulation des groupes de cohomologie supérieurs, nous avons
$$
\dim_k H^1(\mathbf{P}^n_k, \mathcal{K}(m')) = 
- \chi(\mathbf{P}^n_k, \mathcal{K}(m')) +
\dim_k H^0(\mathbf{P}^n_k, \mathcal{K}(m'))
$$
Notons que le $H^0$ a une dimension bornée par la dimension
de $H^0(\mathbf{P}^n_k, \mathcal{O}^{\oplus r}(m'))$ qui
est au maximum $r{n + m' \choose n}$ si $m' > 0$ et zéro
sinon. Enfin, le terme $\chi(\mathbf{P}^n_k,
\mathcal{K}(m'))$ est égal à $r{n + m' \choose n} - P(m')$. Cela
donne une borne du type voulu et achève la démonstration du lemme.
\end{proof}







\section{Frobenius}
\label{section-frobenius}

\noindent
Supposons que $p$ soit un nombre premier. Si $X$ est un schéma, alors nous disons
que `` $X$ a pour caractéristique $p$ '', ou que `` $X$ est de caractéristique $p$
'', ou que `` $X$ est en caractéristique $p$ '' si $p$ est nul dans $\mathcal{O}_X$.

\begin{definition}
\label{definition-absolute-frobenius}
Soit $p$ un nombre premier. Soit $X$ un schéma de caractéristique $p$.
Le {\it Frobenius absolu de $X$} est le morphisme $F_X : X \to X$
donné par l'identité sur l'espace topologique sous-jacent et dont
$F_X^\sharp : \mathcal{O}_X \to \mathcal{O}_X$ est donné par $g \mapsto g^p$.
\end{definition}

\noindent
Cela a du sens car pour tout anneau $A$ de caractéristique $p$,
l'application $F_A : A \to A$, $a \mapsto a^p$ est un
endomorphisme d'anneau qui induit l'identité sur $\Spec(A)$. De
plus, si $A$ est local, alors $F_A$ est un homomorphisme local.
On voit ainsi que le Frobenius absolu de $X$ est
un endomorphisme de $X$ dans la catégorie des schémas. Il s'avère
que le Frobenius absolu définit un endomorphisme du
foncteur identité sur la catégorie des schémas en caractéristique $p$.

\begin{lemma}
\label{lemma-frobenius-endomorphism-identity}
Supposons que $p > 0$ soit un nombre premier.
Supposons que $f : X \to Y$ soit un morphisme de
schémas en caractéristique $p$. Alors le diagramme
$$
\xymatrix{
X \ar[d]_f \ar[r]_{F_X} & X \ar[d]^f \\
Y \ar[r]^{F_Y} & Y
}
$$
est commutatif.
\end{lemma}

\begin{proof}
Cela découle du fait algébrique trivial suivant : si
$\varphi : A \to B$ est un homomorphisme d'anneaux de
caractéristique $p$, alors $\varphi(a^p) = \varphi(a)^p$.
\end{proof}

\begin{lemma}
\label{lemma-frobenius}
Supposons que $p > 0$ soit un nombre premier. Supposons que $X$
soit un schéma en caractéristique $p$. Alors le Frobenius absolu
$F_X : X \to X$ est un homéomorphisme universel, est intégral, et
induit des extensions purement inséparables des corps résiduels.
\end{lemma}

\begin{proof}
Cela découle des résultats correspondants pour l'endomorphisme de
Frobenius $F_A : A \to A$ d'un anneau $A$ de caractéristique $p > 0$.
Voir la discussion dans Algèbre, section
\ref{algebra-section-universal-homeomorphism}, par exemple le lemme \ref{algebra-lemma-p-ring-map}.
\end{proof}

\noindent
Si nous travaillons avec des schémas sur une base fixe,
alors il existe une version relative du morphisme de Frobenius.

\begin{definition}
\label{definition-relative-frobenius}
Supposons que $p > 0$ soit un nombre premier. Supposons que $S$ soit un schéma en
caractéristique $p$. Supposons que $X$ soit un schéma sur $S$. Nous définissons
$$
X^{(p)} = X^{(p/S)} = X \times_{S, F_S} S
$$
vu comme un schéma sur $S$. En appliquant le lemme
\ref{lemma-frobenius-endomorphism-identity}, nous voyons
qu'il existe un unique morphisme $F_{X/S} : X
\longrightarrow X^{(p)}$ sur $S$ s'insérant dans le diagramme commutatif
$$
\xymatrix{
X \ar[rr]_{F_{X/S}} \ar[rrd] \ar@/^1em/[rrrr]^{F_X}
& & X^{(p)} \ar[rr] \ar[d] & & X \ar[d] \\
& & S \ar[rr]^{F_S} & & S
}
$$
où le carré de droite est cartésien. Le morphisme $F_{X/S}$
est appelé le {\it morphisme de Frobenius relatif de $X/S$}.
\end{definition}

\noindent
Remarquez que $X \mapsto X^{(p)}$ est un foncteur ; c'est le foncteur de
changement de base pour le morphisme de Frobenius absolu $F_S : S \to S$. Nous
obtenons les analogues des lemmes précédents pour le morphisme de Frobenius relatif.

\begin{lemma}
\label{lemma-relative-frobenius-endomorphism-identity}
Supposons que $p > 0$ soit un nombre premier. Supposons que
$S$ soit un schéma en caractéristique $p$. Supposons que $f : X
\to Y$ soit un morphisme de schémas sur $S$. Alors le diagramme
$$
\xymatrix{
X \ar[d]_f \ar[r]_{F_{X/S}} & X^{(p)} \ar[d]^{f^{(p)}} \\
Y \ar[r]^{F_{Y/S}} & Y^{(p)}
}
$$
est commutatif.
\end{lemma}

\begin{proof}
Cela découle du lemme
\ref{lemma-frobenius-endomorphism-identity} et des définitions.
\end{proof}

\begin{lemma}
\label{lemma-relative-frobenius}
Supposons que $p > 0$ soit un nombre premier. Supposons que
$S$ soit un schéma en caractéristique $p$. Supposons que $X$
soit un schéma sur $S$. Alors le Frobenius relatif $F_{X/S} : X
\to X^{(p)}$ est un homéomorphisme universel, est intégral, et
induit des extensions purement inséparables des corps résiduels.
\end{lemma}

\begin{proof}
D'après le lemme \ref{lemma-frobenius}, les morphismes $F_X : X
\to X$ et le changement de base $h : X^{(p)} \to X$ de $F_S$ sont
des homéomorphismes universels. Puisque $h \circ F_{X/S} = F_X$,
nous concluons que $F_{X/S}$ est un homéomorphisme universel
(Morphismes, lemme
\ref{morphisms-lemma-permanence-universal-homeomorphism}). D'après les lemmes
\ref{morphisms-lemma-universal-homeomorphism} et
\ref{morphisms-lemma-universally-injective} de Morphismes, nous concluons que $F_{X/S}$ possède également les autres propriétés.
\end{proof}

\begin{lemma}
\label{lemma-relative-frobenius-omega}
Supposons que $p > 0$ soit un nombre premier. Supposons que $S$ soit un schéma en
caractéristique $p$. Supposons que $X$ soit un schéma sur $S$. Alors $\Omega_{X/S} = \Omega_{X/X^{(p)}}$.
\end{lemma}

\begin{proof}
Cela se traduit par le fait algébrique suivant. Supposons
que $A \to B$ soit un homomorphisme d'anneaux de
caractéristique $p$. Posons $B' = B \otimes_{A, F_A} A$ et
considérons l'homomorphisme d'anneaux $F_{B/A} : B' \to B$,
$b \otimes a \mapsto b^pa$. Alors notre assertion est que
$\Omega_{B/A} = \Omega_{B/B'}$. En effet,
$\text{d}(b^pa) = 0$ si $\text{d} : B \to \Omega_{B/A}$ est la
dérivation universelle et donc $\text{d}$ est une dérivation $B'$.
\end{proof}

\begin{lemma}
\label{lemma-relative-frobenius-finite}
Supposons que $p > 0$ soit un nombre premier. Supposons que $S$ soit un
schéma en caractéristique $p$. Supposons que $X$ soit un schéma sur
$S$. Si $X \to S$ est localement de type fini, alors $F_{X/S}$ est fini.
\end{lemma}

\begin{proof}
Cela se traduit par le fait algébrique suivant. Supposons que $A
\to B$ soit un homomorphisme de type fini d'anneaux de
caractéristique $p$. Posons $B' = B \otimes_{A, F_A} A$ et
considérons l'homomorphisme d'anneaux $F_{B/A} : B' \to B$, $b
\otimes a \mapsto b^pa$. Alors notre assertion est que $F_{B/A}$ est
fini. En effet, si $x_1, \ldots, x_n \in B$ sont des générateurs sur
$A$, alors $x_i$ est intégral sur $B'$ parce que $x_i^p = F_{B/A}(x_i
\otimes 1)$. Par conséquent, $F_{B/A} : B' \to B$ est fini d'après le
lemme d'Algèbre \ref{algebra-lemma-characterize-finite-in-terms-of-integral}.
\end{proof}

\begin{lemma}
\label{lemma-geometrically-reduced-p}
Supposons que $k$ soit un corps de caractéristique $p > 0$. Supposons que $X$ soit un
schéma sur $k$. Alors $X$ est géométriquement réduit si et seulement si $X^{(p)}$ est réduit.
\end{lemma}

\begin{proof}
Considérons le Frobenius absolu $F_k : k \to k$. Alors
$F_k(k) = k^p$, en d'autres termes, $F_k : k \to k$ est
s'identifie à l'inclusion de $k$ dans $k^{1/p}$. Ainsi, le
lemme découle du lemme \ref{lemma-geometrically-reduced}.
\end{proof}

\begin{lemma}
\label{lemma-inseparable-deg-p-smooth}
Supposons que $k$ soit un corps de caractéristique $p > 0$. Supposons
que $X$ soit une variété sur $k$. Les énoncés suivants sont équivalents
\begin{enumerate}
\item $X^{(p)}$ est réduit, \item $X$ est
géométriquement réduit, \item il y a un
ouvert non vide $U \subset X$ lisse sur $k$.
\end{enumerate}
Dans ce cas, $X^{(p)}$ est une variété sur $k$ et $F_{X/k} : X
\to X^{(p)}$ est un morphisme dominant fini de degré $p^{\dim(X)}$.
\end{lemma}

\begin{proof}
L'équivalence de (1) et (2) résulte du
lemme \ref{lemma-geometrically-reduced-p}. Nous avons
vu dans le lemme
\ref{lemma-geometrically-reduced-dense-smooth-open} que (2) implique (3). Si
(3) est vrai, alors $U$ est géométriquement réduit (voir
par exemple le lemme
\ref{lemma-geometrically-regular-smooth}) et donc $X$ est
géométriquement réduit par le lemme
\ref{lemma-generic-points-geometrically-reduced}. On voit ainsi que (1), (2) et (3) sont équivalents.

\medskip\noindent
Supposons que (1), (2) et (3) soient vérifiés. Puisque
$F_{X/k}$ est un homéomorphisme (lemme
\ref{lemma-relative-frobenius}), on voit que $X^{(p)}$ est
une variété. Ensuite, $F_{X/k}$ est fini par le lemme
\ref{lemma-relative-frobenius-finite}. Il est dominant car il
est surjectif. Pour calculer le degré (Morphismes, définition
\ref{morphisms-definition-degree}), il suffit de calculer le
degré de $F_{U/k} : U \to U^{(p)}$ (comme $F_{U/k} = F_{X/k}|_U$
par le lemme
\ref{lemma-relative-frobenius-endomorphism-identity}). Après avoir
réduit un peu $U$, nous pouvons supposer qu'il existe un morphisme
étale $h : U \to \mathbf{A}^n_k$, voir Morphismes, lemme
\ref{morphisms-lemma-smooth-etale-over-affine-space}. On a nécessairement $n =
\dim(U)$ parce que $\mathbf{A}^n_k \to \Spec(k)$ est lisse de
dimension relative $n$, le morphisme étale $h$ est lisse de
dimension relative $0$, et $U \to \Spec(k)$ est lisse de dimension
relative $\dim(U)$ et les dimensions relatives s'additionnent
correctement (Morphismes, lemme
\ref{morphisms-lemma-composition-relative-dimension-d}). Remarquons que $h$ est un
morphisme génériquement fini et dominant de variétés, et donc $\deg(h)$ est défini. D'après le
lemme \ref{lemma-relative-frobenius-endomorphism-identity}, nous avons un diagramme commutatif
$$
\xymatrix{
U \ar[rr]_{F_{X/k}} \ar[d]_h & & U^{(p)} \ar[d]^{h^{(p)}} \\
\mathbf{A}^n_k \ar[rr]^{F_{\mathbf{A}^n_k/k}} & &
(\mathbf{A}^n_k)^{(p)}
}
$$
Puisque $h^{(p)}$ est un changement de base de $h$, il
est également étale et il s'ensuit que $h^{(p)}$ est
également un morphisme dominant génériquement fini de
variétés. Le degré de $h^{(p)}$ est le degré de
l'extension $k(X^{(p)})/k((\mathbf{A}^n_k)^{(p)})$, qui est
le même que le degré de l'extension
$k(X)/k(\mathbf{A}^n_k)$ parce que $h^{(p)}$ est le changement
de base de $h$ (un détail technique est omis). Par la multiplicativité des
degrés (Morphismes, lemme
\ref{morphisms-lemma-degree-composition}), il suffit de montrer que le
degré de $F_{\mathbf{A}^n_k/k}$ est $p^n$. Pour le voir, observons que
$(\mathbf{A}^n_k)^{(p)} = \mathbf{A}^n_k$ et que $F_{\mathbf{A}^n_k/k}$
est donné par l'application qui envoie les coordonnées sur leurs puissances $p$.
\end{proof}

\begin{remark}
\label{remark-n-fold-relative-frobenius}
Supposons que $p > 0$ soit un nombre premier. Supposons que $S$ soit un schéma
en caractéristique $p$. Supposons que $X$ soit un schéma sur $S$. Pour $n \geq 1$
$$
X^{(p^n)} = X^{(p^n/S)} = X \times_{S, F_S^n} S
$$
considéré comme un schéma sur $S$. Observons que $X \mapsto
X^{(p^n)}$ est un foncteur. En appliquant le lemme
\ref{lemma-frobenius-endomorphism-identity}, on voit que
$F_{X/S, n} = (F_X^n, \text{id}_S) : X \longrightarrow X^{(p^n)}$
est un morphisme sur $S$ s'insérant dans le diagramme commutatif
$$
\xymatrix{
X \ar[rr]_{F_{X/S, n}} \ar[rrd] \ar@/^1em/[rrrr]^{F_X^n}
& & X^{(p^n)} \ar[rr] \ar[d] & & X \ar[d] \\
& & S \ar[rr]^{F_S^n} & & S
}
$$
où le carré de droite est cartésien. Le
morphisme $F_{X/S, n}$ est parfois appelé le {\it
morphisme de Frobenius relatif itéré $n$ fois de
$X/S$ }. Cela a du sens car nous avons la formule
$$
F_{X/S, n} =
F_{X^{(p^{n - 1})}/S} \circ \ldots \circ F_{X^{(p)}/S} \circ F_{X/S}
$$
ce qui montre que $F_{X/S, n}$ est la composition
de $n$ morphismes de Frobenius relatifs. Puisque nous avons
$$
F_{X^{(p^m)}/S} = F_{X^{(p^{m - 1})}/S}^{(p)} = \ldots = F_{X/S}^{(p^m)}
$$
(détails omis) nous obtenons également que
$$
F_{X/S, n} =
F_{X/S}^{(p^{n - 1})} \circ \ldots \circ F_{X/S}^{(p)} \circ F_{X/S}
$$
\end{remark}




\section{Recollement d'anneaux de dimension un}
\label{section-glueing-dim-1}

\noindent
Cette section contient quelques préliminaires algébriques
pour démontrer qu'un ensemble fini de points de codimension
$1$ d'un schéma séparé est contenu dans un ouvert affine.

\begin{situation}
\label{situation-glue}
On se donne un diagramme commutatif d'anneaux
$$
\xymatrix{
A \ar[r] & K \\
R \ar[u] \ar[r] & B \ar[u]
}
$$
où $K$ est un corps et $A$, $B$ sont des sous-anneaux de $K$ avec
le corps des fractions $K$. Enfin, $R = A \times_K B = A \cap B$.
\end{situation}

\begin{lemma}
\label{lemma-glue-valuation-ring}
Dans la situation \ref{situation-glue}, supposons que $B$ soit un anneau de valuation.
Alors pour chaque élément inversible $u$ de $A$, soit $u \in R$ soit $u^{-1} \in R$.
\end{lemma}

\begin{proof}
En effet, si l'image $c$ de $u$ dans $K$ est
dans $B$, alors $u \in R$. Sinon, $c^{-1} \in
B$ (Algèbre, lemme
\ref{algebra-lemma-valuation-ring-x-or-x-inverse}) et $u^{-1} \in R$.
\end{proof}

\noindent
Le lemme suivant explique la signification de la
condition `` $A \otimes B \to K$ est surjective
'' qui revient assez souvent dans ce qui suit.

\begin{lemma}
\label{lemma-glue-separated}
Dans la situation \ref{situation-glue}, supposons que $A$ soit un anneau
noethérien de dimension $1$. Les assertions suivantes sont équivalentes
\begin{enumerate}
\item $A \otimes B \to K$ n'est pas surjective,
\item il existe un anneau de valuation discrète
$\mathcal{O} \subset K$ contenant à la fois $A$ et $B$.
\end{enumerate}
\end{lemma}

\begin{proof}
Il est clair que (2) implique (1). D'autre part, si $A \otimes B \to
K$ n'est pas surjective, alors l'image $C \subset K$ n'est pas un
corps, donc $C$ possède un idéal maximal non nul $\mathfrak m$.
Choisissons un anneau de valuation $\mathcal{O} \subset K$ dominant
$C_\mathfrak m$. D'après le lemme d'Algèbre
\ref{algebra-lemma-krull-akizuki} appliqué à $A \subset \mathcal{O}$, l'anneau
$\mathcal{O}$ est noethérien. Par conséquent, $\mathcal{O}$ est un anneau de valuation
discrète d'après le lemme d'Algèbre \ref{algebra-lemma-valuation-ring-Noetherian-discrete}.
\end{proof}

\begin{lemma}
\label{lemma-semi-local}
Dans la situation \ref{situation-glue}, supposons
\begin{enumerate}
\item $A$ est un anneau intègre semi-local noethérien de
dimension $1$, \item $B$ est un anneau de valuation discrète,
\end{enumerate}
Alors nous avons les deux possibilités suivantes
\begin{enumerate}
\item [(a)] Si $A^*$ n'est pas contenu dans $R$, alors
$\Spec(A) \to \Spec(R)$ et $\Spec(B) \to \Spec(R)$ sont
des immersions ouvertes recouvrant $\Spec(R)$ et $K = A
\otimes_R B$. \item [(b)] Si $A^*$ est contenu dans $R$,
alors $B$ domine l'un des anneaux locaux de $A$ en un
idéal maximal et $A \otimes B \to K$ n'est pas surjectif.
\end{enumerate}
\end{lemma}

\begin{proof}
L'hypothèse (a) implique qu'il existe une unité $u$ de $A$ dont
l'image dans $K$ appartienne à l'idéal maximal de $B$. Alors $u$
est un non-diviseur de zéro de $R$ et pour chaque $a \in A$ il
existe un $n$ tel que $u^n a \in R$. Il s'ensuit que $A = R_u$.

\medskip\noindent
Supposons que $\mathfrak m_A$ soit le radical de Jacobson de $A$.
Supposons que $x \in \mathfrak m_A$ soit un élément non nul. Puisque
$\dim(A) = 1$, on voit que $K = A_x$. Après avoir remplacé $x$ par $x^n
u^m$ pour certains $n \geq 1$ et $m \in \mathbf{Z}$, on peut supposer que
l'image de $x$ soit une unité de $B$. Pour chaque $b \in B$, le produit $x^nb$
appartient à l'image de $R$ pour un entier $n$ convenable. Ainsi $B = R_x$.

\medskip\noindent
Soit $z \in R$. Si $z \not \in \mathfrak m_A$ et si l'image de $z$ n'appartient pas
à $\mathfrak m_B$, alors $z$ est inversible. Ainsi, $x
+ u$ est inversible dans $R$. Par conséquent, $\Spec(R) = D(x) \cup
D(u)$. Nous avons vu ci-dessus que $D(u) = \Spec(A)$ et $D(x) = \Spec(B)$.

\medskip\noindent
Cas (b). Si $x \in \mathfrak m_A$, alors $1 + x$ est une
unité et donc $1 + x \in R$, c'est-à-dire $x \in R$. Ainsi,
on voit que $\mathfrak m_A \subset R \subset A$. En
fait, dans ce cas, $A$ est intégral sur $R$. En effet,
écrivons $A/\mathfrak m_A = \kappa_1 \times \ldots \times
\kappa_n$ comme un produit de corps. Disons que $x = (c_1,
\ldots, c_r, 0, \ldots, 0)$ est un élément avec $c_i \not = 0$. Alors
$$
x^2 - x(c_1, \ldots, c_r, 1, \ldots, 1)  = 0
$$
Puisque $R$ contient toutes les unités, on voit que
$A/\mathfrak m_A$ est intégral sur l'image de $R$ en lui,
et donc $A$ est intégral sur $R$. Il s'ensuit que $R \subset
A \subset B$ comme $B$ est intégralement clos. De plus, si
$x \in \mathfrak m_A$ est non nul, alors $K = A_x = \bigcup
x^{-n}A = \bigcup x^{-n}R$. Par conséquent $x^{-1} \not \in
B$, c'est-à-dire $x \in \mathfrak m_B$. Nous concluons
$\mathfrak m_A \subset \mathfrak m_B$. Ainsi, $A \cap
\mathfrak m_B$ est un idéal maximal de $A$, ce qui achève la démonstration.
\end{proof}

\begin{lemma}
\label{lemma-semi-local-dimension-one-conductor}
Supposons que $B$ soit un anneau intègre noethérien semi-local de
dimension $1$. Soit $B'$ la clôture
intégrale de $B$ dans son corps des fractions. Alors $B'$
est un anneau de Dedekind semi-local. Supposons que $x$
soit un élément non nul du radical de Jacobson de $B'$. Alors
pour chaque $y \in B'$ il existe un $n$ tel que $x^n y \in B$.
\end{lemma}

\begin{proof}
Supposons que $\mathfrak m_B$ soit le radical de Jacobson
de $B$. La structure de $B'$ résulte de l'Algèbre, lemme
\ref{algebra-lemma-integral-closure-Dedekind}. Étant
donné $x, y \in B'$ comme dans l'énoncé du lemme,
considérons le sous-anneau $B \subset A \subset B'$
engendré par $x$ et $y$. Alors $A$ est fini sur $B$
(Algèbre, lemme
\ref{algebra-lemma-characterize-finite-in-terms-of-integral}).
Puisque le corps des fractions de $B$ et $A$ est le même, nous
voyons que le module fini $A/B$ a son support dans l'ensemble des
points fermés de $B$. Ainsi $\mathfrak m_B^n A \subset B$ pour
un $n$ approprié. De plus, $\Spec(B') \to \Spec(A)$ est
surjectif (Algèbre, lemme
\ref{algebra-lemma-integral-overring-surjective}), donc $A$ est également
semi-local. Il en découle également que $x$ est dans le radical de Jacobson
$\mathfrak m_A$ de $A$. Notons que $\mathfrak m_A = \sqrt{\mathfrak m_B A}$.
Ainsi $x^m y \in \mathfrak m_B A$ pour un certain $m$. Puis $x^{nm} y \in B$.
\end{proof}

\begin{lemma}
\label{lemma-semi-local-both-side}
Dans la situation \ref{situation-glue}, supposons
\begin{enumerate}
\item $A$ est un anneau intègre semi-local noethérien de
dimension $1$, \item $B$ est un anneau intègre semi-local
noethérien de dimension $1$, \item $A \otimes B \to K$ est surjectif.
\end{enumerate}
Alors $\Spec(A) \to \Spec(R)$ et $\Spec(B) \to \Spec(R)$ sont des
immersions ouvertes couvrant $\Spec(R)$ et $K = A \otimes_R B$.
\end{lemma}

\begin{proof}
Cas particulier : $B$ est intégralement clos dans $K$.
Cela signifie que $B$ est un anneau de Dedekind (Algèbre,
lemme \ref{algebra-lemma-characterize-Dedekind}) d'où toutes
ses localisations en des idéaux maximaux sont des anneaux
de valuation discrète. Soient $\mathfrak m_1, \ldots,
\mathfrak m_r$ les idéaux maximaux de $B$. Posons
$$
R_1 = A \times_K B_{\mathfrak m_1}
$$
En observant que $A \otimes_{R_1} B_{\mathfrak m_1} \to K$ est surjectif,
nous concluons d'après le lemme \ref{lemma-semi-local} que $A$ et
$B_{\mathfrak m_1}$ définissent des sous-schémas ouverts couvrant $\Spec(R_1)$
et que $K = A \otimes_{R_1} B_{\mathfrak m_1}$. En particulier, $R_1$ est un
anneau noethérien semi-local de dimension $1$. Par induction, nous définissons
$$
R_{i + 1} = R_i \times_K B_{\mathfrak m_{i + 1}}
$$
pour $i = 1, \ldots, r - 1$. Observons que $R = R_r$ parce
que $B = B_{\mathfrak m_1} \cap \ldots \cap B_{\mathfrak
m_r}$ (voir Algèbre, lemme
\ref{algebra-lemma-normal-domain-intersection-localizations-height-1}).
Il découle de la procédure inductive que $R \to A$ définit
une immersion ouverte $\Spec(A) \to \Spec(R)$. D'autre part,
les idéaux maximaux $\mathfrak n_i$ de $R$ ne se trouvant pas
dans cet ouvert correspondent aux idéaux maximaux $\mathfrak
m_i$ de $B$ et en fait le homomorphisme d'anneaux $R \to B$
définit un isomorphisme $R_{\mathfrak n_i} \to B_{\mathfrak
m_i}$ (détails omis ; indication : à chaque étape nous avons ajouté
exactement un idéal maximal à $\Spec(R_i)$). Il s'ensuit que
$\Spec(B) \to \Spec(R)$ est une immersion ouverte, comme annoncé.

\medskip\noindent
Cas général. Soit $B' \subset K$ la clôture intégrale de $B$.
Voir le lemme \ref{lemma-semi-local-dimension-one-conductor}.
Alors le cas particulier s'applique à $R' = A \times_K B'$.
Choisissons $x \in R'$ qui n'appartient à aucun des idéaux maximaux
de $A$ et appartient à tous les idéaux maximaux de $B'$
(voir Algèbre, lemme \ref{algebra-lemma-chinese-remainder}).
D'après le lemme \ref{lemma-semi-local-dimension-one-conductor},
il existe un entier $n$ tel que $x^n \in R = A \times_K B$.
Remplaçons $x$ par $x^n$ afin que $x \in R$. Pour tout $y \in R'$, il existe
un entier $n$ tel que $x^n y \in R$. D'autre part,
il est clair que $R'_x = A$. Ainsi $R_x = A$.
En échangeant les rôles de $A$ et de $B$, on trouve aussi un $y \in R$
tel que $B = R_y$. Remarquons qu'en inversant à la fois $x$ et $y$, il
ne reste aucun idéal premier hormis $(0)$. Ainsi $K = R_{xy} = R_x \otimes_R R_y$.
Ceci achève la démonstration.
\end{proof}

\begin{lemma}
\label{lemma-glue-a-bunch-of-local-rings}
Soit $K$ un corps. Soient $A_1, \ldots, A_r \subset K$ des anneaux
semi-locaux noethériens de dimension $1$ et de corps des fractions $K$. Si
$A_i \otimes A_j \to K$ est surjectif pour tous $i \not = j$, alors
il existe un anneau intègre semi-local noethérien $A \subset K$
de dimension $1$, contenu dans $A_1, \ldots, A_r$, tel que
\begin{enumerate}
\item $A \to A_i$ induit une immersion ouverte $j_i : \Spec(A_i) \to \Spec(A)$,
\item $\Spec(A)$ est la réunion des ouverts $j_i(\Spec(A_i))$,
\item chaque point fermé de $\Spec(A)$ appartient à exactement un de ces
ouverts.
\end{enumerate}
\end{lemma}

\begin{proof}
On peut en effet prendre $A = A_1 \cap \ldots \cap A_r$. Remarquons d'abord que (3),
une fois (1) et (2) démontrés, résulte de l'hypothèse suivant laquelle
$A_i \otimes A_j \to K$ est surjectif : en effet, si
$\mathfrak m \in j_i(\Spec(A_i)) \cap j_j(\Spec(A_j))$, alors
$A_i \otimes A_j \to K$ se factorise par $A_\mathfrak m$.
Pour démontrer (1) et (2), raisonnons par récurrence sur $r$.
Si $r > 1$, l'hypothèse de récurrence donne (1) et (2) pour
$B = A_2 \cap \ldots \cap A_r$. Le
lemme \ref{lemma-semi-local-both-side} montre alors qu'ils valent pour
$A = A_1 \cap B$.
\end{proof}

\begin{lemma}
\label{lemma-create-globally-generated}
Soit $A$ un anneau intègre de corps des fractions $K$. Soient $B_1, \ldots, B_r \subset K$
des anneaux intègres semi-locaux noethériens de dimension $1$ dont les corps des
fractions sont égaux à $K$. Si $A \otimes B_i \to K$ est surjectif pour $i = 1, \ldots, r$,
alors il existe un $x \in A$ tel que $x^{-1}$ appartienne au radical de Jacobson de
$B_i$ pour $i = 1, \ldots, r$.
\end{lemma}

\begin{proof}
Notons $B_i'$ la clôture intégrale de $B_i$ dans $K$. Supposons qu'il existe un
$x \in A$ non nul tel que $x^{-1}$ appartienne au radical de Jacobson de $B'_i$ pour
$i = 1, \ldots, r$. Alors, d'après le
lemme \ref{lemma-semi-local-dimension-one-conductor},
quitte à remplacer $x$ par une puissance, on obtient $x^{-1} \in B_i$.
Comme $\Spec(B'_i) \to \Spec(B_i)$ est surjectif, il en résulte que
$x^{-1}$ appartient aussi au radical de Jacobson de $B_i$.
On peut donc supposer que chaque $B_i$ est un anneau de Dedekind semi-local.

\medskip\noindent
Si $B_i$ n'est pas local, supprimons $B_i$ de la liste et
remplaçons-le par la famille finie d'anneaux locaux $(B_i)_\mathfrak m$.
On peut ainsi supposer que $B_i$ est un anneau de valuation discrète pour
$i = 1, \ldots, r$.

\medskip\noindent
Soient $v_i : K \to \mathbf{Z}$, $i = 1, \ldots, r$,
les valuations discrètes correspondantes (voir
Algèbre, lemme \ref{algebra-lemma-characterize-Dedekind}).
Nous cherchons un $x \in A$ non nul tel que $v_i(x) < 0$ pour
$i = 1, \ldots, r$. Nous allons procéder par récurrence sur $r$.

\medskip\noindent
Si $r = 1$ et si l'assertion est fausse, alors $A \subset B$ et le morphisme
$A \otimes B \to K$ n'est pas surjectif, ce qui est contradictoire.

\medskip\noindent
Si $r > 1$, l'hypothèse de récurrence fournit un $x \in A$ non nul tel que
$v_i(x) < 0$ pour $i = 1, \ldots, r - 1$. Si $v_r(x) < 0$, le résultat est
acquis; on peut donc supposer que $v_r(x) \geq 0$. Le cas initial fournit un
$y \in A$ non nul tel que $v_r(y) < 0$. Quitte à remplacer $x$ par une
puissance, on peut supposer que $v_i(x) < v_i(y)$ pour $i = 1, \ldots, r - 1$.
Alors $x + y$ est l'élément recherché.
\end{proof}

\begin{lemma}
\label{lemma-power-equal}
Soit $A$ un anneau local noethérien de dimension $1$.
Soit $L = \prod A_\mathfrak p$, le produit portant sur
les idéaux premiers minimaux de $A$. Soient $a_1, a_2 \in \mathfrak m_A$
ayant même image dans $L$. Alors $a_1^n = a_2^n$ pour
un certain $n > 0$.
\end{lemma}

\begin{proof}
Écrivons $a_1 = a_2 + x$. Alors $x$ a une image nulle dans $L$. Par conséquent,
$x$ est un élément nilpotent de $A$, car $\bigcap \mathfrak p$
est le radical de $(0)$, et l'annulateur $I$ de $x$ contient
une puissance de l'idéal maximal, puisque $\mathfrak p \not \in V(I)$
pour tout idéal premier minimal. Écrivons $x^k = 0$ et $\mathfrak m^n \subset I$.
Alors
$$
a_1^{k + n} = a_2^{k + n} + {n + k \choose 1} a_2^{n + k - 1} x +
{n + k \choose 2} a_2^{n + k - 2} x^2 + \ldots +
{n + k \choose k - 1} a_2^{n + 1} x^{k - 1} = a_2^{n + k}
$$
car $a_2 \in \mathfrak m_A$.
\end{proof}

\begin{lemma}
\label{lemma-power-works}
Soit $A$ un anneau local noethérien de dimension $1$.
Soient $L = \prod A_\mathfrak p$ et $I = \bigcap \mathfrak p$,
le produit et l'intersection portant sur
les idéaux premiers minimaux de $A$. Soit $f \in L$ un élément
de la forme $f = i + a$, où $a \in \mathfrak m_A$ et
$i \in IL$. Alors une puissance de $f$ appartient à l'image de $A \to L$.
\end{lemma}

\begin{proof}
Comme $A$ est noethérien, on a $I^t = 0$ pour un certain $t > 0$.
Supposons que $f = a + i$ avec $i \in I^kL$.
Alors $f^n = a^n + na^{n - 1}i \bmod I^{k + 1}L$.
Il suffit donc de montrer que $na^{n - 1}i$ appartient à
l'image de $I^k \to I^kL$ pour un certain $n \gg 0$.
Pour cela, choisissons un $g \in A$ tel que $\mathfrak m_A = \sqrt{(g)}$
(Algèbre, lemme \ref{algebra-lemma-height-1}). Alors $L = A_g$, par exemple d'après
Algèbre, proposition \ref{algebra-proposition-dimension-zero-ring}.
D'autre part, il existe un $n$ tel que $a^n \in (g)$.
On peut donc chasser les dénominateurs des éléments de $L$
en multipliant par une puissance suffisamment grande de $a$.
\end{proof}

\begin{lemma}
\label{lemma-good-intersection}
Soit $A$ un anneau local noethérien de dimension $1$.
Soit $L = \prod A_\mathfrak p$, le produit portant sur
les idéaux premiers minimaux de $A$. Soit $K \to L$ un homomorphisme d'anneaux entier.
Alors il existe $a \in \mathfrak m_A$ et $x \in K$
ayant même image dans $L$ et tels que $\mathfrak m_A = \sqrt{(a)}$.
\end{lemma}

\begin{proof}
D'après le lemme \ref{lemma-power-works}, on peut remplacer $A$ par
$A/(\bigcap \mathfrak p)$ et supposer que $A$ est un anneau réduit
(certains détails sont omis).
On peut aussi remplacer $K$ par l'image de $K \to L$.
Alors $K$ est un anneau réduit. Le morphisme $\Spec(L) \to \Spec(K)$ est
surjectif et fermé (détails omis). Ainsi, $\Spec(K)$ est un espace fini
discret. Il en résulte que $K$ est un produit fini de corps.

\medskip\noindent
Soient $\mathfrak p_j$, $j = 1, \ldots, m$, les idéaux premiers minimaux de $A$.
Notons $L_j$ le corps des fractions de $A_j$, de sorte que
$L = \prod_{j = 1, \ldots, m} L_j$.
Soit $A_j$ la normalisation de $A/\mathfrak p_j$. Alors
$A_j$ est un anneau de Dedekind semi-local possédant au moins
un idéal maximal; voir
Algèbre, lemme \ref{algebra-lemma-integral-closure-Dedekind}.
Soit $n$ la somme des nombres d'idéaux maximaux
de $A_1, \ldots, A_m$. Pour chaque idéal maximal $\mathfrak m \subset A_j$ de ce type,
considérons l'application
$$
v_{\mathfrak m} : L \to L_j \to \mathbf{Z} \cup \{\infty\}
$$
où la seconde flèche est la valuation discrète associée
à l'anneau de valuation discrète $(A_j)_{\mathfrak m}$, prolongée
en envoyant $0$ sur $\infty$. On obtient ainsi $n$ applications
$v_1, \ldots, v_n : L \to \mathbf{Z} \cup \{\infty\}$.
Nous allons trouver un élément $x \in K$ tel que $v_i(x) < 0$
pour tout $i = 1, \ldots, n$.

\medskip\noindent
Affirmons d'abord que, pour chaque $i$, il existe un élément $x \in K$
tel que $v_i(x) < 0$. Supposons en effet que $v_i$ corresponde à
$\mathfrak m \subset A_j$. Si $v_i(x) \geq 0$ pour tout $x \in K$,
alors $K$ s'envoie dans $(A_j)_{\mathfrak m}$ à l'intérieur du corps des fractions
$L_j$ de $A_j$.
L'image de $K$ dans $L_j$ est un corps sur lequel $L_j$ est
algébrique, d'après Algèbre, lemme \ref{algebra-lemma-integral-under-field}.
Avec ce qui précède, cela contredit Algèbre, lemme
\ref{algebra-lemma-valuation-ring-cap-field-finite}.

\medskip\noindent
Supposons avoir trouvé un élément $x \in K$ tel que
$v_1(x) < 0, \ldots, v_r(x) < 0$ pour un certain $r < n$. Si $v_{r + 1}(x) < 0$,
alors $x$ convient pour $r + 1$. Sinon, choisissons un $y \in K$ tel que
$v_{r + 1}(y) < 0$, ce qui est possible par le paragraphe précédent.
Après avoir remplacé $x$ par $x^n$ pour un certain $n > 0$,
on peut supposer que $v_i(x) < v_i(y)$ pour $i = 1, \ldots, r$. Alors
$v_j(x + y) = v_j(x) < 0$ pour $j = 1, \ldots, r$, par les propriétés
des valuations, et de même $v_{r + 1}(x + y) = v_{r + 1}(y) < 0$.
En raisonnant par récurrence, on trouve
$x \in K$ tel que $v_i(x) < 0$ pour $i = 1, \ldots, n$.

\medskip\noindent
En particulier, l'élément $x \in K$ a une projection non nulle
dans chaque facteur de $K$ (rappelons que $K$ est un produit fini de
corps et que, si une composante de $x$ était nulle, l'une
des valeurs $v_i(x)$ serait égale à $\infty$). Ainsi, $x$ est
inversible, et $x^{-1} \in K$ est un élément tel que
$\infty > v_i(x^{-1}) > 0$ pour tout $i$. Il résulte du
lemme \ref{lemma-semi-local-dimension-one-conductor} que,
pour un certain $e < 0$, l'élément $x^e \in K$ s'envoie sur un élément de
$\mathfrak m_A/\mathfrak p_j \subset A/\mathfrak p_j$ pour tout
$j = 1, \ldots, m$. Remarquons que le conoyau de l'application
$\mathfrak m_A \to \prod \mathfrak m_A/\mathfrak p_j$ est annulé
par une puissance de $\mathfrak m_A$. Ainsi, quitte à remplacer
$e$ par un $e$ plus négatif, on trouve un élément $a \in \mathfrak m_A$
dont l'image dans $\mathfrak m_A/\mathfrak p_j$ est égale à
l'image de $x^e$. La paire $(a, x^e)$ satisfait les
conclusions du lemme.
\end{proof}

\begin{lemma}
\label{lemma-localization-semi-local}
Soit $A$ un anneau. Soient $\mathfrak p_1, \ldots, \mathfrak p_r$
un ensemble fini d'idéaux premiers de $A$. Posons
$S = A \setminus \bigcup \mathfrak p_i$. Alors $S$ est un système multiplicatif
et $S^{-1}A$ est un anneau semi-local dont les idéaux maximaux
correspondent aux éléments maximaux de l'ensemble $\{\mathfrak p_i\}$.
\end{lemma}

\begin{proof}
Si $a, b \in A$ et $a, b \in S$, alors $a, b \not \in \mathfrak p_i$,
donc $ab \not \in \mathfrak p_i$, et par suite $ab \in S$. De plus, $1 \in S$.
Ainsi, $S$ est une partie multiplicative de $A$. D'après la description de
$\Spec(S^{-1}A)$ donnée dans
Algèbre, lemme \ref{algebra-lemma-spec-localization},
et d'après
Algèbre, lemme \ref{algebra-lemma-silly},
les idéaux premiers de $S^{-1}A$ correspondent aux idéaux premiers de
$A$ contenus dans l'un des $\mathfrak p_i$.
Les idéaux maximaux de $S^{-1}A$ correspondent donc bijectivement aux
éléments maximaux (pour l'inclusion)\ parmi les éléments de l'ensemble
$\{\mathfrak p_1, \ldots, \mathfrak p_r\}$.
\end{proof}




\section{Schémas noethériens de dimension un}
\label{section-dimension-one}

\noindent
Le résultat principal de cette section affirme qu'un
schéma noethérien séparé de dimension $1$ possède un
faisceau inversible ample. Voir la
proposition \ref{proposition-dim-1-noetherian-separated-has-ample}.

\begin{lemma}
\label{lemma-affine}
Soit $X$ un schéma dont tous les anneaux locaux sont noethériens de dimension
$\leq 1$. Soit $U \subset X$ un ouvert rétrocompact. Notons
$j : U \to X$ le morphisme d'inclusion. Alors $R^pj_*\mathcal{F} = 0$, $p > 0$
pour tout $\mathcal{O}_U$-module quasi-cohérent $\mathcal{F}$.
\end{lemma}

\begin{proof}
Il suffit de vérifier l'annulation de $R^pj_*\mathcal{F}$ sur les germes.
La formation de $R^qj_*$ commute au changement de base plat; voir
Cohomologie des schémas, lemme
\ref{coherent-lemma-flat-base-change-cohomology}.
On peut donc supposer que $X$ est le spectre d'un anneau local noethérien
de dimension $\leq 1$. Dans ce cas, $X$ possède un point fermé
$x$ et un nombre fini d'autres points $x_1, \ldots, x_n$ qui se spécialisent
en $x$, mais non les uns dans les autres (voir
Algèbre, lemme \ref{algebra-lemma-Noetherian-irreducible-components}).
Si $x \in U$, alors $U = X$ et le résultat est clair. Sinon,
$U = \{x_1, \ldots, x_r\}$ pour un certain $r$, après avoir éventuellement renuméroté
les points. Alors $U$ est affine
(Schémas, lemme \ref{schemes-lemma-scheme-finite-discrete-affine}).
Le résultat découle donc de Cohomologie des schémas, lemme
\ref{coherent-lemma-relative-affine-vanishing}.
\end{proof}

\begin{lemma}
\label{lemma-open-in-affine-curve-affine}
Soit $X$ un schéma affine dont tous les anneaux locaux sont noethériens
de dimension $\leq 1$. Alors tout ouvert quasi-compact $U \subset X$ est affine.
\end{lemma}

\begin{proof}
Notons $j : U \to X$ le morphisme d'inclusion. Soit $\mathcal{F}$
un $\mathcal{O}_U$-module quasi-cohérent.
D'après le lemme \ref{lemma-affine}, les images directes supérieures
$R^pj_*\mathcal{F}$ sont nulles. Le $\mathcal{O}_X$-module $j_*\mathcal{F}$
est quasi-cohérent
(Schémas, lemme \ref{schemes-lemma-push-forward-quasi-coherent}).
Ses groupes de cohomologie supérieurs sont donc nuls d'après
Cohomologie des schémas, lemme
\ref{coherent-lemma-quasi-coherent-affine-cohomology-zero}.
La suite spectrale de Leray
Cohomologie, lemme \ref{cohomology-lemma-apply-Leray},
donne $H^p(U, \mathcal{F}) = 0$ pour tout $p > 0$.
Ainsi, $U$ est affine, par exemple d'après
Cohomologie des schémas, lemme
\ref{coherent-lemma-quasi-compact-h1-zero-covering}.
\end{proof}

\begin{lemma}
\label{lemma-complement-codim-1-closed-points}
Soit $X$ un schéma. Soit $U \subset X$ un ouvert. Supposons que
\begin{enumerate}
\item $U$ soit un ouvert rétrocompact de $X$,
\item $X \setminus U$ soit discret, et
\item pour $x \in X \setminus U$, l'anneau local
$\mathcal{O}_{X, x}$ soit noethérien de dimension $\leq 1$.
\end{enumerate}
Alors (1) il existe un $\mathcal{O}_X$-module inversible $\mathcal{L}$
et une section $s$ tels que $U = X_s$, et (2) le morphisme
$\Pic(X) \to \Pic(U)$ est surjectif.
\end{lemma}

\begin{proof}
Écrivons $X \setminus U = \{x_i; i \in I\}$.
Choisissons des ouverts affines $U_i \subset X$ tels que $x_i \in U_i$ et
$x_j \not \in U_i$ pour $j \not = i$. C'est possible par la condition (2).
Écrivons $U_i = \Spec(A_i)$. Soit $\mathfrak m_i \subset A_i$ l'idéal maximal
correspondant à $x_i$. Par notre hypothèse sur les anneaux locaux,
il n'y a qu'un nombre fini d'idéaux premiers
$\mathfrak q \subset \mathfrak m_i$,
$\mathfrak q \not = \mathfrak m_i$ (voir
Algèbre, lemme \ref{algebra-lemma-Noetherian-irreducible-components}).
Le lemme d'évitement des idéaux premiers (Algèbre, lemme
\ref{algebra-lemma-silly}) fournit donc un $f_i \in \mathfrak m_i$
qui n'appartient à aucun de ces idéaux premiers. Alors
$V(f_i) = \{\mathfrak m_i\} \amalg Z_i$ pour un certain sous-ensemble fermé
$Z_i \subset U_i$, car $Z_i$ est un ouvert rétrocompact de
$V(f_i)$ stable par spécialisation; voir
Algèbre, lemme \ref{algebra-lemma-constructible-stable-specialization-closed}.
Quitte à rétrécir $U_i$, on peut supposer que $V(f_i) = \{x_i\}$. Alors
$$
\mathcal{U} : X = U \cup \bigcup U_i
$$
est un recouvrement ouvert de $X$. Considérons le $2$-cocycle à valeurs
dans $\mathcal{O}_X^*$ donné par $f_i$ sur $U \cap U_i$ et par
$f_i/f_j$ sur $U_i \cap U_j$. Il définit un faisceau inversible
$\mathcal{L}$ tel que la section $s$, définie par $1$ sur $U$
et par $f_i$ sur $U_i$, possède les propriétés énoncées dans le lemme.

\medskip\noindent
Soit $\mathcal{N}$ un $\mathcal{O}_U$-module inversible.
Soit $N_i$ le $(A_i)_{f_i}$-module inversible tel que
$\mathcal{N}|_{U \cap U_i}$ soit égal à $\tilde N_i$.
Remarquons que $(A_{\mathfrak m_i})_{f_i}$ est un anneau artinien
(car il s'agit d'un anneau noethérien de dimension zéro; voir
Algèbre, lemme \ref{algebra-lemma-Noetherian-dimension-0}).
C'est donc un produit d'anneaux locaux
(Algèbre, lemme \ref{algebra-lemma-artinian-finite-length}) et
son groupe de Picard est par conséquent trivial. En rétrécissant $U_i$
(c'est-à-dire en remplaçant $A_i$ par $(A_i)_g$ pour un $g \in A_i$,
$g \not \in \mathfrak m_i$),
on peut donc supposer que $N_i = (A_i)_{f_i}$, c'est-à-dire que
$\mathcal{N}|_{U \cap U_i}$ est trivial. Dans ce cas, on voit clairement
comment prolonger $\mathcal{N}$ en un faisceau inversible sur $X$
(en le prolongeant par un module inversible trivial sur chaque $U_i$).
\end{proof}

\begin{lemma}
\label{lemma-find-globally-generated}
Soit $X$ un schéma intègre séparé. Soit $U \subset X$ un ouvert affine non vide
tel que $X \setminus U$ soit un ensemble fini de points
$x_1, \ldots, x_r$ et que $\mathcal{O}_{X, x_i}$ soit noethérien de dimension $1$.
Alors il existe un $\mathcal{O}_X$-module inversible engendré par ses sections globales
$\mathcal{L}$ et une section $s$ tels que $U = X_s$.
\end{lemma}

\begin{proof}
Écrivons $U = \Spec(A)$ et notons $K$ le corps des fonctions rationnelles de $X$.
Posons $B_i = \mathcal{O}_{X, x_i}$ et $\mathfrak m_i = \mathfrak m_{x_i}$.
Comme $x_i \not \in U$, l'ouvert
$U \times_X \Spec(B_i)$ de $\Spec(B_i)$ ne possède qu'un point, c'est-à-dire
$U \times_X \Spec(B_i) = \Spec(K)$.
Puisque $X$ est séparé, $\Spec(K)$ est un sous-schéma fermé
de $U \times \Spec(B_i)$, c'est-à-dire que le morphisme $A \otimes B_i \to K$ est surjectif.
D'après le lemme \ref{lemma-create-globally-generated}, il existe un
$f \in A$ non nul tel que $f^{-1} \in \mathfrak m_i$ pour $i = 1, \ldots, r$.
Choisissons des ouverts $x_i \in U_i \subset X$ tels que $f^{-1} \in \mathcal{O}(U_i)$.
Alors
$$
\mathcal{U} : X = U \cup \bigcup U_i
$$
est un recouvrement ouvert de $X$. Considérons le $2$-cocycle à valeurs
dans $\mathcal{O}_X^*$ donné par $f$ sur $U \cap U_i$ et par
$1$ sur $U_i \cap U_j$. Il définit un faisceau inversible $\mathcal{L}$
muni de deux sections :
\begin{enumerate}
\item une section $s$, définie par $1$ sur $U$
et par $f^{-1}$ sur $U_i$, qui possède les propriétés énoncées dans le lemme, et
\item une section $t$, définie par $f$ sur $U$ et
par $1$ sur $U_i$.
\end{enumerate}
Remarquons que $X_t \supset U_1 \cup \ldots \cup U_r$.
Ainsi, $s, t$ engendrent $\mathcal{L}$, ce qui démontre le lemme.
\end{proof}

\begin{lemma}
\label{lemma-enough-globally-generated-ample}
Soit $X$ un schéma quasi-compact. Si, pour tout $x \in X$,
il existe une paire $(\mathcal{L}, s)$ composée d'un faisceau inversible
$\mathcal{L}$ engendré par ses sections globales et d'une section globale $s$ telle que
$x \in X_s$ et que $X_s$ soit affine, alors $X$ possède un faisceau inversible
ample.
\end{lemma}

\begin{proof}
Comme $X$ est quasi-compact, on peut trouver une famille finie
$(\mathcal{L}_i, s_i)$, $i = 1, \ldots, n$,
de paires telles que $\mathcal{L}_i$ soit engendré par ses sections globales, que $X_{s_i}$
soit affine et que $X = \bigcup X_{s_i}$.
Toujours parce que $X$ est quasi-compact, on peut trouver, pour chaque $i$, une famille
finie de sections $t_{i, j}$ de $\mathcal{L}_i$, $j = 1, \ldots, m_i$,
telle que $X = \bigcup X_{t_{i, j}}$. Posons $t_{i, 0} = s_i$.
Considérons le faisceau inversible
$$
\mathcal{L} = \mathcal{L}_1 
\otimes_{\mathcal{O}_X} \ldots
\otimes_{\mathcal{O}_X} \mathcal{L}_n
$$
et les sections globales
$$
\tau_J = t_{1, j_1} \otimes \ldots \otimes t_{n, j_n}
$$
D'après Propriétés, lemme \ref{properties-lemma-affine-cap-s-open},
l'ouvert $X_{\tau_J}$ est affine dès que $j_i = 0$ pour un certain $i$.
Il est immédiat que ces ouverts recouvrent $X$.
Ainsi, $\mathcal{L}$ est ample par définition.
\end{proof}

\begin{lemma}
\label{lemma-dim-1-noetherian-integral-separated-has-ample}
Soit $X$ un schéma noethérien intègre séparé de dimension $1$.
Alors $X$ possède un faisceau inversible ample.
\end{lemma}

\begin{proof}
Choisissons un recouvrement ouvert affine $X = U_1 \cup \ldots \cup U_n$.
Comme $X$ est noethérien, chacun des ensembles $X \setminus U_i$ est fini.
Le lemme \ref{lemma-find-globally-generated}
fournit donc une paire $(\mathcal{L}_i, s_i)$
composée d'un faisceau inversible $\mathcal{L}_i$ engendré par ses sections globales
et d'une section globale $s_i$ telle que $U_i = X_{s_i}$.
On conclut que $X$ possède un faisceau inversible ample grâce au
lemme \ref{lemma-enough-globally-generated-ample}.
\end{proof}

\begin{lemma}
\label{lemma-surjective-pic-birational-finite}
Soit $f : X \to Y$ un morphisme fini de schémas. Supposons qu'il
existe un ouvert $V \subset Y$ tel que $f^{-1}(V) \to V$ soit un
isomorphisme et que $Y \setminus V$ soit un espace discret. Alors tout
$\mathcal{O}_X$-module inversible est l'image réciproque d'un
$\mathcal{O}_Y$-module inversible.
\end{lemma}

\begin{proof}
Nous utiliserons l'égalité $\Pic(X) = H^1(X, \mathcal{O}_X^*)$; voir
Cohomologie, lemme \ref{cohomology-lemma-h1-invertible}.
Considérons la suite spectrale de Leray associée au faisceau abélien $\mathcal{O}_X^*$
et à $f$; voir Cohomologie, lemme \ref{cohomology-lemma-Leray}.
Considérons le morphisme induit
$$
H^1(X, \mathcal{O}_X^*) \longrightarrow H^0(Y, R^1f_*\mathcal{O}_X^*)
$$
Diviseurs, lemme \ref{divisors-lemma-finite-trivialize-invertible-upstairs},
affirme précisément que ce morphisme est nul. La suite spectrale de Leray donne donc
$H^1(X, \mathcal{O}_X^*) = H^1(Y, f_*\mathcal{O}_X^*)$.
Considérons ensuite le morphisme
$$
f^\sharp : \mathcal{O}_Y^* \longrightarrow f_*\mathcal{O}_X^*
$$
Par hypothèse, le noyau et le conoyau de ce morphisme sont à support
dans le sous-ensemble fermé $T = Y \setminus V$ de $Y$.
Comme $T$ est un espace topologique discret par hypothèse,
les groupes de cohomologie supérieurs de tout faisceau abélien sur $Y$ à support
dans $T$ sont nuls (cela résulte de
Cohomologie, lemme \ref{cohomology-lemma-cohomology-and-closed-immersions},
Modules, lemme \ref{modules-lemma-i-star-exact}, et
du fait que $H^i(T, \mathcal{F}) = 0$ pour tout $i > 0$
et tout faisceau abélien $\mathcal{F}$ sur $T$).
En décomposant le morphisme affiché en suites exactes courtes
$$
0 \to \Ker(f^\sharp) \to \mathcal{O}_Y^* \to \Im(f^\sharp) \to 0,\quad
0 \to \Im(f^\sharp) \to f_*\mathcal{O}_X^* \to \Coker(f^\sharp) \to 0
$$
on conclut d'abord que $H^1(Y, \mathcal{O}_Y^*) \to H^1(Y, \Im(f^\sharp))$
est surjectif, puis que
$H^1(Y, \Im(f^\sharp)) \to H^1(Y, f_*\mathcal{O}_X^*)$ est surjectif.
En combinant tout ce qui précède, on voit que $H^1(Y, \mathcal{O}_Y^*) \to
H^1(X, \mathcal{O}_X^*)$ est surjectif, comme voulu.
\end{proof}

\begin{lemma}
\label{lemma-glue-invertible-sheaves}
Soit $X$ un schéma. Soient $Z_1, \ldots, Z_n \subset X$ des sous-schémas
fermés. Soit $\mathcal{L}_i$ un faisceau inversible sur $Z_i$.
Supposons que
\begin{enumerate}
\item $X$ soit réduit,
\item $X = \bigcup Z_i$ ensemblistement, et
\item $Z_i \cap Z_j$ soit un espace topologique discret pour $i \not = j$.
\end{enumerate}
Alors il existe un faisceau inversible $\mathcal{L}$ sur $X$ dont la restriction
à $Z_i$ est $\mathcal{L}_i$. De plus, si l'on se donne des sections
$s_i \in \Gamma(Z_i, \mathcal{L}_i)$ qui ne s'annulent pas aux
points de $Z_i \cap Z_j$, alors on peut choisir $\mathcal{L}$ de telle sorte
qu'il existe $s \in \Gamma(X, \mathcal{L})$ vérifiant
$s|_{Z_i} = s_i$ pour tout $i$.
\end{lemma}

\begin{proof}
L'existence de $\mathcal{L}$ peut se déduire du
lemme \ref{lemma-surjective-pic-birational-finite},
mais nous donnerons aussi une preuve directe et nous utiliserons
celle-ci pour établir l'assertion concernant les sections.
Posons $T = \bigcup_{i \not = j} Z_i \cap Z_j$. Comme $X$ est réduit, on a
$$
X \setminus T = \bigcup (Z_i \setminus T)
$$
en tant que schémas. L'hypothèse (3) implique que $T$ est une partie discrète de $X$.
Ainsi, pour chaque $t \in T$, il existe un ouvert $U_t \subset X$
tel que $t \in U_t$, mais $t' \not \in U_t$ si $t' \in T$, $t' \not = t$.
En restreignant $U_t$ au besoin, on peut supposer qu'il existe des isomorphismes
$\varphi_{t, i} : \mathcal{L}_i|_{U_t \cap Z_i} \to
\mathcal{O}_{U_t \cap Z_i}$. De plus, pour chaque $i$, choisissons un recouvrement ouvert
$$
Z_i \setminus T = \bigcup\nolimits_j U_{ij}
$$
tel qu'il existe des isomorphismes
$\varphi_{i, j} : \mathcal{L}_i|_{U_{ij}} \cong \mathcal{O}_{U_{ij}}$.
Observons que
$$
\mathcal{U} : X = \bigcup U_t \cup \bigcup U_{ij}
$$
est un recouvrement ouvert de $X$. Nous affirmons que les isomorphismes
$\varphi_{t, i}$ et $\varphi_{i, j}$ permettent de définir un $2$-cocycle à valeurs
dans $\mathcal{O}_X^*$ relativement à ce recouvrement, qui définit $\mathcal{L}$ comme
dans l'énoncé du lemme.

\medskip\noindent
Plus précisément, si $i \not = i'$, alors $U_{i, j} \cap U_{i', j'} = \emptyset$
et il n'y a rien à faire. Sur $U_{i, j} \cap U_{i, j'}$, on a
$\mathcal{O}_X(U_{i, j} \cap U_{i, j'}) =
\mathcal{O}_{Z_i}(U_{i, j} \cap U_{i, j'})$ d'après la première remarque de la preuve.
Ainsi, la fonction de transition de $\mathcal{L}_i$ (plus précisément
$\varphi_{i, j} \circ \varphi_{i, j'}^{-1}$) fournit la valeur de notre
cocycle sur cette intersection.
Sur $U_t \cap U_{i, j}$, on peut procéder de même.
Enfin, si $t \not = t'$, on a
$$
U_t \cap U_{t'} = \coprod (U_t \cap U_{t'}) \cap Z_i
$$
et, de plus, l'intersection $U_t \cap U_{t'} \cap Z_i$ est contenue
dans $Z_i \setminus T$. Par le même raisonnement que précédemment, on voit donc que
$$
\mathcal{O}_X(U_t \cap U_{t'}) =
\prod \mathcal{O}_{Z_i}(U_t \cap U_{t'} \cap Z_i)
$$
et l'on peut utiliser les fonctions de transition de $\mathcal{L}_i$ (plus précisément
$\varphi_{t, i} \circ \varphi_{t', i}^{-1}$) pour définir la valeur de
notre cocycle sur $U_t \cap U_{t'}$. Ceci achève la preuve de l'existence
de $\mathcal{L}$.

\medskip\noindent
Étant données des sections $s_i$ comme dans la dernière assertion du lemme, choisissons,
dans l'argument ci-dessus, $U_t$ de sorte que $s_i|_{U_t \cap Z_i}$ ne s'annule pas et
choisissons $\varphi_{t, i}$ de sorte que $\varphi_{t, i}(s_i|_{U_t \cap Z_i}) = 1$.
Alors l'emploi de $1$ sur $U_t$ et de $\varphi_{i, j}(s_i|_{U_{i, j}})$ sur
$U_{i, j}$ définit une section de $\mathcal{L}$ dont la restriction
est égale à $s_i$ sur $Z_i$.
\end{proof}

\begin{remark}
\label{remark-conductor}
Soit $A$ un anneau réduit. Soient $I, J$ des idéaux de $A$
tels que $V(I) \cup V(J) = \Spec(A)$. Posons $B = A/J$.
Alors $I \to IB$ est un isomorphisme de $A$-modules. En effet, on
a $IB = I + J/J = I/(I \cap J)$, et $I \cap J$ est nul car
$A$ est réduit et $\Spec(A) = V(I) \cup V(J) = V(I \cap J)$.
Ainsi, pour tout $A$-module projectif $P$, on a aussi $IP = I(P/JP)$.
\end{remark}

\begin{lemma}
\label{lemma-dim-1-noetherian-reduced-separated-has-ample}
Soit $X$ un schéma noethérien, réduit, séparé et de dimension $1$.
Alors $X$ possède un faisceau inversible ample.
\end{lemma}

\begin{proof}
Soient $Z_i$, $i = 1, \ldots, n$, les composantes irréductibles de $X$.
Nous les considérons comme des sous-schémas fermés réduits de $X$.
D'après le lemme \ref{lemma-dim-1-noetherian-integral-separated-has-ample},
il existe des faisceaux inversibles amples $\mathcal{L}_i$ sur $Z_i$.
Posons $T = \bigcup_{i \not = j} Z_i \cap Z_j$. Comme $X$ est noethérien
de dimension $1$, l'ensemble $T$ est fini et constitué de points fermés
de $X$. Pour chaque $i$, quitte à remplacer
$\mathcal{L}_i$ par une puissance, on peut choisir $s_i \in \Gamma(Z_i, \mathcal{L}_i)$
tel que $(Z_i)_{s_i}$ soit affine et contienne $T \cap Z_i$; voir
Propriétés, lemme \ref{properties-lemma-ample-finite-set-in-principal-affine}.

\medskip\noindent
D'après le lemme \ref{lemma-glue-invertible-sheaves}, il existe un faisceau inversible
$\mathcal{L}$ sur $X$ et $s \in \Gamma(X, \mathcal{L})$
tels que $(\mathcal{L}, s)|_{Z_i} = (\mathcal{L}_i, s_i)$.
Observons que $X_s$ contient $T$ et coïncide ensemblistement
avec la réunion des sous-schémas fermés affines $(Z_i)_{s_i}$. Il est donc affine d'après
Limites, lemme \ref{limits-lemma-affines-glued-in-closed-affine}.
Pour achever la preuve, il suffit de trouver, pour tout $x \in X$, $x \not \in T$,
un entier $m > 0$ et une section $t \in \Gamma(X, \mathcal{L}^{\otimes m})$
tels que $X_t$ soit affine et que $x \in X_t$. Puisque $x \not \in T$,
on voit que $x \in Z_i$ pour un unique $i$, disons $i = 1$.
Soit $Z \subset X$ le sous-schéma fermé réduit dont l'espace
topologique sous-jacent est $Z_2 \cup \ldots \cup Z_n$.
Soit $\mathcal{I} \subset \mathcal{O}_X$ le faisceau
d'idéaux de $Z$. Notons $\mathcal{I}_1 \subset \mathcal{O}_{Z_1}$
l'image inverse de ce faisceau d'idéaux par le morphisme
d'inclusion $Z_1 \to X$. Observons que
$$
\Gamma(X, \mathcal{I}\mathcal{L}^{\otimes m}) =
\Gamma(Z_1, \mathcal{I}_1 \mathcal{L}_1^{\otimes m})
$$
voir la remarque \ref{remark-conductor}. Il suffit donc de trouver $m > 0$
et $t \in \Gamma(Z_1, \mathcal{I}_1 \mathcal{L}_1^{\otimes m})$
tels que $x \in (Z_1)_t$, cet ouvert étant affine. Puisque $\mathcal{L}_1$ est ample
et que $x$ n'appartient pas à $Z_1 \cap T = V(\mathcal{I}_1)$,
on peut trouver une section
$t_1 \in \Gamma(Z_1, \mathcal{I}_1 \mathcal{L}_1^{\otimes m_1})$
telle que $x \in (Z_1)_{t_1}$; voir Propriétés, proposition
\ref{properties-proposition-characterize-ample}.
Puisque $\mathcal{L}_1$ est ample, on peut trouver une section
$t_2 \in \Gamma(Z_1, \mathcal{L}_1^{\otimes m_2})$
telle que $x \in (Z_1)_{t_2}$ et que $(Z_1)_{t_2}$ soit affine; voir
Propriétés, définition \ref{properties-definition-ample}.
Posons $m = m_1 + m_2$ et $t = t_1 t_2$. Alors
$t \in \Gamma(Z_1, \mathcal{I}_1 \mathcal{L}_1^{\otimes m})$
et $x \in (Z_1)_t$ par construction, tandis que
$(Z_1)_t$ est affine d'après Propriétés, lemme
\ref{properties-lemma-affine-cap-s-open}.
\end{proof}

\begin{lemma}
\label{lemma-lift-line-bundle-from-reduction-dimension-1}
Soit $i : Z \to X$ une immersion fermée de schémas.
Si l'espace topologique sous-jacent à $X$ est noethérien et si
$\dim(X) \leq 1$, alors $\Pic(X) \to \Pic(Z)$ est surjective.
\end{lemma}

\begin{proof}
Considérons la suite exacte courte
$$
0 \to (1 + \mathcal{I}) \cap \mathcal{O}_X^* \to
\mathcal{O}^*_X \to i_*\mathcal{O}^*_Z \to 0
$$
de faisceaux de groupes abéliens sur $X$, où $\mathcal{I}$
est le faisceau quasi-cohérent d'idéaux correspondant à $Z$.
Puisque $\dim(X) \leq 1$, on a $H^2(X, \mathcal{F}) = 0$
pour tout faisceau abélien $\mathcal{F}$; voir
Cohomologie, proposition \ref{cohomology-proposition-vanishing-Noetherian}.
Par conséquent, l'application $H^1(X, \mathcal{O}^*_X) \to H^1(X, i_*\mathcal{O}_Z^*)$
est surjective. D'après
Cohomologie, lemme \ref{cohomology-lemma-cohomology-and-closed-immersions},
on a $H^1(X, i_*\mathcal{O}_Z^*) = H^1(Z, \mathcal{O}_Z^*)$.
Cela démontre le lemme d'après
Cohomologie, lemme \ref{cohomology-lemma-h1-invertible}.
\end{proof}

\begin{proposition}
\label{proposition-dim-1-noetherian-separated-has-ample}
Soit $X$ un schéma noethérien séparé de dimension $1$.
Alors $X$ possède un faisceau inversible ample.
\end{proposition}

\begin{proof}
Soit $Z \subset X$ la réduction de $X$. D'après le
lemme \ref{lemma-dim-1-noetherian-reduced-separated-has-ample},
le schéma $Z$ possède un faisceau inversible ample.
Ainsi, d'après le lemme \ref{lemma-lift-line-bundle-from-reduction-dimension-1},
il existe un $\mathcal{O}_X$-module inversible $\mathcal{L}$
sur $X$ dont la restriction à $Z$ est ample.
Alors $\mathcal{L}$ est ample d'après
Cohomologie des schémas, lemme \ref{coherent-lemma-ample-on-reduction}.
\end{proof}

\begin{remark}
\label{remark-useless-generalization}
En fait, si $X$ est un schéma dont la réduction est un schéma noethérien
séparé de dimension $1$, alors $X$ possède un faisceau inversible
ample. La démonstration est la même que celle de la
proposition \ref{proposition-dim-1-noetherian-separated-has-ample},
à ceci près que l'on utilise
Limites, lemme \ref{limits-lemma-ample-on-reduction},
au lieu de
Cohomologie des schémas, lemme \ref{coherent-lemma-ample-on-reduction}.
\end{remark}

\noindent
Le lemme suivant vaut en fait pour les morphismes séparés quasi-finis,
comme le lecteur peut le constater en utilisant le théorème principal de Zariski
(Compléments sur les morphismes, lemme
\ref{more-morphisms-lemma-quasi-finite-separated-pass-through-finite})
et le lemme \ref{lemma-complement-codim-1-closed-points}.

\begin{lemma}
\label{lemma-surjection-on-pic-quasi-finite}
Soit $f : X \to Y$ un morphisme de schémas. Supposons que
$Y$ soit noethérien de dimension $\leq 1$, que $f$ soit fini et
qu'il existe un ouvert dense $V \subset Y$ tel que
$f^{-1}(V) \to V$ soit une immersion fermée. Alors tout
$\mathcal{O}_X$-module inversible est l'image réciproque d'un $\mathcal{O}_Y$-module inversible.
\end{lemma}

\begin{proof}
Factorisons $f$ en $X \to Z \to Y$, où $Z$ est l'image schématique
de $f$. Alors $X \to Z$ est un isomorphisme au-dessus de $V \cap Z$
et le lemme \ref{lemma-surjective-pic-birational-finite} s'applique.
D'autre part, le
lemme \ref{lemma-lift-line-bundle-from-reduction-dimension-1}
s'applique à $Z \to Y$. Nous omettons certains détails.
\end{proof}








\section{L'invariant delta}
\label{section-delta-invariant}

\noindent
Dans cette section, nous définissons l'invariant $\delta$ d'un point singulier
d'un schéma de Nagata réduit de dimension $1$.

\begin{lemma}
\label{lemma-pre-pre-delta-invariant}
Soit $(A, \mathfrak m)$ un anneau local noethérien de dimension $1$.
Soit $f \in \mathfrak m$. Les conditions suivantes sont équivalentes :
\begin{enumerate}
\item $\mathfrak m = \sqrt{(f)}$,
\item $f$ n'appartient à aucun idéal premier minimal de $A$, et
\item $A_f = \prod_{\mathfrak p\text{ minimal}} A_\mathfrak p$ comme $A$-algèbres.
\end{enumerate}
Un tel $f \in \mathfrak m$ existe. Si $\text{depth}(A) = 1$ (par exemple si
$A$ est réduit), alors (1) -- (3) sont aussi équivalentes à
\begin{enumerate}
\item[(4)] $f$ est un non-diviseur de zéro,
\item[(5)] $A_f$ est l'anneau total des fractions de $A$.
\end{enumerate}
Si $A$ est réduit, alors (1) -- (5) sont aussi équivalentes à
\begin{enumerate}
\item[(6)] $A_f$ est le produit des corps résiduels aux idéaux
premiers minimaux de $A$.
\end{enumerate}
\end{lemma}

\begin{proof}
Le spectre de $A$ ne possède qu'un nombre fini d'idéaux premiers
$\mathfrak p_1, \ldots, \mathfrak p_n$ autres que $\mathfrak m$,
et ceux-ci sont tous minimaux; voir
Algèbre, lemme \ref{algebra-lemma-Noetherian-irreducible-components}.
L'équivalence de (1) et (2) résulte alors de
Algèbre, lemme \ref{algebra-lemma-Zariski-topology}.
Il est clair que (3) implique (2). Réciproquement, si (2) est vraie,
alors le spectre de $A_f$ est la partie
$\{\mathfrak p_1, \ldots, \mathfrak p_n\}$ de $\Spec(A)$
munie de la topologie induite; voir
Algèbre, lemme \ref{algebra-lemma-spec-localization}.
C'est un espace topologique discret fini.
Par conséquent, $A_f = \prod_{\mathfrak p\text{ minimal}} A_\mathfrak p$ d'après
Algèbre, proposition \ref{algebra-proposition-dimension-zero-ring}.
L'existence d'un tel $f$ est affirmée dans
Algèbre, lemme \ref{algebra-lemma-height-1}.

\medskip\noindent
Supposons que $A$ soit de profondeur $1$. (C'est la valeur maximale d'après
Algèbre, lemme \ref{algebra-lemma-bound-depth}, et elle est atteinte si $A$ est réduit d'après
Algèbre, lemme \ref{algebra-lemma-criterion-reduced}.)
Alors $\mathfrak m$ n'est pas un idéal premier associé de $A$.
Tout idéal premier minimal de $A$ est un idéal premier associé
(Algèbre, proposition
\ref{algebra-proposition-minimal-primes-associated-primes}).
Par conséquent, les non-diviseurs de zéro de $A$ sont exactement les éléments
qui n'appartiennent à aucun des idéaux premiers minimaux, d'après
Algèbre, lemme \ref{algebra-lemma-ass-zero-divisors}.
Ainsi, (4) est équivalente à (2).
L'assertion (5) est équivalente à (3) d'après
Algèbre, lemme \ref{algebra-lemma-total-ring-fractions-no-embedded-points}.

\medskip\noindent
Alors $A_\mathfrak p$ est un corps pour tout idéal premier
$\mathfrak p \subset A$ minimal; voir
Algèbre, lemme \ref{algebra-lemma-minimal-prime-reduced-ring}.
Ainsi, (3) est équivalente à (6).
\end{proof}

\begin{lemma}
\label{lemma-pre-delta-invariant}
Soit $(A, \mathfrak m)$ un anneau local de Nagata réduit de dimension $1$.
Soit $A'$ la clôture intégrale de $A$ dans l'anneau total des fractions
de $A$. Alors $A'$ est un anneau de Nagata normal, $A \to A'$ est fini, et
$A'/A$ est un $A$-module de longueur finie.
\end{lemma}

\begin{proof}
L'anneau total des fractions est essentiellement de type fini sur $A$;
ainsi, $A \to A'$ est fini puisque $A$ est de Nagata; voir Algèbre, lemme
\ref{algebra-lemma-nagata-in-reduced-finite-type-finite-integral-closure}.
L'anneau $A'$ est normal, par exemple d'après
Algèbre, lemme \ref{algebra-lemma-characterize-reduced-ring-normal} et
\ref{algebra-lemma-Noetherian-irreducible-components}.
L'anneau $A'$ est de Nagata, par exemple d'après
Algèbre, lemme \ref{algebra-lemma-quasi-finite-over-nagata}.
Choisissons $f \in \mathfrak m$ comme dans le lemme \ref{lemma-pre-pre-delta-invariant}.
Comme $A' \subset A_f$, il est clair que $A_f = A'_f$. Le support du
$A$-module fini $A'/A$ est donc contenu dans $\{\mathfrak m\}$.
Il en résulte qu'il est de longueur finie d'après
Algèbre, lemme \ref{algebra-lemma-support-point}.
\end{proof}

\begin{definition}
\label{definition-delta-invariant-algebra}
Soit $A$ un anneau local de Nagata réduit de dimension $1$.
L'{\it invariant $\delta$ de $A$} est $\text{length}_A(A'/A)$,
où $A'$ est défini comme dans le lemme \ref{lemma-pre-delta-invariant}.
\end{definition}

\noindent
Nous démontrons quelques lemmes sur le comportement de cet invariant.

\begin{lemma}
\label{lemma-delta-invariant-is-zero}
Soit $A$ un anneau local de Nagata réduit de dimension $1$.
L'invariant $\delta$ de $A$ est $0$ si et seulement si
$A$ est un anneau de valuation discrète.
\end{lemma}

\begin{proof}
Si $A$ est un anneau de valuation discrète, alors $A$ est normal et
l'anneau $A'$ est égal à $A$. Réciproquement, si
l'invariant $\delta$ de $A$ est $0$, alors $A$ est intégralement
clos dans son anneau total des fractions, ce qui implique que
$A$ est normal
(Algèbre, lemme \ref{algebra-lemma-characterize-reduced-ring-normal}),
et cela force $A$ à être un anneau de valuation discrète d'après
Algèbre, lemme \ref{algebra-lemma-characterize-dvr}.
\end{proof}

\begin{lemma}
\label{lemma-normalization-same-after-completion}
Soit $A$ un anneau local de Nagata réduit de dimension $1$.
Soit $A \to A'$ comme dans le lemme \ref{lemma-pre-delta-invariant}.
Soient $A^h$, $A^{sh}$, resp.\ $A^\wedge$
l'hensélisé, l'hensélisé strict, resp.\ le complété de $A$.
Alors $A^h$, $A^{sh}$, resp. $A^\wedge$ est un anneau local de Nagata
réduit de dimension $1$, et
$A' \otimes_A A^h$, $A' \otimes_A A^{sh}$, resp. $A' \otimes_A A^\wedge$
est la clôture intégrale de $A^h$, $A^{sh}$, resp.\ $A^\wedge$
dans son anneau total des fractions.
\end{lemma}

\begin{proof}
Observons que $A^\wedge$ est réduit; voir
Compléments d'algèbre, lemme \ref{more-algebra-lemma-completion-reduced}.
Les anneaux $A^h$ et $A^{sh}$ sont réduits d'après
Compléments d'algèbre, lemme \ref{more-algebra-lemma-henselization-reduced}.
Les dimensions de $A$, $A^h$, $A^{sh}$ et $A^\wedge$ sont les mêmes
d'après Compléments d'algèbre, lemmes
\ref{more-algebra-lemma-completion-dimension} et
\ref{more-algebra-lemma-henselization-dimension}.

\medskip\noindent
Rappelons qu'un anneau local noethérien est de Nagata si et seulement si les
fibres formelles de $A$ sont géométriquement réduites; voir
Compléments d'algèbre, lemme \ref{more-algebra-lemma-Nagata-local-ring}.
Cette propriété se transmet à $A^h$ et à $A^{sh}$; voir
les résultats de Compléments d'algèbre, section
\ref{more-algebra-section-properties-formal-fibres},
et en particulier le lemme \ref{more-algebra-lemma-henselization-P-ring}.
Le complété est de Nagata d'après
Algèbre, lemme \ref{algebra-lemma-Noetherian-complete-local-Nagata}.

\medskip\noindent
Passons maintenant à l'assertion sur les clôtures intégrales. Avant de poursuivre,
choisissons $f \in \mathfrak m$ comme dans le
lemme \ref{lemma-pre-pre-delta-invariant}.
Alors l'image de $f$ dans $A^h$, $A^{sh}$ et $A^\wedge$
est manifestement un élément qui vérifie les propriétés (1) -- (6)
dans l'anneau considéré.

\medskip\noindent
Puisque $A \to A'$ est fini, on voit que
$A' \otimes_A A^h$ et $A' \otimes_A A^{sh}$
sont des produits d'anneaux locaux henséliens finis sur $A^h$ et
$A^{sh}$; voir Algèbre, lemme \ref{algebra-lemma-finite-over-henselian}.
Chacun de ces anneaux locaux est l'hensélisé de $A'$ en un
idéal maximal $\mathfrak m' \subset A'$ situé au-dessus de $\mathfrak m$; voir
Algèbre, lemme \ref{algebra-lemma-quasi-finite-henselization} ou
\ref{algebra-lemma-quasi-finite-strict-henselization}.
Ces anneaux locaux sont donc des anneaux intègres normaux d'après
Compléments d'algèbre, lemme \ref{more-algebra-lemma-henselization-normal}.
Il en résulte que $A' \otimes_A A^h$ et $A' \otimes_A A^{sh}$
sont des anneaux normaux. Puisque $A^h \to A' \otimes_A A^h$ et
$A^{sh} \to A' \otimes_A A^{sh}$ sont finis (donc entiers) et puisque
$A' \otimes_A A^h \subset (A^h)_f = Q(A^h)$ et
$A' \otimes_A A^{sh} \subset (A^{sh})_f = Q(A^{sh})$,
on conclut que $A' \otimes_A A^h$ et $A' \otimes_A A^{sh}$ sont
les clôtures intégrales cherchées.

\medskip\noindent
Pour le complété, raisonnons exactement de la même manière. Tout d'abord,
d'après Algèbre, lemme \ref{algebra-lemma-completion-finite-extension}, on a
$$
A' \otimes_A A^\wedge = (A')^\wedge =
\prod\nolimits (A'_{\mathfrak m'})^\wedge
$$
Les anneaux locaux $A'_{\mathfrak m'}$ sont normaux et de dimension $1$
(d'après Algèbre, lemme \ref{algebra-lemma-finite-in-codim-1}, par exemple, ou
la discussion de la
section Algèbre \ref{algebra-section-homomorphism-dimension}).
Ainsi, $A'_{\mathfrak m'}$ est un anneau de valuation discrète; voir
Algèbre, lemme \ref{algebra-lemma-characterize-dvr}.
Par conséquent, $(A'_{\mathfrak m'})^\wedge$ est un anneau de valuation discrète
d'après Compléments d'algèbre, lemme \ref{more-algebra-lemma-completion-dvr}.
Il en résulte que $A' \otimes_A A^\wedge$ est un anneau normal, et
on peut conclure exactement de la même manière que précédemment.
\end{proof}

\begin{lemma}
\label{lemma-delta-same-after-completion}
Soit $A$ un anneau local de Nagata réduit de dimension $1$.
L'invariant $\delta$ de $A$ est égal à
l'invariant $\delta$ de l'hensélisé, de l'hensélisé strict
ou du complété de $A$.
\end{lemma}

\begin{proof}
Traitons le cas du complété $B = A^\wedge$;
les autres cas se démontrent exactement de la même manière.
Soient $A'$, resp.\ $B'$, la clôture intégrale de $A$, resp.\ de $B$,
dans son anneau total des fractions.
Alors $B' = A' \otimes_A B$ d'après le
lemme \ref{lemma-normalization-same-after-completion}.
Par conséquent, $B'/B = (A'/A) \otimes_A B$.
L'égalité résulte maintenant du
Algèbre, lemme \ref{algebra-lemma-pullback-module},
et du fait que $B \otimes_A \kappa_A = \kappa_B$.
\end{proof}

\begin{definition}
\label{definition-delta-invariant}
Soit $k$ un corps. Soit $X$ un $k$-schéma localement algébrique.
Soit $x \in X$ un point tel que $\mathcal{O}_{X, x}$
soit réduit et que $\dim(\mathcal{O}_{X, x}) = 1$.
L'{\it invariant $\delta$ de $X$ en $x$} est
l'invariant $\delta$ de $\mathcal{O}_{X, x}$ défini dans la
définition \ref{definition-delta-invariant-algebra}.
\end{definition}

\noindent
Cette définition a un sens parce que l'anneau local d'un schéma localement
algébrique est de Nagata d'après
Algèbre, proposition \ref{algebra-proposition-ubiquity-nagata}.
Plus généralement, on peut bien entendu poser cette définition
chaque fois que $x \in X$ est un point d'un schéma tel que
l'anneau local $\mathcal{O}_{X, x}$ soit un anneau de Nagata réduit de dimension $1$.
Il résulte du lemme \ref{lemma-delta-same-after-completion}
que l'invariant $\delta$ de $X$ en $x$ est
$$
\delta\text{-invariant of }X\text{ at }x =
\delta\text{-invariant of }\mathcal{O}_{X, x}^h =
\delta\text{-invariant of }\mathcal{O}_{X, x}^\wedge
$$
On en conclut que l'invariant $\delta$ est un invariant
de l'anneau local complet du point.

\begin{lemma}
\label{lemma-delta-invariant-and-change-of-fields}
Soit $k$ un corps. Soit $X$ un $k$-schéma localement algébrique.
Soit $K/k$ une extension de corps et posons $Y = X_K$.
Soit $y \in Y$ d'image $x \in X$.
Supposons que $X$ soit géométriquement réduit en $x$ et que
$\dim(\mathcal{O}_{X, x}) = \dim(\mathcal{O}_{Y, y}) = 1$.
Alors
$$
\delta\text{-invariant of }X\text{ at }x \leq
\delta\text{-invariant of }Y\text{ at }y
$$
\end{lemma}

\begin{proof}
Posons $A = \mathcal{O}_{X, x}$ et $B = \mathcal{O}_{Y, y}$.
D'après le lemme \ref{lemma-geometrically-reduced-at-point},
on voit que $A$ est géométriquement réduit.
Ainsi, $B$ est une localisation de $A \otimes_k K$.
Soit $A \to A'$ comme dans le lemme \ref{lemma-pre-delta-invariant}.
Alors
$$
B' = B \otimes_{(A \otimes_k K)} (A' \otimes_k K)
$$
est fini sur $B$, et $B \to B'$ induit
un isomorphisme sur les anneaux totaux des fractions. En effet, choisissons $f \in \mathfrak m_A$
vérifiant (1) -- (6) du lemme \ref{lemma-pre-pre-delta-invariant};
puisque $\dim(B) = 1$, on voit que $f \in \mathfrak m_B$.
joue le même rôle pour $B$ et l'on voit que $B_f = B'_f$ puisque $A_f = A'_f$.
Soit $B''$ la clôture intégrale de $B$ dans son
anneau total des fractions, comme dans le lemme \ref{lemma-pre-delta-invariant}.
Alors $B' \subset B''$. Ainsi, l'invariant $\delta$ de $Y$ en $y$ est
$\text{length}_B(B''/B)$ et
\begin{align*}
\text{length}_B(B''/B)
& \geq
\text{length}_B(B'/B) \\
& =
\text{length}_B((A'/A) \otimes_A B) \\
& =
\text{length}_B(B/\mathfrak m_A B) \text{length}_A(A'/A)
\end{align*}
d'après Algèbre, lemme \ref{algebra-lemma-pullback-module},
puisque $A \to B$ est plat (étant une localisation de $A \to A \otimes_k K$).
Comme $\text{length}_A(A'/A)$ est l'invariant $\delta$ de $X$ en $x$
et que
$\text{length}_B(B/\mathfrak m_A B) \geq 1$, le lemme est démontré.
\end{proof}

\begin{lemma}
\label{lemma-delta-invariant-and-change-of-fields-better}
Soit $k$ un corps. Soit $X$ un $k$-schéma localement algébrique.
Soit $K/k$ une extension de corps, et posons $Y = X_K$.
Soit $y \in Y$, d'image $x \in X$.
Supposons que les hypothèses (a), (b), (c) du
lemme \ref{lemma-geometrically-normal-in-codim-1}
soient satisfaites pour $x \in X$ et que $\dim(\mathcal{O}_{Y, y}) = 1$.
Alors l'invariant $\delta$
de $X$ en $x$ est l'invariant $\delta$ de $Y$ en $y$.
\end{lemma}

\begin{proof}
Posons $A = \mathcal{O}_{X, x}$ et $B = \mathcal{O}_{Y, y}$.
D'après le lemme \ref{lemma-geometrically-normal-in-codim-1},
$A$ est géométriquement réduit.
Ainsi, $B$ est une localisation de $A \otimes_k K$.
Soit $A \to A'$ comme dans le lemme \ref{lemma-pre-delta-invariant}.
D'après le lemme \ref{lemma-geometrically-normal-in-codim-1},
$A' \otimes_k K$ est normal.
Par conséquent,
$$
B' = B \otimes_{(A \otimes_k K)} (A' \otimes_k K)
$$
est normal, fini sur $B$, et $B \to B'$ induit
un isomorphisme sur les anneaux totaux des fractions. En effet, choisissons $f \in \mathfrak m_A$
satisfaisant (1) -- (6) du lemme \ref{lemma-pre-pre-delta-invariant} ;
comme $\dim(B) = 1$, on a $f \in \mathfrak m_B$, qui
joue le même rôle pour $B$, et l'on a $B_f = B'_f$ puisque $A_f = A'_f$.
Il s'ensuit que $B \to B'$ est comme dans le lemme \ref{lemma-pre-delta-invariant}
pour $B$. Il faut donc montrer que
$\text{length}_A(A'/A) =
\text{length}_B(B'/B) = \text{length}_B((A'/A) \otimes_A B)$.
Comme $A \to B$ est plat (étant une localisation de $A \to A \otimes_k K$)
et que $\mathfrak m_B = \mathfrak m_A B$ (car
$B/\mathfrak m_A B$ est de dimension zéro d'après ce qui précède, et est
une localisation de $K \otimes_k \kappa(x)$, lequel est réduit puisque
$\kappa(x)$ est séparable sur $k$), on conclut par
Algèbre, lemme \ref{algebra-lemma-pullback-module}.
\end{proof}





\section{Le nombre de branches}
\label{section-number-of-branches}

\noindent
Nous avons défini le nombre de branches d'un schéma en un point
dans Propriétés, section \ref{properties-section-unibranch}.

\begin{lemma}
\label{lemma-number-of-branches}
Soit $X$ un schéma. Supposons que tout ouvert quasi-compact de $X$ possède
un nombre fini de composantes irréductibles. Soit $\nu : X^\nu \to X$
le morphisme de normalisation de $X$. Soit $x \in X$.
\begin{enumerate}
\item Le nombre de branches de $X$ en $x$ est le nombre des
antécédents de $x$ dans $X^\nu$.
\item Le nombre de branches géométriques de $X$ en $x$ est
$\sum_{\nu(x^\nu) = x} [\kappa(x^\nu) : \kappa(x)]_s$.
\end{enumerate}
\end{lemma}

\begin{proof}
Remarquons d'abord que l'hypothèse sur $X$ signifie exactement que la
normalisation est définie, voir
Morphismes, définition \ref{morphisms-definition-normalization}.
Alors le germe $A' = (\nu_*\mathcal{O}_{X^\nu})_x$ est la
clôture intégrale de $A = \mathcal{O}_{X, x}$ dans l'anneau total
des fractions de $A_{red}$, voir
Morphismes, lemme \ref{morphisms-lemma-stalk-normalization}.
Puisque $\nu$ est un morphisme intégral,
les points de $X^\nu$ situés au-dessus de $x$ correspondent
aux idéaux premiers de $A'$ situés au-dessus de l'idéal maximal $\mathfrak m$
de $A$. Comme $A \to A'$ est intégral, cela revient aux
idéaux maximaux de $A'$
(Algèbre, lemmes \ref{algebra-lemma-integral-no-inclusion} et
\ref{algebra-lemma-integral-going-up}).
Le lemme découle donc
de son analogue algébrique :
Compléments d'algèbre, lemme \ref{more-algebra-lemma-number-of-branches-1}.
\end{proof}

\begin{lemma}
\label{lemma-geometric-branches-and-change-of-fields}
Soit $k$ un corps. Soit $X$ un $k$-schéma localement algébrique.
Soit $K/k$ une extension de corps. Soit $y \in X_K$ un
point d'image $x$ dans $X$. Alors le nombre de
branches géométriques de $X$ en $x$ est le nombre de branches géométriques
de $X_K$ en $y$.
\end{lemma}

\begin{proof}
Écrivons $Y = X_K$ et soient $X^\nu$, resp.\ $Y^\nu$, les normalisés
de $X$, resp.\ $Y$. Considérons le diagramme commutatif
$$
\xymatrix{
Y^\nu \ar[r] \ar[d] & X^\nu_K \ar[r] \ar[d]_{\nu_K} & X^\nu \ar[d]_\nu \\
Y \ar@{=}[r] & Y \ar[r] & X
}
$$
D'après le lemme \ref{lemma-normalization-and-change-of-fields},
la flèche horizontale supérieure gauche est un homéomorphisme universel.
Elle induit donc des extensions purement inséparables des corps résiduels, voir
Morphismes, lemmes \ref{morphisms-lemma-universal-homeomorphism} et
\ref{morphisms-lemma-universally-injective}.
Ainsi, le nombre de branches géométriques de $Y$ en $y$ est
$\sum_{\nu_K(y') = y} [\kappa(y') : \kappa(y)]_s$
d'après le lemme \ref{lemma-number-of-branches}. De même,
$\sum_{\nu(x') = x} [\kappa(x') : \kappa(x)]_s$ est le nombre de
branches géométriques de $X$ en $x$. À l'aide de Schémas, lemme
\ref{schemes-lemma-points-fibre-product},
notre assertion découle du fait algébrique suivant :
étant données une extension de corps $l/\kappa$ et
une extension algébrique de corps $m/\kappa$, on a
$$
\sum\nolimits_{m \otimes_\kappa l \to m'} [m' : l']_s = [m : \kappa]_s
$$
où la somme porte sur les corps quotients de $m \otimes_\kappa l$.
On peut le démontrer de manière élémentaire, ou appliquer
le lemme \ref{lemma-separably-closed-field-connected-components}
à
$$
\Spec(m \otimes_\kappa l) \times_{\Spec(l)} \Spec(\overline{l}) =
\Spec(m) \otimes_{\Spec(\kappa)} \Spec(\overline{l})
\longrightarrow
\Spec(m) \times_{\Spec(\kappa)} \Spec(\overline{\kappa})
$$
car on peut interpréter $[m : \kappa]_s$ comme le nombre de
composantes connexes du membre de droite, et la somme
$\sum_{m \otimes_\kappa l \to m'} [m' : l']_s$
comme le nombre de composantes connexes du membre de gauche.
\end{proof}

\begin{lemma}
\label{lemma-geometrically-unibranch-and-change-of-fields}
Soit $k$ un corps. Soit $X$ un $k$-schéma localement algébrique.
Soit $K/k$ une extension de corps. Soit $y \in X_K$ un
point d'image $x$ dans $X$. Alors $X$ est géométriquement unibranche
en $x$ si et seulement si $X_K$ est géométriquement unibranche en $y$.
\end{lemma}

\begin{proof}
Cela découle immédiatement du
lemme \ref{lemma-geometric-branches-and-change-of-fields}
et de Compléments d'algèbre, lemme \ref{more-algebra-lemma-number-of-branches-1}.
\end{proof}

\begin{definition}
\label{definition-wedge}
Soient $A$ et $A_i$, $1 \leq i \leq n$, des anneaux locaux. On dit
que {\it $A$ est un bouquet de $A_1, \ldots, A_n$}
s'il existe des isomorphismes
$$
\kappa_{A_1} \to \kappa_{A_2} \to \ldots \to \kappa_{A_n}
$$
et si $A$ est isomorphe à l'anneau constitué des $n$-uplets
$(a_1, \ldots, a_n) \in A_1 \times \ldots \times A_n$ qui ont la
même image dans $\kappa_{A_n}$.
\end{definition}

\noindent
Si l'on se donne un anneau de base $\Lambda$ et que $A$ et $A_i$ sont des $\Lambda$-algèbres,
on exige alors que $\kappa_{A_i} \to \kappa_{A_{i + 1}}$
soient des isomorphismes de $\Lambda$-algèbres,
et que $A$ soit isomorphe, comme $\Lambda$-algèbre, à la $\Lambda$-algèbre
constituée des $n$-uplets
$(a_1, \ldots, a_n) \in A_1 \times \ldots \times A_n$ qui ont la
même image dans $\kappa_{A_n}$. En particulier, si $\Lambda = k$ est un corps
et que les applications $k \to \kappa_{A_i}$ sont des isomorphismes, il existe
un choix unique des isomorphismes
$\kappa_{A_i} \to \kappa_{A_{i + 1}}$, et
on parle souvent du {\it bouquet de $A_1, \ldots, A_n$}.

\begin{lemma}
\label{lemma-delta-number-branches-inequality-sh}
Soit $(A, \mathfrak m)$ un anneau local de Nagata réduit,
strictement hensélien et de dimension $1$. Alors
$$
\delta\text{-invariant of }A \geq \text{number of geometric branches of }A - 1
$$
Si l'égalité a lieu, alors $A$ est un bouquet de $n \geq 1$ anneaux locaux
de valuation discrète strictement henséliens.
\end{lemma}

\begin{proof}
Le nombre de branches géométriques est égal au nombre de branches de $A$
(cela découle immédiatement de
Compléments d'algèbre, définition \ref{more-algebra-definition-number-of-branches}).
Soit $A \to A'$ comme dans le lemme \ref{lemma-pre-delta-invariant}.
Remarquons que le nombre de branches de $A$ est le nombre
d'idéaux maximaux de $A'$, voir
Compléments d'algèbre, lemme \ref{more-algebra-lemma-number-of-branches-1}.
Il existe une surjection
$$
A'/A \longrightarrow
\left(\prod\nolimits_{\mathfrak m'} \kappa(\mathfrak m')\right)/
\kappa(\mathfrak m)
$$
Comme $\dim_{\kappa(\mathfrak m)} \prod \kappa(\mathfrak m')$
est $\geq$ au nombre de branches, l'inégalité est évidente.

\medskip\noindent
Si l'égalité a lieu, alors $\kappa(\mathfrak m') = \kappa(\mathfrak m)$
pour tout $\mathfrak m' \subset A'$, et la flèche affichée ci-dessus
est un isomorphisme. Comme $A$ est hensélien et que
$A \to A'$ est fini, $A'$ est un produit d'anneaux locaux
henséliens, voir Algèbre, lemme \ref{algebra-lemma-finite-over-henselian}.
Les facteurs sont les anneaux locaux $A'_{\mathfrak m'}$ et, comme
$A'$ est normal, ces facteurs sont des anneaux de valuation discrète
(Algèbre, lemme \ref{algebra-lemma-characterize-dvr}).
La flèche affichée étant un isomorphisme,
$A$ est bien le bouquet de ces anneaux locaux.
\end{proof}

\begin{lemma}
\label{lemma-delta-number-branches-inequality}
Soit $(A, \mathfrak m)$ un anneau local de Nagata réduit de dimension $1$. Alors
$$
\delta\text{-invariant of }A \geq \text{number of geometric branches of }A - 1
$$
\end{lemma}

\begin{proof}
On peut remplacer $A$ par l'hensélisé strict de $A$ sans
changer l'invariant $\delta$
(lemme \ref{lemma-delta-same-after-completion}) ni
le nombre de branches géométriques de $A$
(cela découle immédiatement de la définition, voir
Compléments d'algèbre, définition \ref{more-algebra-definition-number-of-branches}).
On peut donc supposer que $A$ est strictement hensélien, puis
appliquer le lemme \ref{lemma-delta-number-branches-inequality-sh}.
\end{proof}






\section{Normalisation des schémas de dimension un}
\label{section-normalization-one-dimensional}

\noindent
Le morphisme de normalisation d'un schéma noethérien de dimension
$1$ possède des propriétés remarquablement bonnes grâce au théorème de Krull--Akizuki.

\begin{lemma}
\label{lemma-normalize-noetherian-dim-1}
Soit $X$ un schéma localement noethérien de dimension $1$.
Soit $\nu : X^\nu \to X$ le morphisme de normalisation. Alors
\begin{enumerate}
\item $\nu$ est intégral, surjectif et induit une bijection
sur les composantes irréductibles,
\item il existe une factorisation $X^\nu \to X_{red} \to X$
et le morphisme $X^\nu \to X_{red}$ est la normalisation
de $X_{red}$,
\item $X^\nu \to X_{red}$ est birationnel,
\item pour tout point fermé $x \in X$, le germe
$(\nu_*\mathcal{O}_{X^\nu})_x$ est la clôture intégrale
de $\mathcal{O}_{X, x}$ dans l'anneau total des fractions
de $(\mathcal{O}_{X, x})_{red} = \mathcal{O}_{X_{red}, x}$,
\item les fibres de $\nu$ sont finies et les extensions de
corps résiduels sont finies,
\item $X^\nu$ est une union disjointe de schémas intègres normaux
localement noethériens, et chaque ouvert affine est le spectre d'un produit
fini d'anneaux de Dedekind.
\end{enumerate}
\end{lemma}

\begin{proof}
Beaucoup de ces résultats sont en fait des propriétés générales du morphisme
de normalisation, voir
Morphismes, lemmes \ref{morphisms-lemma-normalization-reduced},
\ref{morphisms-lemma-stalk-normalization},
\ref{morphisms-lemma-normalization-normal} et
\ref{morphisms-lemma-normalization-birational}.
Ce qui n'est pas clair, c'est que les fibres sont finies,
que les extensions induites des corps résiduels sont finies, et
que $X^\nu$ est localement le spectre d'un
anneau de Dedekind (et est donc noethérien).
Pour le voir, on peut supposer que $X = \Spec(A)$ est
affine, noethérien, de dimension $1$, et que $A$ est réduit.
La description donnée dans
Morphismes, lemme \ref{morphisms-lemma-description-normalization},
permet alors de se ramener au cas où $A$ est un anneau intègre noethérien
de dimension $1$. Dans ce cas, les propriétés voulues
découlent du théorème de Krull--Akizuki sous la forme énoncée dans
Algèbre, lemme \ref{algebra-lemma-integral-closure-Dedekind}.
\end{proof}

\noindent
Bien entendu, le lemme suivant admet une variante lorsque
$X$ n'est pas réduit.

\begin{lemma}
\label{lemma-prepare-delta-invariant}
Soit $X$ un schéma de Nagata réduit de dimension $1$. Soit $\nu : X^\nu \to X$
le morphisme de normalisation. Soit $x \in X$ un point fermé. Alors
\begin{enumerate}
\item $\nu : X^\nu \to X$ est fini, surjectif et birationnel,
\item $\mathcal{O}_X \subset \nu_*\mathcal{O}_{X^\nu}$ et
$\nu_*\mathcal{O}_{X^\nu}/\mathcal{O}_X$ est une somme directe de
faisceaux gratte-ciel, de valeur $\mathcal{Q}_x$ en $x$, laquelle est non nulle
si et seulement si $x$ est un point singulier de $X$,
\item $A' = (\nu_*\mathcal{O}_{X^\nu})_x$ est la clôture intégrale
de $A = \mathcal{O}_{X, x}$ dans son anneau total des fractions,
\item $\mathcal{Q}_x = A'/A$ est de longueur finie, égale à
l'invariant $\delta$ de $X$ en $x$,
\item $A'$ est un anneau semi-local qui est un produit fini
d'anneaux de Dedekind,
\item $A^\wedge$ est un anneau local complet noethérien
réduit, de dimension $1$,
\item $(A')^\wedge$ est la clôture intégrale de $A^\wedge$
dans son anneau total des fractions,
\item $(A')^\wedge$ est un produit fini d'anneaux
de valuation discrète complets, et
\item $A'/A \cong (A')^\wedge/A^\wedge$.
\end{enumerate}
\end{lemma}

\begin{proof}
Nous pouvons utiliser, et utiliserons, tous les résultats du
lemme \ref{lemma-normalize-noetherian-dim-1}.
La finitude de $\nu$ découle de
Morphismes, lemme \ref{morphisms-lemma-nagata-normalization}.
Comme $X$ est réduit, de Nagata, et de dimension $1$, le
l'ensemble des points réguliers est un ouvert dense $U \subset X$ d'après
Compléments d'algèbre, proposition \ref{more-algebra-proposition-ubiquity-J-2}.
Comme un schéma régulier est normal, cela montre que $\nu$
est un isomorphisme au-dessus de $U$.
Puisque $\dim(X) \leq 1$, l'ensemble des points au-dessus desquels $\nu$ n'est pas
un isomorphisme est un ensemble discret de points fermés $x \in X$.
En particulier, on obtient une suite exacte courte
$$
0 \to \mathcal{O}_X \to \nu_*\mathcal{O}_{X^\nu} \to
\bigoplus\nolimits_{x \in X \setminus U} \mathcal{Q}_x \to 0
$$
La description des germes de $\nu_*\mathcal{O}_{X^\nu}$
donnée par le lemme \ref{lemma-normalize-noetherian-dim-1} montre que
$Q_x = A'/A$ est bien de longueur égale à l'invariant $\delta$ de $X$ en $x$.
On a $Q_x \not = 0$ exactement lorsque $x$ est un point singulier, par
exemple d'après le lemme \ref{lemma-delta-invariant-is-zero}.
La description de $A'$ comme produit de domaines de Dedekind semi-locaux découle
elle aussi du lemme \ref{lemma-normalize-noetherian-dim-1}.
La relation entre $A$, $A'$ et $(A')^\wedge$
a été vue dans le lemme \ref{lemma-normalization-same-after-completion}
(et dans sa démonstration).
\end{proof}




\section{Recherche d'ouverts affines}
\label{section-finding-affine-opens}

\noindent
Nous poursuivons la discussion commencée dans
Propriétés, section \ref{properties-section-finding-affine-opens}.
Il s'avère que l'on peut trouver des ouverts affines contenant un ensemble fini donné
de points de codimension $1$ sur un schéma séparé. Voir
la proposition \ref{proposition-finite-set-of-points-of-codim-1-in-affine}.

\medskip\noindent
Nous améliorerons le lemme suivant dans
Descente, lemme \ref{descent-lemma-characterize-open-immersion}.

\begin{lemma}
\label{lemma-characterize-open-immersion}
Soit $f : X \to Y$ un morphisme de schémas. Soit $X^0$ l'ensemble
des points génériques des composantes irréductibles de $X$. Si
\begin{enumerate}
\item $f$ est séparé,
\item il existe un recouvrement ouvert $X = \bigcup U_i$ tel que
$f|_{U_i} : U_i \to Y$ soit une immersion ouverte, et
\item si $\xi, \xi' \in X^0$, $\xi \not = \xi'$, alors $f(\xi) \not = f(\xi')$,
\end{enumerate}
alors $f$ est une immersion ouverte.
\end{lemma}

\begin{proof}
Supposons que $y = f(x) = f(x')$. Choisissons une spécialisation $y_0 \leadsto y$
où $y_0$ est un point générique d'une composante irréductible de $Y$.
Comme $f$ est localement, sur la source, un isomorphisme, on peut choisir des spécialisations
$x_0 \leadsto x$ et $x'_0 \leadsto x'$ qui s'envoient sur $y_0 \leadsto y$.
Remarquons que $x_0, x'_0 \in X^0$. L'hypothèse (3) donne donc $x_0 = x'_0$.
Comme $f$ est séparé, on en conclut que $x = x'$. Ainsi, $f$ est une immersion ouverte.
\end{proof}

\begin{lemma}
\label{lemma-local-isomorphism}
Soit $X \to S$ un morphisme de schémas. Soit $x \in X$ un point
d'image $s \in S$. Si
\begin{enumerate}
\item $\mathcal{O}_{X, x} = \mathcal{O}_{S, s}$,
\item $S$ est réduit,
\item $X \to S$ est de type fini, et
\item $S$ possède un nombre fini de composantes irréductibles,
\end{enumerate}
alors il existe un voisinage ouvert $U$
de $x$ tel que $f|_U$ soit une immersion ouverte.
\end{lemma}

\begin{proof}
On peut supprimer les composantes irréductibles, en nombre fini, de $S$
qui ne contiennent pas $s$. On peut remplacer $S$ par un voisinage ouvert affine
de $s$. On peut remplacer $X$ par un voisinage ouvert affine
de $x$. Écrivons $S = \Spec(A)$ et $X = \Spec(B)$. Soient $\mathfrak q \subset B$,
resp.\ $\mathfrak p \subset A$, les idéaux premiers correspondant à $x$,
resp.\ $s$. Comme $A$ est réduit et que tous les idéaux premiers minimaux de
$A$ sont contenus dans $\mathfrak p$, on a $A \subset A_\mathfrak p$.
Comme $X \to S$ est de type fini, $B$ est de type fini sur $A$.
Soient $b_1, \ldots, b_n \in B$ des éléments qui engendrent $B$ sur $A$.
Puisque $A_\mathfrak p = B_\mathfrak q$, il existe
$f \in A$, $f \not \in \mathfrak p$,
et $a_i \in A$ tels que $b_i$ et $a_i/f$ aient la même image
dans $B_\mathfrak q$. Il existe donc $g \in B$, $g \not \in \mathfrak q$,
tel que $g(fb_i - a_i) = 0$ dans $B$. Il s'ensuit que l'image de
$A_f \to B_{fg}$ contient les images de $b_1, \ldots, b_n$, et
en particulier l'image de $g$.
Choisissons $n \geq 0$ et $f' \in A$ tels que $f'/f^n$ ait
pour image l'image de $g$ dans $B_{fg}$. Puisque $A_\mathfrak p = B_\mathfrak q$,
on a $f' \not \in \mathfrak p$.
On en conclut que $A_{ff'} \to B_{fg}$ est surjective.
Enfin, comme $A_{ff'} \subset A_\mathfrak p = B_\mathfrak q$ (voir ci-dessus),
l'application $A_{ff'} \to B_{fg}$ est injective, donc bijective.
\end{proof}

\begin{lemma}
\label{lemma-points-in-affine}
Soit $f : T \to X$ un morphisme de schémas. Soient $X^0$, resp.\ $T^0$,
les ensembles des points génériques des composantes irréductibles.
Soit $t_1, \ldots, t_m \in T$ un ensemble fini de points
d'images $x_j = f(t_j)$. Si
\begin{enumerate}
\item $T$ est affine,
\item $X$ est quasi-séparé,
\item $X^0$ est fini,
\item $f(T^0) \subset X^0$ et $f : T^0 \to X^0$ est injective, et
\item $\mathcal{O}_{X, x_j} = \mathcal{O}_{T, t_j}$,
\end{enumerate}
alors il existe un ouvert affine de $X$ contenant $x_1, \ldots, x_r$.
\end{lemma}

\begin{proof}
À l'aide de Limites, proposition \ref{limits-proposition-affine},
on se ramène immédiatement au cas où $X$ et $T$ sont réduits.
Nous omettons les détails.

\medskip\noindent
Supposons $X$ et $T$ réduits. On peut écrire $T = \lim_{i \in I} T_i$
comme limite projective de schémas de présentation finie sur $X$
à morphismes de transition affines, voir
Limites, lemme \ref{limits-lemma-relative-approximation}.
Choisissons $i \in I$ tel que $T_i$ soit affine, voir
Limites, lemme \ref{limits-lemma-limit-affine}.
Écrivons $T_i = \Spec(R_i)$ et $T = \Spec(R)$.
Soit $R' \subset R$ l'image de $R_i \to R$.
Alors $T' = \Spec(R')$ est affine, réduit, de type fini sur $X$,
et $T \to T'$ est dominant. Pour $j = 1, \ldots, r$, soit $t'_j \in T'$
l'image de $t_j$. Considérons les homomorphismes d'anneaux locaux
$$
\mathcal{O}_{X, x_j} \to
\mathcal{O}_{T', t'_j} \to
\mathcal{O}_{T, t_j}
$$
Notons $(T')^0$ l'ensemble des points génériques des composantes irréductibles
de $T'$. Soit $\xi \leadsto t'_j$ une spécialisation telle que
$\xi \in (T')^0$. Comme $T \to T'$ est dominant, on peut choisir $\eta \in T^0$ 
d'image $\xi$ (attention : a priori, on ne sait pas si $\eta$ se spécialise
en $t_j$). L'hypothèse (3), appliquée à $\eta$, dit que l'image $\theta$
de $\xi$ dans $X$ correspond à un idéal premier minimal de $\mathcal{O}_{X, x_j}$.
En relevant $\xi$ par l'isomorphisme de (5), on obtient une spécialisation
$\eta' \leadsto t_j$ avec $\eta' \in T^0$ d'image $\theta \leadsto x_j$.
L'injectivité de (4) montre que $\eta = \eta'$. Ainsi,
tout idéal premier minimal de $\mathcal{O}_{T', t'_j}$ est contenu dans
un idéal premier minimal de $\mathcal{O}_{T, t_j}$. On en conclut que
$\mathcal{O}_{T', t'_j} \to \mathcal{O}_{T, t_j}$ est injective,
donc les deux homomorphismes ci-dessus sont des isomorphismes.

\medskip\noindent
D'après le lemme \ref{lemma-local-isomorphism}, il existe un ouvert
$U \subset T'$ contenant tous les points $t'_j$ tel que
$U \to X$ soit un isomorphisme local comme dans le
lemme \ref{lemma-characterize-open-immersion}.
D'après ce lemme, $U \to X$ est une immersion ouverte.
Enfin, d'après
Propriétés, lemme \ref{properties-lemma-ample-finite-set-in-affine},
il existe un ouvert $W \subset U \subset T'$
contenant tous les $t'_j$. L'image de $W$ dans $X$ est
l'ouvert affine recherché.
\end{proof}

\begin{lemma}
\label{lemma-finite-set-codim-1-points-in-affine}
Soit $X$ un schéma intègre séparé. Soit $x_1, \ldots, x_r \in X$
un ensemble fini de points tel que l'anneau local $\mathcal{O}_{X, x_i}$
soit noethérien de dimension $\leq 1$. Alors il existe un sous-schéma
ouvert affine de $X$ contenant tous les $x_1, \ldots, x_r$.
\end{lemma}

\begin{proof}
Soit $K$ le corps des fonctions rationnelles de $X$.
Posons $A_i = \mathcal{O}_{X, x_i}$. Alors $A_i \subset K$ et $K$
est le corps des fractions de $A_i$. Puisque $X$ est séparé et que
$x_i \not = x_j$, il ne peut exister d'anneau de valuation $\mathcal{O} \subset K$
dominant à la fois $A_i$ et $A_j$. En effet, si l'on considère le
diagramme
$$
\xymatrix{
\Spec(\mathcal{O}) \ar[r] \ar[d] & \Spec(A_1) \ar[d] \\
\Spec(A_2) \ar[r] & X
}
$$
et que l'on applique le critère valuatif de séparation
(Schémas, lemme \ref{schemes-lemma-separated-implies-valuative}),
on obtiendrait $x_i = x_j$. Ainsi,
le lemme \ref{lemma-glue-separated} montre
que $A_i \otimes A_j \to K$ est surjectif pour tout $i \not = j$.
D'après le lemme \ref{lemma-glue-a-bunch-of-local-rings},
$A = A_1 \cap \ldots \cap A_r$ est un anneau noethérien
semi-local possédant exactement $r$ idéaux maximaux
$\mathfrak m_1, \ldots, \mathfrak m_r$ tels que $A_i = A_{\mathfrak m_i}$.
De plus,
$$
\Spec(A) = \Spec(A_1) \cup \ldots \cup \Spec(A_r)
$$
est un recouvrement ouvert, et l'intersection de deux membres distincts de ce
recouvrement est $\Spec(K)$. Par conséquent, les morphismes $\Spec(A_i) \to X$
donnés se recollent en un morphisme de schémas
$$
\Spec(A) \longrightarrow X
$$
qui envoie $\mathfrak m_i$ sur $x_i$ et induit des isomorphismes d'anneaux locaux.
Le résultat découle donc du lemme \ref{lemma-points-in-affine}.
\end{proof}

\begin{lemma}
\label{lemma-extra-silly}
Soient $A$ un anneau, $I \subset A$ un idéal,
$\mathfrak p_1, \ldots, \mathfrak p_r$ des idéaux premiers de $A$, et
$\overline{f} \in A/I$ un élément. Si $I \not \subset \mathfrak p_i$
pour tout $i$, alors il existe $f \in A$, $f \not \in \mathfrak p_i$,
qui s'envoie sur $\overline{f}$ dans $A/I$.
\end{lemma}

\begin{proof}
On peut supposer qu'il n'y a aucune relation d'inclusion entre les $\mathfrak p_i$
(quitte à supprimer les idéaux premiers les plus petits). Choisissons d'abord un $f \in A$ relevant
$\overline{f}$. Soit $S$ l'ensemble des $s \in \{1, \ldots, r\}$ tels
que $f \in \mathfrak p_s$. Si $S$ est vide, la démonstration est terminée. Sinon,
considérons l'idéal $J = I \prod_{i \not \in S} \mathfrak p_i$.
Notons que $J$ n'est contenu dans aucun $\mathfrak p_s$ pour $s \in S$,
car il n'y a aucune relation d'inclusion entre les $\mathfrak p_i$ et
que $I$ n'est contenu dans aucun $\mathfrak p_i$.
On peut donc choisir $g \in J$, $g \not \in \mathfrak p_s$ pour
$s \in S$, d'après Algèbre, lemme \ref{algebra-lemma-silly}.
Alors $f + g$ répond aux conditions du lemme.
\end{proof}

\begin{lemma}
\label{lemma-finite-set-codim-1-points-in-affine-per-component}
Soit $X$ un schéma. Soit $T \subset X$ un ensemble fini de points. Supposons que
\begin{enumerate}
\item $X$ possède un nombre fini de composantes irréductibles $Z_1, \ldots, Z_t$, et
\item $Z_i \cap T$ soit contenu dans un ouvert affine du sous-schéma
réduit induit correspondant à $Z_i$.
\end{enumerate}
Alors il existe un sous-schéma ouvert affine de $X$ contenant $T$.
\end{lemma}

\begin{proof}
La proposition \ref{limits-proposition-affine} de Limites
permet de se ramener immédiatement au cas où $X$
est réduit. Nous omettons les détails. Dans la suite de la démonstration, nous
munissons tout sous-ensemble fermé de $X$ de la structure induite de sous-schéma
fermé réduit.

\medskip\noindent
Nous allons montrer par récurrence que l'on peut trouver un ouvert affine
$U \subset Z_1 \cup \ldots \cup Z_r$ contenant
$T \cap (Z_1 \cup \ldots \cup Z_r)$. Pour $r = 1$, cela résulte de l'hypothèse.
Soit $r > 1$, et soit $U \subset Z_1 \cup \ldots \cup Z_{r -  1}$
un ouvert affine contenant $T \cap (Z_1 \cup \ldots \cup Z_{r - 1})$.
Soit $V \subset X_r$ un ouvert affine contenant $T \cap Z_r$ (il existe par
hypothèse). Alors $U \cap V$ contient
$T \cap ( Z_1 \cup \ldots \cup Z_{r - 1} ) \cap Z_r$.
Par conséquent,
$$
\Delta = (U \cap Z_r) \setminus (U \cap V)
$$
ne contient aucun élément de $T$. Notons que $\Delta$ est un sous-ensemble fermé
de $U$. Par évitement des idéaux premiers (Algèbre, lemme \ref{algebra-lemma-silly}),
on peut trouver un ouvert principal $U'$ de $U$ contenant $T \cap U$ et évitant
$\Delta$, c'est-à-dire tel que $U' \cap Z_r \subset U \cap V$.
Après avoir remplacé $U$ par $U'$, on peut supposer que $U \cap V$ est
fermé dans $U$.

\medskip\noindent
Le même raisonnement montre en outre que l'ensemble
$\Delta' = (U \cap (Z_1 \cup \ldots \cup Z_{r - 1})) \setminus (U \cap V)$
ne contient aucun élément de $T$; il existe donc $h \in \mathcal{O}(V)$
tel que $D(h) \subset V$ contienne $T \cap V$ et que
$U \cap D(h) \subset U \cap V$. Comme $U \cap V$ est fermé dans $U$,
on peut appliquer le lemme \ref{lemma-extra-silly}
pour trouver un élément $g \in \mathcal{O}(U)$ dont la restriction à $U \cap V$
est égale à la restriction de $h$ à $U \cap V$ et tel que
$T \cap U \subset D(g)$. On peut alors remplacer $U$ par $D(g)$
et $V$ par $D(h)$, de manière que $U \cap V$ soit
fermé à la fois dans $U$ et dans $V$. Dans ce cas, le schéma $U \cup V$ est
affine d'après Limites, lemme \ref{limits-lemma-affines-glued-in-closed-affine}.
Cela établit l'étape de récurrence, et donc le lemme.
\end{proof}

\noindent
Voici une conséquence des résultats qui précèdent.

\begin{proposition}
\label{proposition-finite-set-of-points-of-codim-1-in-affine}
Soit $X$ un schéma séparé tel que tout ouvert quasi-compact
possède un nombre fini de composantes irréductibles. Soient $x_1, \ldots, x_r \in X$
des points tels que l'anneau local $\mathcal{O}_{X, x_i}$ soit noethérien
de dimension $\leq 1$. Alors il existe un sous-schéma ouvert affine
de $X$ contenant tous les $x_1, \ldots, x_r$.
\end{proposition}

\begin{proof}
On peut remplacer $X$ par un ouvert quasi-compact contenant $x_1, \ldots, x_r$;
on peut donc supposer que $X$ possède un nombre fini de composantes irréductibles.
Le lemme \ref{lemma-finite-set-codim-1-points-in-affine-per-component}
ramène au cas où $X$ est intègre. Ce cas est traité par le
lemme \ref{lemma-finite-set-codim-1-points-in-affine}.
\end{proof}





\section{Courbes}
\label{section-curves}

\noindent
Dans le projet Champs, nous adopterons la définition
suivante d'une courbe.

\begin{definition}
\label{definition-curve}
Soit $k$ un corps. Une {\it courbe} est une variété de dimension $1$ sur $k$.
\end{definition}

\noindent
Deux exemples classiques de courbes sur $k$ sont la droite affine $\mathbf{A}^1_k$
et la droite projective $\mathbf{P}^1_k$. Le schéma $X = \Spec(k[x, y]/(f))$
est une courbe si et seulement si $f \in k[x, y]$ est irréductible.

\medskip\noindent
Notre définition d'une courbe présente les mêmes difficultés que celle d'une
variété; voir la discussion qui suit la définition \ref{definition-variety}.
Elle signifie en outre que toute courbe est munie d'un corps de définition spécifié.
Par exemple, $X = \Spec(\mathbf{C}[x])$ est une courbe sur $\mathbf{C}$,
mais on peut aussi le considérer comme une courbe sur $\mathbf{R}$. Le schéma
$\Spec(\mathbf{Z})$ n'est pas une courbe, bien que les schémas $\Spec(\mathbf{Z})$
et $\mathbf{A}^1_{\mathbf{F}_p}$ se comportent de manière analogue à bien des égards.

\begin{lemma}
\label{lemma-proper-minus-point}
Soit $X$ un schéma séparé irréductible de dimension $> 0$ sur un corps $k$.
Soit $x \in X$ un point fermé. Le sous-schéma ouvert $X \setminus \{x\}$
n'est pas propre sur $k$.
\end{lemma}

\begin{proof}
Puisque $X$ est irréductible, $U = X \setminus \{x\}$ n'est pas fermé dans $X$.
En particulier, l'immersion $U \to X$ n'est pas propre.
D'après Morphismes, lemme \ref{morphisms-lemma-image-proper-scheme-closed}
(on utilise ici que $X$ est séparé), $U \to \Spec(k)$ n'est pas propre non plus.
\end{proof}

\begin{lemma}
\label{lemma-dim-1-quasi-projective}
Soit $X$ un schéma séparé de type fini sur un corps $k$.
Si $\dim(X) \leq 1$, alors $X$ est H-quasi-projectif sur $k$.
\end{lemma}

\begin{proof}
D'après la proposition \ref{proposition-dim-1-noetherian-separated-has-ample},
le schéma $X$ possède un faisceau inversible ample $\mathcal{L}$.
D'après Morphismes, lemme \ref{morphisms-lemma-quasi-projective-finite-type-over-S},
on voit que $X$ est isomorphe à un sous-schéma localement
fermé de $\mathbf{P}^n_k$ au-dessus de $\Spec(k)$. C'est
la définition de H-quasi-projectif sur $k$; voir
Morphismes, définition \ref{morphisms-definition-quasi-projective}.
\end{proof}

\begin{lemma}
\label{lemma-dim-1-proper-projective}
Soit $X$ un schéma propre sur un corps $k$.
Si $\dim(X) \leq 1$, alors $X$ est H-projectif sur $k$.
\end{lemma}

\begin{proof}
D'après le lemme \ref{lemma-dim-1-quasi-projective}, $X$ est un
sous-schéma localement fermé de $\mathbf{P}^n_k$ pour un certain corps $k$.
Puisque $X$ est propre sur $k$, il s'ensuit que $X$ est un sous-schéma fermé
de $\mathbf{P}^n_k$
(Morphismes, lemme \ref{morphisms-lemma-image-proper-scheme-closed}).
\end{proof}

\begin{lemma}
\label{lemma-dim-1-projective-completion}
Soit $X$ un schéma séparé de type fini sur $k$.
Si $\dim(X) \leq 1$, alors il existe une immersion ouverte
$j : X \to \overline{X}$ possédant les propriétés suivantes :
\begin{enumerate}
\item $\overline{X}$ est H-projectif sur $k$, c'est-à-dire que $\overline{X}$
est un sous-schéma fermé de $\mathbf{P}^d_k$ pour un certain $d$,
\item $j(X) \subset \overline{X}$ est dense et schématiquement
dense,
\item $\overline{X} \setminus X = \{x_1, \ldots, x_n\}$
pour certains points fermés $x_i \in \overline{X}$.
\end{enumerate}
\end{lemma}

\begin{proof}
D'après le lemme \ref{lemma-dim-1-quasi-projective}, on peut supposer que $X$ est un
sous-schéma localement fermé de $\mathbf{P}^d_k$ pour un certain $d$. Soit
$\overline{X} \subset \mathbf{P}^d_k$ l'image schématique
de $X \to \mathbf{P}^d_k$; voir Morphismes, définition
\ref{morphisms-definition-scheme-theoretic-image}.
La description donnée dans
Morphismes, lemme \ref{morphisms-lemma-quasi-compact-immersion},
entraîne les propriétés (1) et (2).
On a alors $\dim(X) = 1 \Rightarrow \dim(\overline{X}) = 1$, par exemple en
considérant les points génériques; voir
le lemme \ref{lemma-dimension-locally-algebraic}.
Comme $\overline{X}$ est noethérien, on en déduit que
$\overline{X} \setminus X = \{x_1, \ldots, x_n\}$
est un ensemble fini de points fermés.
\end{proof}

\begin{lemma}
\label{lemma-reduced-dim-1-projective-completion}
Soit $X$ un schéma séparé de type fini sur $k$.
Si $X$ est réduit et $\dim(X) \leq 1$, alors il existe
une immersion ouverte $j : X \to \overline{X}$ telle que
\begin{enumerate}
\item $\overline{X}$ est H-projectif sur $k$, c'est-à-dire que $\overline{X}$
est un sous-schéma fermé de $\mathbf{P}^d_k$ pour un certain $d$,
\item $j(X) \subset \overline{X}$ est dense et schématiquement
dense,
\item $\overline{X} \setminus X = \{x_1, \ldots, x_n\}$
pour certains points fermés $x_i \in \overline{X}$,
\item les anneaux locaux $\mathcal{O}_{\overline{X}, x_i}$
sont des anneaux de valuation discrète pour $i = 1, \ldots, n$.
\end{enumerate}
\end{lemma}

\begin{proof}
Soit $j : X \to \overline{X}$ comme dans le
lemme \ref{lemma-dim-1-projective-completion}.
Considérons la normalisation $X'$ de $\overline{X}$ dans
$X$. D'après le lemme \ref{lemma-relative-normalization-finite},
le morphisme $X' \to \overline{X}$ est fini.
D'après Morphismes, lemme \ref{morphisms-lemma-finite-projective},
$X' \to \overline{X}$ est projectif. D'après Morphismes, lemme
\ref{morphisms-lemma-projective-over-quasi-projective-is-H-projective},
$X' \to \overline{X}$ est H-projectif.
D'après Morphismes, lemme \ref{morphisms-lemma-H-projective-composition},
$X' \to \Spec(k)$ est H-projectif.
Soit $\{x'_1, \ldots, x'_m\} \subset X'$ l'image réciproque
de $\{x_1, \ldots, x_n\} = \overline{X} \setminus X$.
Alors $\dim(\mathcal{O}_{X', x'_i}) = 1$ pour tout $1 \leq i \leq m$.
Par conséquent, les anneaux locaux $\mathcal{O}_{X', x'}$
sont des anneaux de valuation discrète d'après
Morphismes, lemme
\ref{morphisms-lemma-relative-normalization-normal-codim-1}.
Alors $X \to X'$ et $\{x'_1, \ldots, x'_m\}$ répondent aux conditions voulues.
\end{proof}

\begin{lemma}
\label{lemma-dim-1-projective-completion-after-insep}
Soit $X$ un schéma séparé de type fini sur $k$ tel que $\dim(X) \leq 1$.
Alors il existe un diagramme commutatif
$$
\xymatrix{
\overline{Y}_1 \amalg \ldots \amalg \overline{Y}_n \ar[rd] &
Y_1 \amalg \ldots \amalg Y_n \ar[r]_-\nu \ar[d] \ar[l]^j &
X_{k'} \ar[r] \ar[d] &
X \ar[d]^f \\
& \Spec(k'_1) \amalg \ldots \amalg \Spec(k'_n) \ar[r] &
\Spec(k') \ar[r] &
\Spec(k)
}
$$
de schémas possédant les propriétés suivantes :
\begin{enumerate}
\item $k'/k$ est une extension finie purement inséparable de corps,
\item $\nu$ est la normalisation de $X_{k'}$,
\item $j$ est une immersion ouverte d'image dense,
\item $k'_i/k'$ est une extension finie séparable pour $i = 1, \ldots, n$,
\item $\overline{Y}_i$ est lisse, projectif, géométriquement irréductible
et de dimension $\leq 1$ sur $k'_i$.
\end{enumerate}
\end{lemma}

\begin{proof}
Puisque l'on peut remplacer $X$ par sa réduction, on peut et on va supposer $X$ réduit.
Choisissons $X \to \overline{X}$ comme dans le
lemme \ref{lemma-reduced-dim-1-projective-completion}.
Si le lemme est établi pour $\overline{X}$, il en découle
pour $X$ (détails omis).
On peut donc et on va supposer $X$ projectif.

\medskip\noindent
Choisissons une extension finie purement inséparable $k'/k$ telle que la normalisation de
$X_{k'}$ soit géométriquement normale sur $k'$; voir
le lemme \ref{lemma-finite-extension-geometrically-normal}.
Notons $Y = (X_{k'})^\nu$ la normalisation; pour les propriétés
de la normalisation, voir la section \ref{section-normalization}.
Alors $Y$ est géométriquement régulier, car normalité et régularité coïncident
en dimension $\leq 1$; voir
Propriétés, lemme \ref{properties-lemma-normal-dimension-1-regular}.
Ainsi, $Y$ est lisse sur $k'$ d'après
le lemme \ref{lemma-geometrically-regular-smooth}.
Soit $Y = Y_1 \amalg \ldots \amalg Y_n$ la décomposition
de $Y$ en composantes irréductibles.
Posons $k'_i = \Gamma(Y_i, \mathcal{O}_{Y_i})$.
Ce sont des extensions finies séparables de $k'$ d'après
le lemme \ref{lemma-proper-geometrically-reduced-global-sections}.
Le lemme \ref{lemma-baby-stein} achève la démonstration.
\end{proof}

\begin{lemma}
\label{lemma-regular-point-on-curve}
Soit $k$ un corps. Soit $X$ une courbe sur $k$. Soit $x \in X$ un point
fermé. Considérons $x$ comme un sous-schéma fermé (réduit) de $X$, de faisceau
d'idéaux $\mathcal{I}$. Les conditions suivantes sont équivalentes :
\begin{enumerate}
\item $\mathcal{O}_{X, x}$ est régulier,
\item $\mathcal{O}_{X, x}$ est normal,
\item $\mathcal{O}_{X, x}$ est un anneau de valuation discrète,
\item $\mathcal{I}$ est un $\mathcal{O}_X$-module inversible,
\item $x$ est un diviseur de Cartier effectif sur $X$.
\end{enumerate}
Si $k$ est parfait ou si $\kappa(x)$ est séparable sur $k$,
ces conditions sont aussi équivalentes à
\begin{enumerate}
\item[(6)] $X \to \Spec(k)$ est lisse en $x$.
\end{enumerate}
\end{lemma}

\begin{proof}
Puisque $X$ est une courbe, l'anneau local $\mathcal{O}_{X, x}$ est un anneau local
noethérien intègre de dimension $1$ (lemme \ref{lemma-dimension-locally-algebraic}).
Les conditions (4) et (5) sont équivalentes par définition et équivalent au fait que
$\mathcal{I}_x = \mathfrak m_x \subset \mathcal{O}_{X, x}$ soit engendré par un élément
(Diviseurs, lemme \ref{divisors-lemma-effective-Cartier-in-points}).
L'équivalence de (1), (2), (3), (4) et (5) découle donc du
lemme \ref{algebra-lemma-characterize-dvr} d'Algèbre. La dernière assertion
découle du lemme \ref{lemma-dense-smooth-open-variety-over-perfect-field}
lorsque $k$ est parfait. Si $\kappa(x)/k$ est séparable, l'équivalence
découle du lemme \ref{algebra-lemma-separable-smooth} d'Algèbre.
\end{proof}

\begin{remark}
\label{remark-smooth-projective-compactification}
Soit $k$ un corps. Soit $X$ une courbe régulière sur $k$. D'après les
lemmes \ref{lemma-regular-point-on-curve} et
\ref{lemma-reduced-dim-1-projective-completion},
il existe une courbe projective non singulière $\overline{X}$
qui est une compactification de $X$, c'est-à-dire qu'il existe une
immersion ouverte $j : X \to \overline{X}$ dont le complément
est formé d'un nombre fini de points fermés. Si $k$ est parfait,
alors $X$ et $\overline{X}$ sont lisses sur $k$ et
$\overline{X}$ est une compactification projective lisse de $X$.
\end{remark}

\noindent
Observons que si un schéma affine $X$ sur $k$ est propre sur $k$,
alors $X$ est fini sur $k$ (Morphismes, lemme
\ref{morphisms-lemma-finite-proper}) et est donc de
dimension $0$
(Algèbre, lemme \ref{algebra-lemma-finite-dimensional-algebra} et
proposition \ref{algebra-proposition-dimension-zero-ring}).
Ainsi, un schéma de dimension $> 0$ sur $k$ ne peut être à la fois affine et
propre sur $k$. Les deux possibilités du lemme suivant
s'excluent donc mutuellement.

\begin{lemma}
\label{lemma-curve-affine-projective}
Soit $X$ une courbe sur $k$. Alors soit $X$ est un schéma affine, soit $X$
est H-projectif sur $k$.
\end{lemma}

\begin{proof}
Choisissons $X \to \overline{X}$ tel que
$\overline{X} \setminus X = \{x_1, \ldots, x_r\}$ comme dans le
lemme \ref{lemma-reduced-dim-1-projective-completion}.
Alors $\overline{X}$ est également une courbe.
Si $r = 0$, alors $X = \overline{X}$ est H-projectif sur $k$.
On peut donc supposer $r \geq 1$; il s'agit de montrer que
$X$ est affine. D'après le lemme \ref{lemma-open-in-affine-curve-affine},
il suffit de montrer que $\overline{X} \setminus \{x_1\}$ est affine.
Cela nous ramène à l'assertion énoncée au paragraphe suivant.

\medskip\noindent
Soit $X$ une courbe H-projective sur $k$. Soit $x \in X$ un point fermé
tel que $\mathcal{O}_{X, x}$ soit un anneau de valuation discrète. Affirmation :
$U = X \setminus \{x\}$ est affine. D'après le lemme \ref{lemma-regular-point-on-curve},
le point $x$ définit un diviseur de Cartier effectif de $X$.
Pour $n \geq 1$, notons $nx = x + \ldots + x$ la somme de $n$ termes; voir
Diviseurs, définition \ref{divisors-definition-sum-effective-Cartier-divisors}.
Notons $\mathcal{O}_{nx}$ le faisceau structural de $nx$, considéré comme un
module cohérent sur $X$.
Tout module inversible sur le schéma local $nx$ étant trivial,
la première suite exacte courte de
Diviseurs, remarque \ref{divisors-remark-ses-regular-section},
s'écrit
$$
0 \to \mathcal{O}_X
\xrightarrow{1} \mathcal{O}_X(nx) \to \mathcal{O}_{nx} \to 0
$$
dans notre cas. Notons que $\dim_k H^0(X, \mathcal{O}_{nx}) \geq n$.
En effet, d'après le lemme \ref{lemma-chi-tensor-finite}, on a
$H^0(X, \mathcal{O}_{nx}) = \mathcal{O}_{X, x}/(\pi^n)$, où $\pi$ est une uniformisante de
$\mathcal{O}_{X, x}$, et les puissances $\pi^i$ ont des images
$k$-linéairement indépendantes dans $\mathcal{O}_{X, x}/(\pi^n)$
pour $i = 0, 1, \ldots, n - 1$. On a $\dim_k H^1(X, \mathcal{O}_X) < \infty$
d'après Cohomologie des schémas, lemme
\ref{coherent-lemma-proper-over-affine-cohomology-finite}.
Si $n > \dim_k H^1(X, \mathcal{O}_X)$,
la suite exacte longue de cohomologie montre
qu'il existe $s \in \Gamma(X, \mathcal{O}_X(nx))$
qui n'est pas une section de $\mathcal{O}_X$.
Si l'on choisit $n$ minimal avec cette propriété, alors $s$
a pour image un générateur du germe
$\left(\mathcal{O}_X(nx)\right)_x$, faute de quoi il
définirait une section de $\mathcal{O}_X((n - 1)x) \subset \mathcal{O}_X(nx)$.
Pour ce $n$, on en déduit que $s_0 = 1$ et $s_1 = s$
engendrent le module inversible $\mathcal{L} = \mathcal{O}_X(nx)$.

\medskip\noindent
Considérons le morphisme correspondant
$f = \varphi_{\mathcal{L}, (s_0, s_1)} : X \to \mathbf{P}^1_k$
de Constructions, section \ref{constructions-section-projective-space}.
Notons que l'image réciproque de $D_{+}(T_0)$ est $U = X \setminus \{x\}$,
car la section $s_0$ de $\mathcal{L}$ ne s'annule qu'en $x$.
En particulier, $f$ n'est pas constant, c'est-à-dire que $\Im(f)$ contient
plus d'un point. Par conséquent, $f$ envoie nécessairement le point générique
$\eta$ de $X$ sur le point générique de $\mathbf{P}^1_k$.
Ainsi, si $y \in \mathbf{P}^1_k$ est un point fermé, $f^{-1}(\{y\})$
est un fermé de $X$ ne contenant pas $\eta$, et est donc fini.
Enfin, $f$ est propre\footnote{En effet, une variété H-projective
est une variété propre d'après Morphismes, lemme
\ref{morphisms-lemma-projective-is-quasi-projective-proper}.
Un morphisme de variétés dont la source est une variété propre est un morphisme propre
d'après Morphismes, lemme \ref{morphisms-lemma-image-proper-scheme-closed}.}.
D'après Cohomologie des schémas, lemme
\ref{coherent-lemma-proper-finite-fibre-finite-in-neighbourhood}\footnote{On
peut éviter d'utiliser ce lemme, qui repose sur le théorème des fonctions
formelles. En effet, $X$ est projectif, et il suffit donc de montrer
qu'un morphisme propre $f : X \to Y$ à fibres finies entre des schémas quasi-projectifs
sur $k$ est fini. Pour cela, on choisit un ouvert affine de $X$
contenant la fibre de $f$ au-dessus d'un point $y$, en utilisant le fait que tout ensemble fini de
points d'un schéma quasi-projectif sur $k$ est contenu dans un ouvert affine.
Quitte à restreindre $Y$ à un petit voisinage affine de $y$, on se ramène au
cas d'un morphisme propre entre schémas affines. Un tel morphisme est fini d'après
Morphismes, lemme \ref{morphisms-lemma-integral-universally-closed}.},
on conclut que $f$ est fini. Ainsi, $U = f^{-1}(D_{+}(T_0))$
est affine.
\end{proof}

\noindent
Le lemme suivant, combiné au lemme \ref{lemma-proper-minus-point},
montre que, si $X$ est un schéma séparé de dimension $1$ et
de type fini sur $k$, alors $X \setminus Z$ est affine dès que le
sous-ensemble fermé $Z$ rencontre toute composante irréductible de $X$.

\begin{lemma}
\label{lemma-dim-1-nonproper-affine}
Soit $X$ un schéma séparé de type fini sur $k$.
Si $\dim(X) \leq 1$ et si aucune composante irréductible de $X$
n'est propre de dimension $1$, alors $X$ est affine.
\end{lemma}

\begin{proof}
Soit $X = \bigcup X_i$ la décomposition de $X$ en composantes irréductibles.
Nous considérons $X_i$ comme un schéma intègre (muni de la structure
de schéma réduit induite, voir Schémas, définition
\ref{schemes-definition-reduced-induced-scheme}).
En particulier, $X_i$ est un singleton (donc affine) ou une courbe,
donc est affine par le lemme \ref{lemma-curve-affine-projective}.
Alors $\coprod X_i \to X$ est fini surjectif et $\coprod X_i$ est affine.
Nous voyons donc que $X$ est affine par
Cohomologie des schémas, lemme
\ref{coherent-lemma-image-affine-finite-morphism-affine-Noetherian}.
\end{proof}






\section{Degrés sur les courbes}
\label{section-divisors-curves}

\noindent
Nous commençons par définir le degré d'un faisceau inversible et, plus généralement,
celui d'un faisceau localement libre sur un schéma propre
de dimension $1$ sur un corps. Dans la section \ref{section-euler},
nous avons défini la caractéristique d'Euler-Poincaré
d'un faisceau cohérent $\mathcal{F}$ sur un schéma propre $X$ sur un corps
$k$ par la formule
$$
\chi(X, \mathcal{F}) = \sum (-1)^i \dim_k H^i(X, \mathcal{F}).
$$

\begin{definition}
\label{definition-degree-invertible-sheaf}
Soit $k$ un corps, soit $X$ un schéma propre de dimension $\leq 1$
sur $k$, et soit $\mathcal{L}$ un $\mathcal{O}_X$-module inversible.
Le {\it degré} de $\mathcal{L}$ est défini par
$$
\deg(\mathcal{L}) = \chi(X, \mathcal{L}) - \chi(X, \mathcal{O}_X)
$$
Plus généralement, si $\mathcal{E}$ est un faisceau localement libre de rang $n$,
nous définissons le {\it degré} de $\mathcal{E}$ par
$$
\deg(\mathcal{E}) = \chi(X, \mathcal{E}) - n\chi(X, \mathcal{O}_X)
$$
\end{definition}

\noindent
Remarquons que celui-ci dépend du triplet $\mathcal{E}/X/k$.
Si $X$ n'est pas connexe et si $\mathcal{E}$ est localement libre de type fini (mais
pas nécessairement de rang constant), on peut modifier la définition en sommant les degrés
des restrictions de $\mathcal{E}$ aux composantes connexes de $X$.
Si $\mathcal{E}$ est seulement un faisceau cohérent, il existe plusieurs
manières d'étendre la définition\footnote{Si $X$ est une courbe propre
et si $\mathcal{F}$ est un faisceau cohérent sur $X$, on définit souvent
le degré comme $\chi(X, \mathcal{F}) - r\chi(X, \mathcal{O}_X)$,
où $r = \dim_{\kappa(\xi)} \mathcal{F}_\xi$ est le rang de $\mathcal{F}$
au point générique $\xi$ de $X$.}.
Nous montrons dans une série de lemmes que cette définition possède toutes les propriétés
attendues du degré.

\begin{lemma}
\label{lemma-degree-base-change}
Soit $k'/k$ une extension de corps. Soit $X$ un schéma propre de
dimension $\leq 1$ sur $k$. Soit $\mathcal{E}$ un
$\mathcal{O}_X$-module localement libre de rang constant $n$. Alors le degré de
$\mathcal{E}/X/k$ est égal au degré de
$\mathcal{E}_{k'}/X_{k'}/k'$.
\end{lemma}

\begin{proof}
Plus précisément, posons $X_{k'} = X \times_{\Spec(k)} \Spec(k')$.
Soit $\mathcal{E}_{k'} = p^*\mathcal{E}$, où $p : X_{k'} \to X$
est la projection. Par
Cohomologie des schémas, lemme \ref{coherent-lemma-flat-base-change-cohomology},
nous avons
$H^i(X_{k'}, \mathcal{E}_{k'}) = H^i(X, \mathcal{E}) \otimes_k k'$
et
$H^i(X_{k'}, \mathcal{O}_{X_{k'}}) = H^i(X, \mathcal{O}_X) \otimes_k k'$.
Nous voyons donc que les caractéristiques d'Euler-Poincaré sont inchangées, et par conséquent
que le degré est inchangé.
\end{proof}

\begin{lemma}
\label{lemma-degree-additive}
Soit $k$ un corps. Soit $X$ un schéma propre de dimension $\leq 1$
sur $k$. Soit $0 \to \mathcal{E}_1 \to \mathcal{E}_2 \to \mathcal{E}_3 \to 0$
une suite exacte courte de $\mathcal{O}_X$-modules localement libres,
tous de type fini et de rang constant. Alors
$$
\deg(\mathcal{E}_2) = \deg(\mathcal{E}_1) + \deg(\mathcal{E}_3)
$$
\end{lemma}

\begin{proof}
Cela résulte immédiatement de l'additivité des caractéristiques d'Euler-Poincaré
(lemme \ref{lemma-euler-characteristic-additive})
et de l'additivité des rangs.
\end{proof}

\begin{lemma}
\label{lemma-degree-birational-pullback}
Soit $k$ un corps. Soit $f : X' \to X$ un morphisme birationnel de
schémas propres de dimension $\leq 1$ sur $k$. Alors
$$
\deg(f^*\mathcal{E}) = \deg(\mathcal{E})
$$
pour tout faisceau localement libre de type fini et de rang constant. Plus généralement,
il suffit que $f$ induise une bijection entre les composantes irréductibles
de dimension $1$ et des isomorphismes entre les anneaux locaux aux
points génériques correspondants.
\end{lemma}

\begin{proof}
Le morphisme $f$ est propre
(Morphismes, lemme \ref{morphisms-lemma-image-proper-scheme-closed})
et ses fibres sont de dimension $\leq 0$. Ainsi, $f$ est fini
(Cohomologie des schémas, lemme
\ref{coherent-lemma-proper-finite-fibre-finite-in-neighbourhood}).
Par conséquent,
$$
Rf_*f^*\mathcal{E} = f_*f^*\mathcal{E} =
\mathcal{E} \otimes_{\mathcal{O}_X} f_*\mathcal{O}_{X'}
$$
Comme $f$ induit un isomorphisme sur les anneaux locaux aux points génériques de
toutes les composantes irréductibles de dimension $1$, nous voyons que le noyau
et le conoyau dans
$$
0 \to \mathcal{K} \to \mathcal{O}_X \to f_*\mathcal{O}_{X'}
\to \mathcal{Q} \to 0
$$
ont des supports de dimension $\leq 0$. Remarquons que tensoriser cette suite par
$\mathcal{E}$ donne encore une suite exacte puisque $\mathcal{E}$ est localement libre.
Nous obtenons
\begin{align*}
\chi(X, \mathcal{E}) - \chi(X', f^*\mathcal{E})
& =
\chi(X, \mathcal{E}) - \chi(X, f_*f^*\mathcal{E}) \\
& =
\chi(X, \mathcal{E}) - \chi(X, \mathcal{E} \otimes f_*\mathcal{O}_{X'}) \\
& =
\chi(X, \mathcal{K} \otimes \mathcal{E}) -
\chi(X, \mathcal{Q} \otimes \mathcal{E}) \\
& =
n\chi(X, \mathcal{K}) -
n\chi(X, \mathcal{Q}) \\
& =
n\chi(X, \mathcal{O}_X) - n\chi(X, f_*\mathcal{O}_{X'}) \\
& =
n\chi(X, \mathcal{O}_X) - n\chi(X', \mathcal{O}_{X'})
\end{align*}
ce qui prouve l'assertion. La première égalité vaut parce que $f$ est fini, voir
Cohomologie des schémas, lemme \ref{coherent-lemma-relative-affine-cohomology}.
La deuxième égalité résulte de la formule de projection, voir
Cohomologie, lemme \ref{cohomology-lemma-projection-formula}.
La troisième résulte de l'additivité des caractéristiques d'Euler-Poincaré, voir
le lemme \ref{lemma-euler-characteristic-additive}.
La quatrième résulte du lemme \ref{lemma-chi-tensor-finite}.
\end{proof}

\begin{lemma}
\label{lemma-degree-on-proper-curve}
Soit $k$ un corps. Soit $X$ une courbe propre sur $k$ de point générique
$\xi$. Soit $\mathcal{E}$ un $\mathcal{O}_X$-module localement libre de rang $n$
et soit $\mathcal{F}$ un $\mathcal{O}_X$-module cohérent. Alors
$$
\chi(X, \mathcal{E} \otimes \mathcal{F}) =
r \deg(\mathcal{E}) + n \chi(X, \mathcal{F})
$$
où $r = \dim_{\kappa(\xi)} \mathcal{F}_\xi$ est le rang de $\mathcal{F}$.
\end{lemma}

\begin{proof}
Soit $\mathcal{P}$ la propriété d'un faisceau cohérent $\mathcal{F}$
sur $X$ selon laquelle la formule du lemme est satisfaite. Nous affirmons que
les hypothèses (1) et (2) de
Cohomologie des schémas, lemme \ref{coherent-lemma-property},
sont satisfaites pour $\mathcal{P}$. En effet, (1) l'est parce que la caractéristique d'Euler-Poincaré
et le rang $r$ sont additifs dans les suites exactes courtes de faisceaux cohérents.
La condition (2) l'est également : si $Z = X$, nous pouvons prendre
$\mathcal{G} = \mathcal{O}_X$, et $\mathcal{P}(\mathcal{O}_X)$
est vraie par définition du degré. Si $i : Z \to X$ est l'inclusion
d'un point fermé, nous pouvons prendre $\mathcal{G} = i_*\mathcal{O}_Z$,
et $\mathcal{P}$ est satisfaite par le lemme \ref{lemma-chi-tensor-finite}
et par le fait que $r = 0$ dans ce cas.
\end{proof}

\noindent
Soit $k$ un corps. Soit $X$ un schéma de type fini sur $k$ de
dimension $\leq 1$. Soient $C_i \subset X$, $i = 1, \ldots, t$, les
composantes irréductibles de dimension $1$. Nous munissons $C_i$ de la
structure de schéma réduit induite. Soit $\xi_i \in C_i$ le
point générique. La {\it multiplicité de $C_i$ dans $X$} est définie
comme la longueur
$$
m_i = \text{longueur}_{\mathcal{O}_{X, \xi_i}} \mathcal{O}_{X, \xi_i}
$$
Cette définition a un sens parce que $\mathcal{O}_{X, \xi_i}$ est un anneau local
noethérien de dimension zéro, donc est de longueur finie sur lui-même
(Algèbre, proposition \ref{algebra-proposition-dimension-zero-ring}).
Voir Homologie de Chow, section \ref{chow-section-cycle-of-closed-subscheme},
pour de plus amples informations. En fait, le degré d'un faisceau localement libre
ne dépend que de ses restrictions aux composantes irréductibles.

\begin{lemma}
\label{lemma-degree-in-terms-of-components}
Soit $k$ un corps. Soit $X$ un schéma propre de dimension $\leq 1$ sur $k$.
Soit $\mathcal{E}$ un $\mathcal{O}_X$-module localement libre de rang $n$.
Alors
$$
\deg(\mathcal{E}) = \sum m_i \deg(\mathcal{E}|_{C_i})
$$
où $C_i \subset X$, $i = 1, \ldots, t$, sont les composantes irréductibles
de dimension $1$, munies de la structure de schéma réduit induite, et $m_i$ est la
multiplicité de $C_i$ dans $X$.
\end{lemma}

\begin{proof}
Remarquons que l'énoncé a un sens parce que $C_i \to \Spec(k)$
est propre de dimension $1$ ; nous utiliserons ci-dessous le fait que tout sous-schéma fermé
de $X$ est propre sur $k$ par Morphismes, lemmes
\ref{morphisms-lemma-closed-immersion-proper} et
\ref{morphisms-lemma-composition-proper}. Considérons le sous-schéma ouvert
$U_i = X \setminus (\bigcup_{j \not = i} C_j)$ et soit $X_i \subset X$
l'adhérence schématique de $U_i$. Remarquons que $X_i \cap U_i = U_i$
(schématiquement) et que $X_i \cap U_j = \emptyset$
(ensemblistement) pour $i \not = j$ ; cela résulte de la description
de l'adhérence schématique dans
Morphismes, lemme \ref{morphisms-lemma-quasi-compact-immersion}.
Nous pouvons donc appliquer le lemme \ref{lemma-degree-birational-pullback} au morphisme
$X' = \coprod X_i \to X$. Comme il est clair que $C_i \subset X_i$
(schématiquement) et que la multiplicité de $C_i$ dans $X_i$ est égale
à la multiplicité de $C_i$ dans $X$, nous voyons que nous sommes ramenés au cas
étudié dans le paragraphe suivant.

\medskip\noindent
Supposons $X$ irréductible, de point générique $\xi$. Soit
$C = X_{red}$, de multiplicité $m$.
Nous devons montrer que $\deg(\mathcal{E}) = m \deg(\mathcal{E}|_C)$.
Soit $\mathcal{I} \subset \mathcal{O}_X$ l'idéal définissant le sous-schéma fermé
$C$. Soit $e \geq 0$ minimal tel que $\mathcal{I}^{e + 1} = 0$
(Cohomologie des schémas, lemme \ref{coherent-lemma-power-ideal-kills-sheaf}).
Nous raisonnons par récurrence sur $e$. Si $e = 0$, alors $X = C$ et le résultat
est immédiat. Sinon, posons $\mathcal{F} = \mathcal{I}^e$, considéré comme
un $\mathcal{O}_C$-module cohérent (Cohomologie des schémas, lemme
\ref{coherent-lemma-i-star-equivalence}).
Soit $X' \subset X$ le sous-schéma fermé défini par
l'idéal cohérent $\mathcal{I}^e$, et soit $m'$ la multiplicité
de $C$ dans $X'$. En prenant les germes en $\xi$ de la suite exacte courte
$$
0 \to \mathcal{F} \to \mathcal{O}_X \to \mathcal{O}_{X'} \to 0
$$
nous trouvons (en utilisant Algèbre, lemmes \ref{algebra-lemma-length-additive},
\ref{algebra-lemma-dimension-is-length} et
\ref{algebra-lemma-length-independent}) que
$$
m = \text{longueur}_{\mathcal{O}_{X, \xi}} \mathcal{O}_{X, \xi}
= \dim_{\kappa(\xi)} \mathcal{F}_\xi +
\text{longueur}_{\mathcal{O}_{X', \xi}} \mathcal{O}_{X', \xi}
= r + m'
$$
où $r$ est le rang de $\mathcal{F}$ en tant que faisceau cohérent sur $C$.
En tensorisant par $\mathcal{E}$, nous obtenons une suite exacte courte
$$
0 \to \mathcal{E}|_C \otimes \mathcal{F} \to \mathcal{E} \to
\mathcal{E} \otimes \mathcal{O}_{X'} \to 0
$$
Par récurrence, nous avons
$\deg(\mathcal{E}|_{X'}) = m' \deg(\mathcal{E}|_C)$.
Par le lemme \ref{lemma-degree-on-proper-curve}, nous avons
$\chi(\mathcal{E}|_C \otimes \mathcal{F}) =
r \deg(\mathcal{E}|_C) + n \chi(\mathcal{F})$.
En réunissant ces égalités, nous obtenons le résultat.
\end{proof}

\begin{lemma}
\label{lemma-degree-tensor-product}
Soit $k$ un corps, soit $X$ un schéma propre de dimension $\leq 1$
sur $k$, et soient $\mathcal{E}$, $\mathcal{V}$ des
$\mathcal{O}_X$-modules localement libres de type fini et de rang constant. Alors
$$
\deg(\mathcal{E} \otimes \mathcal{V}) =
\text{rang}(\mathcal{E}) \deg(\mathcal{V}) +
\text{rang}(\mathcal{V}) \deg(\mathcal{E})
$$
\end{lemma}

\begin{proof}
Par le lemme \ref{lemma-degree-in-terms-of-components} et un calcul
élémentaire, nous nous ramenons au cas d'une courbe propre.
Ce cas résulte du lemme \ref{lemma-degree-on-proper-curve}.
\end{proof}

\begin{lemma}
\label{lemma-degree-and-det}
Soit $k$ un corps, soit $X$ un schéma propre de dimension $\leq 1$
sur $k$, et soit $\mathcal{E}$ un
$\mathcal{O}_X$-module localement libre de rang $n$. Alors
$$
\deg(\mathcal{E}) = \deg(\wedge^n(\mathcal{E})) = \deg(\det(\mathcal{E}))
$$
\end{lemma}

\begin{proof}
Par le lemme \ref{lemma-degree-in-terms-of-components} et un calcul
élémentaire, nous nous ramenons au cas d'une courbe propre.
Il existe alors une modification $f : X' \to X$ telle que
$f^*\mathcal{E}$ admette une filtration dont les quotients
successifs sont des modules inversibles, voir
Diviseurs, lemme \ref{divisors-lemma-filter-after-modification}.
Par le lemme \ref{lemma-degree-birational-pullback}, nous pouvons travailler sur $X'$.
Nous pouvons donc supposer que nous avons une filtration
$$
0 = \mathcal{E}_0 \subset \mathcal{E}_1 \subset \mathcal{E}_2 \subset
\ldots \subset \mathcal{E}_n = \mathcal{E}
$$
par des $\mathcal{O}_X$-modules localement libres
telle que $\mathcal{L}_i = \mathcal{E}_i/\mathcal{E}_{i - 1}$ soit inversible.
Par Modules, lemme \ref{modules-lemma-det-ses}, et par récurrence, nous trouvons
$\det(\mathcal{E}) = \mathcal{L}_1 \otimes \ldots \otimes \mathcal{L}_n$.
L'égalité résulte donc du lemme \ref{lemma-degree-tensor-product}
et de l'additivité (lemme \ref{lemma-degree-additive}).
\end{proof}

\begin{lemma}
\label{lemma-degree-effective-Cartier-divisor}
Soit $k$ un corps, soit $X$ un schéma propre de dimension $\leq 1$
sur $k$. Soit $D$ un diviseur de Cartier effectif sur $X$.
Alors $D$ est fini sur $\Spec(k)$, de degré
$\deg(D) = \dim_k \Gamma(D, \mathcal{O}_D)$. Pour un faisceau localement libre
$\mathcal{E}$ de rang $n$, nous avons
$$
\deg(\mathcal{E}(D)) = n\deg(D) + \deg(\mathcal{E})
$$
où $\mathcal{E}(D) = \mathcal{E} \otimes_{\mathcal{O}_X} \mathcal{O}_X(D)$.
\end{lemma}

\begin{proof}
Comme $D$ est nulle part dense dans $X$ (Diviseurs, lemme
\ref{divisors-lemma-complement-effective-Cartier-divisor}),
nous voyons que $\dim(D) \leq 0$. Par conséquent, $D$ est fini sur $k$
par le lemme \ref{lemma-algebraic-scheme-dim-0}. Comme $k$ est un corps,
le morphisme $D \to \Spec(k)$ est fini localement libre et possède donc
un degré
(Morphismes, définition \ref{morphisms-definition-finite-locally-free}),
manifestement égal à $\dim_k \Gamma(D, \mathcal{O}_D)$,
comme l'affirme le lemme. Par
Diviseurs, définition
\ref{divisors-definition-invertible-sheaf-effective-Cartier-divisor},
il existe une suite exacte courte
$$
0 \to \mathcal{O}_X \to \mathcal{O}_X(D) \to i_*i^*\mathcal{O}_X(D) \to 0
$$
où $i : D \to X$ est l'immersion fermée. En tensorisant par $\mathcal{E}$,
nous obtenons une suite exacte courte
$$
0 \to \mathcal{E} \to \mathcal{E}(D) \to i_*i^*\mathcal{E}(D) \to 0
$$
L'égalité du lemme résulte de l'additivité de
la caractéristique d'Euler-Poincaré (lemme \ref{lemma-euler-characteristic-additive})
et du lemme \ref{lemma-chi-tensor-finite}.
\end{proof}

\begin{lemma}
\label{lemma-divisible}
Soit $k$ un corps. Soit $X$ un schéma propre sur $k$,
réduit et connexe.
Soit $\kappa = H^0(X, \mathcal{O}_X)$. Alors $\kappa/k$ est une
extension finie de corps et $w = [\kappa : k]$ divise
\begin{enumerate}
\item $\deg(\mathcal{E})$ pour tout $\mathcal{O}_X$-module localement libre
$\mathcal{E}$,
\item $[\kappa(x) : k]$ pour tout point fermé $x \in X$, et
\item $\deg(D)$ pour tout sous-schéma fermé $D \subset X$
de dimension zéro.
\end{enumerate}
\end{lemma}

\begin{proof}
Voir le lemme \ref{lemma-proper-geometrically-reduced-global-sections}
pour les assertions concernant $\kappa$.
Pour tout $\mathcal{O}_X$-module quasi-cohérent, les
$k$-espaces vectoriels $H^i(X, \mathcal{F})$ sont des $\kappa$-espaces vectoriels.
Les divisibilités résultent immédiatement de cette assertion et des
définitions.
\end{proof}

\begin{lemma}
\label{lemma-degree-pullback-map-proper-curves}
Soit $k$ un corps. Soit $f : X \to Y$ un morphisme non constant de
courbes propres sur $k$. Soit $\mathcal{E}$ un
$\mathcal{O}_Y$-module localement libre. Alors
$$
\deg(f^*\mathcal{E}) = \deg(X/Y) \deg(\mathcal{E})
$$
\end{lemma}

\begin{proof}
Le degré de $X$ sur $Y$ est défini dans
Morphismes, définition \ref{morphisms-definition-degree}.
Ainsi, $f_*\mathcal{O}_X$ est un $\mathcal{O}_Y$-module cohérent
de rang $\deg(X/Y)$, c'est-à-dire
$\deg(X/Y) = \dim_{\kappa(\xi)} (f_*\mathcal{O}_X)_\xi$, où $\xi$
est le point générique de $Y$. Nous obtenons donc
\begin{align*}
\chi(X, f^*\mathcal{E})
& =
\chi(Y, f_*f^*\mathcal{E}) \\
& =
\chi(Y, \mathcal{E} \otimes f_*\mathcal{O}_X) \\
& =
\deg(X/Y) \deg(\mathcal{E}) + n \chi(Y, f_*\mathcal{O}_X) \\
& =
\deg(X/Y) \deg(\mathcal{E}) + n \chi(X, \mathcal{O}_X)
\end{align*}
comme voulu. La première égalité vaut parce que $f$ est fini, voir
Cohomologie des schémas, lemme \ref{coherent-lemma-relative-affine-cohomology}.
La deuxième égalité résulte de la formule de projection, voir
Cohomologie, lemme \ref{cohomology-lemma-projection-formula}.
La troisième égalité résulte du lemme \ref{lemma-degree-on-proper-curve}.
\end{proof}

\noindent
Ce qui suit est une conséquence évidente, mais importante, des
résultats ci-dessus sur les degrés.

\begin{lemma}
\label{lemma-check-invertible-sheaf-trivial}
Soit $k$ un corps. Soit $X$ une courbe propre sur $k$.
Soit $\mathcal{L}$ un $\mathcal{O}_X$-module inversible.
\begin{enumerate}
\item Si $\mathcal{L}$ possède une section non nulle, alors
$\deg(\mathcal{L}) \geq 0$.
\item Si $\mathcal{L}$ possède une section non nulle $s$ qui s'annule
en un point, alors $\deg(\mathcal{L}) > 0$.
\item Si $\mathcal{L}$ et $\mathcal{L}^{-1}$ possèdent des sections non nulles, alors
$\mathcal{L} \cong \mathcal{O}_X$.
\item Si $\deg(\mathcal{L}) \leq 0$ et si $\mathcal{L}$ possède une section non nulle,
alors $\mathcal{L} \cong \mathcal{O}_X$.
\item Si $\mathcal{N} \to \mathcal{L}$ est un morphisme non nul de
$\mathcal{O}_X$-modules inversibles, alors $\deg(\mathcal{L}) \geq \deg(\mathcal{N})$,
et, s'il y a égalité, ce morphisme est un isomorphisme.
\end{enumerate}
\end{lemma}

\begin{proof}
Soit $s$ une section non nulle de $\mathcal{L}$. Comme $X$ est une courbe, nous
voyons que $s$ est une section régulière. Il existe donc un diviseur de
Cartier effectif $D \subset X$ et un isomorphisme
$\mathcal{L} \to \mathcal{O}_X(D)$ envoyant $s$ sur la section canonique
$1$ de $\mathcal{O}_X(D)$, voir
Diviseurs, lemme \ref{divisors-lemma-characterize-OD}.
Alors $\deg(\mathcal{L}) = \deg(D)$ par
le lemme \ref{lemma-degree-effective-Cartier-divisor}.
Comme $\deg(D) \geq 0$ et que l'égalité $= 0$ a lieu
si et seulement si $D = \emptyset$, cela prouve (1) et (2).
Dans le cas (3), nous voyons que $\deg(\mathcal{L}) = 0$ et que
$D = \emptyset$. De même pour (4). Pour prouver (5), appliquons
(1) et (4) au faisceau inversible
$$
\mathcal{L} \otimes_{\mathcal{O}_X} \mathcal{N}^{\otimes -1} =
\SheafHom_{\mathcal{O}_X}(\mathcal{N}, \mathcal{L})
$$
dont le degré est $\deg(\mathcal{L}) - \deg(\mathcal{N})$
par le lemme \ref{lemma-degree-tensor-product}.
\end{proof}

\begin{lemma}
\label{lemma-no-sections-dual-nef}
Soit $k$ un corps. Soit $X$ un schéma propre sur $k$,
réduit, connexe et équidimensionnel de dimension $1$.
Soit $\mathcal{L}$ un $\mathcal{O}_X$-module inversible.
Si $\deg(\mathcal{L}|_C) \leq 0$ pour toute composante irréductible $C$
de $X$, alors ou bien $H^0(X, \mathcal{L}) = 0$, ou bien
$\mathcal{L} \cong \mathcal{O}_X$.
\end{lemma}

\begin{proof}
Soit $s \in H^0(X, \mathcal{L})$ non nul. Comme $X$ est réduit, il existe
une composante irréductible $C$ de $X$ telle que $s|_C \not = 0$.
Mais si $s|_C$ est non nul, alors
$s$ ne s'annule nulle part sur $C$ par
le lemme \ref{lemma-check-invertible-sheaf-trivial}.
Cela implique à son tour que $s$ ne s'annule nulle part sur
toute composante irréductible de $X$ rencontrant $C$.
Comme $X$ est connexe, nous en concluons que $s$
ne s'annule nulle part, et le lemme en résulte.
\end{proof}

\begin{lemma}
\label{lemma-ample-curve}
Soit $k$ un corps. Soit $X$ une courbe propre sur $k$.
Soit $\mathcal{L}$ un $\mathcal{O}_X$-module inversible.
Alors $\mathcal{L}$ est ample si et seulement si $\deg(\mathcal{L}) > 0$.
\end{lemma}

\begin{proof}
Si $\mathcal{L}$ est ample, il existe un $n > 0$ et une section
$s \in H^0(X, \mathcal{L}^{\otimes n})$ tels que $X_s$ soit affine. Comme
$X$ n'est pas affine (sinon, d'après
Morphismes, lemme \ref{morphisms-lemma-finite-proper},
$X$ serait fini), on voit que $s$ s'annule en un point.
Ainsi, $\deg(\mathcal{L}^{\otimes n}) > 0$ d'après le
lemme \ref{lemma-check-invertible-sheaf-trivial}.
D'après le lemme \ref{lemma-degree-tensor-product},
on conclut que $\deg(\mathcal{L}) = 1/n\deg(\mathcal{L}^{\otimes n}) > 0$.

\medskip\noindent
Supposons que $\deg(\mathcal{L}) > 0$. Alors
$$
\dim_k H^0(X, \mathcal{L}^{\otimes n}) \geq \chi(X, \mathcal{L}^n)
= n\deg(\mathcal{L}) + \chi(X, \mathcal{O}_X)
$$
croît linéairement en $n$. Pour toute famille finie de points
fermés $x_1, \ldots, x_t$ de $X$, on peut donc trouver un $n$ tel que
$\dim_k H^0(X, \mathcal{L}^{\otimes n}) > \sum \dim_k \kappa(x_i)$.
(Rappelons que, d'après le théorème des zéros de Hilbert, les extensions
$\kappa(x_i)/k$ sont finies, voir par exemple
Morphismes, lemme \ref{morphisms-lemma-closed-point-fibre-locally-finite-type}.)
On peut donc trouver une section non nulle $s \in H^0(X, \mathcal{L}^{\otimes n})$
qui s'annule en $x_1, \ldots, x_t$. En particulier, si l'on choisit
$x_1, \ldots, x_t$ de telle sorte que $X \setminus \{x_1, \ldots, x_t\}$
soit affine, alors $X_s$ est également affine
(par exemple d'après Propriétés, lemme \ref{properties-lemma-affine-cap-s-open},
bien que, si l'on choisit l'ensemble fini de telle sorte que
$\mathcal{L}|_{X \setminus \{x_1, \ldots, x_t\}}$ soit trivial,
ce soit immédiat). On peut donc trouver un $n > 0$
et une section non nulle $s \in H^0(X, \mathcal{L}^{\otimes n})$ telle
que $X_s$ soit affine.

\medskip\noindent
Nous allons montrer que, pour tout faisceau quasi-cohérent d'idéaux $\mathcal{I}$,
il existe un $m > 0$ tel que
$H^1(X, \mathcal{I} \otimes \mathcal{L}^{\otimes m})$ soit nul.
Cela achèvera la démonstration d'après
Cohomologie des schémas, lemme \ref{coherent-lemma-vanshing-gives-ample}.
Pour cela, considérons les applications
$$
\mathcal{I} \xrightarrow{s}
\mathcal{I} \otimes \mathcal{L}^{\otimes n} \xrightarrow{s}
\mathcal{I} \otimes \mathcal{L}^{\otimes 2n} \xrightarrow{s} \ldots
$$
Puisque $\mathcal{I}$ est sans torsion, ces applications sont injectives et
sont des isomorphismes sur $X_s$ ; leurs conoyaux ont donc un $H^1$ nul
(par exemple d'après Cohomologie des schémas, lemme
\ref{coherent-lemma-coherent-support-dimension-0}).
On en conclut que les applications d'espaces vectoriels
$$
H^1(X, \mathcal{I}) \to
H^1(X, \mathcal{I} \otimes \mathcal{L}^{\otimes n}) \to
H^1(X, \mathcal{I} \otimes \mathcal{L}^{\otimes 2n}) \to \ldots
$$
sont surjectives. D'autre part, $H^1(X, \mathcal{I})$ est de dimension
finie, et tout élément a une image nulle à partir d'un certain rang d'après
Cohomologie des schémas, lemme
\ref{coherent-lemma-section-affine-open-kills-classes}.
Il existe donc un $e > 0$ tel que
$H^1(X, \mathcal{I} \otimes \mathcal{L}^{\otimes en})$ soit nul.
Cela achève la démonstration.
\end{proof}

\begin{lemma}
\label{lemma-ampleness-in-terms-of-degrees-components}
Soit $k$ un corps. Soit $X$ un schéma propre de dimension $\leq 1$
sur $k$. Soit $\mathcal{L}$ un $\mathcal{O}_X$-module inversible.
Soient $C_i \subset X$, $i = 1, \ldots, t$, les composantes irréductibles
de dimension $1$. Les conditions suivantes sont équivalentes :
\begin{enumerate}
\item $\mathcal{L}$ est ample, et
\item $\deg(\mathcal{L}|_{C_i}) > 0$ pour $i = 1, \ldots, t$.
\end{enumerate}
\end{lemma}

\begin{proof}
Soient $x_1, \ldots, x_r \in X$ les points fermés isolés.
Considérons $x_i = \Spec(\kappa(x_i))$ comme un schéma.
Considérons le morphisme de schémas
$$
f : C_1 \amalg \ldots \amalg C_t \amalg x_1 \amalg \ldots \amalg x_r
\longrightarrow X
$$
C'est un morphisme fini surjectif entre des schémas propres sur $k$
(détails omis). Ainsi, $\mathcal{L}$ est ample si et seulement si
$f^*\mathcal{L}$ est ample (Cohomologie des schémas, lemme
\ref{coherent-lemma-surjective-finite-morphism-ample}).
On conclut donc par le lemme \ref{lemma-ample-curve}.
\end{proof}

\begin{lemma}
\label{lemma-general-degree-g-line-bundle}
Soit $k$ un corps algébriquement clos. Soit $X$ une courbe propre sur $k$.
Il existe alors
\begin{enumerate}
\item un $\mathcal{O}_X$-module inversible $\mathcal{L}$ tel que
$\dim_k H^0(X, \mathcal{L}) = 1$ et $H^1(X, \mathcal{L}) = 0$, et
\item un $\mathcal{O}_X$-module inversible $\mathcal{N}$ tel que
$\dim_k H^0(X, \mathcal{N}) = 0$ et $H^1(X, \mathcal{N}) = 0$.
\end{enumerate}
\end{lemma}

\begin{proof}
Choisissons une immersion fermée $i : X \to \mathbf{P}^n_k$
(lemme \ref{lemma-dim-1-proper-projective}).
En posant $\mathcal{L} = i^*\mathcal{O}_{\mathbf{P}^n}(d)$
pour $d \gg 0$, on voit qu'il existe un faisceau inversible
$\mathcal{L}$ tel que $H^0(X, \mathcal{L}) \not = 0$ et
$H^1(X, \mathcal{L}) = 0$ (voir
Cohomologie des schémas, lemme \ref{coherent-lemma-vanshing-gives-ample},
pour l'annulation, et les références qui y sont données pour la non-annulation).
Nous allons achever la démonstration de (1) par récurrence descendante sur
$t = \dim_k H^0(X, \mathcal{L})$. Le cas initial $t = 1$ est immédiat.
Supposons que $t > 1$.

\medskip\noindent
Soit $U \subset X$ l'ouvert non vide des points non singuliers
étudié dans le lemme \ref{lemma-dense-smooth-open-variety-over-perfect-field}.
Soit $s \in H^0(X, \mathcal{L})$ non nul. Il existe un point
fermé $x \in U$ tel que $s$ ne s'annule pas en $x$. Soit $\mathcal{I}$
le faisceau d'idéaux de $i : x \to X$ comme dans le
lemme \ref{lemma-regular-point-on-curve}. Considérons la
suite exacte courte
$$
0 \to \mathcal{I} \otimes_{\mathcal{O}_X} \mathcal{L} \to
\mathcal{L} \to i_*i^*\mathcal{L} \to 0
$$
Remarquons que $H^0(X, i_*i^*\mathcal{L}) = H^0(x, i^*\mathcal{L})$
est de dimension $1$, puisque $x$ est un point $k$-rationnel ($k$ est
algébriquement clos). Comme $s$ ne s'annule pas en $x$, on en conclut que
$$
H^0(X, \mathcal{L}) \longrightarrow H^0(X, i_*i^*\mathcal{L})
$$
est surjective. Par conséquent,
$\dim_k H^0(X, \mathcal{I} \otimes_{\mathcal{O}_X} \mathcal{L}) = t - 1$.
Enfin, la suite exacte longue de cohomologie montre également que
$H^1(X, \mathcal{I} \otimes_{\mathcal{O}_X} \mathcal{L}) = 0$,
ce qui achève la démonstration de l'étape de récurrence.

\medskip\noindent
Pour obtenir un faisceau inversible comme en (2), prenons un faisceau inversible
$\mathcal{L}$ comme en (1) et appliquons l'argument du paragraphe précédent
une fois encore.
\end{proof}

\begin{lemma}
\label{lemma-vanishing-degree-2g-and-1-line-bundle}
Soit $k$ un corps algébriquement clos. Soit $X$ une courbe propre sur $k$.
Posons $g = \dim_k H^1(X, \mathcal{O}_X)$. Pour tout
$\mathcal{O}_X$-module inversible $\mathcal{L}$ tel que $\deg(\mathcal{L}) \geq 2g - 1$,
on a $H^1(X, \mathcal{L}) = 0$.
\end{lemma}

\begin{proof}
Soit $\mathcal{N}$ le module inversible obtenu dans
le lemme \ref{lemma-general-degree-g-line-bundle}, partie (2).
Le degré de $\mathcal{N}$ est
$\chi(X, \mathcal{N}) - \chi(X, \mathcal{O}_X) = 0 - (1 - g) = g - 1$.
Le degré de $\mathcal{L} \otimes \mathcal{N}^{\otimes - 1}$
est donc $\deg(\mathcal{L}) - (g - 1) \geq g$.
Par conséquent,
$\chi(X, \mathcal{L} \otimes \mathcal{N}^{\otimes -1}) \geq g + 1 - g = 1$.
Il existe ainsi une section globale non nulle $s$ dont le schéma des zéros est un
diviseur de Cartier effectif $D$ de degré $\deg(\mathcal{L}) - (g - 1)$.
On obtient une suite exacte courte
$$
0 \to \mathcal{N} \xrightarrow{s} \mathcal{L} \to i_*(\mathcal{L}|_D) \to 0
$$
où $i : D \to X$ est le morphisme d'inclusion. On en conclut que
$H^0(X, \mathcal{L})$ s'applique isomorphiquement sur $H^0(D, \mathcal{L}|_D)$,
qui est de dimension $\deg(\mathcal{L}) - (g - 1)$. Le résultat découle
de la définition du degré.
\end{proof}






\section{Intersections numériques}
\label{section-num}

\noindent
Dans cette section, nous allons exploiter la caractéristique d'Euler-Poincaré des
faisceaux cohérents sur les schémas propres afin d'obtenir des nombres d'intersection
numériques pour les modules inversibles. Notre principal outil sera le
lemme suivant.

\begin{lemma}
\label{lemma-numerical-polynomial-from-euler}
Soit $k$ un corps. Soit $X$ un schéma propre sur $k$. Soit $\mathcal{F}$
un $\mathcal{O}_X$-module cohérent. Soient
$\mathcal{L}_1, \ldots, \mathcal{L}_r$ des $\mathcal{O}_X$-modules inversibles.
L'application
$$
(n_1, \ldots, n_r) \longmapsto
\chi(X, \mathcal{F} \otimes
\mathcal{L}_1^{\otimes n_1} \otimes \ldots \otimes
\mathcal{L}_r^{\otimes n_r})
$$
est un polynôme numérique en $n_1, \ldots, n_r$, dont le degré total est au
plus la dimension du support de $\mathcal{F}$.
\end{lemma}

\begin{proof}
Nous procédons par récurrence sur $\dim(\text{Supp}(\mathcal{F}))$.
Si ce nombre est nul, la fonction est constante de valeur
$\dim_k \Gamma(X, \mathcal{F})$ d'après le lemme \ref{lemma-chi-tensor-finite}.
Supposons que $\dim(\text{Supp}(\mathcal{F})) > 0$.

\medskip\noindent
Si $\mathcal{F}$ a des points associés immergés, on peut considérer
la suite exacte courte
$0 \to \mathcal{K} \to \mathcal{F} \to \mathcal{F}' \to 0$
construite dans Diviseurs, lemme \ref{divisors-lemma-remove-embedded-points}.
Comme la dimension du support de $\mathcal{K}$ est strictement plus petite,
le résultat vaut pour $\mathcal{K}$ par l'hypothèse de récurrence, avec
un degré total strictement plus petit.
Par additivité de la caractéristique d'Euler-Poincaré
(lemme \ref{lemma-euler-characteristic-additive}),
il suffit d'établir le résultat pour $\mathcal{F}'$. On peut donc supposer que
$\mathcal{F}$ n'a pas de points associés immergés.

\medskip\noindent
Si $i : Z \to X$ est une immersion fermée et $\mathcal{F} = i_*\mathcal{G}$,
le résultat pour $X$, $\mathcal{F}$,
$\mathcal{L}_1, \ldots, \mathcal{L}_r$ équivaut au résultat
pour $Z$, $\mathcal{G}$, $i^*\mathcal{L}_1, \ldots, i^*\mathcal{L}_r$
(car les groupes de cohomologie coïncident, voir
Cohomologie des schémas, lemme \ref{coherent-lemma-relative-affine-cohomology}).
En appliquant Diviseurs, lemme \ref{divisors-lemma-no-embedded-points-endos},
on peut supposer que $X$ n'a pas de composantes immergées et que
$X = \text{Supp}(\mathcal{F})$.

\medskip\noindent
Choisissons une section méromorphe régulière $s$ de $\mathcal{L}_1$, voir
Diviseurs, lemme \ref{divisors-lemma-regular-meromorphic-section-exists}.
Soit $\mathcal{I} \subset \mathcal{O}_X$ l'idéal des
dénominateurs de $s$, et considérons les applications
$$
\mathcal{I}\mathcal{F} \to \mathcal{F},\quad
\mathcal{I}\mathcal{F} \to \mathcal{F} \otimes \mathcal{L}_1
$$
de Diviseurs, lemme \ref{divisors-lemma-make-maps-regular-section}.
Elles sont injectives et ont pour conoyaux $\mathcal{Q}$, $\mathcal{Q}'$,
à support dans des sous-schémas fermés rares de $X = \text{Supp}(\mathcal{F})$.
Le produit tensoriel par le module inversible
$\mathcal{L}_1^{\otimes n_1} \otimes \ldots \otimes \mathcal{L}_r^{\otimes n_r}$
est exact ; par conséquent, en utilisant de nouveau l'additivité,
on voit que
\begin{align*}
&\chi(X, \mathcal{F} \otimes \mathcal{L}_1^{\otimes n_1} \otimes \ldots \otimes
\mathcal{L}_r^{\otimes n_r}) -
\chi(X, \mathcal{F} \otimes \mathcal{L}_1^{\otimes n_1 + 1}
\otimes \ldots \otimes \mathcal{L}_r^{\otimes n_r}) \\
& =
\chi(\mathcal{Q} \otimes \mathcal{L}_1^{\otimes n_1} \otimes \ldots \otimes
\mathcal{L}_r^{\otimes n_r}) -
\chi(\mathcal{Q}' \otimes \mathcal{L}_1^{\otimes n_1} \otimes \ldots \otimes
\mathcal{L}_r^{\otimes n_r})
\end{align*}
Ainsi, la fonction $P(n_1, \ldots, n_r)$ du lemme
a la propriété que
$$
P(n_1 + 1, n_2, \ldots, n_r) - P(n_1, \ldots, n_r)
$$
est un polynôme numérique dont le degré total est $<$ la dimension
du support de $\mathcal{F}$. Bien entendu, par symétrie, il en va
de même de
$$
P(n_1, \ldots, n_{i - 1}, n_i + 1, n_{i + 1}, \ldots, n_r)
- P(n_1, \ldots, n_r)
$$
pour tout $i \in \{1, \ldots, r\}$. Un argument arithmétique simple montre
que $P$ est un polynôme numérique de degré total au plus
$\dim(\text{Supp}(\mathcal{F}))$.
\end{proof}

\noindent
Le lemme suivant montre, en gros, que le coefficient dominant ne dépend que
de la longueur du module cohérent aux points génériques de son
support.

\begin{lemma}
\label{lemma-numerical-polynomial-leading-term}
Soit $k$ un corps. Soit $X$ un schéma propre sur $k$. Soit
$\mathcal{F}$ un $\mathcal{O}_X$-module cohérent. Soient
$\mathcal{L}_1, \ldots, \mathcal{L}_r$ des $\mathcal{O}_X$-modules inversibles.
Posons $d = \dim(\text{Supp}(\mathcal{F}))$.
Soient $Z_i \subset X$ les composantes irréductibles
de $\text{Supp}(\mathcal{F})$ de dimension $d$. Soit $\xi_i \in Z_i$
le point générique, et posons
$m_i = \text{longueur}_{\mathcal{O}_{X, \xi_i}}(\mathcal{F}_{\xi_i})$.
Alors
$$
\chi(X, \mathcal{F} \otimes \mathcal{L}_1^{\otimes n_1} \otimes \ldots \otimes
\mathcal{L}_r^{\otimes n_r}) -
\sum\nolimits_i
m_i\ \chi(Z_i, \mathcal{L}_1^{\otimes n_1} \otimes \ldots \otimes
\mathcal{L}_r^{\otimes n_r}|_{Z_i})
$$
est un polynôme numérique en $n_1, \ldots, n_r$ de degré total $< d$.
\end{lemma}

\begin{proof}
Considérons les couples $(\xi , Z)$ où $Z \subset X$ est un
sous-schéma fermé intègre de dimension $d$ et où $\xi$ est son point générique.
Alors le $\mathcal{O}_{X, \xi}$-module de type fini $\mathcal{F}_\xi$
a son support contenu dans $\{\xi\}$ ; sa longueur
$m_Z = \text{longueur}_{\mathcal{O}_{X, \xi}}(\mathcal{F}_\xi)$
est donc finie (Algèbre, lemme \ref{algebra-lemma-support-point})
et est nulle sauf si $Z = Z_i$ pour un certain $i$. L'expression
du lemme peut donc s'écrire
$$
E(\mathcal{F}) =
\chi(X, \mathcal{F} \otimes \mathcal{L}_1^{\otimes n_1} \otimes \ldots \otimes
\mathcal{L}_r^{\otimes n_r}) -
\sum\nolimits
m_Z\ \chi(Z, \mathcal{L}_1^{\otimes n_1} \otimes \ldots \otimes
\mathcal{L}_r^{\otimes n_r}|_Z)
$$
où la somme porte sur les sous-schémas fermés intègres $Z \subset X$
de dimension $d$. L'application $\mathcal{F} \mapsto E(\mathcal{F})$
est additive pour les suites exactes courtes
$0 \to \mathcal{F} \to \mathcal{F}' \to \mathcal{F}'' \to 0$
de $\mathcal{O}_X$-modules cohérents dont le support est de dimension
$\leq d$. Cela résulte de l'additivité des caractéristiques d'Euler-Poincaré
(lemme \ref{lemma-euler-characteristic-additive})
et de l'additivité des longueurs
(Algèbre, lemme \ref{algebra-lemma-length-additive}).
Appliquons Cohomologie des schémas, lemme \ref{coherent-lemma-coherent-filter},
pour obtenir une filtration
$$
0 = \mathcal{F}_0 \subset \mathcal{F}_1 \subset
\ldots \subset \mathcal{F}_m = \mathcal{F}
$$
par des sous-faisceaux cohérents tels que, pour chaque $j = 1, \ldots, m$,
il existe un sous-schéma fermé intègre $V_j \subset X$
et un faisceau d'idéaux non nul $\mathcal{I}_j \subset \mathcal{O}_{V_j}$
tels que
$$
\mathcal{F}_j/\mathcal{F}_{j - 1}
\cong (V_j \to X)_* \mathcal{I}_j
$$
Il en résulte que $V_j \subset \text{Supp}(\mathcal{F})$ et
donc que $\dim(V_j) \leq d$.
Par l'additivité relevée ci-dessus, il suffit
d'établir le résultat pour chacun des sous-quotients
$\mathcal{F}_j/\mathcal{F}_{j - 1}$. Il suffit donc d'établir
le résultat lorsque $\mathcal{F} = (V \to X)_*\mathcal{I}$, où
$V \subset X$ est un sous-schéma fermé intègre de dimension $\leq d$
et $\mathcal{I} \subset \mathcal{O}_V$ est un faisceau cohérent
d'idéaux non nul. Si $\dim(V) < d$, et plus généralement pour $\mathcal{F}$
dont le support est de dimension $< d$, le premier terme
de $E(\mathcal{F})$ est de degré total $< d$ d'après le
lemme \ref{lemma-numerical-polynomial-from-euler},
et le second terme est nul. Si $\dim(V) = d$, on peut utiliser la
suite exacte courte
$$
0 \to (V \to X)_*\mathcal{I} \to (V \to X)_*\mathcal{O}_V
\to (V \to X)_*(\mathcal{O}_V/\mathcal{I}) \to 0
$$
Le résultat vaut pour le faisceau du milieu, car
le seul $Z$ qui figure dans la somme est $Z = V$,
avec $m_Z = 1$, et parce que
$$
H^i(X, ((V \to X)_*\mathcal{O}_V) \otimes 
 \mathcal{L}_1^{\otimes n_1} \otimes \ldots \otimes
\mathcal{L}_r^{\otimes n_r}) =
H^i(V,  \mathcal{L}_1^{\otimes n_1} \otimes \ldots \otimes
\mathcal{L}_r^{\otimes n_r}|_V)
$$
d'après la formule de projection
(Cohomologie, section \ref{cohomology-section-projection-formula}) et
Cohomologie des schémas, lemme
\ref{coherent-lemma-relative-affine-cohomology} ;
dans ce cas, on a donc en fait $E(\mathcal{F}) = 0$.
Le résultat vaut pour le faisceau de droite, car son support
est de dimension $< d$. Il vaut donc pour le faisceau
de gauche, ce qui démontre le lemme.
\end{proof}

\begin{definition}
\label{definition-intersection-number}
Soit $k$ un corps. Soit $X$ un schéma propre sur $k$. Soit
$i : Z \to X$ un sous-schéma fermé de dimension $d$. Soient
$\mathcal{L}_1, \ldots, \mathcal{L}_d$ des
$\mathcal{O}_X$-modules inversibles. On définit le {\it nombre d'intersection}
$(\mathcal{L}_1 \cdots \mathcal{L}_d \cdot Z)$
comme le coefficient du monôme $n_1 \ldots n_d$ dans le polynôme numérique
$$
\chi(X, i_*\mathcal{O}_Z \otimes \mathcal{L}_1^{\otimes n_1} \otimes
\ldots \otimes \mathcal{L}_d^{\otimes n_d}) =
\chi(Z, \mathcal{L}_1^{\otimes n_1} \otimes
\ldots \otimes \mathcal{L}_d^{\otimes n_d}|_Z)
$$
Dans le cas particulier
où $\mathcal{L}_1 = \ldots = \mathcal{L}_d = \mathcal{L}$,
on écrit $(\mathcal{L}^d \cdot Z)$.
\end{definition}

\noindent
L'égalité affichée dans la définition résulte de
la formule de projection
(Cohomologie, section \ref{cohomology-section-projection-formula}) et de
Cohomologie des schémas, lemme
\ref{coherent-lemma-relative-affine-cohomology}.
Nous allons établir quelques lemmes sur ces nombres d'intersection.

\begin{lemma}
\label{lemma-intersection-number-integer}
Dans la situation de la définition \ref{definition-intersection-number},
le nombre d'intersection
$(\mathcal{L}_1 \cdots \mathcal{L}_d \cdot Z)$
est un entier.
\end{lemma}

\begin{proof}
Tout polynôme numérique de degré $e$ en $n_1, \ldots, n_d$ s'écrit
de manière unique comme combinaison $\mathbf{Z}$-linéaire des fonctions
${n_1 \choose k_1}{n_2 \choose k_2} \ldots {n_d \choose k_d}$ avec
$k_1 + \ldots + k_d \leq e$. Appliquons cela avec $e = d$.
Nous laissons les détails en exercice.
\end{proof}

\begin{lemma}
\label{lemma-intersection-number-additive}
Dans la situation de la définition \ref{definition-intersection-number},
le nombre d'intersection
$(\mathcal{L}_1 \cdots \mathcal{L}_d \cdot Z)$
est additif : si $\mathcal{L}_i = \mathcal{L}_i' \otimes \mathcal{L}_i''$,
alors
$$
(\mathcal{L}_1 \cdots \mathcal{L}_i \cdots \mathcal{L}_d \cdot Z) =
(\mathcal{L}_1 \cdots \mathcal{L}_i' \cdots \mathcal{L}_d \cdot Z) +
(\mathcal{L}_1 \cdots \mathcal{L}_i'' \cdots \mathcal{L}_d \cdot Z)
$$
\end{lemma}

\begin{proof}
Cela résulte de ce que, d'après le lemme \ref{lemma-numerical-polynomial-from-euler},
la fonction
$$
(n_1, \ldots, n_{i - 1}, n_i', n_i'', n_{i + 1}, \ldots, n_d)
\mapsto
\chi(Z, \mathcal{L}_1^{\otimes n_1} \otimes
\ldots \otimes (\mathcal{L}_i')^{\otimes n_i'} \otimes
(\mathcal{L}_i'')^{\otimes n_i''} \otimes \ldots \otimes
\mathcal{L}_d^{\otimes n_d}|_Z)
$$
est un polynôme numérique de degré total au plus $d$ en $d + 1$ variables.
\end{proof}

\begin{lemma}
\label{lemma-intersection-number-in-terms-of-components}
Dans la situation de la définition \ref{definition-intersection-number},
soient $Z_i \subset Z$ les composantes irréductibles de dimension $d$. Posons
$m_i = \text{longueur}_{\mathcal{O}_{X, \xi_i}}(\mathcal{O}_{Z, \xi_i})$,
où $\xi_i \in Z_i$ est le point générique. Alors
$$
(\mathcal{L}_1 \cdots \mathcal{L}_d \cdot Z) =
\sum m_i(\mathcal{L}_1 \cdots \mathcal{L}_d \cdot Z_i)
$$
\end{lemma}

\begin{proof}
Cela résulte immédiatement du lemme \ref{lemma-numerical-polynomial-leading-term}
et des définitions.
\end{proof}

\begin{lemma}
\label{lemma-intersection-number-and-pullback}
Soit $k$ un corps. Soit $f : Y \to X$ un morphisme de schémas propres sur $k$.
Soit $Z \subset Y$ un sous-schéma fermé intègre de dimension $d$, et soient
$\mathcal{L}_1, \ldots, \mathcal{L}_d$ des $\mathcal{O}_X$-modules inversibles.
Alors
$$
(f^*\mathcal{L}_1 \cdots f^*\mathcal{L}_d \cdot Z) =
\deg(f|_Z : Z \to f(Z)) (\mathcal{L}_1 \cdots \mathcal{L}_d \cdot f(Z))
$$
où $\deg(Z \to f(Z))$ est défini comme dans
Morphismes, définition \ref{morphisms-definition-degree},
ou vaut $0$ si $\dim(f(Z)) < d$.
\end{lemma}

\begin{proof}
Le membre de gauche est calculé à l'aide du coefficient de $n_1 \ldots n_d$
dans la fonction
$$
\chi(Y, \mathcal{O}_Z \otimes f^*\mathcal{L}_1^{\otimes n_1} \otimes
\ldots \otimes f^*\mathcal{L}_d^{\otimes n_d}) =
\sum (-1)^i
\chi(X, R^if_*\mathcal{O}_Z \otimes
\mathcal{L}_1^{\otimes n_1} \otimes \ldots \otimes
\mathcal{L}_d^{\otimes n_d})
$$
L'égalité résulte du lemme \ref{lemma-euler-characteristic-morphism}
et de la formule de projection
(Cohomologie, lemme \ref{cohomology-lemma-projection-formula}).
Si $f(Z)$ est de dimension $< d$, alors le membre de droite
est un polynôme de degré total $< d$ d'après le
lemme \ref{lemma-numerical-polynomial-from-euler},
et le résultat est établi. Supposons que $\dim(f(Z)) = d$. Soit
$\xi \in f(Z)$ le point générique. D'après la
théorie de la dimension (voir les lemmes \ref{lemma-dimension-locally-algebraic} et
\ref{lemma-dimension-fibres-locally-algebraic}),
le point générique de $Z$ est l'unique point de $Z$ envoyé sur $\xi$.
Alors $f : Z \to f(Z)$ est fini au-dessus d'un ouvert non vide de $f(Z)$, voir
Morphismes, lemme \ref{morphisms-lemma-generically-finite}.
Ainsi $\deg(f : Z \to f(Z))$ est défini et est en fait égal
à la longueur de la fibre de $f_*\mathcal{O}_Z$ en $\xi$, considérée
comme module sur $\mathcal{O}_{X, \xi}$. De plus, la fibre de
$R^if_*\mathcal{O}_X$ en $\xi$ est nulle pour $i > 0$, car
nous venons de voir que $f|_Z$ est fini dans un voisinage de $\xi$
(de sorte que Cohomologie des Schémas, lemme
\ref{coherent-lemma-finite-pushforward-coherent} donne cette annulation).
Ainsi, les termes $\chi(X, R^if_*\mathcal{O}_Z \otimes
\mathcal{L}_1^{\otimes n_1} \otimes \ldots \otimes
\mathcal{L}_d^{\otimes n_d})$ avec $i > 0$ sont de degré total
$< d$, et
$$
\chi(X, f_*\mathcal{O}_Z \otimes
\mathcal{L}_1^{\otimes n_1} \otimes \ldots \otimes
\mathcal{L}_d^{\otimes n_d})
=
\deg(f : Z \to f(Z)) \chi(f(Z),
\mathcal{L}_1^{\otimes n_1} \otimes \ldots \otimes
\mathcal{L}_d^{\otimes n_d}|_{f(Z)})
$$
modulo un polynôme de degré total $< d$, d'après le
lemme \ref{lemma-numerical-polynomial-leading-term}.
Le résultat voulu en résulte.
\end{proof}

\begin{lemma}
\label{lemma-numerical-intersection-effective-Cartier-divisor}
Soit $k$ un corps. Soit $X$ propre sur $k$. Soit $Z \subset X$ un
sous-schéma fermé de dimension $d$. Soient $\mathcal{L}_1, \ldots, \mathcal{L}_d$
des $\mathcal{O}_X$-modules inversibles. Supposons qu'il existe un
diviseur de Cartier effectif $D \subset Z$ tel que
$\mathcal{L}_1|_Z \cong \mathcal{O}_Z(D)$. Alors
$$
(\mathcal{L}_1 \cdots \mathcal{L}_d \cdot Z) =
(\mathcal{L}_2 \cdots \mathcal{L}_d \cdot D)
$$
\end{lemma}

\begin{proof}
On peut remplacer $X$ par $Z$ et $\mathcal{L}_i$ par $\mathcal{L}_i|_Z$.
On peut donc supposer que $X = Z$ et $\mathcal{L}_1 = \mathcal{O}_X(D)$.
Alors $\mathcal{L}_1^{-1}$ est le faisceau d'idéaux de $D$ et l'on peut
considérer la suite exacte courte
$$
0 \to \mathcal{L}_1^{\otimes -1} \to \mathcal{O}_X \to \mathcal{O}_D \to 0
$$
Posons
$P(n_1, \ldots, n_d) =
\chi(X, \mathcal{L}_1^{\otimes n_1} \otimes \ldots \otimes
\mathcal{L}_d^{\otimes n_d})$
et
$Q(n_1, \ldots, n_d) =
\chi(D, \mathcal{L}_1^{\otimes n_1} \otimes \ldots \otimes
\mathcal{L}_d^{\otimes n_d}|_D)$.
L'additivité donne alors
$$
P(n_1, \ldots, n_d) - P(n_1 - 1, n_2, \ldots, n_d) =
Q(n_1, \ldots, n_d)
$$
Comme le degré total de $P$ est au plus $d$, on voit que
le coefficient de $n_1 \ldots n_d$ dans $P$ est égal au coefficient
de $n_2 \ldots n_d$ dans $Q$.
\end{proof}

\begin{lemma}
\label{lemma-ample-positive}
Soit $k$ un corps. Soit $X$ propre sur $k$. Soit $Z \subset X$ un
sous-schéma fermé de dimension $d$. Si $\mathcal{L}_1, \ldots, \mathcal{L}_d$
sont amples, alors $(\mathcal{L}_1 \cdots \mathcal{L}_d \cdot Z)$ est strictement positif.
\end{lemma}

\begin{proof}
Raisonnons par récurrence sur $d$. Le cas $d = 0$
résulte du lemme \ref{lemma-chi-tensor-finite}. Supposons $d > 0$.
D'après le lemme \ref{lemma-intersection-number-in-terms-of-components},
on peut supposer que $Z$ est un sous-schéma fermé intègre.
En fait, on peut remplacer $X$ par $Z$ et $\mathcal{L}_i$
par $\mathcal{L}_i|_Z$ pour se ramener au cas où $Z = X$ est une
variété propre de dimension $d$.
D'après le lemme \ref{lemma-intersection-number-additive},
on peut remplacer $\mathcal{L}_1$ par une puissance tensorielle positive.
On peut donc supposer qu'il existe une section non nulle
$s \in \Gamma(X, \mathcal{L}_1)$
telle que $X_s$ soit affine (on utilise ici la définition d'un
faisceau inversible ample, voir
Propriétés, définition \ref{properties-definition-ample}).
Remarquons que $X$ n'est pas affine, car être propre et affine
implique être fini (Morphismes, lemme \ref{morphisms-lemma-finite-proper}),
ce qui contredit $d > 0$. Il s'ensuit que $s$ possède un schéma des zéros non vide
$Z(s) \subset X$. Comme $X$ est une variété, $s$ est une section régulière
de $\mathcal{L}_1$ ; ainsi $Z(s)$ est un diviseur de Cartier effectif,
donc $Z(s)$ est de codimension $1$ dans $X$, et
par conséquent $Z(s)$ est de dimension $d - 1$ (on utilise ici les résultats de
Diviseurs, sections \ref{divisors-section-effective-Cartier-divisors},
\ref{divisors-section-effective-Cartier-invertible} et
\ref{divisors-section-Noetherian-effective-Cartier},
ainsi que la théorie de la dimension comme dans le lemme \ref{lemma-dimension-locally-algebraic}).
D'après le lemme \ref{lemma-numerical-intersection-effective-Cartier-divisor},
on a
$$
(\mathcal{L}_1 \cdots \mathcal{L}_d \cdot X) =
(\mathcal{L}_2 \cdots \mathcal{L}_d \cdot Z(s))
$$
Le membre de droite est strictement positif par récurrence, ce qui achève la preuve.
\end{proof}

\begin{definition}
\label{definition-degree}
Soit $k$ un corps. Soit $X$ un schéma propre sur $k$. Soit
$\mathcal{L}$ un $\mathcal{O}_X$-module inversible ample.
Pour tout sous-schéma fermé $Z$, son {\it degré relativement à
$\mathcal{L}$}, noté $\deg_\mathcal{L}(Z)$, est
le nombre d'intersection $(\mathcal{L}^d \cdot Z)$,
où $d = \dim(Z)$.
\end{definition}

\noindent
D'après le lemme \ref{lemma-ample-positive}, le degré d'un sous-schéma est toujours un
entier strictement positif. Remarquons que $\deg_\mathcal{L}(Z) = d$ si et seulement si
$$
\chi(Z, \mathcal{L}^{\otimes n}|_Z) = \frac{d}{\dim(Z)!} n^{\dim(Z)} + l.o.t
$$
ce que l'on voit en utilisant l'identité
$$
(n_1 + \ldots + n_{\dim(Z)})^{\dim(Z)} =
\dim(Z)!\ n_1 \ldots n_{\dim(Z)} + \text{autres termes}
$$

\begin{lemma}
\label{lemma-degree-finite-morphism-in-terms-degrees}
Soit $k$ un corps. Soit $f : Y \to X$ un morphisme fini
dominant de variétés propres sur $k$. Soit $\mathcal{L}$
un $\mathcal{O}_X$-module inversible ample.
Alors
$$
\deg_{f^*\mathcal{L}}(Y) = \deg(f) \deg_\mathcal{L}(X)
$$
où $\deg(f)$ est défini dans
Morphismes, définition \ref{morphisms-definition-degree}.
\end{lemma}

\begin{proof}
L'énoncé a un sens, car $f^*\mathcal{L}$ est ample d'après
Morphismes, lemme \ref{morphisms-lemma-pullback-ample-tensor-relatively-ample}.
Cela étant, le résultat est un cas particulier du
lemme \ref{lemma-intersection-number-and-pullback}.
\end{proof}

\noindent
Enfin, nous mettons en relation le nombre d'intersection avec une courbe et la notion
de degré des modules inversibles sur les courbes introduite dans la
section \ref{section-divisors-curves}.

\begin{lemma}
\label{lemma-intersection-numbers-and-degrees-on-curves}
Soit $k$ un corps. Soit $X$ un schéma propre sur $k$.
Soit $Z \subset X$ un sous-schéma fermé de dimension $\leq 1$.
Soit $\mathcal{L}$ un $\mathcal{O}_X$-module inversible.
Alors
$$
(\mathcal{L} \cdot Z) = \deg(\mathcal{L}|_Z)
$$
où $\deg(\mathcal{L}|_Z)$ est défini dans la
définition \ref{definition-degree-invertible-sheaf}.
Si $\mathcal{L}$ est ample, alors
$\deg_\mathcal{L}(Z) = \deg(\mathcal{L}|_Z)$.
\end{lemma}

\begin{proof}
Cela résulte du fait que la fonction
$n \mapsto \chi(Z, \mathcal{L}|_Z^{\otimes n})$ est de degré $1$
et que son coefficient dominant est donc la différence de valeurs consécutives.
\end{proof}

\begin{proposition}[Riemann-Roch asymptotique]
\label{proposition-asymptotic-riemann-roch}
Soit $k$ un corps. Soit $X$ un schéma propre sur $k$ de dimension $d$.
Soit $\mathcal{L}$ un $\mathcal{O}_X$-module inversible ample.
Alors
$$
\dim_k \Gamma(X, \mathcal{L}^{\otimes n}) \sim c n^d + l.o.t.
$$
où $c = \deg_\mathcal{L}(X)/d!$ est une constante strictement positive.
\end{proposition}

\begin{proof}
Cela résulte des définitions,
du lemme \ref{lemma-ample-positive} et de l'annulation
de la cohomologie en degrés supérieurs donnée dans
Cohomologie des Schémas, lemme \ref{coherent-lemma-vanshing-gives-ample}.
\end{proof}






\section{Dimension de plongement}
\label{section-embedding-dimension}

\noindent
Il existe plusieurs façons de définir la dimension de plongement, mais, pour
les points fermés des schémas algébriques sur des corps algébriquement clos,
toutes ces définitions sont équivalentes à la suivante.

\begin{definition}
\label{definition-embed-dim}
Soit $k$ un corps algébriquement clos. Soit $X$ un
$k$-schéma localement algébrique et soit $x \in X$ un point fermé. La
{\it dimension de plongement de $X$ en $x$} est
$\dim_k \mathfrak m_x/\mathfrak m_x^2$.
\end{definition}

\noindent
Quelques faits sur la dimension de plongement.
Soient $k, X, x$ comme dans la définition \ref{definition-embed-dim}.
\begin{enumerate}
\item La dimension de plongement de $X$ en $x$ est
la dimension de
l'espace tangent $T_{X/k, x}$ (définition \ref{definition-tangent-space})
comme $k$-espace vectoriel.
\item La dimension de plongement de $X$ en $x$ est le plus petit entier
$d \geq 0$ tel qu'il existe une surjection
$$
k[[x_1, \ldots, x_d]] \longrightarrow \mathcal{O}_{X, x}^\wedge
$$
de $k$-algèbres.
\item La dimension de plongement de $X$ en $x$ est le plus petit entier
$d \geq 0$ tel qu'il existe un voisinage ouvert $U \subset X$
de $x$ et une immersion fermée $U \to Y$, où $Y$ est une variété lisse
de dimension $d$ sur $k$.
\item La dimension de plongement de $X$ en $x$ est le plus petit entier
$d \geq 0$ tel qu'il existe un voisinage ouvert $U \subset X$
de $x$ et un morphisme non ramifié $U \to \mathbf{A}^d_k$.
\item Si l'on se donne une immersion fermée $X \to Y$ avec $Y$ lisse
sur $k$, alors la dimension de plongement de $X$ en $x$ est le plus petit
entier $d \geq 0$ tel qu'il existe un sous-schéma fermé $Z \subset Y$
tel que $X \subset Z$, que $Z \to \Spec(k)$ soit lisse en $x$ et que
$\dim_x(Z) = d$.
\end{enumerate}
Si ces faits nous sont nécessaires, nous formulerons un énoncé précis et en donnerons
une démonstration.

\medskip\noindent
Corps de base non algébriquement clos ou points non fermés.
Soit $k$ un corps et soit $X$ un $k$-schéma localement algébrique.
Si $x \in X$ est un point, plusieurs choix sont possibles pour la
dimension de plongement de $X$ en $x$. On peut en effet prendre
\begin{enumerate}
\item $\dim_{\kappa(x)}(\mathfrak m_x/\mathfrak m_x^2)$,
\item $\dim_{\kappa(x)}(T_{X/k, x}) =
\dim_{\kappa(x)}(\Omega_{X/k, x} \otimes_{\mathcal{O}_{X, x}} \kappa(x))$
(lemme \ref{lemma-tangent-space-cotangent-space}),
\item le plus petit entier $d \geq 0$ tel qu'il existe un
voisinage ouvert $U \subset X$ de $x$ et une immersion fermée
$U \to Y$, où $Y$ est une variété lisse de dimension $d$ sur $k$.
\end{enumerate}
En caractéristique zéro, (1) $=$ (2) si $x$ est un point fermé ;
plus généralement, cette égalité vaut si $\kappa(x)$ est algébrique séparable
sur $k$, voir le
lemme \ref{lemma-tangent-space-rational-point}.
La définition géométrique (3) semble correspondre le mieux à
l'intuition géométrique qu'évoque l'expression
``dimension de plongement''.
Puisque l'on peut montrer que (3) et (2) définissent le même nombre
(cela résulte du lemme \ref{lemma-immersion-into-affine})
c'est ce que nous utiliserons.
Dans notre terminologie, nous préciserons que nous
prenons la dimension de plongement relativement au corps de base.

\begin{definition}
\label{definition-embedding-dimension}
Soit $k$ un corps. Soit $X$ un $k$-schéma localement algébrique.
Soit $x \in X$ un point. La {\it dimension de plongement de $X/k$ en $x$}
est $\dim_{\kappa(x)}(T_{X/k, x})$.
\end{definition}

\noindent
Si $(A, \mathfrak m, \kappa)$ est un anneau local noethérien,
la {\it dimension de plongement de $A$} est parfois définie comme la dimension de
$\mathfrak m/\mathfrak m^2$ sur $\kappa$. Nous avons vu ci-dessus que,
si $A$ est donné comme algèbre sur un corps $k$, il peut être préférable
d'utiliser $\dim_\kappa(\Omega_{A/k} \otimes_A \kappa)$. Appelons
cette quantité la {\it dimension de plongement de $A/k$}.
Avec cette terminologie, on a
$$
\begin{aligned}
\text{dimension de plongement de }X/k\text{ en }x
&= \text{dimension de plongement de }\mathcal{O}_{X, x}/k \\
&= \text{dimension de plongement de }\mathcal{O}_{X, x}^\wedge/k
\end{aligned}
$$
si $k, X, x$ sont comme dans la définition \ref{definition-embedding-dimension}.






\section{Théorèmes de Bertini}
\label{section-bertini}

\noindent
Dans cette section, nous démontrons des résultats de la forme suivante : étant donnée une variété
projective lisse $X$ sur un corps $k$, il existe un diviseur ample $H \subset X$
qui est lisse.

\begin{lemma}
\label{lemma-generate-over-complement}
Soit $k$ un corps. Soit $X$ un schéma propre sur $k$. Soit $\mathcal{L}$
un $\mathcal{O}_X$-module inversible ample. Soit $Z \subset X$ un
sous-schéma fermé. Alors il existe un entier $n_0$ tel que, pour tout
$n \geq n_0$, le noyau $V_n$ de
$\Gamma(X, \mathcal{L}^{\otimes n}) \to \Gamma(Z, \mathcal{L}^{\otimes n}|_Z)$
engendre $\mathcal{L}^{\otimes n}|_{X \setminus Z}$ et
le morphisme canonique
$$
X \setminus Z \longrightarrow \mathbf{P}(V_n)
$$
est une immersion de schémas sur $k$.
\end{lemma}

\begin{proof}
Soit $\mathcal{I} \subset \mathcal{O}_X$ le faisceau quasi-cohérent d'idéaux de
$Z$. Remarquons que, via l'inclusion
$\mathcal{I} \otimes_{\mathcal{O}_X} \mathcal{L}^{\otimes n} \subset
\mathcal{L}^{\otimes n}$, on a
$V_n = \Gamma(X, \mathcal{I} \otimes_{\mathcal{O}_X} \mathcal{L}^{\otimes n})$.
Choisissons $n_1$ tel que, pour $n \geq n_1$, le faisceau
$\mathcal{I} \otimes \mathcal{L}^{\otimes n}$ soit engendré par ses sections globales, voir
Propriétés, proposition \ref{properties-proposition-characterize-ample}.
Il en résulte que $V_n$ engendre $\mathcal{L}^{\otimes n}|_{X \setminus Z}$
pour $n \geq n_1$.

\medskip\noindent
Pour $n \geq n_1$, notons
$\psi_n : V_n \to
\Gamma(X \setminus Z, \mathcal{L}^{\otimes n}|_{X \setminus Z})$
l'application de restriction. On obtient un morphisme canonique
$$
\varphi = \varphi_{\mathcal{L}^{\otimes n}|_{X \setminus Z}, \psi_n} :
X \setminus Z
\longrightarrow
\mathbf{P}(V_n)
$$
d'après Constructions, exemple \ref{constructions-example-projective-space}.
Choisissons $n_2$ tel que, pour tout $n \geq n_2$, le faisceau inversible
$\mathcal{L}^{\otimes n}$ soit très ample sur $X$.
Nous affirmons que $n_0 = n_1 + n_2$ convient.

\medskip\noindent
Preuve de l'affirmation. Supposons $n \geq n_0$ et écrivons $n = n_1 + n'$.
Pour $x \in X \setminus Z$, on peut choisir $s_1 \in V_1$ ne
s'annulant pas en $x$. Posons $V' = \Gamma(X, \mathcal{L}^{\otimes n'})$.
Par le choix de $n$ et de $n'$, le morphisme correspondant
$\varphi' : X \to \mathbf{P}(V')$ est une immersion fermée. Ainsi, si l'on choisit
$s' \in \Gamma(X, \mathcal{L}^{\otimes n'})$ ne s'annulant pas en $x$,
alors $X_{s'} = (\varphi')^{-1}(D_+(s'))$ (voir
Constructions, lemme \ref{constructions-lemma-invertible-map-into-proj})
est affine et $X_{s'} \to D_+(s')$ est une immersion fermée.
Alors $s = s_1 \otimes s' \in V_n$ ne s'annule pas en $x$.
Si $D_+(s) \subset \mathbf{P}(V_n)$ désigne
l'ouvert affine correspondant de notre espace projectif, alors
$\varphi^{-1}(D_+(s)) = X_s \subset X \setminus Z$ (voir la référence ci-dessus).
L'ouvert $X_s = X_{s'} \cap X_{s_1}$ est affine, voir
Propriétés, lemme \ref{properties-lemma-affine-cap-s-open}.
Considérons l'homomorphisme d'anneaux
$$
\text{Sym}(V)_{(s)} \longrightarrow \mathcal{O}_X(X_s)
$$
qui définit le morphisme $X_s \to D_+(s)$. Comme $X_{s'} \to D_+(s')$
est une immersion fermée, les images des éléments
$$
\frac{s_1 \otimes t'}{s_1 \otimes s'}
$$
où $t' \in V'$, engendrent l'image de
$\mathcal{O}_X(X_{s'}) \to \mathcal{O}_X(X_s)$.
Puisque $X_s \to X_{s'}$ est une immersion ouverte,
il s'ensuit que $X_s \to D_+(s)$ est une immersion de schémas affines
(voir ci-dessous). Ainsi, $\varphi_n$ est une immersion d'après
Morphismes, lemme \ref{morphisms-lemma-check-immersion}.

\medskip\noindent
Soient $a : A' \to A$ et $c : B \to A$ des homomorphismes d'anneaux tels que
$\Spec(a)$ soit une immersion et $\Im(a) \subset \Im(c)$.
Posons $B' = A' \times_A B$, avec les projections $b : B' \to B$
et $c' : B' \to A'$. Par hypothèse, $c'$
est surjectif, de sorte que $\Spec(c')$ est une immersion fermée.
Par conséquent, $\Spec(c') \circ \Spec(a)$ est une immersion
(Schémas, lemme \ref{schemes-lemma-composition-immersion}).
Alors $\Spec(c)$ est nécessairement une immersion, car il factorise
l'immersion $\Spec(c') \circ \Spec(a) = \Spec(b) \circ \Spec(c)$, voir
Morphismes, lemme \ref{morphisms-lemma-immersion-permanence}.
\end{proof}

\begin{situation}
\label{situation-family-divisors}
Soit $k$ un corps,
soit $X$ un schéma sur $k$,
soit $\mathcal{L}$ un $\mathcal{O}_X$-module inversible,
soit $V$ un $k$-espace vectoriel de dimension finie, et
soit $\psi : V \to \Gamma(X, \mathcal{L})$ une application $k$-linéaire.
Supposons $\dim(V) = r$ et fixons une base $v_1, \ldots, v_r$ de $V$.
On obtient alors un ``diviseur universel''
$$
H_{univ} = Z(s_{univ}) \subset  \mathbf{A}^r \times_k X
$$
qui est le schéma des zéros
(Diviseurs, définition \ref{divisors-definition-zero-scheme-s}) de la section
$$
s_{univ} = \sum\nolimits_{i = 1, \ldots, r} x_i \psi(v_i) \in
\Gamma(\mathbf{A}^r \times_k X, \text{pr}_2^*\mathcal{L})
$$
Pour une extension de corps $k'/k$, les $k'$-points $v \in \mathbf{A}^r_k(k')$
correspondent aux vecteurs $(a_1, \ldots, a_r)$ d'éléments de $k'$.
Ainsi, on peut d'une part considérer $v$ comme l'élément
$v = \sum_{i = 1, \ldots, r} a_i v_i \in V \otimes_k k'$
et, d'autre part, associer à $v$ la section
$$
\psi(v) = \sum\nolimits_{i = 1, \ldots, r} a_i \psi(v_i) \in
\Gamma(X_{k'}, \mathcal{L}|_{X_{k'}})
$$
Avec ces notations, il est clair que la fibre de $H_{univ}$ au-dessus de
$v \in V \otimes k'$ est le schéma des zéros de $\psi(v)$.
En formule :
$$
H_v = H_{univ, v} = Z(\psi(v))
$$
Nous désignerons cette valeur commune par $H_v$, comme indiqué. Enfin, dans
cette situation, soit $P$ une propriété des vecteurs $v \in V \otimes_k k'$
pour une extension de corps arbitraire $k'/k$\footnote{Par exemple, nous pourrions
considérer la condition que $H_v$ soit lisse sur $k'$, ou géométriquement
irréductible sur $k'$.}. Nous disons que $P$ {\it est vérifiée pour un vecteur général}
$v \in V \otimes_k k'$ s'il existe un ouvert de Zariski non vide
$U \subset \mathbf{A}^r_k$ tel que, si $v$ correspond à un $k'$-point
de $U$, quelle que soit l'extension $k'/k$, alors $P(v)$ est vérifiée.
\end{situation}

\begin{lemma}
\label{lemma-bertini}
Dans la situation \ref{situation-family-divisors}, supposons que
\begin{enumerate}
\item $X$ soit lisse sur $k$,
\item l'image de $\psi : V \to \Gamma(X, \mathcal{L})$
engendre $\mathcal{L}$,
\item le morphisme correspondant
$\varphi_{\mathcal{L}, \psi} : X \to \mathbf{P}(V)$ soit
une immersion.
\end{enumerate}
Alors, pour un vecteur général $v \in V \otimes_k k'$,
le schéma $H_v$ est lisse sur $k'$.
\end{lemma}

\begin{proof}
(Nous observons que $X$ est séparé et de type fini comme sous-schéma
localement fermé d'un espace projectif.)
Utilisons les notations introduites avant l'énoncé du lemme.
Considérons les projections
$$
\xymatrix{
\mathbf{A}^r_k \times_k X \ar[d] &
H_{univ} \ar[l] \ar[ld]^p \ar[r] \ar[rd]_q &
\mathbf{A}^r_k \times_k X \ar[d] \\
X & &
\mathbf{A}^r_k
}
$$
Soit $\Sigma \subset H_{univ}$ le lieu singulier du morphisme
$q : H_{univ} \to \mathbf{A}^r_k$, c'est-à-dire l'ensemble des points où
$q$ n'est pas lisse. Alors $\Sigma$ est fermé, car le lieu lisse
d'un morphisme est ouvert par définition. Comme toute fibre d'un
morphisme lisse est lisse, il suffit de montrer que $q(\Sigma)$
est contenu dans un sous-ensemble fermé strict de $\mathbf{A}^r_k$.
Puisque $\Sigma$ (muni de la structure de schéma induite réduite) est un
schéma de type fini sur $k$, il suffit de montrer que $\dim(\Sigma) < r$.
Cela résulte du lemme \ref{lemma-dimension-fibres-locally-algebraic}.
Comme les dimensions ne changent pas lorsqu'on remplace $k$ par un corps plus grand
(Morphismes, lemme \ref{morphisms-lemma-dimension-fibre-after-base-change}),
nous pouvons supposer, et nous supposons désormais, que $k$ est algébriquement clos.
D'après la théorie de la dimension (lemme \ref{lemma-dimension-fibres-locally-algebraic}),
il suffit de montrer que, pour tout point fermé $x \in X \setminus Z$,
$p^{-1}(\{x\}) \cap \Sigma$ est de dimension $< r - \dim(X')$,
où $X'$ est l'unique composante irréductible de $X$ contenant $x$.
Comme $X$ est lisse sur $k$ et que $x$ est un point fermé, on a
$\dim(X') = \dim \mathfrak m_x/\mathfrak m_x^2$
(Morphismes, lemme \ref{morphisms-lemma-smooth-omega-finite-locally-free}
et Algèbre, lemme \ref{algebra-lemma-rank-omega}).
Il suffit donc que
$$
\dim p^{-1}(x) \cap \Sigma < r - \dim \mathfrak m_x/\mathfrak m_x^2
$$
pour tout point fermé $x \in X$.

\medskip\noindent
Puisque les sections globales provenant de $V$ engendrent $\mathcal{L}$, pour toute composante irréductible
$X'$ de $X$, il existe un ouvert de Zariski non vide de $\mathbf{A}^r$ tel que les
fibres de $q$ au-dessus de cet ouvert ne contiennent pas $X'$. (Par exemple, si $x' \in X'$
est un point fermé, on peut prendre l'ouvert correspondant aux
vecteurs $v \in V$ tels que $\psi(v)$ ne s'annule pas en $x'$. Cet
ouvert est le complémentaire d'un hyperplan dans $\mathbf{A}^r_k$.)
Soit $U \subset \mathbf{A}^r$ l'intersection (non vide) de ces ouverts.
Alors les fibres de $q^{-1}(U) \to U$ sont des diviseurs de Cartier effectifs
sur les fibres de $U \times_k X \to U$ (car une section non nulle
d'un module inversible sur un schéma intègre est une section régulière). Par suite, le
morphisme $q^{-1}(U) \to U$ est plat d'après
Diviseurs, lemme \ref{divisors-lemma-fibre-Cartier}.
Ainsi, pour tout point fermé $x \in X$ et tout $v \in V = \mathbf{A}^r_k(k)$,
si $(x, v) \in H_{univ}$, c'est-à-dire si $x \in H_v$,
alors $q$ est lisse en $(x, v)$ si et seulement si la fibre
$H_v$ est lisse en $x$, voir
Morphismes, lemme \ref{morphisms-lemma-smooth-at-point}.

\medskip\noindent
Considérons l'image $\psi(v)_x$ dans la fibre $\mathcal{L}_x$
de la section correspondant à $v \in V$. On a
$$
x \in H_v \Leftrightarrow \psi(v)_x \in \mathfrak m_x\mathcal{L}_x
$$
Si cela est vrai, alors on a
$$
H_v\text{ est singulier en }x \Leftrightarrow
\psi(v)_x \in \mathfrak m_x^2\mathcal{L}_x
$$
En effet, $\psi(v)_x$ n'appartient pas à $\mathfrak m_x^2\mathcal{L}_x$
$\Leftrightarrow$
l'équation locale de $H_v \subset X$ en $x$ n'appartient pas
à $\mathfrak m_x^2$
$\Leftrightarrow$
$\mathcal{O}_{H_v, x}$ est régulier
(Algèbre, lemme \ref{algebra-lemma-regular-ring-CM})
$\Leftrightarrow$
$H_v$ est lisse en $x$ sur $k$
(Algèbre, lemme \ref{algebra-lemma-separable-smooth}).
On conclut que les points fermés de $p^{-1}(x) \cap \Sigma$ correspondent
aux $v \in V$ tels que $\psi(v)_x \in \mathfrak m_x^2\mathcal{L}_x$.
Cependant, comme $\varphi_{\mathcal{L}, \psi}$ est une immersion,
l'application
$$
V \longrightarrow \mathcal{L}_x/\mathfrak m_x^2\mathcal{L}_x
$$
est surjective (un petit détail est omis). D'après ce qui précède, les points fermés
du lieu $p^{-1}(x) \cap \Sigma$, considéré comme un sous-espace de $V$,
forment le noyau de cette application, lequel a donc pour dimension
$r - \dim \mathfrak m_x/\mathfrak m_x^2 - 1$, comme voulu.
\end{proof}






\section{Enriques-Severi-Zariski}
\label{section-vanishing-negative}

\noindent
Dans cette section, nous démontrons des résultats de la forme suivante : tensoriser par
un module inversible « très négatif » annule la cohomologie en
petits degrés. Nous en déduisons aussi la connexité
d'une section par une hypersurface d'un schéma propre normal de
dimension $\geq 2$.

\begin{lemma}
\label{lemma-vanishin-h0-negative}
Soit $k$ un corps. Soit $X$ un schéma propre sur $k$. Soit $\mathcal{L}$
un $\mathcal{O}_X$-module inversible ample. Soit $\mathcal{F}$ un
$\mathcal{O}_X$-module cohérent. Si $\text{Ass}(\mathcal{F})$ ne
contient aucun point fermé, alors
$\Gamma(X, \mathcal{F} \otimes_{\mathcal{O}_X} \mathcal{L}^{\otimes n}) = 0$
pour $n \ll 0$.
\end{lemma}

\begin{proof}
Pour un $\mathcal{O}_X$-module cohérent $\mathcal{F}$,
notons $\mathcal{P}(\mathcal{F})$ la propriété suivante : il
existe un $n_0 \in \mathbf{Z}$ tel que, pour $n \leq n_0$,
toute section $s$ de
$\mathcal{F} \otimes_{\mathcal{O}_X} \mathcal{L}^{\otimes n}$
ait un support constitué uniquement de points fermés.
Comme $\text{Ass}(\mathcal{F}) =
\text{Ass}(\mathcal{F} \otimes_{\mathcal{O}_X} \mathcal{L}^{\otimes n})$,
on voit qu'il suffit de démontrer que $\mathcal{P}$
est satisfaite par tout module cohérent sur $X$.
Pour ce faire, nous allons démontrer que les conditions (1), (2) et (3) de
Cohomologie des schémas, lemme
\ref{coherent-lemma-property-higher-rank-cohomological},
sont satisfaites.

\medskip\noindent
Pour vérifier la condition (1), supposons que
$$
0 \to \mathcal{F}_1 \to \mathcal{F} \to \mathcal{F}_2 \to 0
$$
soit une suite exacte courte de $\mathcal{O}_X$-modules cohérents
telle que $\mathcal{P}$ soit satisfaite par $\mathcal{F}_i$, $i = 1, 2$.
Soient $n_1, n_2$ les bornes ainsi obtenues.
Soit $\mathcal{F}'_2 \subset \mathcal{F}_2$ le plus grand
sous-module cohérent dont le support est un ensemble fini de points
fermés. Soit $\mathcal{I} \subset \mathcal{O}_X$ l'annulateur
de $\mathcal{F}'_2$. Comme $\mathcal{L}$ est ample, il existe un $e > 0$
tel que $\mathcal{I} \otimes_{\mathcal{O}_X} \mathcal{L}^{\otimes e}$
soit engendré par ses sections globales. Posons $n_0 = \min(n_2, n_1 - e)$. Soit $n \leq n_0$
et soit $t$ une section globale de
$\mathcal{F} \otimes \mathcal{L}^{\otimes n}$.
L'image de $t$ dans $\mathcal{F}_2 \otimes \mathcal{L}^{\otimes n}$
appartient à $\mathcal{F}'_2 \otimes \mathcal{L}^{\otimes n}$ puisque
$n \leq n_2$. Par conséquent, pour tout
$s \in \Gamma(X, \mathcal{I} \otimes_{\mathcal{O}_X} \mathcal{L}^{\otimes e})$,
le produit $t \otimes s$ appartient à
$\mathcal{F}_1 \otimes \mathcal{L}^{\otimes n + e}$.
Ainsi, le support de $t \otimes s$ est contenu dans
l'ensemble fini des points fermés de $\text{Ass}(\mathcal{F}_1)$
puisque $n + e \leq n_1$. Par notre choix de $e$,
pour tout point on peut choisir $s$ qui y soit inversible, sauf s'il
appartient au support de $\mathcal{F}'_2$ ;
on conclut que le support de $t$ est contenu dans la réunion de
l'ensemble fini des points fermés de $\text{Ass}(\mathcal{F}_1)$ et de
l'ensemble fini des points fermés de $\text{Ass}(\mathcal{F}_2)$.
Ceci achève la démonstration de la condition (1).

\medskip\noindent
La condition (2) est immédiate.

\medskip\noindent
Pour la condition (3), choisissons $\mathcal{G} = \mathcal{O}_Z$.
Dans ce cas, si $Z$ est un point fermé de $X$, il n'y a
rien à démontrer. Si $\dim(Z) > 0$, nous allons montrer que
$\Gamma(Z, \mathcal{L}^{\otimes n}|_Z) = 0$
pour $n < 0$. En effet, soit $s$ une section non nulle d'une puissance négative
de $\mathcal{L}|_Z$. Choisissons une section non nulle $t$ d'une puissance positive
de $\mathcal{L}|_Z$ (ce choix est possible puisque $\mathcal{L}$ est ample, voir
Propriétés, proposition \ref{properties-proposition-characterize-ample}).
Alors $s^{\deg(t)} \otimes t^{\deg(s)}$
est une section globale non nulle de $\mathcal{O}_Z$
(puisque $Z$ est intègre), et donc
une unité (lemme \ref{lemma-proper-geometrically-reduced-global-sections}).
Cela implique que $t$ est une section trivialisante
d'une puissance positive de $\mathcal{L}$.
Ainsi, la fonction $n \mapsto \dim_k \Gamma(X, \mathcal{L}^{\otimes n})$
est bornée sur un ensemble infini d'entiers positifs, ce qui contredit
le théorème de Riemann-Roch asymptotique
(proposition \ref{proposition-asymptotic-riemann-roch}),
puisque $\dim(Z) > 0$.
\end{proof}

\begin{lemma}[Enriques-Severi-Zariski]
\label{lemma-vanishin-h1-negative}
Soit $k$ un corps. Soit $X$ un schéma propre sur $k$. Soit $\mathcal{L}$
un $\mathcal{O}_X$-module inversible ample. Soit $\mathcal{F}$ un
$\mathcal{O}_X$-module cohérent. Supposons que,
pour tout point fermé $x \in X$, on ait $\text{profondeur}(\mathcal{F}_x) \geq 2$.
Alors
$H^1(X, \mathcal{F} \otimes_{\mathcal{O}_X} \mathcal{L}^{\otimes m}) = 0$
pour $m \ll 0$.
\end{lemma}

\begin{proof}
Choisissons une immersion fermée $i : X \to \mathbf{P}^n_k$ telle que
$i^*\mathcal{O}(1) \cong \mathcal{L}^{\otimes e}$ pour un certain $e > 0$
(voir Morphismes, lemme
\ref{morphisms-lemma-finite-type-over-affine-ample-very-ample}).
Il suffit alors de démontrer le lemme pour
$$
\mathcal{G} =
i_*(\mathcal{F} \oplus \mathcal{F} \otimes \mathcal{L} \oplus \ldots
\oplus \mathcal{F} \otimes \mathcal{L}^{\otimes e - 1})
\quad\text{et}\quad \mathcal{O}(1)
$$
sur $\mathbf{P}^n_k$. En effet, on a
$$
H^1(\mathbf{P}^n_k, \mathcal{G}(m)) =
\bigoplus\nolimits_{j = 0, \ldots, e - 1}
H^1(X, \mathcal{F} \otimes \mathcal{L}^{\otimes j + me})
$$
d'après
Cohomologie des schémas, lemme \ref{coherent-lemma-relative-affine-cohomology}.
De plus, si $y \in \mathbf{P}^n_k$ est un point fermé, alors
$\text{profondeur}(\mathcal{G}_y) = \infty$ si $y \not \in i(X)$
et $\text{profondeur}(\mathcal{G}_y) = \text{profondeur}(\mathcal{F}_x)$
si $y = i(x)$ puisque, dans ce cas,
$\mathcal{G}_y \cong \mathcal{F}_x^{\oplus e}$ comme module sur
$\mathcal{O}_{\mathbf{P}^n_k, x}$ et que l'on peut, par exemple, utiliser
Algèbre, lemme \ref{algebra-lemma-depth-goes-down-finite},
pour obtenir l'égalité.

\medskip\noindent
Supposons que $X = \mathbf{P}^n_k$, que $\mathcal{L} = \mathcal{O}(1)$ et que
$k$ soit infini. Choisissons $s \in H^0(\mathbf{P}^1_k, \mathcal{O}(1))$
qui détermine une suite exacte
$$
0 \to \mathcal{F}(-1) \xrightarrow{s} \mathcal{F} \to \mathcal{G} \to 0
$$
comme au lemme \ref{lemma-exact-sequence-induction}. Puisque le morphisme
$\mathcal{F}(-1) \to \mathcal{F}$ est, localement sur des ouverts affines, donné par
la multiplication par un non-diviseur de zéro dans $\mathcal{F}$,
on voit que, pour tout point fermé $x \in \mathbf{P}^n_k$, on a
$\text{profondeur}(\mathcal{G}_x) \geq 1$, voir
Algèbre, lemme \ref{algebra-lemma-depth-drops-by-one}.
Ainsi, d'après le lemme \ref{lemma-vanishin-h0-negative},
on a $H^0(\mathcal{G}(m)) = 0$ pour $m \ll 0$.
En considérant la suite exacte longue de cohomologie après tensorisation
(voir remarque \ref{remark-exact-sequence-induction-cohomology}),
on constate que la suite de nombres
$$
\dim H^1(\mathbf{P}^n_k, \mathcal{F}(m))
$$
se stabilise pour $m \leq m_0$, pour un certain entier $m_0$.
Soit $N$ la dimension commune de ces
espaces pour $m \leq m_0$. Il faut montrer que $N = 0$.

\medskip\noindent
Pour $d > 0$ et $m \leq m_0$, considérons l'application bilinéaire
$$
H^0(\mathbf{P}^n_k, \mathcal{O}(d)) \times
H^1(\mathbf{P}^n_k, \mathcal{F}(m - d))
\longrightarrow
H^1(\mathbf{P}^n_k, \mathcal{F}(m))
$$
D'après l'algèbre linéaire, il existe un sous-espace de codimension $\leq N^2$ que l'on note
$V_m \subset H^0(\mathbf{P}^n_k, \mathcal{O}(d))$ et tel que
la multiplication par tout $s' \in V_m$ annule
$H^1(\mathbf{P}^n_k, \mathcal{F}(m - d))$.
Notons que, pour $m' < m \leq m_0$, le diagramme
$$
\xymatrix{
H^0(\mathbf{P}^n_k, \mathcal{O}(d)) \times
H^1(\mathbf{P}^n_k, \mathcal{F}(m' - d)) \ar[r] \ar[d]^{1 \times s^{m' - m}} &
H^1(\mathbf{P}^n_k, \mathcal{F}(m')) \ar[d]^{s^{m' - m}}\\
H^0(\mathbf{P}^n_k, \mathcal{O}(d)) \times
H^1(\mathbf{P}^n_k, \mathcal{F}(m - d)) \ar[r] &
H^1(\mathbf{P}^n_k, \mathcal{F}(m))
}
$$
est commutatif, et les flèches verticales sont des isomorphismes. Ainsi, $V_m = V$ est
indépendant de $m \leq m_0$. Pour $x \in \text{Ass}(\mathcal{F})$,
posons $Z = \overline{\{x\}}$. Pour $d$ suffisamment grand, l'application linéaire
$$
H^0(\mathbf{P}^n_k, \mathcal{O}(d)) \to H^0(Z, \mathcal{O}(d)|_Z)
$$
est de rang $> N^2$ puisque $\dim(Z) \geq 1$ (cela résulte, par exemple,
du théorème de Riemann-Roch asymptotique et de l'amplitude de $\mathcal{O}(1)$ ; détails
omis). On peut donc trouver $s' \in V$ tel que $s'$ ne s'annule
en aucun point associé de $\mathcal{F}$ (on utilise le fait que l'ensemble
des points associés est fini). On obtient alors
$$
0 \to \mathcal{F}(-d) \xrightarrow{s'} \mathcal{F} \to \mathcal{G}' \to 0
$$
et, comme précédemment, on conclut que la multiplication par $s'$
sur $H^1(\mathbf{P}^n_k, \mathcal{F}(m - d))$ est injective
pour $m \ll 0$. Cela contredit le choix de $s'$, à moins que
$N = 0$, comme voulu.

\medskip\noindent
Il reste à traiter le cas où $k$ est fini.
Dans ce cas, soit $K/k$ une extension algébrique infinie.
Notons $\mathcal{F}_K$ et $\mathcal{L}_K$ les images réciproques
de $\mathcal{F}$ et $\mathcal{L}$ sur $X_K = \Spec(K) \times_{\Spec(k)} X$.
On a
$$
H^1(X_K, \mathcal{F}_K \otimes \mathcal{L}_K^{\otimes m}) =
H^1(X, \mathcal{F} \otimes \mathcal{L}^{\otimes m}) \otimes_k K
$$
d'après Cohomologie des schémas, lemme
\ref{coherent-lemma-flat-base-change-cohomology}.
D'autre part, un point fermé $x_K$ de $X_K$ a pour image un point fermé
$x$ de $X$ puisque $K/k$ est une extension algébrique.
L'homomorphisme d'anneaux $\mathcal{O}_{X, x} \to \mathcal{O}_{X_K, x_K}$
est plat (lemme \ref{lemma-change-fields-flat}). Par conséquent, on a
$$
\text{profondeur}(\mathcal{F}_{x_K}) =
\text{profondeur}(\mathcal{F}_x \otimes_{\mathcal{O}_{X, x}}
\mathcal{O}_{X_K, x_K}) \geq
\text{profondeur}(\mathcal{F}_x)
$$
d'après Algèbre, lemme \ref{algebra-lemma-apply-grothendieck-module}
(en fait, on a égalité ici, mais nous n'en avons pas besoin).
Le résultat sur $k$ découle donc
du résultat sur le corps infini $K$, ce qui achève la démonstration.
\end{proof}

\begin{lemma}
\label{lemma-connectedness-ample-divisor}
Soit $k$ un corps. Soit $X$ un schéma propre sur $k$. Soit $\mathcal{L}$
un $\mathcal{O}_X$-module inversible ample. Soit
$s \in \Gamma(X, \mathcal{L})$. Supposons que
\begin{enumerate}
\item $s$ soit une section régulière
(Diviseurs, définition \ref{divisors-definition-regular-section}),
\item pour tout point fermé $x \in X$, on ait
$\text{profondeur}(\mathcal{O}_{X, x}) \geq 2$, et
\item $X$ soit connexe.
\end{enumerate}
Alors le schéma des zéros $Z(s)$ de $s$ est connexe.
\end{lemma}

\begin{proof}
Puisque $s$ est une section régulière, il en est de même de
$s^n \in \Gamma(X, \mathcal{L}^{\otimes n})$ pour tout $n > 1$.
De plus, le morphisme d'inclusion $Z(s) \to Z(s^n)$ induit une bijection
sur les espaces topologiques sous-jacents. Par conséquent, si $Z(s)$ est non connexe,
$Z(s^n)$ l'est aussi. Considérons maintenant la suite exacte courte canonique
$$
0 \to \mathcal{L}^{\otimes -n} \xrightarrow{s^n}
\mathcal{O}_X \to \mathcal{O}_{Z(s^n)} \to 0
$$
Considérons la $k$-algèbre $R_n = \Gamma(X, \mathcal{O}_{Z(s^n)})$.
Si $Z(s)$ est non connexe, c'est-à-dire si $Z(s^n)$ est non connexe,
alors soit $R_n$ est nul lorsque $Z(s^n) = \emptyset$, soit
$R_n$ contient un idempotent non trivial lorsque $Z(s^n) = U \amalg V$, avec
$U, V \subset Z(s^n)$ ouverts et non vides (le lecteur pourra consulter le
lemme \ref{lemma-proper-geometrically-reduced-global-sections}).
Ainsi, l'homomorphisme $\Gamma(X, \mathcal{O}_X) \to R_n$ ne peut être un isomorphisme.
Il en résulte que l'un des deux groupes $H^0(X, \mathcal{L}^{\otimes -n})$ et
$H^1(X, \mathcal{L}^{\otimes -n})$ est non nul pour une infinité
d'entiers positifs $n$. Cela contredit le lemme \ref{lemma-vanishin-h0-negative} ou
\ref{lemma-vanishin-h1-negative},
ce qui achève la démonstration.
\end{proof}








\begin{multicols}{2}[\section{Autres chapitres}]
\noindent
Préliminaires
\begin{enumerate}
\item \hyperref[introduction-section-phantom]{Introduction}
\item \hyperref[conventions-section-phantom]{Conventions}
\item \hyperref[sets-section-phantom]{Théorie des ensembles}
\item \hyperref[categories-section-phantom]{Catégories}
\item \hyperref[topology-section-phantom]{Topologie}
\item \hyperref[sheaves-section-phantom]{Faisceaux sur les espaces}
\item \hyperref[sites-section-phantom]{Sites et faisceaux}
\item \hyperref[stacks-section-phantom]{Champs}
\item \hyperref[fields-section-phantom]{Corps}
\item \hyperref[algebra-section-phantom]{Algèbre commutative}
\item \hyperref[brauer-section-phantom]{Groupes de Brauer}
\item \hyperref[homology-section-phantom]{Algèbre homologique}
\item \hyperref[derived-section-phantom]{Catégories dérivées}
\item \hyperref[simplicial-section-phantom]{Méthodes simpliciales}
\item \hyperref[more-algebra-section-phantom]{Compléments d'algèbre}
\item \hyperref[smoothing-section-phantom]{Lissage des morphismes d'anneaux}
\item \hyperref[modules-section-phantom]{Faisceaux de modules}
\item \hyperref[sites-modules-section-phantom]{Modules sur les sites}
\item \hyperref[injectives-section-phantom]{Objets injectifs}
\item \hyperref[cohomology-section-phantom]{Cohomologie des faisceaux}
\item \hyperref[sites-cohomology-section-phantom]{Cohomologie sur les sites}
\item \hyperref[dga-section-phantom]{Algèbre différentielle graduée}
\item \hyperref[dpa-section-phantom]{Algèbre à puissances divisées}
\item \hyperref[sdga-section-phantom]{Faisceaux différentiels gradués}
\item \hyperref[hypercovering-section-phantom]{Hyperrecouvrements}
\end{enumerate}
Schémas
\begin{enumerate}
\setcounter{enumi}{25}
\item \hyperref[schemes-section-phantom]{Schémas}
\item \hyperref[constructions-section-phantom]{Constructions de schémas}
\item \hyperref[properties-section-phantom]{Propriétés des schémas}
\item \hyperref[morphisms-section-phantom]{Morphismes de schémas}
\item \hyperref[coherent-section-phantom]{Cohomologie des schémas}
\item \hyperref[divisors-section-phantom]{Diviseurs}
\item \hyperref[limits-section-phantom]{Limites de schémas}
\item \hyperref[varieties-section-phantom]{Variétés}
\item \hyperref[topologies-section-phantom]{Topologies sur les schémas}
\item \hyperref[descent-section-phantom]{Descente}
\item \hyperref[perfect-section-phantom]{Catégories dérivées de schémas}
\item \hyperref[more-morphisms-section-phantom]{Compléments sur les morphismes}
\item \hyperref[flat-section-phantom]{Compléments sur la platitude}
\item \hyperref[groupoids-section-phantom]{Schémas en groupoïdes}
\item \hyperref[more-groupoids-section-phantom]{Compléments sur les schémas en groupoïdes}
\item \hyperref[etale-section-phantom]{Morphismes étales de schémas}
\end{enumerate}
Thèmes de théorie des schémas
\begin{enumerate}
\setcounter{enumi}{41}
\item \hyperref[chow-section-phantom]{Homologie de Chow}
\item \hyperref[intersection-section-phantom]{Théorie de l'intersection}
\item \hyperref[pic-section-phantom]{Schémas de Picard des courbes}
\item \hyperref[weil-section-phantom]{Théories cohomologiques de Weil}
\item \hyperref[adequate-section-phantom]{Modules adéquats}
\item \hyperref[dualizing-section-phantom]{Complexes dualisants}
\item \hyperref[duality-section-phantom]{Dualité pour les schémas}
\item \hyperref[discriminant-section-phantom]{Discriminants et différentes}
\item \hyperref[derham-section-phantom]{Cohomologie de de Rham}
\item \hyperref[local-cohomology-section-phantom]{Cohomologie locale}
\item \hyperref[algebraization-section-phantom]{Géométrie algébrique et formelle}
\item \hyperref[curves-section-phantom]{Courbes algébriques}
\item \hyperref[resolve-section-phantom]{Résolution des surfaces}
\item \hyperref[models-section-phantom]{Réduction semi-stable}
\item \hyperref[functors-section-phantom]{Foncteurs et morphismes}
\item \hyperref[equiv-section-phantom]{Catégories dérivées de variétés}
\item \hyperref[pione-section-phantom]{Groupes fondamentaux de schémas}
\item \hyperref[etale-cohomology-section-phantom]{Cohomologie étale}
\item \hyperref[crystalline-section-phantom]{Cohomologie cristalline}
\item \hyperref[proetale-section-phantom]{Cohomologie pro-étale}
\item \hyperref[relative-cycles-section-phantom]{Cycles relatifs}
\item \hyperref[more-etale-section-phantom]{Compléments de cohomologie étale}
\item \hyperref[trace-section-phantom]{La formule des traces}
\end{enumerate}
Espaces algébriques
\begin{enumerate}
\setcounter{enumi}{64}
\item \hyperref[spaces-section-phantom]{Espaces algébriques}
\item \hyperref[spaces-properties-section-phantom]{Propriétés des espaces algébriques}
\item \hyperref[spaces-morphisms-section-phantom]{Morphismes d'espaces algébriques}
\item \hyperref[decent-spaces-section-phantom]{Espaces algébriques décents}
\item \hyperref[spaces-cohomology-section-phantom]{Cohomologie des espaces algébriques}
\item \hyperref[spaces-limits-section-phantom]{Limites d'espaces algébriques}
\item \hyperref[spaces-divisors-section-phantom]{Diviseurs sur les espaces algébriques}
\item \hyperref[spaces-over-fields-section-phantom]{Espaces algébriques sur des corps}
\item \hyperref[spaces-topologies-section-phantom]{Topologies sur les espaces algébriques}
\item \hyperref[spaces-descent-section-phantom]{Descente et espaces algébriques}
\item \hyperref[spaces-perfect-section-phantom]{Catégories dérivées d'espaces}
\item \hyperref[spaces-more-morphisms-section-phantom]{Compléments sur les morphismes d'espaces}
\item \hyperref[spaces-flat-section-phantom]{Platitude sur les espaces algébriques}
\item \hyperref[spaces-groupoids-section-phantom]{Groupoïdes dans les espaces algébriques}
\item \hyperref[spaces-more-groupoids-section-phantom]{Compléments sur les groupoïdes dans les espaces}
\item \hyperref[bootstrap-section-phantom]{Amorçage}
\item \hyperref[spaces-pushouts-section-phantom]{Sommes amalgamées d'espaces algébriques}
\end{enumerate}
Thèmes de géométrie
\begin{enumerate}
\setcounter{enumi}{81}
\item \hyperref[spaces-chow-section-phantom]{Groupes de Chow des espaces}
\item \hyperref[groupoids-quotients-section-phantom]{Quotients de groupoïdes}
\item \hyperref[spaces-more-cohomology-section-phantom]{Compléments sur la cohomologie des espaces}
\item \hyperref[spaces-simplicial-section-phantom]{Espaces simpliciaux}
\item \hyperref[spaces-duality-section-phantom]{Dualité pour les espaces}
\item \hyperref[formal-spaces-section-phantom]{Espaces algébriques formels}
\item \hyperref[restricted-section-phantom]{Algébrisation des espaces formels}
\item \hyperref[spaces-resolve-section-phantom]{Résolution des surfaces revisitée}
\end{enumerate}
Théorie des déformations
\begin{enumerate}
\setcounter{enumi}{89}
\item \hyperref[formal-defos-section-phantom]{Théorie formelle des déformations}
\item \hyperref[defos-section-phantom]{Théorie des déformations}
\item \hyperref[cotangent-section-phantom]{Le complexe cotangent}
\item \hyperref[examples-defos-section-phantom]{Problèmes de déformation}
\end{enumerate}
Champs algébriques
\begin{enumerate}
\setcounter{enumi}{93}
\item \hyperref[algebraic-section-phantom]{Champs algébriques}
\item \hyperref[examples-stacks-section-phantom]{Exemples de champs}
\item \hyperref[stacks-sheaves-section-phantom]{Faisceaux sur les champs algébriques}
\item \hyperref[criteria-section-phantom]{Critères de représentabilité}
\item \hyperref[artin-section-phantom]{Axiomes d'Artin}
\item \hyperref[quot-section-phantom]{Espaces de Quot et de Hilbert}
\item \hyperref[stacks-properties-section-phantom]{Propriétés des champs algébriques}
\item \hyperref[stacks-morphisms-section-phantom]{Morphismes de champs algébriques}
\item \hyperref[stacks-limits-section-phantom]{Limites de champs algébriques}
\item \hyperref[stacks-cohomology-section-phantom]{Cohomologie des champs algébriques}
\item \hyperref[stacks-perfect-section-phantom]{Catégories dérivées de champs}
\item \hyperref[stacks-introduction-section-phantom]{Introduction aux champs algébriques}
\item \hyperref[stacks-more-morphisms-section-phantom]{Compléments sur les morphismes de champs}
\item \hyperref[stacks-geometry-section-phantom]{La géométrie des champs}
\end{enumerate}
Thèmes de théorie des espaces de modules
\begin{enumerate}
\setcounter{enumi}{107}
\item \hyperref[moduli-section-phantom]{Champs de modules}
\item \hyperref[moduli-curves-section-phantom]{Espaces de modules de courbes}
\end{enumerate}
Divers
\begin{enumerate}
\setcounter{enumi}{109}
\item \hyperref[examples-section-phantom]{Exemples}
\item \hyperref[exercises-section-phantom]{Exercices}
\item \hyperref[guide-section-phantom]{Guide bibliographique}
\item \hyperref[desirables-section-phantom]{Desiderata}
\item \hyperref[coding-section-phantom]{Style de programmation}
\item \hyperref[obsolete-section-phantom]{Obsolète}
\item \hyperref[fdl-section-phantom]{Licence de documentation libre GNU}
\item \hyperref[index-section-phantom]{Index généré automatiquement}
\end{enumerate}
\end{multicols}


\bibliography{my}
\bibliographystyle{amsalpha}

\end{document}
